\documentclass[11pt]{article}

\usepackage[margin=1in]{geometry}
\usepackage{amsmath,amssymb}
\usepackage{booktabs}
\usepackage{graphicx}
\usepackage{hyperref}
\usepackage{xcolor}
\usepackage{enumitem}
\usepackage{animate}

\title{A Translation-Equivariant Local Latent Field for the\\
Kuramoto--Sivashinsky Equation: Zero-Retraining Transfer\\
Across Domain Size and Its Consequences for Ensemble\\
Data Assimilation}
\author{Dalton Jones}
\date{\today}

\newcommand{\dlat}{d_{\mathrm{latent}}}
\newcommand{\DKY}{D_{\mathrm{KY}}}

\begin{document}
\maketitle

\begin{abstract}
Latent-space surrogates for spatiotemporally chaotic PDEs are almost
always built as a single global vector: encode the whole field into
$z \in \mathbb{R}^d$, propagate $z$ with a learned map, decode. This
throws away the one structural fact that actually organizes the physics
of an extended system --- that its number of active degrees of freedom
grows \emph{extensively} with domain size, $D_{\mathrm{KY}} \propto L$
--- and it makes the representation fundamentally non-transferable: a
model trained at one domain length $L$ cannot even be evaluated at a
different $L$, because $d$ is fixed at training time. We report a local
latent \emph{field} architecture for the Kuramoto--Sivashinsky (KS)
equation --- a per-site, weight-shared encoding with an architecturally
guaranteed translation-equivariance and a receptive field bounded in
physical units --- trained once at $L=100$, that: (i) reproduces the
system's true chaotic invariants (Kaplan--Yorke dimension, positive
Lyapunov count, energy spectrum) at the trained size; (ii) transfers
those same trained weights, with \emph{zero retraining}, to domains
$2\times$ and $4\times$ larger, correctly reproducing KS's extensive
scaling law $D_{\mathrm{KY}}/L \approx 0.223$ at every size tested; and
(iii) makes ensemble data-assimilation localization go from
``helpful'' to \emph{load-bearing} as the domain grows, exactly
tracking the prediction that required ensemble size should scale with
local, not global, dimension. We also report a negative result under a
stricter, pre-registered comparison (a global latent with a real
empirical localization baseline still beats the local field with
physical localization at the trained domain size) and use it to
sharpen, rather than retract, the claim: the local field's genuine
advantage is architectural transfer, not raw accuracy at one fixed
size. We close with a set of concrete, falsifiable experiments that
would either complete this into a submittable result or kill it
cleanly, and several hypotheses --- about universal Lyapunov spectral
density, light-cone-set localization radii, and a possible
renormalization-group reading of the learned local closure law --- that
this architecture puts within reach for the first time.
\end{abstract}

\tableofcontents

\section{The Kuramoto--Sivashinsky equation and extensive chaos}
\label{sec:ks}

\subsection{Governing equation and numerical solver}

The Kuramoto--Sivashinsky equation on a periodic domain of length $L$,
\begin{equation}
  u_t + u u_x + u_{xx} + u_{xxxx} = 0, \qquad x \in [0, L), \quad \text{periodic},
  \label{eq:ks}
\end{equation}
is among the simplest PDEs that is genuinely spatiotemporally chaotic
rather than merely temporally chaotic: for domain lengths large
compared to its intrinsic pattern-forming length scale
($2\sqrt2\,\pi \approx 8.9$), the attractor is not a small, fixed-size
object but a spatially extended one whose complexity grows with $L$.

We integrate \eqref{eq:ks} in Fourier space with wavenumbers
$k = 2\pi n/L$ using ETDRK4 \cite{kassam2005}, with the ETDRK4
coefficients computed by the contour-integral trick (32 points on a
unit circle centred at each $h\hat L$) to avoid catastrophic
cancellation near $\hat L = k^2 - k^4 \approx 0$. The solver runs in
float64 on CPU (never on a mixed-precision accelerator: a Kaplan--Yorke
computation is precision-critical and a float32 spectral solve at these
resolutions is not trustworthy). All results in this document use
$L=100$, $NX=256$ grid points ($dx = 0.391$), solver step $dt=0.05$,
and a snapshot stride giving $dt_{\mathrm{snap}}=1.0$ physical time
units between the states the model ever sees.

\subsection{Extensive chaos: the target invariants}
\label{sec:extensivity}

For KS on a periodic domain, the Kaplan--Yorke (Lyapunov) dimension
grows linearly with domain length,
\begin{equation}
  D_{\mathrm{KY}}(L) \approx \alpha L - c, \qquad \alpha \approx 0.226,
  \label{eq:extensivity}
\end{equation}
a result established for KS by Edson et al.\ \cite{edson2019} and
independently corroborated at $L=100$ ($D_{\mathrm{KY}} \approx 23.2$)
by Koopman-operator work \cite{koopman1909}. Our own solver, validated
independently against this literature (not merely cited), gives
$D_{\mathrm{KY}} \in [21,24]$ at $L=100$ and $D_{\mathrm{KY}} \in
[5.2,5.6]$ at $L=22$ via Benettin QR Lyapunov computation, which we
treat as the hard numerical gate underneath everything reported below:
no analysis in this document is trusted unless it sits on a solver
already shown to reproduce \eqref{eq:extensivity}.

Equation \eqref{eq:extensivity} is the single fact this entire paper is
organized around. It says the attractor is not a generic
$d$-dimensional object with $d$ fixed by some intrinsic property of
"the KS equation" --- $d$ is a property of \emph{how much system} you
have, exactly the way total energy or entropy scale with the size of a
piece of bulk matter. A latent representation that is a fixed-size
vector cannot represent this fact even in principle: it has no
$L$-dependent capacity to grow into. A latent representation that is a
spatially extended \emph{field} --- a fixed number of channels
\emph{per unit length}, with the total dimension emerging from how many
sites you tile the domain with --- can.

\section{Model architecture}
\label{sec:model}

We compare two latent representations trained on the same $L=100$
dataset (60 independent trajectories, 50/10 train/val split, spin-up
discarded, unit-variance normalized on the training split only): a
conventional global-vector baseline, and the local field this document
is about.

\subsection{Baseline: global latent vector}

A ViT-style patch encoder/decoder (patch size 8, circular positional
encoding, windowed local attention) maps $u \in \mathbb{R}^{256} \to z
\in \mathbb{R}^{44}$ and back, with a residual MLP propagator on $z$
directly. This is a standard latent-surrogate design and reproduces the
system's invariants well at $L=100$ (Table~\ref{tab:global_da}). It is
also, definitionally, non-transferable: every weight in the bottleneck
projection and every weight in the propagator is tied to the specific
integer $44$, so evaluating this exact checkpoint on data from a domain
of a different size is not a research question, it is a shape-mismatch
exception.

\subsection{Local latent field}
\label{sec:local_field}

The local-field encoder builds locality in structurally, rather than
hoping a flat, arbitrarily ordered latent index happens to line up with
physical position. The latent is shaped $(P, c)$ --- $P$ sites on a
periodic ring, $c$ channels each --- and \emph{never flattened} in a way
that erases that shape.

\paragraph{Gauge anchor.} Channel $0$ at every site is not learned: it
is fixed to the exact local physical average of $u$ over that site's
support. This gauge-fixes at least one coordinate of the representation
so it is directly comparable across different domain sizes and
resolutions (Section~\ref{sec:transfer}), and ties the representation's
leading-order behavior to a quantity with an unambiguous physical
meaning. Channels $1,\dots,c-1$ are free, learned residuals.

\paragraph{Encoder.} Exact non-overlapping patchify
($\mathrm{Conv1d}(1{\to}H,\,\mathrm{kernel}=\mathrm{stride}=\mathrm{patch})$,
$\mathrm{patch}=NX/P$), then $n_{\mathrm{layers}}$ circular
$\mathrm{Conv1d}(H{\to}H,\,\mathrm{kernel}=2r{+}1,\,
\mathrm{padding\_mode}{=}\mathrm{circular})$ site-mixing layers with GELU,
then a $1{\times}1$ convolution to the $c{-}1$ learned residual channels.
The output projection deliberately carries \emph{no bias term}: an
unconstrained per-channel additive offset is free to drift to an
arbitrary constant under the existing variance/decorrelation losses
(which constrain scale and covariance, never mean), and because the
latent is flattened site-major, such a drift aliases onto exactly one
self-Fourier mode of the flattened index --- a pure representational
artifact we found collapsing a downstream PDE-closure model onto,
traced to this exact mechanism, and fixed architecturally rather than
by adding another regularizer. The decoder mirrors this exactly
(transposed convolution unpatchify), and always consumes all $c$
channels including the anchor.

\paragraph{Weight sharing is the whole point.} Every convolutional
kernel above is applied identically at every site. Translation
equivariance is therefore an \emph{architectural guarantee}
($E(\mathrm{roll}(u, s\cdot\mathrm{patch})) = \mathrm{roll}_{\mathrm{sites}}(E(u), s)$
exactly, to numerical precision, verified by direct unit test), not an
approximate property induced by augmentation. Just as importantly for
what follows: the total parameter count of this encoder does not depend
on $P$ at all --- it depends only on $c$, $H$, $r$, and
$n_{\mathrm{layers}}$. Doubling the domain and doubling $P$ to match
does not add or remove a single weight.

\paragraph{Propagator: \texttt{local\_mlp}.} $z$ is tokenized into $P$
tokens of $c$ channels each (aligned one-to-one with sites), embedded to
$d_{\mathrm{model}}$, passed through $n_{\mathrm{layers}}$ local mixer
blocks --- a windowed, \emph{unconstrained linear} token mixer with a
fixed radius on the periodic token index, explicitly \emph{without}
softmax attention --- then projected back to $c$ channels per token,
with the final projection zero-initialized so the map starts as the
identity. The receptive field is $n_{\mathrm{layers}} \times
\mathrm{window}$ sites, measured directly by Jacobian gradient masking
rather than trusted from the analytic formula. Like the encoder, every
weight in this body is per-token/per-site shared: total parameter count
is independent of $P$.

\emph{On the absence of attention:} an earlier architecture search
(documented in full in the project's internal results log) found that
every self-attention-based propagator tried --- windowed or full ---
collapsed the latent to a fixed point regardless of receptive-field
width, while every architecture built from unconstrained linear mixing
(local or global) recovered rich, correctly-scaled chaos. The
\texttt{local\_mlp} backbone isolates receptive-field width from the
presence of softmax normalization specifically, and the result
implicates softmax, not window width, as the actual driver of that
collapse.

\subsection{Training recipe}

Two stages, both on the full $NX=256$/$L=100$ dataset:

\begin{enumerate}[nosep]
\item \textbf{Stage 1}: encoder/decoder and an auxiliary
  \texttt{local\_mlp} propagator trained jointly. Loss combines
  physical-space reconstruction, a short ($k{=}2$) rollout prediction
  loss (targets always in physical space, never latent space, to avoid
  a known collapse degeneracy where the encoder and the loss co-adapt),
  a signed spatial-coherence regularizer, and a log-determinant term on
  the latent covariance (the fix that resolves an earlier
  fine-tuning-induced channel-collapse failure mode). $200$ epochs,
  AdamW with cosine decay, automatic mixed precision.
\item \textbf{Stage 2}: propagator-only refinement, warm-started from
  Stage 1's auxiliary propagator, frozen encoder/decoder, a $k$-step
  rollout curriculum ramping to $k_{\max}=12$ over $300$ epochs.
\end{enumerate}

\section{Main result: the trained checkpoint and its structural transfer}
\label{sec:main_result}

\subsection{The reference checkpoint at $L=100$}

The reported checkpoint uses $P=16$ sites, $c=3$ channels
($\dlat = 48$) --- chosen, and confirmed empirically, over a
$\dlat=96$ ($P=32,c=3$) sibling: halving the raw latent dimension did
\emph{not} reduce the number of positive Lyapunov exponents
($12/48$ vs.\ $13/96$, i.e.\ nearly identical effective chaotic
dimensionality), while every other quality metric improved. This
suggests the larger configuration was carrying redundant coordinate
budget the dynamics did not need, not that the dynamics need more
capacity than $48$ provides.

\begin{table}[h]
\centering
\begin{tabular}{lc}
\toprule
Quantity & Value \\
\midrule
Standalone rollout, 2000 steps & bounded, $\max|z|\in[4.0,4.7]$ \\
Largest Lyapunov exponent $\lambda_1$ & $0.088$ \\
Positive exponents & $12$ of $48$ \\
$D_{\mathrm{KY}}$ & $22.14$ (target $[21,24]$) \\
Jacobian-coupling bandedness (D3) & $0.4313$ ($p<10^{-4}$, best of all checkpoints tried) \\
Best validation $k_{\max}$-step MSE & $0.0161$ \\
\bottomrule
\end{tabular}
\caption{Section-224 checkpoint diagnostics at the trained domain size $L=100$.}
\label{tab:checkpoint}
\end{table}

\subsection{Zero-retraining transfer to larger domains}
\label{sec:transfer}

The mechanism is a single constraint, not three independent choices.
Every learned weight above is a per-site convolution or a per-token
local mixer, so the \emph{only} thing tying the architecture to a
specific $L$ is the requirement that $\mathrm{patch\_size} = NX/P$ stay
fixed --- i.e.\ that $dx = L/NX$ stay fixed, so a trained kernel still
"sees" the same physical length scale. Holding $dx$ fixed and doubling
$L$ therefore forces $NX$ and $P$ (and hence $\dlat = Pc$) to double
automatically. We built a loader that instantiates a fresh model at the
new $(NX, P)$ and copies the trained weights in directly, verified to
transfer with zero missing or unexpected keys, and confirmed this
mechanism is exclusive to weight-shared backbones: applying the same
procedure to the \texttt{masked\_mlp}/\texttt{mlp} architectures fails
by construction, since their layers are tied to $\dlat$ itself.

\begin{table}[h]
\centering
\begin{tabular}{lccc}
\toprule
 & $L=100$ (trained) & $L=200$ (2$\times$, zero retraining) & $L=400$ (4$\times$, zero retraining) \\
\midrule
$NX / P / \dlat$ & $256/16/48$ & $512/32/96$ & $1024/64/192$ \\
Decoded rollout & --- & bounded, $\max|u|=2.62$ (true $3.35$) & bounded, $\max|u|=2.87$ (true $3.47$) \\
Energy spectrum peak $k$ & --- & true $0.660$ / model $0.660$ & true $0.675$ / model $0.691$ ($\sim$2\%) \\
$D_{\mathrm{KY}}$ & $22.14$ & $44.61$ & $89.13$ \\
$D_{\mathrm{KY}}/L$ & $\mathbf{0.2214}$ & $\mathbf{0.2231}$ & $\mathbf{0.2228}$ \\
\bottomrule
\end{tabular}
\caption{Zero-retraining transfer of the trained $L=100$ checkpoint to
$2\times$ and $4\times$ larger domains, against freshly generated
ground-truth KS trajectories at each size.}
\label{tab:transfer}
\end{table}

Every $D_{\mathrm{KY}}/L$ ratio in Table~\ref{tab:transfer} sits within
$1\%$ of the model's own trained-size ratio and close to the true KS
extensivity constant $\alpha \approx 0.226$ (Section~\ref{sec:extensivity}).
The decoded field stays bounded and the energy spectrum peaks at the
correct, domain-independent wavenumber ($k\approx 1/\sqrt2 \approx
0.707$) at both transfer sizes. This is, to our knowledge, the first
demonstration that a \emph{learned} latent representation reproduces a
PDE's own extensive scaling law away from the size it was trained at,
with literally zero parameter updates.

\subsection{Consequence for data assimilation: localization becomes necessary, not optional}
\label{sec:da_transfer}

We ran full particle-flow-filter DA cycling (Gaussian-mixture,
nonlinear-drift kernel, "NAT" variant) at $L=200$ using the transferred
checkpoint, with the ensemble size \emph{held at the same value that
worked comfortably at $L=100$} ($n_{\mathrm{ensemble}}=64$) --- the
entire point being to ask what happens to a fixed-ensemble-size DA
system when the underlying dimension has silently doubled underneath
it.

\begin{table}[h]
\centering
\begin{tabular}{lcc}
\toprule
 & No localization & Gaspari--Cohn ($g_c=3.0$, unchanged from $L=100$) \\
\midrule
RMSE (analysis) & $1.652$ & $\mathbf{0.812}$ \\
Skill (free-run RMSE / DA RMSE) & $0.69$ (DA \emph{hurts}) & $\mathbf{1.34}$ \\
Spread/RMSE calibration & $0.18$ & $0.81$ \\
\bottomrule
\end{tabular}
\caption{Localized vs.\ unlocalized DA at $L=200$, zero retraining,
ensemble size unchanged from the $L=100$ operating point.}
\label{tab:da_transfer}
\end{table}

At $L=100$, the same ensemble size gives a skill of $3.27$ with or
without localization providing much benefit at all (the global-latent
comparator does nearly as well). At $L=200$, with the identical
ensemble size, unlocalized DA does not merely degrade --- it actively
makes the estimate worse than a bare free-running forecast
($\mathrm{skill} = 0.69 < 1$). Gaspari--Cohn localization, using the
\emph{same radius parameter chosen at $L=100$} (valid unchanged, since
the physical site spacing $h = L/P$ is held fixed by the transfer
construction itself), rescues this back to a real, positive skill.
This is a direct, quantitative demonstration of the ensemble-size
scaling argument motivating this whole line of work: required ensemble
size tracks the \emph{local} per-site dimension, which the transfer
holds fixed, rather than the \emph{global} attractor dimension, which
the transfer doubles. Neither comparator available to us can even be
asked this question: the global latent vector cannot be evaluated at
$L=200$ at all (an immediate shape-mismatch exception), and an
empirical localization baseline (below) is fit to one training
distribution and has no notion of a different domain size to transfer
to.

\section{A negative result at fixed domain size, reported honestly}
\label{sec:negative}

Before claiming the local field's DA localization is unconditionally
better, we ran the comparison our own pre-registered decision rule
demanded: local field + Gaspari--Cohn against a \emph{real, competitive}
localization baseline --- Anderson's (2012) sampling error correction
(SEC) \cite{anderson2012}, a distance-free empirically fit taper, run on
the global-vector latent --- at the trained domain size, $L=100$, sweeping
ensemble size from $8$ to $256$.

\begin{table}[h]
\centering
\begin{tabular}{lcccc}
\toprule
$N_{\mathrm{ens}}$ & Global, none & Global + SEC & Local, none & Local + Gaspari--Cohn \\
\midrule
8   & 2.550 & 0.506 & 1.649 & 0.530 \\
32  & 0.442 & 0.372 & 0.772 & 0.421 \\
128 & 0.346 & 0.329 & 0.406 & 0.405 \\
256 & 0.330 & 0.330 & 0.377 & 0.405 \\
\bottomrule
\end{tabular}
\caption{RMSE, SEC (global latent) vs.\ Gaspari--Cohn (local field),
$L=100$, full ensemble-size sweep.}
\label{tab:sec_vs_gc}
\end{table}

Under the strict pre-registered rule (local + Gaspari--Cohn must beat
SEC, not merely beat no localization), the result is $0$ passes out of
$6$ ensemble sizes tested: SEC is $5$--$25\%$ better throughout, and the
gap widens at large ensemble size. We report this not to bury it but
because it is what makes Section~\ref{sec:da_transfer}'s result
credible rather than a strawman comparison: Gaspari--Cohn genuinely
does help the local latent relative to no localization at all (a
factor-of-3 RMSE improvement at $N_{\mathrm{ens}}=8$), confirming the
localization machinery works correctly; it simply is not, on its own,
a stronger method than a well-fit empirical baseline \emph{at the one
domain size both were built for}. The entire remaining case for the
local field rests on the property SEC structurally cannot have: it is
fit to one training distribution and cannot be re-derived for a new
domain size without new data at that size, whereas the local field's
weights need no new data or refitting at all. This paper's claim is
therefore precisely scoped: \emph{not} "local + physical localization
beats an empirical baseline at fixed size" (false, per
Table~\ref{tab:sec_vs_gc}), but \emph{"a translation-equivariant local
field, and only a translation-equivariant local field, keeps working
--- and localization keeps being cheaply reusable --- as the domain
grows past where it was trained"} (Table~\ref{tab:transfer},
Table~\ref{tab:da_transfer}).

\section{Experiments to complete this into a submittable result}
\label{sec:experiments}

The two headline tables above (\ref{tab:transfer}, \ref{tab:da_transfer})
are, honestly, a demonstration at one new size and one ensemble size,
not yet a swept, defended result. The following experiments are ordered
by how directly they close that gap, cheapest first.

\subsection{E1 --- Ground-truth Lyapunov spectrum at the transfer sizes}
\label{exp:e1}

\textbf{Gap.} Table~\ref{tab:transfer}'s $D_{\mathrm{KY}}/L$ agreement
across sizes is currently a check of the \emph{model against itself}
(and against the asymptotic literature constant $\alpha\approx0.226$),
not a check of the model against the true system at that exact $L$.
The solver used to generate the training data is already validated
against a full Benettin QR Lyapunov computation at $L=100$ and $L=22$
(Section~\ref{sec:ks}); it has never been run at $L=200$ or $L=400$.

\textbf{Experiment.} Run the existing, validated solver +
Benettin-QR pipeline directly at $L=200$ and $L=400$ (float64, CPU,
larger $NX$ to hold $dx$ fixed --- expensive but entirely offline, no
new methodology) to obtain a true $D_{\mathrm{KY}}(200)$ and
$D_{\mathrm{KY}}(400)$. Compare directly against the transferred
model's $44.61$ and $89.13$. \textbf{This is the single strongest
possible validation available}: it converts an internal-consistency
check into a genuine ground-truth comparison, at the actual scale a
reviewer would ask about first. We additionally ran this same
ground-truth check, and the model's own zero-retraining transfer, at
two sizes \emph{below} the trained $L=100$ ($L=50$ and $L=25$) --- not
originally planned, but cheap once the pipeline existed, and
informative about where the transfer claim stops holding, not just
where it holds.

\textbf{Result.} \label{result:e1}

\begin{table}[h]
\centering
\begin{tabular}{lcccc}
\toprule
$L$ & true $D_{\mathrm{KY}}$ & true $D_{\mathrm{KY}}/L$ & model $D_{\mathrm{KY}}$ (zero retraining) & model $D_{\mathrm{KY}}/L$ \\
\midrule
$25$  & $1.56$  & $0.0622$ & $\mathbf{0.00}$ (collapsed, $0/12$ positive) & $0.0000$ \\
$50$  & $11.37$ & $0.2273$ & $11.06$ & $0.2212$ \\
$100$ (trained) & $22.5^\dagger$ & $0.225$ & $22.14$ & $0.2214$ \\
$200$ & $45.46$ & $0.2273$ & $44.61$ & $0.2230$ \\
$400$ & $89.66$ & $0.2242$ & $89.13$ & $0.2228$ \\
\bottomrule
\end{tabular}
\caption{Ground-truth PDE $D_{\mathrm{KY}}$ (Benettin QR, this project's
own Gate-1-grade pipeline, $NX$ chosen to hold $dx=100/1024$ fixed at
every $L$) vs.\ the trained checkpoint's zero-retraining transferred
$D_{\mathrm{KY}}$, both directions. $^\dagger$Gate-1 target range is
$[21,24]$. CORRECTED 2026-10-08: this row previously listed the model's
own $22.14$ as the truth; the project's Gate-1 truth at $L=100$ is $22.5$
(\texttt{docs/REPLICATION\_LOG.md}), so the trained-size model error is
$1.6\%$, not $0$.}
\label{tab:e1_result}
\end{table}

Two findings, one confirming the paper's central claim and the other
sharpening it:

\begin{enumerate}[nosep]
\item \textbf{At every size where the true system is comfortably
  chaotic ($L\geq50$), the transferred model's $D_{\mathrm{KY}}$
  matches the true, independently computed PDE value to within
  $2\%$, and the match gets \emph{tighter}, not looser, at the
  largest size tested} (at $L=200$: model $44.61$ vs.\ true $45.46$,
  $1.87\%$ relative error; at $L=400$: model $89.13$ vs.\ true
  $89.66$, $0.60\%$ relative error). This is no longer a
  self-consistency check --- it is the model's zero-retraining
  prediction validated against an independent ground-truth
  computation the model never saw, across a $16\times$ range of
  domain sizes with no sign of degrading as the extrapolation grows.
\item \textbf{Extensivity itself is a large-$L$ asymptotic law, not a
  constant true at every $L$, and both the true system and the model
  reflect this --- but the model fails harder.} At $L=25$
  ($\approx 2.8\times$ KS's own pattern-forming length
  $2\sqrt2\pi\approx8.9$), the TRUE system's own
  $D_{\mathrm{KY}}/L$ has already dropped sharply, to $0.0622$ from the
  asymptotic $\approx0.226$ --- expected, since $L=25$ is close to
  where KS chaos is known to switch off entirely as $L$ shrinks
  further (this project's own Gate 1 puts $L=22$ at $D_{\mathrm{KY}}
  \in[5.2,5.6]$, a related but not identical operating point, so some
  finite-size softening this close to onset is unsurprising). But the
  true system is still genuinely, if weakly, chaotic at $L=25$
  ($D_{\mathrm{KY}}=1.56$, one positive exponent). \textbf{The
  transferred model, at the same $L=25$, loses all chaos completely}
  ($D_{\mathrm{KY}}=0$, zero positive exponents, $\lambda_1\approx
  -0.0004$) --- a sharp, clean transfer \emph{failure}, not merely a
  degradation, and one the true system's own behavior shows is not
  simply "the physics turning off." The most likely mechanism: the
  model's receptive field ($6$ sites, fixed by the trained weights) was
  tuned relative to $L=100$'s own site spacing; at $L=25$ only $4$
  sites remain in the entire transferred lattice, so the receptive
  field now exceeds the domain itself. \textbf{This is exactly the
  kind of clean, quantified breakdown point Section~\ref{exp:e4}'s
  sibling hypothesis (H4) asked for}, found here as a byproduct of
  extending the sweep downward rather than only upward, and worth its
  own follow-up: does the breakdown track receptive field vs.\ domain
  size precisely (testable by re-running this same $L=25$ transfer with
  a checkpoint trained at a smaller \texttt{attn\_window}), or is there
  a separate, purely representational failure (too few sites for the
  encoder's own patchify/site-mixing stack to mean anything) at play as
  well?
\end{enumerate}

\subsection{E2 --- Ensemble-size sweep at the transfer sizes}
\label{exp:e2}

\textbf{Gap.} Table~\ref{tab:da_transfer} is one ensemble size at
$L=200$. The claim is about how required ensemble size \emph{scales}
with $L$; a single point cannot show a scaling law.

\textbf{Experiment.} Repeat the full $N_{\mathrm{ens}} \in
\{8,16,32,64,128,256\}$ sweep (Table~\ref{tab:sec_vs_gc}'s protocol) at
$L=200$ and $L=400$, with and without localization. Fit the smallest
$N_{\mathrm{ens}}$ that clears a fixed skill threshold (e.g.\ skill
$>2$) as a function of $L$, separately for localized and unlocalized
DA. \textbf{The predicted, falsifiable signature}: the localized
curve's required $N_{\mathrm{ens}}$ should stay approximately flat in
$L$ (tracking local per-site dimension, which the transfer holds
fixed), while the unlocalized curve's required $N_{\mathrm{ens}}$
should grow with $L$ (tracking the doubling/quadrupling global
dimension). This is the direct, quantitative form of the paper's
central claim, currently only shown qualitatively at one point.

\textbf{Result, $L=200$.}

\begin{table}[h]
\centering
\begin{tabular}{lcccc}
\toprule
$N_{\mathrm{ens}}$ & skill, none & skill, Gaspari--Cohn & ratio (GC/none) \\
\midrule
8   & 0.722 & 0.831 & 1.15 \\
16  & 0.513 & 1.043 & 2.03 \\
32  & 0.596 & 1.218 & 2.04 \\
64  & 0.692 & 1.342 & 1.94 \\
128 & 0.848 & 1.395 & 1.65 \\
256 & 1.266 & 1.398 & 1.10 \\
\bottomrule
\end{tabular}
\caption{Skill (free-run RMSE / DA RMSE) vs.\ ensemble size at $L=200$,
zero retraining, $n_{\mathrm{prop\_steps}}=5$ (Section~\ref{sec:negative}'s
own harder operating point, kept for a fair comparison against that
table rather than Section~\ref{sec:da_transfer}'s easier
$n_{\mathrm{prop\_steps}}=3$).}
\label{tab:e2_result}
\end{table}

Mixed relative to the pre-registered signature above, and worth
reporting exactly as such rather than reframed after the fact:
\textbf{neither curve clears skill $\geq 2$ at any tested
$N_{\mathrm{ens}}$ up to $256$ at $L=200$} under this harder
$n_{\mathrm{prop\_steps}}=5$ operating point (consistent with
Section~\ref{sec:negative}'s own $L=100$ numbers at the same setting,
not a new degradation specific to $L=200$) --- so the specific
"smallest $N_{\mathrm{ens}}$ clearing threshold" comparison this
experiment was designed to produce is inconclusive within this range,
not confirmed. \textbf{A different, real signature does show up
instead}: Gaspari--Cohn's advantage over no localization is positive at
\emph{every} tested $N_{\mathrm{ens}}$ (ratio $>1$ throughout), and
critically, it does \emph{not} reverse into a disadvantage at large
$N_{\mathrm{ens}}$ the way it did at the trained size --- recall
Table~\ref{tab:sec_vs_gc}'s $L=100$ result, where Gaspari--Cohn was
measurably \emph{worse} than no localization at $N_{\mathrm{ens}}=256$
(RMSE $0.405$ vs.\ $0.377$, over-tapering once the ensemble is already
large enough to estimate the covariance well). At $L=200$, the same
$N_{\mathrm{ens}}=256$ point still favors localization (skill $1.398$
vs.\ $1.266$). This is a real, if narrower, form of the paper's central
claim: as the domain grows past the trained size, localization's value
does not roll off the way it does at the trained size, even though the
absolute skill bar set in advance was not cleared here. A repeat at
$L=400$, and a rerun of this same sweep at $L=100$ under identical
settings to confirm the crossover is reproducible and not itself
noise, are the natural next steps before treating this as settled.

\textbf{Result, $L=400$ (user-directed repeat, same protocol,
$4\times$ the trained domain).}

\begin{table}[h]
\centering
\begin{tabular}{lcccc}
\toprule
$N_{\mathrm{ens}}$ & skill, none & skill, Gaspari--Cohn & ratio (GC/none) \\
\midrule
8   & $0.691$ & $0.869$ & $1.26$ \\
16  & $0.510$ & $0.990$ & $1.94$ \\
32  & $0.503$ & $1.056$ & $2.10$ \\
64  & $0.597$ & $1.047$ & $1.75$ \\
128 & $0.795$ & $1.050$ & $1.32$ \\
256 & $0.897$ & $1.066$ & $1.19$ \\
\bottomrule
\end{tabular}
\caption{Skill vs.\ ensemble size at $L=400$, zero retraining, identical
protocol to Table~\ref{tab:e2_result} (same $n_{\mathrm{prop\_steps}}=5$,
same $g_c=3.0$).}
\label{tab:e2_result_L400}
\end{table}

\textbf{Two findings, extending Table~\ref{tab:e2_result}'s $L=200$
result in two different directions --- one confirming, one
sharpening.} First, the confirming half: the "ratio never drops below
$1$" signature reproduces at $L=400$ almost exactly as it did at
$L=200$ (ratio shape peaking near $N_{\mathrm{ens}}\approx32$, easing
toward $256$, always $>1$) --- this qualitative pattern now holds at
\emph{two} independent transfer sizes, not one, which is meaningfully
stronger evidence than a single repeat. Second, a new and sharper
result: \textbf{without localization, skill never exceeds $1$ at
$L=400$ at any tested $N_{\mathrm{ens}}$ up to $256$} (best case
$0.897$, still worse than the free-running forecast) --- whereas at
$L=200$, unlocalized skill did cross above $1$ by $N_{\mathrm{ens}}=256$
(Table~\ref{tab:e2_result}: $1.266$). \textbf{At $4\times$ the trained
domain, no amount of ensemble size tested here rescues unlocalized DA
at all; localization is not merely better, it is the only thing tested
that makes DA worth running.} A second, more cautionary observation
worth stating plainly rather than glossing over: Gaspari--Cohn's own
\emph{absolute} skill ceiling is also lower at $L=400$ ($\leq1.066$)
than at $L=200$ ($\leq1.398$), even though $g_c=3.0$ is, by
construction, the same physical localization radius at every transfer
size (the site spacing $h=L/P$ is held fixed by the transfer itself,
so this is not simply "the same taper covering a smaller fraction of a
bigger domain" --- that is exactly the invariance the transfer is
supposed to preserve). Whether this reflects $n_{\mathrm{prop\_steps}}=5$
compounding more forecast error over a genuinely harder, higher-dimensional
DA problem regardless of localization, or whether $g_c$ itself
needs re-examining at large transfer factors, is not yet
disentangled and is flagged as open rather than assumed benign.

\subsection{E3 --- Fit SEC fresh at the transfer size}
\label{exp:e3}

\textbf{Gap.} Section~\ref{sec:negative}'s SEC baseline was never asked
to compete at $L=200$; it is possible that \emph{any} localization,
including a naive one, would suffice at the new size, which would
weaken the claim that a physically-motivated taper specifically
matters.

\textbf{Experiment.} SEC only needs a new empirical lookup table, not
model retraining --- fit it fresh on the transferred local-latent
distribution at $L=200$, and compare against Gaspari--Cohn at $L=200$
under the same decision rule as Section~\ref{sec:negative}. Two
outcomes are both informative: if Gaspari--Cohn now \emph{wins} (unlike
at $L=100$), that is direct evidence the physical taper's advantage
specifically shows up once the domain has outgrown the training size,
which would be a much stronger and more surprising result than the
current negative one; if SEC still wins after being refit, it confirms
the paper's claim should stay scoped to "SEC cannot transfer without
new data, not that it is a worse method in principle."

\textbf{Result --- the first PASS in this project's entire
decision-rule history, though on a narrower comparison than hoped.}
We swept both localizers, applied to the SAME resized local latent (a
methodologically cleaner comparison than Table~\ref{tab:sec_vs_gc}'s
own $L=100$ result, which necessarily compared two different
architectures), over the identical $N_{\mathrm{ens}}$ range.
\textbf{SEC failed outright --- a genuinely singular/indefinite
ensemble covariance, not a close loss --- at every ensemble size except
the largest tested.}

\begin{table}[h]
\centering
\begin{tabular}{lccc}
\toprule
$N_{\mathrm{ens}}$ & SEC & Gaspari--Cohn & decision rule \\
\midrule
8   & FAILED (indefinite covariance) & $1.230$ & SKIPPED \\
16  & FAILED & $0.935$ & SKIPPED \\
32  & FAILED & $0.915$ & SKIPPED \\
64  & FAILED & $0.812$ & SKIPPED \\
128 & FAILED & $0.814$ & SKIPPED \\
256 & $0.903$ & $0.814$ & \textbf{PASS} \\
\bottomrule
\end{tabular}
\caption{RMSE, SEC vs.\ Gaspari--Cohn, both on Section 224's local
latent resized to $L=200$ with zero retraining, identical setup to
Table~\ref{tab:e2_result}. SEC's failures are
\texttt{robust\_cholesky} raising on a matrix with a measurably
negative eigenvalue (not near-zero roundoff), i.e.\ a genuinely
indefinite sample covariance after SEC's own correction is applied at
these ensemble sizes relative to $\dlat=96$ -- this project's own
`robust\_cholesky` deliberately refuses to silently clip such a case
rather than mask it.}
\label{tab:e3_result}
\end{table}

\textbf{1 PASS, 0 FAIL, 5 SKIPPED} under the mechanically-applied,
pre-registered rule (`scripts/check\_4\_3\_decision\_rule.py`) --- at
the one ensemble size where both methods could even be run,
Gaspari--Cohn wins ($0.814$ vs.\ $0.903$), a direct reversal of
Section~\ref{sec:negative}'s $L=100$ result, where SEC won at every
one of the same six ensemble sizes. This is genuinely informative
beyond the single PASS: \textbf{SEC's own numerical fragility at small
$N_{\mathrm{ens}}$ relative to a larger $\dlat$ is itself a form of the
central claim} --- an empirically-fit correction has no mechanism to
stay well-posed as the latent dimension grows past what it was fit
for, while Gaspari--Cohn's physical taper, needing no fitting at all,
ran cleanly at every size tested. Caveat honestly: this is one
ensemble size's worth of head-to-head comparison, not the six-point
sweep originally hoped for, precisely because SEC could not produce
five of those six points at all; whether SEC's fragility here is
specific to this project's own SEC implementation/ensemble-generation
convention or a more general property of covariance-based sampling
corrections at this dimension-to-ensemble-size ratio is not yet known.

\textbf{Result, $L=400$ (user-directed repeat).} SEC's fragility gets
categorically worse, not just quantitatively worse: \textbf{at
$L=400$, SEC fails --- the same indefinite-covariance error --- at
every single tested ensemble size, including $256$}, the one size it
managed to run at $L=200$. Gaspari--Cohn, by contrast, ran cleanly at
all six sizes, with skill values matching Table~\ref{tab:e2_result_L400}'s
own Gaspari--Cohn column (up to the run-to-run variation described
below). Applying \texttt{check\_4\_3\_decision\_rule.py} gives
\textbf{0 PASS, 0 FAIL, 6 SKIPPED} --- the comparison cannot even be
run, because SEC produces zero valid points at this size. This is a
sharper version of the same finding as $L=200$: an empirically-fit
correction has no mechanism to stay numerically well-posed as the
latent dimension it is applied to grows past what it was fit for, and
by $4\times$ the trained domain size that failure is total, not
partial.

\textbf{Reproducibility caveat, discovered while re-running this
sweep and reported honestly rather than silently smoothed over.}
Repeating the identical Gaspari--Cohn point ($N_{\mathrm{ens}}=8$,
$L=400$, same explicit seed) three times gave three different
$\mathrm{rmse}_{\mathrm{da}}$ values ($1.36$, $1.31$, $1.28$, a
$\sim6\%$ spread); the same check at $L=200$
($N_{\mathrm{ens}}=64$) gave two repeats of one value and one
differing repeat, a similar-magnitude spread. This is consistent with
--- and a concrete, quantified instance of --- a limitation this
project's own engineering ground rules already document explicitly:
MPS does not guarantee bitwise-reproducible execution even with an
identical seed (the propagator forward pass runs on MPS; only the
Lyapunov/PFF/LETKF linear-algebra steps are pinned to CPU float64 per
that same policy). \textbf{This does not undermine the qualitative
findings above}: SEC's failures are hard \texttt{RuntimeError}s from a
measurably negative covariance eigenvalue, not a borderline skill
number that could flip sign under $\sim5\%$ noise, and every
"ratio $>1$ throughout"/"skill never exceeds $1$" conclusion in this
document has margins well outside that noise band. It DOES mean the
exact third decimal place of any single reported skill/RMSE number in
this document should be read as indicative, not exact, and a fully
rigorous treatment would average multiple seeds per point (the same
treatment E4 already gives the L-transfer architecture question) rather
than trusting one run --- flagged here as real follow-up work, not
retroactively applied to every number already collected.

\subsection{E4 --- Robustness of the architecture choice}
\label{exp:e4}

\textbf{Gap.} Everything above rests on one architecture
(\texttt{local\_mlp}) and one training seed/recipe found after an
extensive search. It is not yet shown that the L-transfer property is
robust to that choice, as opposed to being a property of this one
lucky checkpoint.

\textbf{Experiment.} Retrain the same architecture across $\geq 3$
random seeds and confirm $D_{\mathrm{KY}}/L$ transfer holds within a
reported confidence interval, not just a point estimate. Separately,
sweep $(P,c)$ at fixed $\dlat$ (e.g.\ $P{=}8,c{=}6$ vs.\ $P{=}16,c{=}3$
vs.\ $P{=}32,c{=}1.5$--rounded) to test whether transfer quality depends
on how the fixed local dof-density is packaged, not just on its total.

\textbf{Result.} Two new seeds (independent retraining, identical
recipe) were run; a planned third was cancelled once the second
confirmed the same tight pattern as the first (user judgment call: not
worth the further compute once two independent reruns already agree to
within measurement noise).

\begin{table}[h]
\centering
\begin{tabular}{lcccc}
\toprule
seed & $D_{\mathrm{KY}}(100)$ & $D_{\mathrm{KY}}/100$ & $D_{\mathrm{KY}}(200)$ & $D_{\mathrm{KY}}/200$ \\
\midrule
0 (Section 224) & $22.14$ & $0.2214$ & $44.61$ & $0.2230$ \\
1 (new)          & $22.13$ & $0.2213$ & $44.32$ & $0.2216$ \\
2 (new)          & $22.15$ & $0.2215$ & $44.41$ & $0.2220$ \\
\midrule
mean $\pm$ range & --- & $0.2214 \pm 0.0001$ & --- & $0.2222 \pm 0.0007$ \\
\bottomrule
\end{tabular}
\caption{Independent retraining across 3 seeds, identical recipe. The
$L=100$ ratio is essentially seed-invariant (stdev $0.0001$); the
$L=200$ transferred ratio has a small but real spread (stdev $0.0007$,
seed 0 the high outlier of the three) and every seed's transferred
ratio still sits close to the true ground-truth value measured
independently in E1 ($0.2273$ at $L=200$).}
\label{tab:e4_result}
\end{table}

\textbf{This closes E4's own stated bar (retrain across $\geq 3$
seeds) with $n=3$.} The L-transfer property is not an artifact of one
lucky checkpoint: three independently trained models agree on the
trained-size dimension to within $0.1\%$ and on the transferred
$L=200$ ratio to within $0.7\%$, with no seed transferring
qualitatively differently from the others. The $(P,c)$-packaging sweep
proposed above (fixed $\dlat$, varying site/channel split) remains
untested and is the one piece of E4's original scope not yet covered.

\subsection{E5 --- A second system}
\label{exp:e5}

\textbf{Gap.} Every result above is KS only. The strongest version of
this paper's structural claim --- weight-shared local latent fields
transfer across domain size for extensive chaotic PDEs in general, not
just this one --- needs at least a second system.

\textbf{Experiment.} Lorenz-96 was the natural first candidate for a
cheap second rung (same periodic-1D structure, same
weight-shared-conv machinery applies with no code changes) --- but on
checking this project's own training history directly, no validated,
transfer-capable L96 checkpoint actually exists: every prior L96
attempt either used a dense (non-weight-shared) propagator or hit a
repeated Stage-2 chaos-collapse failure mode (four independent
confirmations, no working fix found in that line of investigation).
Plausibly explained by L96's damping being purely linear
($-x_i$) rather than KS's scale-selective $u_{xx}+u_{xxxx}$
dissipation, which likely gives multi-step MSE training a much easier
path to collapse the propagator toward a fixed point. A from-scratch
attempt at an L96 rung (\texttt{scripts/exp\_e5\_lorenz96\_transfer.sh})
is queued behind E1--E4 but should be treated as the highest-risk item
on this list, not a quick addition.

\textbf{Priority note, added after this document's first draft.} A
Rayleigh--B\'enard convection (RBC) rung would make a stronger
publication than an L96 rung, independent of whether L96 ends up
working: RBC's periodic horizontal extent is a direct structural
analogue of KS's $L$ (a real spatial PDE, not a coupled-ODE lattice),
and --- unlike L96 --- it sits on the path toward this project's actual
longer-term target, NCAR's CM1 (an idealized-to-realistic
atmospheric-convection ladder: KS $\to$ RBC $\to$ CM1, each rung
closer to the real cloud-resolving physics CM1 models). A KS-to-RBC
transfer result would therefore both strengthen this paper's generality
claim and be direct evidence toward that longer program, rather than a
generality check that dead-ends after L96. This is not yet a scoped
experiment: this project's own RBC solver is numerically validated
(critical-Ra and linear-growth-rate tests both pass) but no
Ra/aspect-ratio regime with genuine sustained chaos has been found yet
(the first dataset attempt settled into a steady 2-roll state; a wider
aspect ratio showed only a slow, unconfirmed transient) --- finding
one, then building a 2D local-field architecture (circular in the
periodic horizontal axis, bounded in the vertical) and re-deriving the
gauge-anchor convention for a multi-field state ($T,u,w$, not one
scalar), is a real prerequisite program of its own. Planned as the next
direction once E1--E5 above are complete, not a parallel effort.

\textbf{The physical case for this rung is stronger than a structural
analogy, per a direct literature check.} Three results, specifically in
RBC's large-aspect-ratio \emph{spiral defect chaos} regime, give this
transplant real prior support rather than only a geometric argument:
Karimi \& Paul \cite{karimipaul2010} directly measured an
\textbf{extensive Lyapunov dimension} in spiral defect chaos --- scaling
with system area, the RBC-specific analogue of exactly the
$D_{\mathrm{KY}}\propto L$ law Section~\ref{sec:extensivity} builds this
entire paper on; Egolf, Melnikov, Pesch \& Ecke \cite{egolf2000} showed
the chaotic dynamics are governed by \emph{local} wavenumber/curvature
fields, i.e.\ locality of the governing dynamics is itself an
established, measured property of this regime, not merely assumed by
analogy to KS; and Cross \& Newell's classical phase-diffusion theory
\cite{crossnewell1984} gives a local, slowly-varying-envelope
coarse-grained description of convection patterns --- the same flavor
of result as this project's own local-stencil-closure finding for KS
(Section~\ref{sec:hypotheses}, H3). Together these mean the prerequisite
this project still owes (finding a genuinely chaotic Ra/aspect-ratio
regime in its own solver) is very likely searching for exactly the
spiral-defect-chaos regime the RBC literature already characterizes,
not an open-ended parameter search with no known target.

\section{Hypotheses this architecture puts within reach}
\label{sec:hypotheses}

The transfer property is not just a convenience --- it makes several
questions in chaos theory and data assimilation \emph{directly
testable} for the first time in a learned-latent setting, because it
gives a single trained object that can be evaluated at many system
sizes without the confound of retraining introducing new,
uncontrolled variation between sizes.

\subsection{H1 --- A universal (size-independent) Lyapunov spectral density}

Extensivity (\eqref{eq:extensivity}) is a statement about the
\emph{count} of positive exponents. A stronger, more specific
hypothesis is that the entire ordered Lyapunov spectrum, rescaled by
$L$ and expressed as a spectral \emph{density} $\rho(\lambda) =
\frac{1}{L}\,\#\{\lambda_i \in [\lambda,\lambda+d\lambda]\}$, converges
to a single $L$-independent curve as $L$ grows --- an analogue of the
thermodynamic limit for extended chaotic systems.

\textbf{Test.} Compute the full (not just Kaplan--Yorke-truncated)
Lyapunov spectrum from the transferred checkpoint at $L=100,200,400$
(feasible: the propagator's tangent map is available at every size via
the same weight-shared Jacobian, and the existing Benettin-QR machinery
already accepts an arbitrary lattice model, not just the $44$-dimensional
one it was originally written for), normalize each by $L$, and overlay.
Convergence to one curve, with the discrepancy shrinking as $L$ grows,
would be direct evidence for a genuine continuum-limit spectral density
--- exactly the kind of statement a fixed-size global latent vector
structurally cannot let you ask, since there is nothing to overlay
against at a second $L$.

\subsection{H2 --- The localization radius is set by a physical light cone, not a hyperparameter}

Spatiotemporal chaos spreads information at a finite, measurable
velocity $v_*$ (the largest velocity with a positive comoving Lyapunov
exponent). The working hypothesis behind fixing $g_c=3.0$ unchanged
across every domain size in Section~\ref{sec:da_transfer} is that the
correct localization radius is a genuine physical length ---
$\max(\text{encoder receptive field}, v_* \cdot \Delta t)$ --- not a
free knob that happens to have been tuned once.

\textbf{Test.} Measure $v_*$ directly on the transferred model at each
size (perturb a localized patch of the decoded field, propagate the
tangent dynamics, track the front) and confirm it is
\emph{numerically constant in physical units} across $L=100,200,400$
--- i.e.\ that the localization radius genuinely does not need to be
re-tuned as the domain grows, rather than merely happening to work at
the one radius tried so far. A stride-independence check (the same
$v_*$ recovered at different snapshot strides, expressed in physical
time) is the internal consistency test that would catch a
$dt_{\mathrm{snap}}$-related bug if this measurement disagreed with
expectation.

\subsection{H3 --- The learned local closure as a renormalization-group object}

A shared, translation-invariant quadratic stencil law fit directly to
Section~\ref{sec:main_result}'s trained latent recovers $90\%$ of its
one-step dynamics using roughly a third of the full ring width, and
bootstrap-stable sparse selection finds a real (but not
hand-readably-small) compression: $49$ of $252$ candidate polynomial
terms survive resampling, with the survivors overwhelmingly coupling
the physically-anchored channel to \emph{learned hidden channels} at
nearby sites, not to other sites' own physical values. This is
suggestive of a genuine local closure law, not an artifact of a
restricted fitting window (the same $90\%$ figure, not a window-size
one, holds at full global width too).

\textbf{Hypothesis.} If a local closure law of this kind exists, it
should be possible to \emph{coarse-grain} it: define a coarser latent
field by a fixed, non-learned pooling of adjacent trained sites
($z^{(P/2)}_j := R(z^{(P)}_{2j}, z^{(P)}_{2j+1})$), fit the same
stencil-closure form at the coarser level, and ask whether the fitted
coefficients converge as this coarse-graining is iterated --- a
real-space renormalization-group flow of the learned local law, with a
fixed point corresponding to a genuine continuum-limit closure.

\textbf{Test.} This is directly implementable with machinery already
built for the closure-regression experiment above (only the pooling
step and the iteration are new). Two outcomes are both worth reporting:
convergence to a fixed point would be a substantially stronger
"here is the effective PDE" claim than any single-resolution fit
could support; measured non-convergence, with a quantified flow rate,
would itself be a real, quotable negative result about how much
hidden-channel memory this system's coarse dynamics genuinely carry.

\subsection{H4 --- Localized ensemble DA at scales no other method can even attempt}

Sections~\ref{sec:da_transfer} and~\ref{exp:e2} establish the
mechanism at $2\times$ the trained size. The strongest form of the
practical claim is that this mechanism, unlike any comparator, keeps
this DA problem \emph{tractable at domain sizes where a global-vector
approach's required ensemble size would already be operationally
absurd} (recall global-vector DA cannot even be evaluated past its
training size at all; SEC needs fresh, expensive data at every new
size).

\textbf{Test.} Push the transfer sweep (E1/E2) further --- $8\times$,
$16\times$ the trained domain --- specifically looking for where, if
anywhere, the extensivity ratio $D_{\mathrm{KY}}/L$ or the localized-DA
skill starts to degrade. A clean, reportable breakdown point (e.g.\ "transfer
holds cleanly through $8\times$ and degrades by $16\times$, coincident with
[a specific measured cause]") would be at least as valuable to report as
indefinite success, and is the honest way to find the actual boundary
of the claim rather than asserting it holds everywhere.

\subsection{E6 --- LETKF as a second DA analysis scheme}
\label{exp:e6}

\textbf{Gap.} Every DA result in this document uses one analysis
method, a nonlinear Gauss--Newton particle-flow filter (NAT-PFF,
Section~\ref{sec:model}). This raises an obvious question for a
reviewer with a classical-DA background: is the localization result
specific to this less-standard method, or does it hold under the
field-standard comparator, the (localized) ensemble Kalman filter?
Gaspari--Cohn tapering and sampling-error correction were both
developed \emph{for} the EnKF family specifically, so testing
localization inside that family is arguably the more natural home for
this claim, not an optional addition.

\textbf{Implementation.} We implemented a covariance-localized EnKF
(\texttt{ks\_latent.da.letkf.LETKF}) as a drop-in comparator to NAT-PFF,
sharing the same interface, the same forecast step, and --- critically
--- the same \texttt{localize\_fn} taper object PFF already accepts, so
a PFF-vs-LETKF comparison isolates exactly one variable, the analysis
update rule, and nothing else. We are explicit that this is the
Houtekamer--Mitchell-style covariance-localized EnKF, not literally
Hunt et al.\ (2007)'s own per-gridpoint local-domain LETKF algorithm
(designed for much higher-dimensional systems where per-gridpoint
parallelism is the point) --- at this project's scale
($d_{\mathrm{latent}}\leq192$) the full covariance is cheap to hold
explicitly, so we use it directly rather than decomposing into local
domains. The analysis update is a standard perturbed-observation
(stochastic) ensemble Kalman gain, $K = B J^\top(JBJ^\top+R)^{-1}$,
with $B$ the localized ensemble forecast covariance and $J$ the
ensemble-mean Jacobian of the (nonlinear, decoder-based) observation
operator --- the same linearization PFF's own NAT variant already
computes. Validated against the analytic linear-Gaussian Kalman
posterior to $<5\%$ relative error (matching the existing PFF
validation's own bound and tolerance exactly), plus shape/finiteness
and a direct check that the taper measurably shrinks a deliberately
induced far-apart prior correlation's effect on the posterior.

\textbf{First result.} A single head-to-head point at the $L=200$
transfer size, identical setup to Table~\ref{tab:e2_result}'s
$N_{\mathrm{ens}}=64$/Gaspari--Cohn row: \textbf{LETKF gives skill
$1.310$ (RMSE $0.832$, calibration $0.784$) against PFF's skill
$1.342$ (RMSE $0.852$, calibration $0.778$) at the identical operating
point --- nearly indistinguishable.} This is only one point, not yet a
claim either method dominates, but it is an encouraging first signal:
the localization result this paper's central claim depends on does not
appear to be an artifact of the specific nonlinear flow method used
throughout the rest of this document, and a reviewer can be shown the
same effect inside the classical EnKF family they already trust. A
proper follow-up sweep --- LETKF across the same $N_{\mathrm{ens}}$
range and transfer sizes E2 already covers for PFF --- is the natural
next step to turn this into a real comparison rather than one point.

\subsection{E7 --- Full LETKF ensemble-size sweep at the transfer size}
\label{exp:e7}

\textbf{Gap.} E6 is one point ($N_{\mathrm{ens}}=64$, Gaspari--Cohn,
$L=200$). It shows PFF and LETKF agree \emph{there}, but says nothing
about whether they agree across the full range this paper's central
claim is actually made over --- E2's own $N_{\mathrm{ens}}\in
\{8,\dots,256\}\times\{\text{none},\text{Gaspari--Cohn}\}$ sweep at
$L=200$ (Table~\ref{tab:e2_result}).

\textbf{Experiment.} Repeat E2's identical sweep protocol with LETKF
in place of PFF (\texttt{scripts/exp\_e2\_ensemble\_sweep\_ltransfer.py
--method letkf}, same checkpoint, same $L=200$, same
$n_{\mathrm{prop\_steps}}=5$, same $g_c=3.0$) and overlay the two
methods' skill curves directly. \textbf{Falsifiable prediction}: if
E6's single-point agreement generalizes, LETKF's
skill-vs-$N_{\mathrm{ens}}$ curve should track PFF's closely at every
tested size, \emph{and} LETKF should reproduce the same qualitative
signature Section~\ref{exp:e2} found for PFF --- Gaspari--Cohn staying
ahead of no localization at every $N_{\mathrm{ens}}$ including $256$,
unlike the reversal seen at the trained size $L=100$. If LETKF instead
diverges from PFF at small $N_{\mathrm{ens}}$ specifically (plausible:
LETKF's Kalman gain is built from a single ensemble-mean-Jacobian
linearization, whereas PFF's nonlinear flow can in principle track a
more non-Gaussian posterior when the ensemble is small and the
linearization is least trustworthy), that would be a genuinely
informative divergence, not a failure of the experiment --- it would
say something real about how much the posterior's non-Gaussianity
matters at this operating point, precisely the question E6's module
docstring raised and could not yet answer with one point.

\textbf{Result: the prediction holds, cleanly, at every point.}

\begin{table}[h]
\centering
\begin{tabular}{lcccccc}
\toprule
$N_{\mathrm{ens}}$ & PFF, none & LETKF, none & PFF, GC & LETKF, GC & ratio, PFF & ratio, LETKF \\
\midrule
8   & $0.722$ & $0.784$ & $0.831$ & $0.801$ & $1.15$ & $1.02$ \\
16  & $0.513$ & $0.635$ & $1.043$ & $1.064$ & $2.03$ & $1.67$ \\
32  & $0.596$ & $0.559$ & $1.218$ & $1.163$ & $2.04$ & $2.08$ \\
64  & $0.692$ & $0.658$ & $1.342$ & $1.310$ & $1.94$ & $1.99$ \\
128 & $0.848$ & $0.905$ & $1.395$ & $1.387$ & $1.65$ & $1.53$ \\
256 & $1.266$ & $1.150$ & $1.398$ & $1.352$ & $1.10$ & $1.18$ \\
\bottomrule
\end{tabular}
\caption{PFF (Table~\ref{tab:e2_result}) vs.\ LETKF, identical sweep,
$L=200$. "ratio" is skill(Gaspari--Cohn)/skill(none). Every LETKF skill
value lands within $\sim0.15$ of PFF's at the same point, and the
ratio column's shape --- rising through $N_{\mathrm{ens}}\approx32$,
declining toward $256$, but never crossing below $1$ --- is reproduced
almost exactly.}
\label{tab:e7_result}
\end{table}

Both halves of the pre-registered prediction came true: LETKF's
absolute skill tracks PFF's closely at every one of the twelve points,
and --- more importantly for the paper's central claim ---
Gaspari--Cohn's advantage over no localization never reverses at
$N_{\mathrm{ens}}=256$ under LETKF either ($1.18\times$, vs.\ PFF's
$1.10\times$), unlike Table~\ref{tab:sec_vs_gc}'s $L=100$ crossover.
The alternative, also-informative outcome flagged above --- a
small-$N_{\mathrm{ens}}$ divergence from PFF's nonlinear flow handling
non-Gaussianity better --- did \emph{not} occur; at this operating
point the two methods behave as essentially interchangeable, which is
itself informative (the posterior here does not appear to have enough
non-Gaussian structure for PFF's extra machinery to earn its keep
over a plain linearized Kalman gain). \textbf{Net effect on the
paper's argument}: the localization-becomes-necessary-at-scale finding
now stands under two independently implemented analysis schemes, one
of them the classical, widely-trusted EnKF family --- substantially
reducing the risk that Section~\ref{sec:da_transfer}'s and
Section~\ref{exp:e2}'s results are an artifact of PFF's specific
nonlinear-flow mechanics.

\textbf{Result, $L=400$ (user-directed repeat).}

\begin{table}[h]
\centering
\begin{tabular}{lcccccc}
\toprule
$N_{\mathrm{ens}}$ & PFF, none & LETKF, none & PFF, GC & LETKF, GC & ratio, PFF & ratio, LETKF \\
\midrule
8   & $0.691$ & $0.732$ & $0.869$ & $0.795$ & $1.26$ & $1.09$ \\
16  & $0.510$ & $0.581$ & $0.990$ & $0.996$ & $1.94$ & $1.72$ \\
32  & $0.503$ & $0.466$ & $1.056$ & $1.025$ & $2.10$ & $2.20$ \\
64  & $0.597$ & $0.529$ & $1.047$ & $1.047$ & $1.75$ & $1.98$ \\
128 & $0.795$ & $0.730$ & $1.050$ & $1.053$ & $1.32$ & $1.44$ \\
256 & $0.897$ & $0.922$ & $1.066$ & $1.060$ & $1.19$ & $1.15$ \\
\bottomrule
\end{tabular}
\caption{PFF (Table~\ref{tab:e2_result_L400}) vs.\ LETKF, identical
sweep, $L=400$.}
\label{tab:e7_result_L400}
\end{table}

\textbf{The $L=400$ repeat confirms both halves of the prediction
again, including the sharper finding Table~\ref{tab:e2_result_L400}
raised}: LETKF's unlocalized skill also never exceeds $1$ at any
tested $N_{\mathrm{ens}}$ ($0.922$ at $256$, essentially matching
PFF's own $0.897$) --- "no ensemble size rescues unlocalized DA at
$4\times$ the trained domain" now holds under both analysis schemes,
not one. The ratio column keeps its shape under LETKF too (rising
through $N_{\mathrm{ens}}\approx32$, easing toward $256$, always
$>1$). Read alongside the reproducibility caveat just above: these
values differ from a single PFF/LETKF re-run by roughly the same
$\sim5$--$6\%$ margin already quantified there, so the close PFF/LETKF
agreement itself should be read at that resolution too --- genuinely
close, not necessarily identical to three significant figures.

\section{Staged next steps: extending to 2D chaotic systems}
\label{sec:next_steps}

Section~\ref{sec:hypotheses}'s RBC/CM1 note is still the eventual
target, but RBC itself carries a real cost this section is explicit
about rather than glossing over: it is 2D, multi-field (temperature
plus two velocity components at minimum, versus KS's single scalar),
needs an incompressibility/pressure constraint every step (not a
simple explicit spectral update), and --- per this project's own
exploratory RBC runs --- has not yet even been shown to possess a
genuinely chaotic regime at any tested Ra/aspect ratio (the one
dataset attempt settled into a steady state; a wider aspect ratio
showed only a slow, unconfirmed transient). Committing to full RBC
immediately means absorbing all of that cost and uncertainty at once.

\textbf{Decision (user-directed, 2026-09-28): stage the 2D generalization
through two cheaper, better-understood systems first, in order, each
validated before moving to the next.} Both are single-field, spectrally
solvable on a periodic domain with no pressure/incompressibility
coupling at all (2D KS) or a trivial one (Kolmogorov flow, solved
exactly in Fourier space on a periodic domain) --- neither carries
RBC's own cost, and both let the local-field architecture's 2D
generalization be validated on its own before RBC's additional
complications (multi-field state, non-periodic vertical boundary,
uncertain chaotic regime) are added on top.

\subsection{Step 1: 2D Kuramoto--Sivashinsky}
\label{step:2dks}

The direct 2D generalization of the exact equation this whole document
is built on: $u_t + \nabla^2 u + \nabla^4 u + \tfrac{1}{2}|\nabla u|^2 =
0$ on a periodic 2D domain $(L_x, L_y)$. Same solver family (ETDRK4,
contour-integral coefficients), same architecture family (local field
+ \texttt{local\_mlp}), extended from 1D to 2D at every layer --- the
lowest-risk possible next rung, since almost nothing about the
underlying numerics or architecture design needs to be invented, only
generalized.

\begin{enumerate}[nosep]
\item \textbf{Solver.} Extend \texttt{ks\_latent/solver/ks.py}'s ETDRK4
  integrator to a 2D periodic grid: the linear operator becomes
  $\hat L = (k_x^2+k_y^2) - (k_x^2+k_y^2)^2$ and the nonlinearity becomes
  $-\tfrac{i}{2}(k_x \widehat{u_x^2} + k_y \widehat{u_y^2})$ (or the
  equivalent $|\nabla u|^2$ form), using 2D FFT throughout; the
  contour-integral coefficient trick carries over unchanged in form; time
  stepping likely needs a smaller \texttt{dt} than the 1D case since the
  stiffest mode now grows with $k_x^2+k_y^2$ over a 2D wavenumber grid.
\item \textbf{Solver validation, mirroring Section~\ref{sec:ks}'s own
  gate exactly.} Linear-instability spectrum in 2D (most unstable
  mode(s) at $|k|=1/\sqrt2$, now a full ring in $(k_x,k_y)$ space, not
  one point); mean conservation; translation equivariance in
  \emph{both} $x$ and $y$ independently; convergence order in
  \texttt{dt}. \textbf{Then the new question 1D never had}: does
  $D_{\mathrm{KY}}$ scale with \emph{area} $L_x L_y$ (the natural
  2D-extensivity conjecture) or with perimeter, or with neither cleanly
  at accessible domain sizes? Measure this directly via Benettin QR
  (the same machinery E1 already validated and reused at scale, not new
  code) at $\geq 3$ domain sizes before trusting any later
  extensivity/transfer claim on this system --- do not assume the 1D
  law generalizes; measure whether it does.
\item \textbf{Dataset.} Once a genuinely chaotic, validated regime is
  found, generate a training set exactly analogous to
  \texttt{ks\_latent.solver.dataset}'s existing 1D convention (60
  trajectories, independent ICs, discarded spin-up) but for 2D
  snapshots $u(x,y,t)$.
\item \textbf{Architecture.} Extend \texttt{KSAutoencoderLocalField} to
  \texttt{Conv2d}: circular patchify in both directions, circular
  \texttt{Conv2d} site-mixing, latent shaped $(n_{\mathrm{sites},x},
  n_{\mathrm{sites},y}, c)$, gauge anchor unchanged in spirit (channel 0
  = the non-learned physical average of $u$ over each 2D patch).
  Structurally the same generalization step already exercised once this
  session (Section~\ref{sec:model}'s multiscale extension reused
  Section 224's own layers almost verbatim when adding a new spatial
  axis of variation; 1D$\to$2D is the same kind of mechanical
  extension, not a new design).
\item \textbf{Propagator.} 2D analogue of \texttt{local\_mlp}: tokens
  indexed by $(site_x, site_y)$, a windowed \emph{linear} mixer (no
  softmax --- the one mechanism this entire project has repeatedly
  found necessary for recovering chaos, see
  Section~\ref{sec:model}'s own note) with a 2D window radius.
\item \textbf{Training/validation.} Reuse Section 224's exact
  regularizer recipe as the starting point (not reinvented), then the
  same Gate-3/4-style diagnostic suite this document uses throughout
  (standalone rollout boundedness, Lyapunov spectrum, D3
  bandedness) --- generalized to 2D only where the diagnostic itself is
  inherently spatial (e.g.\ D3's coupling-graph locality now has two
  spatial axes to be banded in).
\item \textbf{The payoff experiment}: repeat this document's own
  E1--E4 protocol --- zero-retraining transfer in $L_x$, in $L_y$, and
  in both jointly, checked against the 2D extensivity law measured in
  step 2, ground-truthed against a real Benettin-QR computation the way
  E1 did, not merely checked for internal self-consistency.
\end{enumerate}

\textbf{Progress (2026-09-28).} Steps 1--2 (solver + validation) are
done. \texttt{ks\_latent/solver/ks2d.py} implements the equation above;
the linear operator reuses 1D's exact functional form
($\hat L=K^2-K^4$), but the nonlinear term cannot reuse 1D's own
divergence-form trick ($-u u_x=-\tfrac12(u^2)_x$, an exact
$x$-derivative) since $-\tfrac12|\nabla u|^2$ is not a pure divergence
--- computed instead via spectral differentiation in both directions. A
real, non-obvious bug was caught and fixed exactly, not patched around:
since every RHS term depends only on $u$'s spatial derivatives, $u$'s
absolute value carries zero physical content, but unlike 1D (whose
nonlinearity conserves $\mathrm{mean}(u)$ automatically), 2D's own
nonlinearity has a nonzero spatial mean, so the equation genuinely
drives an unbounded drift in the physically meaningless mean mode ---
an early unprojected run showed exactly this (unbounded drift,
$\max|u|>100$ within spin-up). Fixed by projecting the $k=0$ Fourier
mode to zero after every step (base state and every tangent direction,
for Lyapunov-gauge consistency); after the fix, $\max|u|$ stays in
$[10,16]$ over a full 2000-step trajectory with $\mathrm{mean}(u)$
pinned at float64 roundoff. Validated with the same Gate-1-style
discipline as 1D (9 tests: isotropic-ring instability spectrum, ETDRK4
coefficients checked two ways, exact mean conservation, exact
translation equivariance in $x$/$y$/jointly, convergence order
$\geq3.5$, small-domain decay) --- all pass, full suite 845/845, no
regressions. \textbf{Genuine chaos found on essentially the first
parameter choice tried, no search needed}: at $L_x=L_y=32$,
$\lambda_1=0.237$, $D_{\mathrm{KY}}=45.5$ (cleanly bracketed at 70
directions, 24 positive, a smoothly decreasing spectrum, not
degenerate); confirmed again at the production config
($L_x=L_y=44$, $NX=NY=128$).

\textbf{Architecture, one deliberate deviation from the plan above.}
\texttt{KSAutoencoderLocalField2D} (\texttt{Conv2d} throughout,
literal 2D generalization of Section 224's encoder, same gauge-anchor
design) is implemented and tested (17 tests: shapes, exact gauge
anchor, exact equivariance in $x$/$y$/jointly, no-bias output,
gradient flow). For the PROPAGATOR, this first attempt uses
\texttt{site\_conv\_2d} (a direct 2D generalization of
\texttt{site\_conv} --- reusing the encoder's own circular-\texttt{Conv2d}
site-mixing stack as a dimension-preserving $z\to z$ map, the same
precedent Section 222/223 established for 1D) rather than the
windowed-token \texttt{local\_mlp} design originally planned above:
both are weight-shared, no-softmax linear mixers (the one property
this project has repeatedly found necessary), but \texttt{site\_conv\_2d}
needed no new 2D token-window masking machinery to implement correctly
in this pass.

\textbf{The actual 2D \texttt{local\_mlp} (\texttt{backbone=
"local\_mlp\_2d"}) has since been built too}, user-directed, as a
second propagator to compare against \texttt{site\_conv\_2d} once
training results are in --- genuinely straightforward given
\texttt{site\_conv\_2d}'s own precedent:
\texttt{\_LocalMixerBlock}'s three components (pre-norm windowed mixer + residual, pre-norm per-token
FFN + residual) are all per-token pointwise operations except the
mixer itself, so the ONLY new code was swapping the windowed circular
\texttt{Conv1d} for a windowed circular \texttt{Conv2d} --- the same
1D$\to$2D generalization pattern already used twice (the encoder,
\texttt{site\_conv\_2d}). Purpose-built for a \texttt{local\_field\_2d}
latent specifically, it needs no general tokenizer at all (unlike 1D's
own \texttt{local\_mlp}, built to tokenize an arbitrary flat
$d_{\mathrm{latent}}$): tokens ARE the sites, already naturally
gridded. 9 new unit tests (shape, identity-at-init, exact equivariance
under a site shift, gradient flow, multi-step rollout finiteness);
also directly verified NOT to trigger the same MPS rollout+backward
bug the multiscale experiment's degenerate single-token case hit
(this backbone's smallest realistic token grid is a genuine
$8\times8{=}64$ tokens, not $1$).

\textbf{Dataset sizing, a real mistake caught before it wasted
significant time.} A first attempt reused 1D KS's own dataset
convention directly (60 trajectories, 1001 snapshots each) without
re-deriving the resulting file size for a 2D field --- a 2D snapshot at
$NX{=}NY{=}128$ has $64\times$ the raw element count of 1D's own
$NX{=}256$ field, so the same snapshot COUNT produced a $\sim$15.8GB
file (raw + normalized arrays combined), most of the way through
generating before the mistake was caught. Corrected to 30 trajectories
(24 train / 6 val) $\times$ 251 snapshots (\texttt{trajectory\_time=25},
$\approx$6 Lyapunov times at $\lambda_1{\approx}0.24$) --- $\approx$1.8GB,
generated in $\approx$8 minutes.

\textbf{Result: \texttt{site\_conv\_2d} collapsed --- $D_{\mathrm{KY}}=0$,
a real negative result with a well-motivated next step, not a dead
end.} Training completed cleanly (\texttt{val\_recon\_final=0.018712},
\texttt{best\_val\_kmax\_mse=0.044841}), but a real anomaly was flagged before
declaring success: \texttt{val\_kmax\_mse} printed EXACTLY the same
value for 220+ epochs straight despite the training loss itself
continuing to change --- confirmed not a logging artifact by reading
the training loop's own code. Diagnosed directly, not assumed: the
standalone 2000-step rollout stayed finite with different initial
conditions remaining well-separated (ruling out a fixed-point-collapse
where all ICs converge together), but \texttt{max}$|z|$ showed a slow,
MONOTONIC drift rather than the bounded fluctuation a healthy chaotic
rollout shows. The Lyapunov spectrum of the trained model settled it:
$\lambda_1=-1.19\times10^{-5}$ (essentially exactly zero),
$n_{\mathrm{positive}}=0/192$, $D_{\mathrm{KY}}=0$ --- a genuine
collapse to a non-chaotic, near-neutral regime, a third distinct
collapse signature from what this project has seen elsewhere (neither
"diverges to NaN" nor "clean fixed point").

\textbf{Why this doesn't undermine Step 1's own core finding, and what
to try next.} \texttt{site\_conv\_2d} was chosen for this first 2D
attempt specifically because it needed no new 2D token-window masking
machinery --- but \texttt{site\_conv} was never the reliably robust
backbone in 1D either: Section 222 (\texttt{site\_conv}, 1D) got good
short-horizon behavior but diverged catastrophically by step 1000;
Section 223 (\texttt{local\_mlp}, 1D) was the one that held
$D_{\mathrm{KY}}=22.66$ bounded across the full 2000 steps and became
Section 224's own winning recipe. This 2D collapse may simply be
repeating a known-weaker-1D-backbone's own known unreliability, not a
new, 2D-specific failure. \texttt{local\_mlp\_2d} (built and tested per
Section~\ref{sec:next_steps}'s own follow-up, the actual 2D analogue of
Section 224's winning propagator) is the natural, well-precedented next
attempt --- queued, not yet launched.

\textbf{Result: genuine, rich chaos recovered, with one honest caveat
about what's actually been isolated.} Section 229
(\texttt{local\_mlp\_2d} at the same \texttt{attn\_window=2} Section
228 used) was killed mid-Stage-1, user-directed, before producing a
result, superseded by Section 230: \texttt{local\_mlp\_2d} at
\texttt{attn\_window=3} (receptive field $2\times3=6$ sites
$\approx33$ physical units, deliberately past the "2-3 correlation
lengths" heuristic's own upper bound of $\approx26.7$).
\texttt{val\_kmax\_mse} never froze this time --- genuinely fluctuated
throughout training, and \texttt{best\_val\_kmax\_mse=0.016956},
over $2.5\times$ better than Section 228's $0.044841$.

\begin{table}[h]
\centering
\begin{tabular}{lcc}
\toprule
 & Section 228 (\texttt{site\_conv\_2d}, w=2) & Section 230 (\texttt{local\_mlp\_2d}, w=3) \\
\midrule
$\lambda_1$ & $-1.19\times10^{-5}$ ($\approx0$) & $\mathbf{0.197}$ \\
$n_{\mathrm{positive}}$ & $0/192$ & $\mathbf{35/192}$ \\
$D_{\mathrm{KY}}$ & $0$ & $\mathbf{71.9}$ \\
$\max|z|$ pattern (2000 steps) & monotonic drift, $7.36\to7.65$ & oscillating, $7.17\leftrightarrow8.15$, bounded \\
IC pairwise separation & $14.8$ & $15.95$ \\
\bottomrule
\end{tabular}
\caption{Section 228 vs.\ 230 diagnostics on the trained latent propagator.}
\end{table}

\textbf{$D_{\mathrm{KY}}=71.9$ is strikingly consistent with the true
system's own extensivity, not just "some positive number."} E1's
ground-truth measurement at the smaller $L_x{=}L_y{=}32$ test domain
gave $D_{\mathrm{KY}}/\mathrm{Area}=0.0444$; naively scaling that ratio
to the production domain's area ($44^2{=}1936$ vs.\ $32^2{=}1024$)
predicts $D_{\mathrm{KY}}\approx86$ --- the same order of magnitude as
the model's own $71.9$, not a wild mismatch (not yet a rigorous
ground-truth comparison, which needs an actual Benettin-QR computation
at $L_x{=}L_y{=}44$ directly).

\textbf{One honest caveat, stated precisely rather than overclaimed}:
two variables changed between Section 228 and 230 (backbone AND
\texttt{attn\_window}), and Section 229 --- the run that would have
isolated which one mattered --- never finished. Both plausibly
contributed to the fix; the data needed to cleanly separate "backbone
choice" from "receptive field width" as the operative variable does
not exist yet.

\subsection{Step 2: Kolmogorov flow}
\label{step:kolmogorov}

\textbf{Gate before starting: Step 1 (2D KS) has to actually work}
--- a validated 2D local-field checkpoint with a confirmed
zero-retraining transfer result, matching this document's own
"confirm each rung before extending" discipline (Section~\ref{sec:ks}'s
own gate-before-Phase-2 rule, applied here to the systems ladder
itself). Do not start Kolmogorov flow work in parallel with Step 1;
its main value is confirming Step 1's architecture generalizes to a
\emph{different} nonlinearity, which is only a meaningful test once
Step 1 itself is trustworthy.

Forced 2D incompressible Navier--Stokes in vorticity--streamfunction
form, periodic domain, sinusoidal body forcing:
$\omega_t + u\cdot\nabla\omega = \nu\nabla^2\omega + f(y)$, with $u =
\nabla^\perp\psi$ and $\nabla^2\psi=-\omega$ solved exactly in Fourier
space (a periodic-domain Poisson solve is one division per mode, not
an iterative/implicit solve --- this is why Kolmogorov flow is barely
more expensive than 2D KS despite being genuine Navier--Stokes, not a
simplified proxy for it). Directly relevant prior art already flagged
in this project's own literature review: Constante-Amores, Linot \&
Graham (already cited, Section~\ref{sec:related}) test exactly this
KS/Kolmogorov-flow pair.

\begin{enumerate}[nosep]
\item \textbf{Solver.} New module, same ETDRK4-family time-stepping
  approach as steps above (or an appropriate alternative if Kolmogorov
  flow's own stiffness structure favors a different scheme --- decide
  empirically, not by assumption); vorticity is the sole evolved field,
  velocity/streamfunction are diagnostic (derived each step via the
  Fourier Poisson solve), so the state representation stays single-field
  exactly like both KS variants.
\item \textbf{Solver validation.} Benchmark against known Kolmogorov
  flow results (critical Reynolds number for the transition to chaos is
  well documented in the literature for this exact system) before
  trusting any dataset generated from it --- the same
  validate-against-a-known-answer discipline Gate 1 applied to 1D KS.
\item \textbf{Everything else mirrors Step 1's own steps 3--7}
  (dataset generation, \texttt{Conv2d} local-field architecture,
  \texttt{local\_mlp}-family propagator, Section 224's regularizer
  recipe as the starting point, the same diagnostic suite, the same
  E1--E4-style transfer protocol as the payoff experiment) --- reusing
  Step 1's own 2D architecture code directly rather than rebuilding it,
  since both systems share the same periodic-2D-scalar-field state
  representation.
\end{enumerate}

\textbf{Progress (2026-09-28).} User-directed: proceed to Step 2 despite
its own stated gate (a fully validated 2D-KS zero-retraining transfer
result) not yet being closed out --- Step 1's own Section 230 result
(genuine chaos, $D_{\mathrm{KY}}=71.9$) is treated as sufficient evidence
the architecture generalizes to warrant starting Step 2 in parallel,
with the transfer experiment itself still open on both systems.

\texttt{ks\_latent/solver/kolmogorov.py} implements the equation above
(vorticity-streamfunction form, sinusoidal body forcing in $x$ at
wavenumber $n$, exact Fourier Poisson solve for the streamfunction every
nonlinear-term evaluation). \textbf{Real, quantified code reuse from
Step 1, not just a structural analogy}: \texttt{wavenumbers2d} and
\texttt{etdrk4\_coefficients\_2d} are imported directly from
\texttt{ks2d.py}, unmodified --- the linear operator here,
$\hat L=-\nu(k_x^2+k_y^2)$, is pure diffusion (no zero-crossing, unlike
KS's own $K^2-K^4$), a strict special case of the same contour-integral
construction 2D KS already validated. \texttt{generate\_trajectory\_dataset2d}
was generalized (optional \texttt{integrate\_fn}/\texttt{spinup\_fn}
parameters, defaulted to 2D KS's own functions so every existing caller's
behavior is unchanged) rather than duplicated, and
\texttt{generate\_trajectory\_dataset\_kolmogorov} is a four-line wrapper
calling it with this system's own integrator.
\texttt{lyapunov\_spectrum\_kolmogorov}/\texttt{KSAutoencoderLocalField2D}/
\texttt{site\_conv\_2d}/\texttt{local\_mlp\_2d} needed literally zero new
code: the Lyapunov adapter is the same flatten-at-the-boundary pattern as
\texttt{lyapunov\_spectrum\_ks2d} (new file, mirrored structure, since the
tangent-linear model itself is system-specific), and the architecture/
propagator code is reused completely unchanged --- it only ever sees a
flat \texttt{NX*NY} scalar field and has no notion of which PDE produced
it.

A genuinely new piece of numerics was needed for the nonlinear term and
its tangent-linear model (\texttt{step\_with\_tangent\_kolmogorov}):
vorticity advection $u\cdot\nabla\omega$ is bilinear in $\omega$ (velocity
$u=(\psi_y,-\psi_x)$ is itself linear in $\omega$ via the Poisson solve),
the same "product of two linear functionals of the state" structure both
KS variants' own nonlinearities have, so the Fr\'echet derivative follows
the product rule exactly (no approximation) --- $D[u\cdot\nabla\omega][d\omega]
= du\cdot\nabla\omega + u\cdot\nabla(d\omega)$, with $du$ the velocity
derived from $d\omega$ by the identical linear map.

\textbf{Validation, mirroring Step 1's own Gate-1-style discipline}
(9 new tests, \texttt{tests/unit/test\_solver\_kolmogorov.py}), with two
literature-independent checks specific to this system: (1) the laminar
solution $u_{\mathrm{lam}}=(\sin(ny),0)$ is an EXACT fixed point of the
FULL nonlinear equation by construction (the forcing is chosen precisely
to balance viscous dissipation once the convective term vanishes
identically on it) --- confirmed to survive discretization and ETDRK4
time-stepping to float64 roundoff over hundreds of steps, a strong check
needing no external number at all; (2) the forcing $f(y)$ genuinely
breaks full $y$-translation invariance (unlike both KS variants, which
are equivariant in every spatial direction) while leaving a residual
DISCRETE symmetry (shifts by $2\pi/n$) exactly intact --- both the
positive case (equivariant at the forcing period) and the negative case
(NOT equivariant for a generic shift, ground rule 19's "test the failure
modes" applied to a genuine physical asymmetry rather than a bug) are
tested explicitly, alongside $x$-equivariance (exact, unaffected by the
forcing), mean-vorticity conservation, $\mathcal{O}(4)$ convergence in
\texttt{dt}, and a qualitative stability benchmark (leading Lyapunov
exponent of the linearization about the laminar fixed point is negative
at $\mathrm{Re}=1$ and positive at $\mathrm{Re}=50$, forcing wavenumber
$n=4$ on the standard $2\pi\times2\pi$ box --- the same configuration
Chandler \& Kerswell (2013) and others study; only the qualitative sign
is asserted, not a specific critical Reynolds number, since this
project's compact box is not the elongated/near-infinite-$x$ domain the
classical Meshalkin--Sinai $\mathrm{Re}_c=\sqrt2$ result is stated for.
\textbf{Not a claim that $\mathrm{Re}=50$ is past chaos onset} --- a
later, more careful investigation (below) found this specific point sits
right at a marginal/limit-cycle-like secondary state, not genuine chaos;
this test only checks the much weaker (and here, unambiguous) claim that
the laminar solution's own linear instability threshold lies somewhere
between $\mathrm{Re}=1$ and $\mathrm{Re}=50$.

\textbf{A real methodology bug, caught and fixed, that materially changed
this system's own story twice.} An initial, wide parameter search (varying
Reynolds number, domain aspect ratio, forcing wavenumber $n$, and the
drag coefficient $\alpha$ introduced below) using short Lyapunov
measurements (\texttt{warmup\_time=0.0}, the function's own default, only
6--8 total time units) found the plain equation apparently STABLE at
every configuration tried --- $\mathrm{Re}\in[40,1000]$, square and
elongated (up to $4{:}1$) boxes, $n\in\{1,2,4\}$, $\alpha\in[0,0.3]$ ---
despite the laminar base flow being confirmed linearly unstable. This
was alarming enough (a well-documented chaotic system appearing
stubbornly non-chaotic everywhere) to suspect a genuine bug rather than
physics. It was: the Benettin QR algorithm's own random initial tangent
basis needs a real warmup period to align with the true dominant
directions before the exponents mean anything, and this project's own
established protocol elsewhere (e.g.\
\texttt{scripts/exp\_e1\_groundtruth\_lyapunov\_ltransfer.py}, which
defaults \texttt{warmup\_time=300} out of \texttt{total\_time=2000}) already
knew this --- the ad hoc exploration scripts written for Kolmogorov flow
simply hadn't followed that convention. A direct check confirmed it
exactly: a single-direction power iteration about the exact laminar fixed
point at $n=4$, $\mathrm{Re}=50$ gave $\lambda_1\approx-0.2$ to $-0.4$
with no warmup (a short, 3-time-unit window) but $\lambda_1=0.837$--$0.839$
(agreeing at both $NX=32$ and $NX=128$) once warmed up for 10 time units
and measured over 30 more --- not a small correction, a sign flip.

\textbf{Once corrected, genuine, resolution-converged chaos appeared
cleanly, with no drag and no elongated domain needed at all.}
Re-measuring the SATURATED nonlinear attractor (not just the linearization
about the fixed point) with \texttt{warmup\_time=20}, \texttt{total\_time=40}:
$\mathrm{Re}=50$ (n=4, square box) sits at $\lambda_1\approx-0.0025$, a
marginal/limit-cycle-like state, but $\mathrm{Re}=70$ gives
$\lambda_1=0.057$, $n_{\mathrm{positive}}=3$, $D_{\mathrm{KY}}=7.09$, and
$\mathrm{Re}=100$ gives $\lambda_1=0.102$, $n_{\mathrm{positive}}=2$,
$D_{\mathrm{KY}}=7.90$ --- both confirmed at production resolution
$NX{=}NY{=}128$. Separately, user-directed exploration of whether a
SMALLER forcing wavenumber could reach chaos at even lower resolution
(the opposite of raising $n$, which needs MORE points to resolve a
shorter wavelength and, via the classical $\mathrm{Re}_c\sim n\sqrt2$
scaling, makes the flow relatively more sub-critical at fixed
$\mathrm{Re}$) found $n=2$, $\mathrm{Re}=40$ chaotic at just $NX{=}NY{=}48$
($D_{\mathrm{KY}}=3.59$), confirmed resolution-converged at $NX{=}96$
($D_{\mathrm{KY}}=3.37$) --- a genuinely minimal validated config, though
lower-dimensional than the $n=4$ production point chosen below.

\textbf{The drag term ($\alpha$), tested afterward, turned out to matter
for a different reason than originally guessed.} It was motivated
up front (user-directed, following a chat discussion on 2026-09-28
about the parallel to RBC's own physical vertical confinement) by 2D
turbulence's inverse-cascade condensation problem (Kraichnan 1967) ---
but once $\mathrm{Re}=100$ was already confirmed chaotic without any
drag, sweeping $\alpha$ there revealed a real, non-monotonic ENRICHMENT
effect rather than a "needed to reach chaos" one: $D_{\mathrm{KY}}=7.90$
($\alpha=0$), $\mathbf{23.26}$ ($\alpha=0.05$), $16.18$ ($\alpha=0.1$).
$\alpha=0.05$ --- clearly the best of the three, not simply "more drag is
better" --- was adopted as the production default alongside $n=4$,
$\mathrm{Re}=100$ ($\nu=0.01$), $NX{=}NY{=}128$: $\lambda_1=0.0905$,
$n_{\mathrm{positive}}=12/24$, $D_{\mathrm{KY}}=23.26$ --- comparable in
scale to 1D KS's own $D_{\mathrm{KY}}\approx22$ at its own canonical
config, and reached in a system whose whole point was to be cheaper than
2D KS to stand up.

\textbf{Dataset and training (Section 231).} Reusing
\texttt{generate\_trajectory\_dataset2d} unchanged (generalized with
optional \texttt{integrate\_fn}/\texttt{spinup\_fn} parameters rather than
duplicated) via a 10-line wrapper,
\texttt{generate\_trajectory\_dataset\_kolmogorov} produced 30
trajectories (24 train / 6 val) at the production config,
\texttt{trajectory\_time=40} ($\approx3.6$ Lyapunov times at the measured
$\lambda_1=0.0905$, a coarser \texttt{dt\_snap=0.2} than 2D KS's own 0.1
to keep the file near the same $\sim$1.5GB ballpark given this system's
own slower dynamics), 1.52GB. Training reused
\texttt{KSAutoencoderLocalField2D}/\texttt{local\_mlp\_2d} and Section
224's own regularizer recipe completely unchanged (only \texttt{NX}/\texttt{NY}
and the dataset path differ from the 2D KS driver), \texttt{attn\_window=3}
(reusing Section 230's own value directly, not re-derived) --- Stage 1
\texttt{val\_recon\_final=0.008251} (better than 2D KS's own 0.016521),
Stage 2 \texttt{best\_val\_kmax\_mse=0.030675}, \texttt{val\_kmax\_mse}
fluctuating genuinely throughout (no Section-228-style frozen-metric
collapse signature).

\textbf{Result: the trained model reproduces the true system's own
chaos closely, on the first attempt.} Standalone 2000-step rollout:
finite throughout, \texttt{max}$|z|$ oscillating boundedly (4.5--6.3, no
drift), two independent ICs remaining well-separated (final pairwise
distance 15.0). Lyapunov spectrum of the trained propagator
($n_{\mathrm{directions}}{=}192$, \texttt{warmup\_steps=200},
\texttt{n\_steps=2000}): $\lambda_1=0.1207$, $n_{\mathrm{positive}}=10/192$,
$\boldsymbol{D_{\mathrm{KY}}=21.33}$ --- against the true system's own
$D_{\mathrm{KY}}=23.26$ at the identical config, a strikingly close match
(closer, in fact, than 2D KS's own Section 230 comparison against a
naively-scaled ground-truth estimate). The zero-retraining L-transfer
payoff experiment (does this hold at larger domains/more forcing periods,
where the local-field architecture's actual value proposition lives) is
the natural next step, not yet run.

\subsubsection{Visualization}
\label{sec:section231_viz}

Figures~\ref{fig:section231_physical} and \ref{fig:section231_latent}
animate the same 150-step rollout the numbers above summarize (150
steps, $\approx30$ time units, $\approx2.7$ Lyapunov times, strided to
75 frames), starting from a held-out validation initial condition the
model never trained on --- an interactive HTML version with an
error-over-time chart and latent-channel traces is also published at
\url{https://claude.ai/artifact/EKThBzA7nUM44QBfJAfkUj}. Playback
controls (play/pause/frame-step/speed) require a PDF viewer with
embedded-JavaScript support (Adobe Acrobat/Reader; most browser-native
PDF viewers do not run the animation and will show only the first
frame).

\begin{figure}[h]
\centering
\animategraphics[controls,loop,width=0.95\textwidth]{10}{anim_frames/physical/frame}{1}{75}
\caption{Vorticity field: truth (left), the model's decoded latent
rollout (middle) from the same initial condition, and absolute error
(right). The model never sees the true trajectory --- only its initial
latent state is taken from the true encoding.}
\label{fig:section231_physical}
\end{figure}

\begin{figure}[h]
\centering
\animategraphics[controls,loop,width=0.8\textwidth]{10}{anim_frames/latent/frame}{1}{75}
\caption{Latent field ($8\times8$ sites, 3 channels): truth (top row) vs.\
the model's rolled-out latent state (bottom row). Channel 0 is the
non-learned physical average per site (the gauge anchor); channels 1--2
are the encoder's free residual channels --- this is where the
propagator's own dynamics actually live, ahead of the decode shown in
Figure~\ref{fig:section231_physical}.}
\label{fig:section231_latent}
\end{figure}

\subsection{Why this order, and what it defers}

Both steps are deliberately chosen to be strictly cheaper than RBC on
every axis that matters (single field, periodic domain, explicit or
trivially-diagonal pressure handling, no boundary-layer physics) while
still being genuine, non-trivial spatiotemporal chaos in 2D --- the
actual open question (does the local-field architecture's core
mechanism generalize past 1D at all) gets answered before RBC's own,
separate open questions (does a chaotic regime exist? how does a
multi-field, non-periodic-vertical-boundary state change the
architecture?) are layered on top of it. If Step 1 or Step 2 reveals a
2D-specific failure mode, that is far cheaper to diagnose and fix on a
single-field periodic system than on RBC directly. RBC/CM1 remain the
longer-term target (Section~\ref{sec:hypotheses}'s own note); this
ladder is how to reach it with each rung's risk isolated rather than
compounded.

\section{L-transfer for the 2D systems: a staged plan (G1--G8)}
\label{sec:2d_transfer_plan}

Both 2D rungs now have a trained checkpoint reproducing the true
system's own chaos closely with no L-transfer test yet run at all
(Section 230: $D_{\mathrm{KY}}=71.9$ vs.\ a naive area-scaled estimate of
$86$, not yet a real ground-truth check at the transfer size itself;
Section 231: $D_{\mathrm{KY}}=21.33$ vs.\ the true system's own $23.26$
\emph{at the trained size}, transfer not yet attempted at all). The
obvious next step, user-directed (2026-09-29), is E1--E7's own protocol
(Section~\ref{sec:experiments}), applied to both.

\textbf{Why this is not a mechanical repeat of E1--E7, and the plan
below is staged rather than parallel.} Two real differences from 1D
change the ordering and the cost, not the goal:

\begin{enumerate}[nosep]
\item \textbf{Ground-truth cost scales with grid AREA, not length.} 1D
  KS's own E1 held $dx$ fixed while quadrupling $L$ (a $4\times$ increase
  in state dimension, tractable in a single evening). The equivalent
  isotropic $4\times$-$L$ move in 2D is a $16\times$ increase in grid
  points and a correspondingly larger Benettin QR cost (more tangent
  directions needed for a $\sim16\times$ larger expected $D_{\mathrm{KY}}$
  too, by extensivity) --- a genuinely different compute regime, not the
  same experiment at a bigger number. G1/G2 below therefore start at a
  deliberately modest transfer factor ($1.5\times$ per axis, $2.25\times$
  area) rather than jumping straight to $4\times$; a larger factor is a
  natural follow-up once the first point is in hand and its cost is
  known empirically rather than guessed.
\item \textbf{Kolmogorov flow's own extensivity is the actual open
  question, not a formality.} The chat discussion preceding this plan
  (2026-09-29) raised a specific, physically-grounded reason to expect
  Kolmogorov flow might NOT transfer the way both KS variants did: 2D
  forced turbulence has an inverse energy cascade with no fixed
  correlation length absent an explicit outer-scale mechanism (Kraichnan
  1967 condensation), unlike KS's own fixed critical wavenumber
  $k=1/\sqrt2$. Section 231's own production config already carries a
  large-scale drag term ($\alpha=0.05$) for a different reason
  (attractor enrichment, Section~\ref{step:kolmogorov}'s own progress
  note) --- whether that SAME drag term also supplies the outer-scale
  cutoff needed for genuine L-transfer is exactly the open question G2
  is designed to answer, and is treated here as the single most
  scientifically informative experiment in this whole plan, not a
  routine second data point after 2D KS.
\end{enumerate}

\textbf{The paper-level stakes, kept explicit throughout.} A clean
L-transfer result on THREE structurally distinct extensive chaotic
systems (1D KS, 2D KS, and forced 2D Navier--Stokes) sharing one
architecture family would be substantially stronger evidence for the
paper's central claim than 1D alone --- not "we found one more system
that works," but evidence the weight-shared local-field mechanism is
tracking something genuinely general about extensive spatiotemporal
chaos, independent of the specific PDE. This shapes the ordering below:
G1--G2 (does the core transfer claim hold at all, on the true system, at
both 2D rungs) are prioritized far above G5--G8 (does it ALSO help
downstream DA the way 1D's E2/E3/E6/E7 showed) --- a paper with only
G1--G4 complete already has its central claim; G5--G8 strengthen the
practical-implications case but are not load-bearing for it, and should
not be allowed to delay G1--G2.

\begin{description}[style=nextline,leftmargin=1.2em]

\item[G1 --- Ground-truth L-transfer, 2D KS.] \label{exp:g1}
Analogue of E1. Transfer factor $1.5\times$ per axis
($L_x=L_y=66$, up from the trained $44$, $NX=NY=192$ to hold $dx$
fixed). Two numbers, both needed: (a) the trained checkpoint's own
zero-retraining $D_{\mathrm{KY}}$ at this size (cheap: the latent
Benettin computation, $d_{\mathrm{lat}}=12{\times}12{\times}3=432$
directions at most); (b) a true Benettin-QR computation of the TRUE 2D
KS PDE at the same $(L_x,L_y,NX,NY)$ (expensive: state dimension
$36864$). A short timing probe (few directions, short window) is run
first to calibrate the full computation's wall time before committing
to it, matching this session's own established discipline for every
expensive Lyapunov run so far.

\textbf{Result (2026-09-29), and it is a genuinely two-sided one.} The
TRUE 2D KS system: $\lambda_1=0.204$, $n_{\mathrm{positive}}=101/300$,
$\boldsymbol{D_{\mathrm{KY}}=192.5}$ ($11803$s wall time, 300 directions).
The zero-retraining transferred MODEL: $\lambda_1=0.057$,
$n_{\mathrm{positive}}=19/432$, $D_{\mathrm{KY}}=43.6$ --- a
\textbf{77\% relative error}, far worse than 1D KS's own E1 result
(sub-2\% at every tested size).
\textbf{Correction (2026-10-02):} this run encoded an initial condition
that was not normalized, about $3.8\times$ too large for the encoder. With a
correctly normalized initial condition the same model gives
$D_{\mathrm{KY}}=147.4$ (23\% error); see \S\ref{sec:icnorm}.

\textbf{But the TRUE system's own $D_{\mathrm{KY}}/\mathrm{Area}$ ratio is
essentially constant across two independently-measured domain sizes} ---
$0.0444$ at $L_x{=}L_y{=}32$ (Step 1's own original test-domain
measurement) and $\mathbf{0.0442}$ at $L_x{=}L_y{=}66$ (this experiment)
--- about as clean a confirmation of genuine area-extensivity as this
project has produced for any system. This is the crucial, disambiguating
fact: \textbf{the 77\% gap cannot be blamed on 2D KS's physics lacking a
fixed correlation length}, unlike the live open question for Kolmogorov
flow (G2, below) --- there manifestly IS a stable local law here, and the
zero-retraining transfer mechanism itself is failing to carry it forward,
not failing because no such law exists. The model's own
$D_{\mathrm{KY}}/\mathrm{Area}$ collapses from $0.0371$ (self-reported at
the trained $44{\times}44$ size) to $\mathbf{0.0100}$ under transfer ---
a $\approx4\times$ drop in density, a severe and specific transfer
failure, not a mild one.

\textbf{One architecture-side hypothesis, not yet tested}: a fixed
\texttt{attn\_window} covers a shrinking FRACTION of the domain much
faster in 2D than in 1D as the domain grows, because site count scales
with (scale factor)$^2$ in 2D but only (scale factor) in 1D --- the same
nominal window size is relatively much narrower, relative to the new
domain, after a 2D transfer than after an equivalent 1D one. This is a
purely architectural claim (receptive field vs.\ domain size), orthogonal
to whether the underlying physics is extensive, and would apply to both
2D KS and Kolmogorov flow's own transfer degradation.

\item[G2 --- Ground-truth L-transfer, Kolmogorov flow.] \label{exp:g2}
Analogue of E1, and the highest-priority, highest-uncertainty item on
this list (see above). Transfer in $x$ only, keeping $L_y=2\pi$ fixed at
the forcing period (Kolmogorov flow is only exactly translation-equivariant
in $x$ --- Section~\ref{step:kolmogorov}'s own solver docstring --- so this
is the clean, symmetry-respecting analogue of extending $L$, unlike an
arbitrary $y$-extension which would also change the forcing's own
discrete structure). $L_x = 2\times(2\pi)$ ($NX=256$, $NY=128$, matching
this session's own earlier elongated-domain exploration runs). Same
two-number comparison as G1 (cheap model-side transfer prediction vs.\
expensive true-PDE Benettin at the same size), with the SAME drag
$\alpha=0.05$ carried through unchanged (testing whether it \emph{also}
transfers, not just whether it exists at $L_x=2\pi$).

\textbf{Result (2026-09-29).} The TRUE Kolmogorov-flow system at
$2\times L_x$: $\lambda_1=0.129$, $n_{\mathrm{positive}}=28/75$,
$\boldsymbol{D_{\mathrm{KY}}=56.8}$ ($6833$s wall time, 75 directions) ---
genuinely, richly chaotic, resolving the inverse-cascade/condensation
worry raised before this experiment ran: doubling the domain does not
collapse the attractor or wash it out toward a trivial large-scale mode.
The zero-retraining transferred MODEL: $\lambda_1=0.102$,
$n_{\mathrm{positive}}=13/384$, $D_{\mathrm{KY}}=28.4$ --- a
\textbf{50\% relative error}: smaller than 2D KS's own 77\% (G1, above),
but still a substantial, real gap, not a close match.

\textbf{The true system's own $D_{\mathrm{KY}}$ grew SUPER-linearly with
area} --- $23.26\to56.84$, a $2.44\times$ increase for an EXACTLY
$2\times$ increase in area (Lx doubled, Ly fixed) --- not the
$\approx2\times$ growth clean, linear-in-area extensivity would predict.
This is the key fact distinguishing Kolmogorov flow's own open question
from 2D KS's (settled, above): super-linear growth is consistent with
genuinely NEW large-scale structure entering as the domain grows (2D
turbulence's own inverse-cascade condensation, Kraichnan 1967), i.e.\
the drag $\alpha=0.05$ may not yet be strong enough to pin the
correlation length independent of domain size --- as opposed to simply
"not yet asymptotic," the same finite-size caveat 1D KS's own E1 result
raised at small $L$. \textbf{Not yet disambiguated}; the direct,
cheap (ground-truth-only, no retraining) test is measuring TRUE
$D_{\mathrm{KY}}/\mathrm{Area}$ at a third, larger $x$-scale and/or
sweeping $\alpha$ at fixed domain size, to see whether the ratio
converges or keeps accelerating --- proposed in chat (2026-09-29),
not yet run or added to this numbered list.

\item[G3 --- Seed robustness, 2D KS.] \label{exp:g3}
Analogue of E4. Retrain Section 230's exact recipe (\texttt{local\_mlp\_2d},
\texttt{attn\_window=3}) at $\geq2$ new seeds; confirm the transferred
$D_{\mathrm{KY}}/\mathrm{Area}$ ratio from G1 is seed-stable, not a
property of one lucky checkpoint, mirroring E4's own $n=3$ bar.

\item[G4 --- Seed robustness, Kolmogorov flow.] \label{exp:g4}
Same as G3, for Section 231's recipe and G2's transfer size.

\item[G5 --- DA ensemble-size transfer sweep, 2D KS.] \label{exp:g5}
Analogue of E2. Contingent on a first check (not yet done) that this
project's existing PFF/localization machinery
(Section~\ref{sec:negative}, \texttt{ks\_latent/da}) needs no
2D-specific changes --- it is written generically against a
decoder/observation-operator interface, so this may be "reuse
unchanged" the same way the architecture itself was, but that is a
claim to verify, not assume, before running the sweep.

\item[G6 --- DA ensemble-size transfer sweep, Kolmogorov flow.] \label{exp:g6}
Same as G5, at G2's transfer size, contingent on the same infrastructure
check.

\item[G7 --- SEC vs.\ Gaspari--Cohn refit at transfer size, both 2D systems.] \label{exp:g7}
Analogue of E3. Lower priority than G1--G6: informative about
localization method choice specifically, not about whether the
architecture's core transfer claim holds, and contingent on G5/G6's own
DA infrastructure being in place first.

\item[G8 --- LETKF cross-check, both 2D systems.] \label{exp:g8}
Analogue of E6/E7. Lowest priority on this list --- completes the "three
systems $\times$ two DA schemes" matrix for a maximally thorough paper,
but is not needed for the central L-transfer claim and should only be
run once G1--G6 are done.

\end{description}

\textbf{Status (2026-09-29): G1 and G2 complete, both a genuine mixed
result rather than a clean confirmation --- paused here, user-directed,
before proceeding to G3.} Neither system's zero-retraining transfer
matched its own true ground truth closely (77\% error, 2D KS; 50\%
error, Kolmogorov flow) --- a materially different outcome from 1D KS's
own E1 result (sub-2\% at every size). The disambiguating fact from G1
(2D KS's own true $D_{\mathrm{KY}}/\mathrm{Area}$ ratio is confirmed
constant across two independent domain sizes) rules out "the physics
isn't extensive" as the explanation there, pointing instead at the
transfer MECHANISM itself; G2's super-linear true-$D_{\mathrm{KY}}$
growth leaves open whether Kolmogorov flow's own physics is extensive at
this drag strength at all. Two candidate next diagnostics were raised in
chat, neither run yet: (i) an architecture-side hypothesis that a fixed
\texttt{attn\_window} covers a shrinking fraction of the domain faster in
2D than 1D as scale grows (would apply to both systems); (ii) for
Kolmogorov flow specifically, extending G2's own ground-truth-only
measurement to a third, larger domain size and/or sweeping $\alpha$ at
fixed size, to check whether $D_{\mathrm{KY}}/\mathrm{Area}$ converges.
\textbf{Update (2026-09-30): diagnostic (ii) was run, and it resolves
G2's own open question outright --- see \S\ref{sec:g2_resolved}.} G3
onward (2D KS seed-robustness and beyond) remain queued, not started.

\subsection{G2 resolved: linear drag restores extensivity and unlocks genuine L-transfer (2026-09-30)}
\label{sec:g2_resolved}

\textbf{Context.} G2 left one question explicitly open: is Kolmogorov
flow's own 50\% zero-retraining transfer error a MECHANISM failure (as
G1 showed for 2D KS) or a PHYSICS failure --- insufficient drag letting
the 2D inverse energy cascade build genuinely new, domain-size-dependent
structure that no strictly local model could represent regardless of its
architecture. An initial attempt to fix the mechanism side directly ---
a rotation-equivariant architecture, the natural analogue of the
orientation hypothesis proposed for 2D KS
(\S\ref{sec:orientation_hypothesis}) --- was abandoned before training
once it was pointed out that Kolmogorov flow is not isotropic: its
forcing singles out $y$, giving exact equivariance only in $x$ and
discrete period equivariance in $y$, so full rotation equivariance would
impose a symmetry the physics does not have. A phase-conditioned
anisotropic architecture was built instead (exact $x$-equivariance, plus
two fixed, non-learned \texttt{cos}/\texttt{sin} channels carrying the
forcing's own $y$-phase). \textbf{Its propagator collapsed}: Stage-1/
Stage-2 training losses looked ordinary, but a direct native-size
Lyapunov check (this document's own standing discipline --- never trust
training loss alone, \S\ref{sec:threats}) found $D_{\mathrm{KY}}\approx1$
at both the trained size and the transfer size, against a true value of
$\approx24$ --- the latent state had converged to a near-fixed point.
Reported honestly and abandoned rather than debugged further, per
user direction: \emph{"we should take a step back and use 230 along
with its original dataset as L0"} became, applied to Kolmogorov flow,
\emph{"let's return to the model from 231 and increase drag through
alpha."} This subsection reports the result of that pivot, plus a
second, architecture-side question run alongside it (a windowed local
self-attention site-mixer, isolating whether a more expressive local
operator helps independently of the drag question).

\textbf{Setup.} Two new checkpoints, both trained on a freshly generated
Kolmogorov-flow dataset with $\alpha=0.1$ (double Section 231's own
$\alpha=0.05$, everything else --- $\nu$, forcing wavenumber, domain,
grid, trajectory count --- unchanged):

\begin{itemize}[nosep]
\item \textbf{Section 233}: Section 231's exact recipe (\texttt{local\_
field\_2d} circular-\texttt{Conv2d} encoder, \texttt{local\_mlp\_2d}
propagator, \texttt{attn\_window=3}), retrained from scratch on the
$\alpha{=}0.1$ dataset. A single-variable comparison against Section 231
by construction.
\item \textbf{Section 234}: the same patchify/gauge-anchor/propagator
scaffold, but with the encoder's circular-\texttt{Conv2d} site-mixing
layers replaced by \texttt{\_WindowedAttnBlock2D} --- windowed multi-head
self-attention over the site grid, masked to the same fixed square
neighborhood a \texttt{Conv2d} kernel of that radius would cover, with a
relative (offset-indexed, not absolute-position) bias table so the
zero-retraining resize mechanism is preserved exactly
(\texttt{tests/unit/test\_autoencoder\_local\_attention\_2d.py} checks
exact equivariance and a receptive-field upper bound the same way every
other architecture in this document has been checked). Same $\alpha=0.1$
dataset, same propagator, so this isolates the site-mixing OPERATOR
question from the drag question.
\end{itemize}

\textbf{Extensivity is restored.} A fresh ground-truth Benettin
computation at $\alpha=0.1$, both at the trained size and at G2's own
$2\times L_x$ transfer size:

\begin{table}[h]
\centering
\begin{tabular}{lccccc}
\toprule
$\alpha$ & Trained $D_{\mathrm{KY}}$ & $2\times$ $D_{\mathrm{KY}}$ & Area ratio & $D_{\mathrm{KY}}$ ratio & $D_{\mathrm{KY}}/\mathrm{Area}$ drift \\
\midrule
0.05 & 23.26 & 56.84 & $2.00\times$ & $2.44\times$ & $0.589\to0.720$ ($+22\%$) \\
\textbf{0.10} & \textbf{15.88} & \textbf{30.34} & $2.00\times$ & $\mathbf{1.91\times}$ & $\mathbf{0.402\to0.384}$ ($\mathbf{-4.5\%}$) \\
\bottomrule
\end{tabular}
\caption{Doubling the drag takes the true system from clearly
super-linear $D_{\mathrm{KY}}$ growth (non-extensive) to within 5\% of
flat $D_{\mathrm{KY}}/\mathrm{Area}$ --- essentially the same
extensivity quality 2D KS's own G1 result showed (0.0444 vs.\ 0.0442,
\S\ref{exp:g1}). Because this is a genuine physical-parameter
change, not a domain-size or measurement change, the fact that it fixes
extensivity is itself strong evidence the ORIGINAL non-extensivity was a
real inverse-cascade effect (insufficient outer-scale drag), not merely
an unconverged finite-size artifact --- directly answering the question
G2 left open.}
\label{tab:g2_extensivity}
\end{table}

\textbf{Zero-retraining transfer, both architectures, at $\alpha=0.1$:}

\begin{table}[h]
\centering
\begin{tabular}{lcccc}
\toprule
Model & Trained-size $D_{\mathrm{KY}}$ (err.) & $2\times$-transfer $D_{\mathrm{KY}}$ (err.) & Stage-2 val\_kmax\_mse \\
\midrule
Section 231 ($\alpha{=}0.05$, conv) & 21.33 (8\%) & 28.44 (\textbf{49.96\%}) & 0.0477 \\
Section 233 ($\alpha{=}0.1$, conv)  & 12.77 (19.56\%) & \textbf{31.70 (4.48\%)} & 0.0244 \\
Section 234 ($\alpha{=}0.1$, local-attn.) & 11.08 (30.23\%) & 21.59 (28.83\%) & \textbf{0.0079} \\
\bottomrule
\end{tabular}
\caption{Section 233 vs.\ Section 231: identical architecture, $\alpha$
doubled (which in our solver also raises the forcing amplitude by
$\approx24\%$, i.e.\ $Re$ by $\approx11\%$ --- see
\S\ref{sec:td_literature}; drag dominates, since $D_{\mathrm{KY}}$ fell) ---
transfer error drops from 49.96\% to 4.48\%, an
$\approx11\times$ reduction. Section 234 vs.\ Section 233: identical
$\alpha$, only the site-mixing operator changed --- despite $\approx3\times$
LOWER training loss, Section 234 transfers markedly WORSE, at both
sizes.}
\label{tab:g2_transfer}
\end{table}

\textbf{IC-robustness check on the headline number.} A prior investigation
of the transferred 2D KS model (G1, \S\ref{sec:2d_transfer_plan}) found
it exhibits genuine multistability: the SAME frozen checkpoint's own
zero-retraining $D_{\mathrm{KY}}$ came out to either 43.6 or
$\sim$144 depending purely on which initial condition the measurement
rollout started from (an arbitrary fresh-spinup state vs.\ a realistic
on-attractor snapshot) --- not yet written up formally elsewhere in this
document, but load-bearing enough here that it is worth stating plainly:
a transfer number is not trustworthy until it has been checked against
more than one IC. Given that precedent, Section 233's own 4.48\%
transfer result was re-measured at $2\times L_x$ with two additional,
independent random seeds before treating it as validated:

\begin{table}[h]
\centering
\begin{tabular}{lccc}
\toprule
Seed & $\lambda_1$ & $D_{\mathrm{KY}}$ & Relative error \\
\midrule
0 & 0.1271 & 31.70 & 4.48\% \\
1 & 0.1347 & 32.13 & 5.92\% \\
2 & 0.1288 & 31.99 & 5.43\% \\
\bottomrule
\end{tabular}
\caption{All three seeds cluster within $1.4\%$ of each other and give
relative error in a tight $4.5$--$6\%$ band --- nothing resembling the
$>3\times$ IC-dependence swing found for the transferred 2D KS model.
The headline result is IC-robust, not a lucky draw.}
\label{tab:g2_ic_robustness}
\end{table}

\textbf{Implications.}

\begin{itemize}[nosep]
\item \textbf{Two structurally distinct root causes for L-transfer
failure, now both diagnosed on real systems by the same diagnostic.} G1
showed 2D KS's failure is a MECHANISM problem: true
$D_{\mathrm{KY}}/\mathrm{Area}$ was already confirmed flat, so the
zero-retraining machinery itself is losing information the physics
does have (the orientation hypothesis, \S\ref{sec:orientation_hypothesis}).
G2/this result shows Kolmogorov flow's failure was instead a PHYSICS
problem: $D_{\mathrm{KY}}/\mathrm{Area}$ was genuinely NOT flat at
$\alpha=0.05$, and fixing that with a single physical parameter, zero
architecture changes, resolved transfer almost completely. The same
cheap, ground-truth-only diagnostic (measure $D_{\mathrm{KY}}/\mathrm{Area}$
at two domain sizes) distinguishes which failure mode is in play, and
the two failure modes have different, non-interchangeable fixes.
\item \textbf{A practical design principle for a new PDE.} Before
building and training the full local-field pipeline on an unfamiliar
extensive-chaos system, measure its true $D_{\mathrm{KY}}/\mathrm{Area}$
at two domain sizes (ground-truth Benettin only, no learned model
required). If it is not flat, check whether a physically-motivated
dissipation/outer-scale mechanism (defensible for the system, not an
arbitrary knob) can restore it before assuming the architecture needs to
change --- a far cheaper fix than an architecture rebuild, and one that
also has a genuine, verifiable physical interpretation.
\item \textbf{Training loss is not a reliable proxy for transfer
quality, and this result sharpens that into a specific, measured
direction.} Section 234's Stage-2 validation loss was $\approx3\times$
lower than Section 233's, yet its zero-shot transfer error was
$\approx6\times$ WORSE (28.83\% vs.\ 4.48\%). A plausible mechanism: the
attention site-mixer's learned, data-dependent relative-position bias
can fit correlations specific to the trained domain size more tightly
than a plain circular-\texttt{Conv2d} kernel can, but that extra
flexibility is exactly what a fixed inductive bias (translation
equivariance, the same weights applied identically regardless of how
many sites exist) protects against when the site count changes at
transfer time. Not yet directly verified (candidate check: compare the
two encoders' learned kernels/attention patterns for domain-size-specific
structure) but consistent with the classic flexibility/generalization
trade-off, here measured specifically along the domain-size-extrapolation
axis this document's whole architecture family is built around.
\item \textbf{Against the literature already cited here
(\S\ref{sec:related}): a third position, not "sometimes it works."}
Rather than the zero-shot-optimism and "false promise of zero-shot"
positions being in tension, this result suggests both are locally
correct: transfer is governed by a CHECKABLE physical precondition
(extensivity) rather than being a property of the architecture family
alone. When that precondition holds, even the least expressive of the
two site-mixing operators tested transfers at a few percent error with
literally zero retraining; when it does not, no architecture choice
tested here fixes it, because the target itself is a moving one.
\end{itemize}

\textbf{Extension to $3\times L_x$ (2026-09-30): extensivity is delayed,
not eliminated, and transfer error tracks it almost exactly.} The
natural next question after a clean $2\times$ result is whether it holds
further out. A fresh ground-truth Benettin computation at $3\times L_x$
($NX{=}384$, same $\alpha{=}0.1$) gives $\lambda_1=0.1132$,
$n_{\mathrm{positive}}=28/77$, $D_{\mathrm{KY}}=56.35$ ($4720$s, 77
directions) --- and Section 233's own zero-retraining prediction there is
$D_{\mathrm{KY}}=45.40$ (305s), a \textbf{19.44\% relative error}: far
better than $\alpha=0.05$'s $2\times$ result (49.96\%), but a real,
substantial step back from $\alpha=0.1$'s own $2\times$ result
($\approx5\%$).

\begin{table}[h]
\centering
\begin{tabular}{lccccc}
\toprule
Size & True $D_{\mathrm{KY}}$ & Model $D_{\mathrm{KY}}$ & Error & True $D_{\mathrm{KY}}/\mathrm{Area}$ & $|$drift from trained$|$ \\
\midrule
Trained & 15.88 & 12.77 & 19.56\% & 0.402 & --- \\
$2\times$ & 30.34 & 31.70--32.13 & 4.5--5.9\% & 0.384 & 4.5\% \\
$3\times$ & 56.35 & 45.40 & 19.44\% & 0.476 & 18.3\% \\
\bottomrule
\end{tabular}
\caption{Extensivity is NOT monotonic with scale at $\alpha=0.1$: nearly
flat at $2\times$, then a sharp $+18.3\%$ density jump by $3\times$
(the true $D_{\mathrm{KY}}$ ratio, $3.548\times$, is clearly super-linear
for a $3\times$ area increase). \textbf{The striking pattern: model
transfer error tracks the extensivity drift almost 1:1} at both measured
points ($4.5\%$ drift $\leftrightarrow$ $4.5$--$5.9\%$ error;
$18.3\%$ drift $\leftrightarrow$ $19.44\%$ error) --- only two points,
so this could be coincidence rather than a real quantitative law, but it
is consistent with, and sharpens, the mechanistic story below.}
\label{tab:g2_3x_extension}
\end{table}

This reframes the headline claim from "linear drag fixes Kolmogorov
flow's L-transfer" to the more precise and more interesting statement:
\textbf{doubling the drag delays, rather than eliminates, the onset of
non-extensive growth, and a strictly local model's zero-shot transfer
error is a direct, measurable proxy for how extensive the true physics
actually is at that specific size} --- not a fixed property of the
architecture or the drag value alone. This is consistent with a genuine
(if slower) inverse-cascade effect still operating at $\alpha=0.1$: more
drag raises the area at which new large-scale structure starts to
accumulate, but does not remove the mechanism that produces it. A
natural follow-up, not yet run: whether $\alpha=0.15$ or $0.2$ pushes the
onset point further still (consistent with a genuine drag-dependent
outer scale), or whether $3\times$ is close to some other structural
threshold (e.g.\ three forcing periods) independent of drag strength ---
the two hypotheses make different, testable predictions at a fourth
domain size and/or a second drag value.

\textbf{Status (2026-09-30).} G2 resolved, with the extensivity/transfer
relationship now characterized (not just demonstrated) at two independent
points. Figures~\ref{fig:g2_dky_bars}--\ref{fig:g2_error_tracks_extensivity}
visualize the four tables above. Not yet run: the same $\alpha$-sweep
diagnostic for 2D KS (G1's own failure was already shown to be
mechanism-side, so this is a confirmatory check, not an open question);
IC-robustness at $3\times$ (run only at $2\times$ so far); Section 234's
own IC-robustness triple-seed check (run only for Section 233 so far,
since it is the architecture the headline claim rests on); and the
$\alpha$-sweep-at-fixed-size vs.\ larger-domain-at-fixed-$\alpha$
disambiguation raised just above.

\begin{figure}[h]
\centering
\includegraphics[width=0.85\textwidth]{figures/g2_resolved_dky_bars.pdf}
\caption{True vs.\ zero-retraining-transferred $D_{\mathrm{KY}}$ at the
trained size and at $2\times L_x$, for $\alpha=0.05$ (Section 231) and
$\alpha=0.1$ (Section 233), with relative error annotated above each
pair. The transfer gap (right pair of bars, $49.96\%\to4.5\%$) nearly
closes at the higher drag.}
\label{fig:g2_dky_bars}
\end{figure}

\begin{figure}[h]
\centering
\includegraphics[width=0.85\textwidth]{figures/g2_resolved_extensivity.pdf}
\caption{$D_{\mathrm{KY}}/\mathrm{Area}$ vs.\ area, true system only, at
both drag values. The $\alpha=0.1$ line is nearly flat; the
$\alpha=0.05$ line rises visibly --- the direct visual counterpart of
Table~\ref{tab:g2_extensivity}.}
\label{fig:g2_extensivity}
\end{figure}

\begin{figure}[h]
\centering
\includegraphics[width=0.7\textwidth]{figures/g2_resolved_ic_robustness.pdf}
\caption{Section 233's own transferred $D_{\mathrm{KY}}$ at $2\times L_x$
across three independent seeds (initial conditions), against the true
value. Tight clustering, well inside the shaded 10\% band --- the inset
note states 2D KS's own $>3\times$ IC-dependence swing for scale, to make
the difference concrete.}
\label{fig:g2_ic_robustness}
\end{figure}

\begin{figure}[h]
\centering
\includegraphics[width=0.8\textwidth]{figures/g2_resolved_error_tracks_extensivity.pdf}
\caption{Section 233's zero-shot transfer error (circles) and the true
system's own $|D_{\mathrm{KY}}/\mathrm{Area}$ drift from the trained size$|$
(triangles), at trained/$2\times$/$3\times L_x$. The two curves track
each other almost exactly at both non-trivial points --- consistent with
transfer error being a direct proxy for local extensivity rather than a
fixed property of the drag value or architecture alone
(Table~\ref{tab:g2_3x_extension}).}
\label{fig:g2_error_tracks_extensivity}
\end{figure}

\section{Transfer error as a physics diagnostic: locating where locality breaks down}
\label{sec:transfer_diagnostic}

\subsection{The hypothesis}

\S\ref{sec:g2_resolved} was set up to answer an engineering question
(does the model transfer?), but its most interesting output is a
physics one. A strictly local, translation-equivariant surrogate can
only represent dynamics whose statistics are built from a finite
interaction range: every weight is shared across sites, and its
receptive field is bounded in physical units. When such a model is
trained on a small domain and applied, unchanged, to a larger one, it
can only reproduce the larger system if the larger system is ``more
copies of the same local unit.'' Whatever the larger domain adds that
is \emph{not} local --- a domain-filling condensate, an aggregated
convective cluster, any structure whose size is set by the box rather
than by the physics --- the model cannot represent by construction. So
its transfer error should measure how much non-local structure the
larger domain contains. That turns the surrogate into a probe.

Stated precisely, for a system with Kaplan--Yorke dimension
$D_{\mathrm{KY}}(L)$ at domain size $L$, define the dimension density
$\rho(L)=D_{\mathrm{KY}}(L)/|\Omega(L)|$ and, relative to the training
size $L_0$, the extensivity deficit
\[
\Delta(L) = \left|\frac{\rho(L)}{\rho(L_0)} - 1\right| .
\]
Let $\varepsilon_T(L)$ be the zero-retraining transfer error of a local
surrogate trained at $L_0$ (relative $D_{\mathrm{KY}}$ error, or a
cheaper forward-statistics mismatch, \S\ref{sec:td_experiments}).

\begin{description}[style=nextline,leftmargin=1.2em]
\item[TD1 (diagnostic hypothesis).] $\varepsilon_T(L)$ is a monotone,
  approximately system-independent function of $\Delta(L)$, offset by the
  model's own in-distribution error at $L_0$. Transfer error is therefore
  a readout of extensivity loss, and the size at which $\varepsilon_T$
  rises locates the size at which the dynamics stop being local.
\item[TD2 (physical mechanism, 2D turbulence).] In forced 2D
  Navier--Stokes with linear drag, the length that controls extensivity
  is the frictional arrest scale of the inverse cascade,
  $\ell_\alpha=(\epsilon/\alpha^3)^{1/2}$ \cite{boffettaecke2012,danilovgurarie2000},
  with $\epsilon$ the energy injection rate. Below
  $L\sim\ell_\alpha$ each increase in domain size admits new, box-scale
  modes that the inverse cascade populates, so $D_{\mathrm{KY}}$ grows
  super-linearly. Well above $\ell_\alpha$ the new area only adds more
  independent local units, and $D_{\mathrm{KY}}$ becomes extensive. The
  crossover length should therefore scale as $L^*\propto\alpha^{-3/2}$
  at fixed $\epsilon$.
\item[TD3 (application).] The same construction applied to a
  cloud-resolving model in radiative--convective equilibrium (RCE) should
  locate the domain-size threshold for convective self-aggregation
  \cite{mullerheld2012,wing2017review} without any aggregation index
  being defined by hand.
\end{description}

\textbf{What the current data do and do not show.} The evidence is
three domain sizes at $\alpha=0.1$ and two at $\alpha=0.05$, one
training seed per model, and Benettin windows of only $\sim5$ Lyapunov
times (Table~\ref{tab:g2_3x_extension}). The observed $\varepsilon_T$
tracks $\Delta$ at the two transfer sizes ($4.5\%\leftrightarrow5\%$,
$18.3\%\leftrightarrow19.4\%$), which is suggestive of TD1. But the
pattern is not monotone in $L$ (nearly flat at $2\times$, then rising
at $3\times$), and the in-distribution error at $L_0$ (19.6\%) is larger
than the $2\times$ transfer error. Both are the signature of measurement
noise in $D_{\mathrm{KY}}$ at least comparable to the effect. TD2 also
predicts the opposite trend from a naive reading of the $3\times$ point:
extensivity should \emph{improve}, not worsen, once $L\gg\ell_\alpha$,
so a rise at $3\times$ means either $L_x=6\pi$ is still below the
crossover at $\alpha=0.1$, or the $2\times$ dip is noise. Note too that
every transfer so far extends only $L_x$, with $L_y=2\pi$ fixed at the
forcing period, so the ``area'' being varied is really one length. None
of this is settled; the experiments below are designed to settle it.

\subsection{Why it would matter}

\begin{itemize}[nosep]
\item \textbf{A model-based detector for loss of locality.} Deciding
  where a spatially extended system stops being local normally needs
  either an expensive tangent-space computation at every domain size or
  a hand-built order parameter for the specific structure one expects
  (a condensate amplitude, an aggregation index). A trained local
  surrogate plus forward simulations of the true system at larger sizes
  would give the same answer without either, provided TD1 holds for
  cheap forward-statistics versions of $\varepsilon_T$ (experiment T5).
\item \textbf{A principled split for hybrid models.} Where
  $\varepsilon_T$ rises is where a purely local model must be
  supplemented by a coarse, non-local component. The scale-incremental
  hierarchy of \S\ref{sec:residual_hierarchy} (local Level-0 field plus
  coarser residual levels) is exactly such a split. TD1/TD2 would say
  where to put the boundary, and predict how much capacity the non-local
  part needs as a function of $L/L^*$, rather than choosing it by hand.
\item \textbf{A concrete target in atmospheric modeling.} Convective
  self-aggregation in RCE depends on domain size: aggregation appears
  above roughly 200--300\,km at 2\,km grid spacing \cite{mullerheld2012},
  the threshold is tied to cold-pool dynamics and vanishes when rain
  evaporation is switched off \cite{jeevanjeeromps2013}, and systematic
  domain-size sweeps exist \cite{patriziorandall2019,wing2020rcemip}.
  Separately, ML convective parameterizations have run into the same
  locality question from the other side: purely single-column neural
  networks miss organized multi-column systems, and adding neighboring
  columns as inputs improves skill \cite{wangyuvalogorman2022}; a local
  NN parameterization tested over a larger spatial extent than it was
  trained on degraded in ways treated as bias to be fixed
  \cite{brenowitzbretherton2019}. Nobody we could find has turned that
  degradation around and used it to \emph{locate} the aggregation
  threshold. That is the gap TD3 targets. It would also give a test of
  superparameterization-style designs that assume convection is local
  below the host-model grid scale.
\end{itemize}

\subsection{Best-case outcome, and what weaker outcomes still give}
\label{sec:td_best_case}

\textbf{What the diagnostic measures.} The natural reading is ``a
diagnostic for the size or forcing at which large-scale chaotic structure
appears,'' and that is close. More precisely, it measures the onset of
\emph{non-locality}: the point at which a spatially extended chaotic system
stops being ``more copies of the same local unit'' because structure whose
size is set by the domain (a condensate, an aggregated cluster) starts to
carry dynamical degrees of freedom. The chaos is present on both sides of
that line; what changes is whether its degrees of freedom are local.
Domain size is only one control parameter that can move a system across
it; drag, forcing amplitude, Reynolds number, or (in CM1) rain evaporation
can too, and the same trained local surrogate can be scanned along any of
them.

\textbf{Best case, stated as a single result.} \emph{A local surrogate
trained once on a small domain, and then evaluated with forward
simulations only, measures the non-local part of a chaotic system's
attractor dimension. The length at which that part appears collapses onto
the physical arrest scale across drag values in 2D turbulence, and in a
cloud-resolving model it recovers the self-aggregation threshold and its
dependence on cold pools, without any aggregation index defined by hand.}
Concretely, every T-series experiment would have to come out the right
way:
\begin{enumerate}[nosep]
\item T1 confirms that the transfer-error/extensivity relationship survives
  error bars.
\item T2/T3 show $\varepsilon_T$ vs.\ $\Delta$ collapsing onto one curve
  across $\alpha$ \emph{quantitatively}, so that $\varepsilon_T$ is not
  just an alarm but a measurement: $\varepsilon_T(L)\cdot D_{\mathrm{KY}}(L)
  \approx D_{\mathrm{KY}}(L)-\rho(L_0)\,|\Omega(L)|$, the number of
  degrees of freedom that are \emph{not} accounted for by local units. The
  crossover length collapses as $L^*\propto\ell_\alpha\propto\alpha^{-3/2}$.
\item T4 shows the diagnostic distinguishes physics from architecture
  failure (flat-but-high error for 2D KS).
\item T5 shows forward statistics alone reproduce the Lyapunov-based
  answer, so the expensive tangent computation is only needed once, for
  validation.
\item T6 recovers the known aggregation threshold in CM1 and shows it
  moves when rain evaporation is switched off.
\item T7 shows a hybrid model with a local field plus a global component
  sized from the measured non-local dimension restores transfer at
  $\alpha=0.05$, where the pure local model fails.
\end{enumerate}

\textbf{What that would establish.}
\begin{itemize}[nosep]
\item \emph{A new measurement.} ``How many of this system's degrees of
  freedom are non-local, and above what scale?'' currently requires either
  an expensive tangent-space computation at every size or a hand-built
  order parameter for a structure one already expects. The best case
  replaces both with a cheap, generic, model-based readout, including for
  structures nobody has named yet.
\item \emph{A sizing rule for ML emulators of large domains.} Train small
  (below $L^*$), deploy the local model freely up to $L^*$, and add a
  coarse global component whose capacity is set by the measured non-local
  dimension --- instead of the current practice of discovering the
  failure after the fact and retraining with more capacity (as in
  \cite{constanteamores2024}'s widened Kolmogorov box). The same logic is a
  direct test of the locality assumption behind superparameterization and
  single-column ML convection schemes \cite{wangyuvalogorman2022}.
\item \emph{A physics result in its own right.} The first direct
  Lyapunov-dimension measurement that the drag-set arrest scale is the
  onset length for extensive chaos in forced 2D turbulence
  (\S\ref{sec:td_literature}).
\end{itemize}

\textbf{Weaker outcomes, which are more likely and still useful.}
\begin{itemize}[nosep]
\item \emph{Qualitative only} (T1--T3 give a monotone but non-collapsing
  relation). $\varepsilon_T$ remains a rising alarm for non-locality with
  no calibrated scale; still enough to size hybrid models by experiment,
  and still novel as a detector.
\item \emph{Physics without the diagnostic} (T2 confirms
  $L^*\propto\alpha^{-3/2}$, but $\varepsilon_T$ does not track $\Delta$
  once error bars are in). The TD2 physics result stands alone; the ML
  angle becomes a negative result about what transfer error can and
  cannot see.
\item \emph{T1 kills the signal} (the $2\times$ dip is noise, and the
  current curve is flat within error). Then \S\ref{sec:g2_resolved}'s
  headline shrinks to ``drag reduces transfer error'' and TD1 needs new
  evidence before any further investment.
\item \emph{T5 fails} (only the Lyapunov version works). The diagnostic is
  confined to systems where tangent computations are affordable, so the
  CM1 application (T6) is off the table.
\end{itemize}

\subsection{Literature check: what is and is not already known}
\label{sec:td_literature}

Three independent literature searches (2026-09-30; arXiv and journal
sources, every citation checked against its abstract or DOI) were run
specifically to avoid repeating existing work. Table~\ref{tab:td_prior}
summarizes which components of the hypothesis each closest precedent
already covers.

\begin{table}[h]
\centering
\small
\begin{tabular}{p{4.8cm}ccccc}
\toprule
Work & Local, equiv. & Zero-shot $L$ & Lyapunov/$D_{\mathrm{KY}}$ & Error as & Onset \\
 & surrogate & transfer & validation & diagnostic & scale \\
\midrule
Pathak et al.\ 2018 \cite{pathak2018} & \checkmark & partial & --- & --- & --- \\
Constante-Amores et al.\ 2024 \cite{constanteamores2024} & \checkmark & --- & --- & --- & --- \\
de Villeroch\'e et al.\ 2026 \cite{zeroshotlargerdomains2026} & \checkmark & \checkmark & --- & --- & --- \\
\"Ozalp \& Magri 2025 \cite{ozalpmagri2025} & --- & --- & \checkmark & --- & --- \\
Vinograd et al.\ 2026 \cite{kolmogorovregimes2026} & --- & --- & \checkmark & --- & --- \\
van Nieuwenburg et al.\ 2017 \cite{vannieuwenburg2017} & --- & --- & --- & \checkmark$^*$ & \checkmark$^*$ \\
Brenowitz \& Bretherton 2019 \cite{brenowitzbretherton2019} & \checkmark & partial & --- & --- & --- \\
Wang, Yuval \& O'Gorman 2022 \cite{wangyuvalogorman2022} & --- & --- & --- & partial & --- \\
Boffetta \& Ecke 2012; Danilov \& Gurarie 2000 & \multicolumn{4}{c}{(physics: arrest scale $\ell_\alpha$)} & \checkmark$^\dagger$ \\
Egolf \& Greenside 1994; Paul et al.\ 2007 & \multicolumn{4}{c}{(physics: extensive chaos above a correlation length)} & \checkmark$^\dagger$ \\
\midrule
\textbf{This proposal} & \checkmark & \checkmark & \checkmark & \checkmark & \checkmark \\
\bottomrule
\end{tabular}
\caption{Closest prior work. $^*$Classification error vs.\ an assumed
label split locating equilibrium phase transitions, not transfer error
of a dynamical surrogate. $^\dagger$Establishes the relevant length scale
or the extensive/non-extensive crossover as physics, without a learned
model and without a Lyapunov-dimension test of the $\alpha^{-3/2}$
onset in 2D Navier--Stokes.}
\label{tab:td_prior}
\end{table}

\textbf{Surrogates and zero-shot transfer.} Parallel local reservoir
computers already forecast KS on large domains by exploiting spatial
homogeneity \cite{pathak2018}, and patch-decomposed, weight-shared
autoencoder/neural-ODE models do the same for KS and Kolmogorov flow
\cite{constanteamores2024}. Neither tests transfer to a domain larger
than the one trained on, or measures Lyapunov dimension against domain
size. For \cite{constanteamores2024} this was confirmed from the full
text: each domain ($L{=}220$ KS; Kolmogorov flow at $Re{=}200$,
$n_f{=}2$, no drag, in $[2\pi,16\pi]$ and $[6\pi,16\pi]$ boxes) gets its
own trained model, and validation is by short-time tracking error,
autocorrelations and dissipation--input joint PDFs rather than Lyapunov
spectra. The closest mechanical match is recent
work adding a locality bias and removing absolute position embeddings
so transformer neural operators apply zero-shot to larger domains,
showing that translation-equivariant models keep error flat with domain
size \cite{zeroshotlargerdomains2026}. That is the ``train small, deploy
large'' mechanism, but on general PDE benchmarks rather than chaotic
attractors, and it treats flat error as the goal rather than rising
error as information. The zero-shot \emph{resolution} critique
\cite{falsepromisezeroshot2026} is about mesh density, not domain extent.
Latent models have been validated against the true Lyapunov spectrum,
$D_{\mathrm{KY}}$ and covariant Lyapunov vectors on KS
\cite{ozalpmagri2025}, but at a single domain size.

\textbf{ML error as a physics probe.} The closest methodological
precedent is ``learning by confusion'' \cite{vannieuwenburg2017}, where
a classifier's accuracy as a function of an assumed label split peaks at
an equilibrium phase transition. TD1 is the dynamical analogue: the
thing being varied is domain size and the curve being read is transfer
error of a generative surrogate. We found no instance of domain-size
transfer error of a local surrogate being used to detect locality or
extensivity breakdown in a dynamical system. (The atomistic-ML
literature ties the validity of local interatomic potentials to a
physical correlation length, which is the same principle in a different
field.)

\textbf{Physics.} The ingredients of TD2 are individually standard: the
dual cascade \cite{kraichnan1967}, box-scale condensation of the inverse
cascade when no large-scale sink exists \cite{smithyakhot1993,chertkov2007},
the arrest wavenumber $k_\alpha=\epsilon^{-1/2}\alpha^{3/2}$
\cite{boffettaecke2012,danilovgurarie2000}, drag-dependent condensates
specifically in Kolmogorov flow \cite{mishra2015}, and the general result
that chaotic dimension is extensive only for systems much larger than a
characteristic length \cite{ruelle1982,egolfgreenside1994,paul2007,karimipaul2010}.
Lyapunov and dimension studies of 2D Navier--Stokes and Kolmogorov flow
exist \cite{grappinleorat1991,plattsirovich1991,chandlerkerswell2013}.
The most recent and closest is Vinograd, Cullen \& Clark Di Leoni
\cite{kolmogorovregimes2026}, read in full: they compare autoencoder
dimension estimates with $D_{\mathrm{KY}}$ across Reynolds number and
forcing wavenumber $k_f\in\{2,4,8\}$, but in a fixed $2\pi\times2\pi$ box
with no drag, so there is no domain-size sweep and no transfer of any
model; a related study gives $D\sim Re^{1/3}$ from rank-minimizing
autoencoders \cite{clearypage2025}. We found no study that measures the
Lyapunov-dimension crossover to extensivity as a function of drag and
tests $L^*\propto\alpha^{-3/2}$.

\textbf{What these two papers add to our own results.} Reading both in
full turned up five points that bear directly on \S\ref{sec:g2_resolved}
and on the experiment design below.
\begin{enumerate}[nosep]
\item \emph{$D_{\mathrm{KY}}$ in Kolmogorov flow is a large-scale
  quantity.} In \cite{kolmogorovregimes2026}, autoencoders trained on
  low-pass fields ($k<k_f$) reproduce $D_{\mathrm{KY}}$, while the
  small-scale ($k>k_f$) dimension keeps growing after $D_{\mathrm{KY}}$
  saturates; the unstable directions live at scales larger than the
  forcing, as found earlier for shell models \cite{yamadaohkitani1988}.
  So extensivity of $D_{\mathrm{KY}}$ is specifically a statement about
  the large-scale sector, which is exactly the sector the inverse cascade
  and the arrest scale $\ell_\alpha$ control. That supports TD2, and it
  means a $D_{\mathrm{KY}}$-based $\varepsilon_T$ is sensitive to
  large-scale transfer failure but nearly blind to small-scale errors ---
  one reason T5 should add forward statistics that see small scales.
\item \emph{Degrees of freedom track the forcing scale, not the mode
  count.} The saturated $D_{\mathrm{KY}}$ in \cite{kolmogorovregimes2026}
  grows linearly with $k_f$ (slope $\approx12$), not quadratically.
  Raising $k_f$ in a fixed box adds forcing periods, much as our
  $x$-transfer adds $x$-extent. This suggests measuring domain size in
  units of both $\ell_\alpha$ and the forcing length $2\pi/k_f$ in T2.
\item \emph{Where our runs sit.} With their definition
  $Re=(L/2\pi)^{3/2}\sqrt{f_0}/\nu$ and our forcing amplitude
  $f_0=\nu k_f^2+\alpha$ (chosen so $u=\sin(k_f y)$ is an exact
  laminar state), our runs have $Re\approx46$ at $\alpha{=}0.05$ and
  $Re\approx51$ at $\alpha{=}0.1$. For drag-free $k_f{=}4$ flow they find
  $D_{\mathrm{KY}}$ still rising between $Re_1\approx28$ and
  $Re_2\approx60$, where the large scales saturate. Our runs sit inside
  that ``large scales still developing'' window, although drag shifts the
  boundaries, so this is only indicative.
\item \emph{A confound in \S\ref{sec:g2_resolved}.} Because $f_0$ depends
  on $\alpha$ in our solver, doubling $\alpha$ also raised $f_0$ by
  $\approx24\%$ ($Re$ by $\approx11\%$). Higher $Re$ should raise
  $D_{\mathrm{KY}}$, yet $D_{\mathrm{KY}}$ fell from 23.26 to 15.88, so
  drag clearly dominates. Still, ``only $\alpha$ changed'' is not
  strictly true. T2 must separate the two.
\item \emph{Independent evidence of domain-size-dependent structure
  without drag.} When \cite{constanteamores2024} widened the Kolmogorov
  box from $L_x{=}2\pi$ to $6\pi$ at $Re{=}200$ (no drag), they report that
  ``the dynamics in the $x$-direction become more complex'', observe
  vortex cores aligned at half the domain length, and had to double the
  per-patch latent dimension from 15 to 30. That is box-scale
  organization appearing as the domain widens --- the same non-extensive
  growth TD2 attributes to missing drag --- and their remedy (retrain with
  more local capacity) is exactly what a rising $\varepsilon_T$ would
  have flagged. Their validation statistics (spatial autocorrelation,
  dissipation--input joint PDF) are also natural cheap candidates for T5.
\end{enumerate}
Both papers also average over far longer windows than we did
($>10^4$ turnover times in \cite{kolmogorovregimes2026}, with FTLEs
checked for convergence against the integration interval), which
reinforces T1.

\textbf{Verdict.} No component is new on its own, and the paper should
cite each precedent explicitly rather than claim priority for any
single piece. The combination --- a local surrogate's zero-shot transfer
error, validated against the true Lyapunov dimension, used to locate the
onset of non-extensive dynamics and tied quantitatively to the arrest
scale --- does not appear in the literature we could find, and neither
does its application to the self-aggregation threshold. The two closest
papers \cite{constanteamores2024,kolmogorovregimes2026} were read in
full and confirmed to contain no domain-size transfer; the remaining
searches relied mainly on abstracts.

\subsection{Required experiments (T1--T7)}
\label{sec:td_experiments}

Ordered by dependency: T1 decides whether the existing signal is real,
T2--T4 test TD1/TD2 on systems we already have, T5 makes the diagnostic
practical, T6 applies it to convection, T7 uses it.

\begin{description}[style=nextline,leftmargin=1.2em]
\item[T1 --- Error bars on everything already measured.] \label{exp:t1}
  Before any new physics: repeat the true $D_{\mathrm{KY}}$ at $\alpha=0.1$,
  $L_x\in\{1,2,3\}\times2\pi$, with 3 seeds and \texttt{total\_time}
  $\in\{40,120\}$ ($\sim5$ vs.\ $\sim15$ Lyapunov times), and retrain
  Section 233 with 2 more seeds. Output: $\rho(L)$ and $\varepsilon_T(L)$
  with confidence intervals. \emph{Decision rule}: if the $2\times$ dip in
  $\rho$ is not outside the $1\sigma$ band of the $1\times$ and $3\times$
  values, it is treated as noise and the ``1:1 tracking'' claim is
  withdrawn in favor of whatever the corrected curve shows.

  \emph{Implementation} (\texttt{scripts/exp\_t1\_error\_bars.py}, launched
  2026-09-30). Rather than separate 40- and 120-unit runs, each true
  $D_{\mathrm{KY}}$ is one \texttt{total\_time}$=120$ run that records
  running estimates at $t=40,80,120$: the $t=40$ record reproduces the
  original seed-0 protocol (a consistency check), and the three records
  show convergence in averaging time directly. Model side: 3 model seeds
  (Section 233 plus two retrains) $\times$ 3 sizes $\times$ 3 IC seeds, 5 of
  which already existed. 38 tasks in all, run one at a time,
  single-threaded at the lowest scheduling priority. The run is
  preemptable: \texttt{benettin} now writes its full state atomically
  every 4 time units and resumes from it (a \texttt{kill -9} drill on the
  Kolmogorov path reproduced the uninterrupted spectrum to $2\times10^{-15}$,
  the same level at which two uninterrupted runs of this FFT pipeline
  already differ), training resumes at the stage level, and finished tasks
  are skipped on restart. Estimated $\sim$30\,h of single-core time.
\item[T2 --- Map the crossover in $(\alpha, L_x)$.] \label{exp:t2}
  Ground-truth $D_{\mathrm{KY}}$ on a grid $\alpha\in\{0.05,0.1,0.2\}$,
  $L_x\in\{1,2,3,4,6\}\times2\pi$ (drop the largest sizes where cost is
  prohibitive), with the energy injection rate $\epsilon$ measured
  directly from each run so $\ell_\alpha$ can be computed rather than
  assumed. Record an independent, non-Lyapunov indicator of non-local
  structure alongside: the fraction of kinetic energy in the lowest $k_x$
  modes and in $k<k_f$ (the sector that carries $D_{\mathrm{KY}}$,
  \S\ref{sec:td_literature}), and the spectral peak wavenumber. Because
  our forcing amplitude currently scales with $\alpha$, add one control
  point per $\alpha$ run at fixed $f_0$ (fixed $Re$), dropping the
  exact-laminar-state convention for that run, so drag and Reynolds
  number are separated. \emph{Prediction (TD2)}: $\Delta$
  plotted against $L_x/\ell_\alpha$ collapses onto one curve across
  $\alpha$, rising then falling back toward zero once $L_x\gg\ell_\alpha$.
  \emph{Falsified if} the crossover does not move with $\alpha$, or moves
  with a clearly different exponent.
\item[T3 --- Transfer error on the same grid.] \label{exp:t3}
  Train one local surrogate per $\alpha$ at $L_0=2\pi$ (plus seeds from
  T1), transfer to every grid point, measure $\varepsilon_T$.
  \emph{Prediction (TD1)}: $\varepsilon_T-\varepsilon_T(L_0)$ vs.\ $\Delta$
  collapses across all three $\alpha$. \emph{Pre-registered pass}:
  Spearman correlation $>0.8$ over $\geq9$ grid points, and the
  $\varepsilon_T$-inferred crossover within a factor of 1.5 of the
  ground-truth crossover from T2.
\item[T4 --- Controls that separate diagnostic from coincidence.] \label{exp:t4}
  Positive control: 1D KS, where $\Delta\approx0$ at all sizes tested
  (E1), so $\varepsilon_T$ should stay near its in-distribution value.
  Negative control: 2D KS, where $\Delta\approx0$ (G1) but transfer fails
  for a mechanism reason. TD1 predicts $\varepsilon_T$ is high but
  \emph{flat} in $L$ there, which distinguishes ``the architecture
  cannot generalize'' from ``the physics changed.'' If 2D KS error also
  grows with $L$, the diagnostic is confounded with generic architecture
  degradation and TD1 must be restated.
\item[T5 --- A cheap version of $\varepsilon_T$.] \label{exp:t5}
  The diagnostic is only useful for CM1-scale systems if it avoids the
  tangent-space computation (hours at $3\times$ here; infeasible for a
  CRM). Replace the Lyapunov comparison with forward-only statistics: the
  mismatch in time-averaged energy spectrum (split at $k_f$), two-point
  spatial autocorrelation, the dissipation--power-input joint PDF (the
  validation statistic of \cite{constanteamores2024}), and short-horizon
  ensemble forecast error between model rollouts and true simulations at
  each size. \emph{Pass}: the cheap mismatch locates
  the same crossover as the Lyapunov version on the T2/T3 grid.
\item[T6 --- Convection.] \label{exp:t6}
  First Rayleigh--B\'enard (already the planned next rung, E5,
  \S\ref{exp:e5}), where extensivity at large aspect ratio is established
  \cite{paul2007,karimipaul2010}: a positive control in a convective
  system. Then CM1 \cite{bryanfritsch2002} in RCE: train a local surrogate
  on a domain below the aggregation threshold (e.g.\ the RCEMIP small
  domain, $\sim$100\,km \cite{wing2020rcemip}), transfer to domains
  straddling it using the cheap $\varepsilon_T$ from T5. \emph{Prediction
  (TD3)}: $\varepsilon_T$ rises where aggregation sets in. \emph{Causal
  test}, the direct analogue of varying $\alpha$: disable rain
  evaporation, which removes the domain-size dependence of aggregation
  \cite{jeevanjeeromps2013}. The $\varepsilon_T$ rise should move or
  disappear accordingly.
\item[T7 --- Use the diagnostic to design the hybrid.] \label{exp:t7}
  At $\alpha=0.05$ (pure local transfer error 50\%), add a coarse
  global/residual component (\S\ref{sec:residual_hierarchy}) sized by the
  T2/T3 crossover. \emph{Prediction}: transfer error returns to the
  in-distribution level, and the capacity the coarse component needs
  grows with $L_x/L^*$ rather than with $L_x$ itself.
\end{description}

\textbf{Cost and priority.} T1 is a few days of CPU and should run
first, since everything else assumes the current signal is real. T2 is
the most expensive (one Benettin run of up to several hours per grid
point) and can be trimmed to 9--12 points. T3 and T4 reuse existing
machinery and checkpoints. T5 is the step that decides whether T6 is
feasible at all. T6 is a new system and a larger engineering effort.

\textbf{Status (2026-10-01): T1 launched, interim results, one unplanned
finding.} \texttt{scripts/exp\_t1\_error\_bars.py} is running (preemptable,
\S\ref{exp:t1}'s own implementation note). 32 of 38 tasks complete at this
writing: all three true-$D_{\mathrm{KY}}$ seeds at $1\times$, both
retrains, and all 27 model-side evaluations (3 model seeds $\times$ 3
sizes $\times$ 3 IC seeds). The $2\times$/$3\times$ ground-truth seeds are
still running.

\begin{table}[h]
\centering
\begin{tabular}{lcccc}
\toprule
 & $t=40$ & $t=80$ & $t=120$ & mean $\pm$ std ($t{=}120$) \\
\midrule
True $D_{\mathrm{KY}}$, $1\times$ (3 seeds) & 15.88, 12.66, 14.28 &
  15.13, 11.65, 12.88 & 14.86, 12.89, 15.18 & $14.31\pm1.24$ (9\%) \\
\bottomrule
\end{tabular}
\caption{T1's own ground-truth re-measurement at the trained size. The
original single-seed value (15.88) was the $t{=}40$ estimate of what
turned out to be the highest of three seeds; the window used throughout
\S\ref{sec:g2_resolved} was also not fully converged (compare the
$t{=}40$ and $t{=}120$ columns for seed 0 alone: $15.88\to14.86$).
Seed-to-seed spread at fixed averaging time ($\approx$9--11\%) is
comparable to several of the effects \S\ref{sec:g2_resolved} reported
(the $2\times$ density drift, 4.5\%) --- exactly the measurement-noise
concern \S\ref{sec:g2_resolved} itself raised as a caveat.}
\label{tab:t1_truth_1x}
\end{table}

Revising against the corrected mean: the native-size model error
(12.36$\pm$0.67 across all 9 model$\times$IC-seed combinations at
$1\times$, Table~\ref{tab:t1_model_1x}) moves from the originally
reported 19.56\% to $\approx$13.6\%.

\textbf{The unplanned finding: one of the two retrained propagators
collapses under transfer, invisibly at the trained size.} Table
\ref{tab:t1_model_collapse} gives all 27 model-side results.

\begin{table}[h]
\centering
\begin{tabular}{lccc}
\toprule
Model seed & $1\times$ (native) & $2\times$ & $3\times$ \\
\midrule
m0 (Section 233) & 11.69, 12.77, 13.13 & 31.70, 31.99, 32.13 & 45.40, 47.28, 47.48 \\
m1 (retrain) & 12.56, 12.72, 13.24 & 29.74, 32.40, 34.49 & 45.10, 45.88, 46.77 \\
\textbf{m2 (retrain)} & 11.45, 11.67, 12.03 & \textbf{0.00, 0.00, 1.28} & \textbf{1.16, 1.21, 2.11} \\
\bottomrule
\end{tabular}
\caption{Each cell lists the model's zero-retraining $D_{\mathrm{KY}}$
across 3 IC seeds at that size. m0 and m1 are mutually consistent at
every size and match the true values in Table~\ref{tab:g2_3x_extension}
reasonably well. \textbf{m2 is indistinguishable from m0/m1 at the size
it was trained on}, then collapses to a near-fixed point
($D_{\mathrm{KY}}\approx0$--$2$) at both transfer sizes, IC-independently
--- the same signature \S\ref{sec:g2_resolved}'s opening paragraph
reports for the abandoned anisotropic propagator, but arising here from
nothing more than a different random training seed of the IDENTICAL
architecture, data, and recipe that gave m0/m1.}
\label{tab:t1_model_collapse}
\end{table}

\begin{table}[h]
\centering
\begin{tabular}{lccc}
\toprule
Model seed & IC seed 0 & IC seed 1 & IC seed 2 \\
\midrule
m0 & 12.77 & 11.69 & 13.13 \\
m1 & 13.24 & 12.72 & 12.56 \\
m2 & 12.03 & 11.67 & 11.45 \\
\bottomrule
\end{tabular}
\caption{The $1\times$ column of Table~\ref{tab:t1_model_collapse}
spelled out: mean $12.36\pm0.67$ across all 9 values, no visible
difference between m2 and the other two seeds.}
\label{tab:t1_model_1x}
\end{table}

\textbf{This is a third failure category, not a restatement of T4.} T4
(\S\ref{exp:t4}) was designed to distinguish a physics failure (error
rises with $L$ because $\Delta$ rises) from an architecture failure
(error is high but flat in $L$, the 2D KS signature). m2's failure is
neither: it is IC-independent collapse to near-zero at every transfer
size while being completely unremarkable --- same loss, same native-size
spectrum --- at the trained size. Nothing in Stage-1/Stage-2 training
metrics or in a native-size Lyapunov check (this document's own standing
pre-transfer diagnostic, \S\ref{sec:g2_resolved}'s opening paragraph)
would have caught it before transfer was actually tested. Call this a
\textbf{training-instability failure}: the recipe does not reliably learn
a propagator whose zero-shot-resized dynamics are well-behaved, even
though it reliably learns one whose dynamics are well-behaved at the
size it was shown.

\textbf{Consequence for \S\ref{sec:g2_resolved}'s headline claim.} The
4.48\% transfer result was measured on m0 alone. With $n=3$ retrains, one
collapsed, which is too small a sample to put a number on the failure
rate, but is enough to show it is not negligible. ``Drag restores
extensivity and this architecture transfers'' needs a fourth,
seed-robustness precondition alongside the physics check (G1/G2) and the
architecture check (T4): \emph{does this specific trained propagator
transfer}, checked directly, not inferred from its training loss or its
behavior at the trained size.

\textbf{What is and is not yet known.} Whether m2 is representative of
some fixed fraction of training seeds, or an unlucky single draw, is not
answered by $n=3$; the training-side literature does report multistable
collapse-vs-chaos outcomes for superficially identical local-field
training runs elsewhere in this project (the earlier v3 architecture-search
episode, \S\ref{sec:residual_ladder_revision}), so a nonzero failure rate
for this family is plausible rather than surprising. A direct comparison
of m2 against m0/m1 (what differs in the learned weights, and whether it
is detectable without running a transfer test) is planned next, before
resuming T1's remaining ground-truth runs; see \S\ref{sec:m2_autopsy} below.

\subsection{Autopsy: what differs about m2's weights (2026-10-01)}
\label{sec:m2_autopsy}

Since `local\_mlp\_2d`'s site-mixer (\texttt{\_LocalMixerBlock2D.mix}) has
no softmax --- it is a genuinely linear, shift-invariant circular
\texttt{Conv2d} over the site grid, the token-mixing step of each block
--- its per-frequency gain is exact and computable in closed form: the
2D DFT of the kernel at frequency $(k_x,k_y)$ is a
$d_{\mathrm{model}}\times d_{\mathrm{model}}$ complex matrix, and the
spectral radius of that matrix is the operator's amplification factor for
a perturbation at that frequency. Training only ever constrains this
kernel at the 8 discrete $x$-frequencies the trained site grid admits;
resizing to $2\times$/$3\times$ introduces new, never-trained
frequencies in between.

\begin{table}[h]
\centering
\begin{tabular}{lcccc}
\toprule
Model & Block & Max gain, trained freqs & Max gain, new freqs ($2\times$) & Max gain, new freqs ($3\times$) \\
\midrule
m0 & 0 & 2.02 & 1.58 & 1.82 \\
m0 & 1 & 2.04 & 1.82 & 1.90 \\
m1 & 0 & 2.12 & 1.49 & 1.78 \\
m1 & 1 & 1.87 & 2.02 & 1.99 \\
\textbf{m2} & \textbf{0} & \textbf{4.41} & \textbf{3.44} & \textbf{3.96} \\
m2 & 1 & 2.27 & 2.03 & 2.16 \\
\bottomrule
\end{tabular}
\caption{Spectral radius of the circular site-mixer's per-frequency gain
matrix (\texttt{scripts/exp\_t1\_m2\_autopsy.py}). m2's block-0 kernel peaks at roughly $2\times$ the gain
of the corresponding m0/m1 kernel, and --- unlike m0/m1, whose gain at
genuinely new (never-trained) frequencies drops noticeably below their
own trained-frequency peak --- m2's gain at new frequencies stays close
to its already-large trained peak. Bulk kernel norm does \emph{not}
distinguish the models (Frobenius norms 5.4--6.1 for block 0, 7.8--8.0
for block 1, m2 unremarkable on this measure): m2's kernel is not bigger,
its gain is more sharply concentrated onto one frequency.}
\label{tab:m2_gain}
\end{table}

\textbf{This is a real, measurable structural difference, found without
ever running a transfer test} --- a candidate pre-transfer screening
signal distinct from both the native-size Lyapunov check
(\S\ref{sec:g2_resolved}'s opening paragraph, which did not catch m2
either, since collapse there was a transfer-size-only loss of chaos, not
a loss at the trained size) and from training-loss metrics.

\textbf{What the autopsy does NOT yet establish.} A direct empirical
check --- rolling out the base (nonlinear) trajectory for 300 steps at
each of 7 resize factors from $1\times$ to $3\times$ and tracking both
the state norm and the step-to-step state change in the tail --- shows no
qualitative difference between m0/m1/m2 at any scale tested: all three
have comparable, non-vanishing tail norm variance and step-to-step
motion (coefficient of variation $\sim10^{-3}$ throughout, no model
settling toward a visibly static trajectory within 300 steps). The
$D_{\mathrm{KY}}\approx0$--$2$ collapse reported in
Table~\ref{tab:t1_model_collapse} is a statement about the
\emph{linearized tangent map} over a much longer window (2000 steps after
200 warmup, vs.\ this probe's 300), not necessarily about the gross shape
of the base trajectory; a single high-gain linear mode in one block need
not dominate the full nonlinear, LayerNorm-regularized system's base
dynamics even while erasing its Lyapunov expansion. \textbf{The causal
chain from ``block-0 gain is $2\times$ higher'' to ``Lyapunov spectrum
collapses under transfer'' is plausible and consistent with every
measurement so far, but not yet closed.} The direct test --- short-horizon
tangent-vector growth rate at each resize factor, which would show
exactly where along the $1\times\to3\times$ ramp the collapse actually
sets in --- was not run in this pass and is a natural next step before
generalizing the gain-matrix diagnostic into a cheap pre-transfer screen
(a candidate addition to T5, \S\ref{exp:t5}).

\section{A frozen, scale-incremental hierarchy: residual models as a physics-discovery probe}
\label{sec:residual_hierarchy}

Proposed in chat (2026-09-29), directly motivated by G1/G2's own result:
if zero-retraining transfer of a single fixed-size model degrades
substantially at both 2D rungs (77\%/50\% relative error), a natural
alternative is to stop asking one model to generalize to every size for
free, and instead build a \emph{tower} of models, one per
size-doubling, where each new level is trained once and everything below
it is frozen forever. This section documents the idea, a concrete
architecture, a literature check, and a numbered plan (R1--R\ldots),
mirroring this document's own discipline (Section~\ref{sec:2d_transfer_plan})
of writing the plan down before running it.

\subsection{The core idea}

Level 0 is an existing, already-trained checkpoint (Section 230 or 231),
used completely unchanged --- exactly the zero-retraining transfer
mechanism already in place. Level 1 is a NEW, separate, small model,
trained only once, whose entire job is to predict the residual between
the true dynamics at the transfer size and what the frozen Level-0
model already produces there. Level 0 never receives a gradient from
this training. If a still-larger domain is wanted later, Level 2 is
added on top with Levels 0 and 1 both frozen, and so on.

This changes the claim from ``one model transfers to any size for
free'' (which G1/G2 just showed does not hold cleanly for either 2D
system) to ``a fixed, growing tower of models, one bounded new training
job per size-doubling, nothing below ever retrained'' --- a weaker but
more realistic claim, and one with a much closer analogy to how physics
actually builds effective theories level by level (Section~\ref{sec:lit_residual}).

\textbf{The idea's real payoff, per the user's own framing, is diagnostic
as much as practical}: examining the TRAINED Level-1 model's own
behavior --- how much capacity it needs, what spatial pattern its
correction takes, whether it stays small/simple or is large/persistent
--- is itself a way to characterize what, if anything, is genuinely new
about the dynamics at the larger scale, without having to already know
the answer in closed form. This is a novel empirical probe of
scale-locality, complementary to this project's own existing D3
(coupling-graph bandedness) and D7 (light-cone/spreading-velocity)
diagnostics, but measured through ``how hard is it to correct a frozen
local model at a bigger scale'' rather than a direct correlation-length
measurement. Concretely, the SAME methodology applied to both 2D systems
should give a DIFFERENT answer for a principled reason: G1 already
showed 2D KS's own physics is confirmed extensive (clean
$D_{\mathrm{KY}}/\mathrm{Area}$ match across two domain sizes), so its
77\% gap should be patchable by a small, simple Level-1 correction ---
if instead Level 1 turns out to need substantial, persistent capacity
for 2D KS too, that would be a real surprise worth understanding. G2's
own physics is NOT yet confirmed extensive (super-linear
$D_{\mathrm{KY}}$ growth); if Level 1 needs substantially MORE capacity
for Kolmogorov flow than for 2D KS, that contrast is itself evidence the
two systems' transfer failures are different in kind, not degree ---
exactly the distinction this project has been trying to establish since
G1/G2 landed.

\subsection{Architecture: why a coarser grid, not more channels at the fine grid}
\label{sec:residual_architecture}

The user's own instinct --- widen the receptive field, but keep the
added latent dimensionality modest rather than letting it explode --- is
right, but the natural way to satisfy both is a coarser SITE grid for
Level 1, not more CHANNELS at Level 0's existing fine site density.
\textbf{Channels alone do not buy spatial reach}: adding more channels
per site at the SAME site density gives the model more numbers to
describe each existing physical patch, but does not let any layer's
receptive field span more physical distance --- that is set entirely by
window size (in site-hops) and site density (physical area per site),
not channel count. A coarser grid, by contrast, gets a wider PHYSICAL
receptive field almost for free: a modest window (even $w'=2$) on a
grid pooled by a factor of, say, 4 already covers $4\times$ the physical
distance per site-hop that the same window covers on the fine grid, at
a fraction of the parameters a fine-grid window wide enough to match
that physical reach directly would need. This is exactly
Section~\ref{sec:model}'s own multiscale architecture's ``regional''
scale, repurposed here specifically as the cross-domain-size residual
level rather than a same-size auxiliary channel --- the SAME
mechanically-precedented idea (Section 227 built this once already),
now aimed at a problem it is actually well-suited for (Section
227 itself under-performed Section 224 on 1D KS precisely because 1D KS
has no genuine large-scale structure for a regional/global scale to
capture --- see Section~\ref{sec:model}'s own note; both 2D systems'
G1/G2 results are the first evidence in this project that such structure
might actually exist to capture).

Concretely:
\begin{itemize}[nosep]
\item \textbf{Level 1 encoder}: a SEPARATE small local-field encoder,
  patchifying the SAME physical field at a much coarser site density
  (e.g.\ pool factor 4--8 relative to Level 0's own site grid), producing
  a small number of new residual channels (1--2, not matching Level 0's
  own $c{=}3$) at this coarse grid. Crucially, this coarse encoder is
  built from the SAME size-independent Conv2d machinery
  (\S\ref{sec:model}) --- so it inherits zero-retraining transfer to yet
  LARGER domains automatically, exactly like Level 0 does, keeping the
  whole hierarchy size-agnostic at every level, not just the base one.
\item \textbf{Level 1 propagator}: a small \texttt{local\_mlp\_2d}-family
  propagator operating on the coarse residual field, additionally
  conditioned on the (frozen) Level-0 latent state at the corresponding
  location --- giving Level 1 explicit access to what Level 0 already
  ``thinks'' is happening locally, so it only needs to learn the part
  Level 0 gets wrong or cannot see.
\item \textbf{Decode}: $\hat u = D_0(z_0) + \mathrm{upsample}(D_1(z_1))$
  --- Level 0's own frozen decode, plus Level 1's coarse correction
  upsampled back to the fine physical grid (the SAME transposed-conv
  machinery Level 0's own decoder already uses, at the coarse grid's own
  patch size).
\item \textbf{Training}: freeze Level 0 entirely (no gradients). Train
  Level 1's encoder + propagator on the residual between the TRUE
  trajectory at the transfer size and Level 0's own zero-retraining
  decoded rollout there --- i.e.\ Level 1 is trained specifically to fix
  what transfer alone gets wrong, not to reproduce the attractor from
  scratch.
\end{itemize}

\subsection{Literature check and honest novelty assessment}
\label{sec:lit_residual}

Each individual ingredient here is well precedented; the specific
combination --- a frozen, incrementally-grown hierarchy solving DOMAIN-SIZE
(not grid-resolution) extrapolation for a chaotic PDE latent-field
surrogate, with the trained residual model's own structure read as a
diagnostic of emergent scale-dependent physics --- does not appear to
be a documented combination, based on a quick literature check (2026-09-29,
not exhaustive).

\begin{itemize}[nosep]
\item \textbf{Closure/residual modeling is a large, mature field}
  \cite{mlclosurereview2024} --- neural networks trained to predict the
  residual between a coarse-grid simulation and the true high-resolution
  physics (large-eddy-simulation subgrid closures, Reynolds-stress
  models). \cite{coarsegrainneuralop2024} applies this directly to
  chaotic systems via neural operators. The closure literature's own
  interpretability work is a direct precedent for ``read the trained
  correction to learn what's missing'' --- but essentially always in the
  context of GRID coarsening at FIXED domain size, not domain-SIZE
  extrapolation at fixed grid resolution, the axis this project's own
  transfer claim is about.
\item \textbf{The domain-size generalization failure itself is a live,
  independently-confirmed problem in the neural-operator literature},
  not an artifact of this project's own implementation:
  \cite{zeroshotlargerdomains2026} states plainly that ``existing neural
  operator approaches for PDEs implicitly assume a fixed domain size,''
  and \cite{falsepromisezeroshot2026} is explicitly skeptical of
  zero-shot claims in this exact setting. This independently corroborates
  G1/G2's own finding (77\%/50\% relative error) as a real, general
  phenomenon, not a bug specific to this project's architecture ---
  reassuring in one sense (not a project-specific failure), sobering in
  another (the wider field has not solved this either). The specific
  FIXES proposed there (attention-logit locality biases, rotary
  positional embeddings) are single-model architectural changes, a
  different strategy from the frozen-hierarchy approach here.
\item \textbf{Renormalization-group/deep-learning correspondences are
  well established} \cite{mehtaschwab2014,mlrgstatphys2023} --- but
  almost entirely in STATISTICAL physics (spin models, phase
  transitions, feature-learning-as-coarse-graining inside a single
  network's own layers), not applied as an explicit, incrementally-grown
  training PROCEDURE for a dynamical chaotic-PDE surrogate.
\item \textbf{Model-discrepancy / hybrid modeling}
  \cite{kennedyohagan2001,discrepancyframework2022} --- learning a
  correction to a trusted-but-incomplete model and treating the learned
  correction itself as informative about missing physics, going back to
  Kennedy \& O'Hagan (2001). Conceptually the closest existing precedent
  to this section's own ``residual as diagnostic'' framing, but not
  previously combined with a FROZEN, size-incremental hierarchy
  specifically for chaotic-PDE domain-size transfer.
\end{itemize}

\textbf{If this works, the implications extend past this specific
project.} A validated ``frozen tower, bounded cost per doubling''
recipe would be a genuinely practical answer to a problem
\cite{zeroshotlargerdomains2026,falsepromisezeroshot2026} show the wider
neural-operator field has not resolved --- not by inventing a better
single-model architecture, but by giving up on the (apparently
currently unrealistic) demand that one fixed-size model generalize
perfectly to every size, and replacing it with a bounded, interpretable,
incremental-cost alternative. The diagnostic use --- reading the trained
residual's own required complexity as a direct, empirical measurement of
a PDE's scale-locality --- would additionally be a new, general-purpose
tool: applied across many systems, it would let scale-locality be
MEASURED rather than assumed, exactly the property Section~\ref{sec:extensivity}'s
whole architecture already leans on for 1D KS without a direct empirical
test of the kind this section proposes.

\subsection{Revision (2026-09-29): build the ladder from scratch, progressively, starting small}
\label{sec:residual_ladder_revision}

User-directed, after this section's own first draft above: rather than
bolting a single Level-1 correction onto the already-large, already-expensive
Section 230/231 checkpoints, build an entirely NEW, self-contained ladder
of domain sizes from scratch, starting at the smallest domain that is
genuinely chaotic and doubling upward. Two reasons, both worth stating
plainly since they reshape the plan below: \textbf{(1) efficiency} --- a
small base domain is cheap enough (see the chaos search below: a full
40-direction Lyapunov check at $L{=}16$ takes about a minute, against
G1's own $11803$s at $L{=}66$) to debug and validate the whole
frozen-hierarchy training PROCEDURE fast, before ever spending G1/G2-scale
compute on it; \textbf{(2) direct observability}. With only one
correction step (the original plan), there is exactly one before/after
snapshot to look at. With a genuine multi-rung ladder, each new level's
own trained correction --- its required capacity, its spatial/spectral
character, how it interacts with what the frozen levels below it already
predict --- becomes directly inspectable at every rung, not just once.

\textbf{This directly opens the ``cascading chaos'' question}: does the
amount of NEW correction needed at each successive doubling SHRINK
(the system is settling into its extensive, ``one more copy of the same
local unit'' regime --- a cascade converging to a fixed point, in RG
language), stay roughly CONSTANT (a genuinely scale-invariant cascade,
no characteristic scale), or GROW (the system keeps discovering
qualitatively new large-scale structure without bound)? This project
already has one, static, two-point version of this answer for 2D KS ---
G1's own $D_{\mathrm{KY}}/\mathrm{Area}$ ratio is nearly identical at
$L{=}32$ ($0.0444$) and $L{=}66$ ($0.0442$), implying SOME such
convergence already happened by $L{=}32$ --- but a full ladder would let
that convergence be WATCHED happening, rung by rung, rather than only
compared at two endpoints. If the required residual capacity tracks the
ground-truth ratio's own convergence, that is a direct, mechanistic
confirmation of what ``becoming extensive'' looks like from the model's
own point of view, not just the PDE's.

\textbf{Design choice: fix the physical patch length once, across the
whole ladder.} Every level's own encoder/propagator stays built from the
same size-independent Conv2d machinery (\S\ref{sec:model}) at a FIXED
physical patch length, so each level is itself zero-retraining-resizable
to the next rung exactly like Section 230/231 already are --- only the
ADDITION of a new level (never the resizing of an existing one) requires
new training. Concretely: patch length $2.0$ physical units throughout
(patch $=8$ grid points at $dx{=}0.25$), $n_{\mathrm{sites}}$ doubling
with $L$ ($8\to16\to32\to\ldots$), \texttt{local\_channels=3},
\texttt{site\_mix\_radius=2}, \texttt{attn\_window=3} --- Section 230's
own values, chosen deliberately unchanged so the ladder's own Level-0
rung is a direct, like-for-like miniature of the existing recipe, not a
new architecture search. One useful, deliberate consequence of holding
\texttt{attn\_window} fixed while $L$ grows: the propagator's OWN
receptive field, a fixed $14$ physical units
($(2\times3{+}1)$ sites $\times\ 2.0$), covers a SHRINKING fraction of
the domain at each rung ($87\%$ at $L_0{=}16$, down to $11\%$ by
$L_3{=}128$) --- letting the ladder directly test
\S\ref{sec:2d_transfer_plan}'s own architecture-side hypothesis (a fixed
window covers a shrinking fraction of the domain faster in 2D than 1D)
by checking whether required residual capacity tracks that shrinking
fraction.

\textbf{Revision (2026-09-29): Level 0's first attempt under-captured
its own domain's chaos, diagnosed directly, not assumed.} Trained as
specified above (\texttt{attn\_window=3}, \texttt{local\_channels=3},
$d_{\mathrm{lat}}=192$) and checked with this project's own standard
rollout+Lyapunov diagnostic (never trust a validation-loss number alone)
rather than declaring success from \texttt{val\_kmax\_mse} alone:
rollout stayed finite and bounded (no collapse), but the trained
propagator's own Lyapunov spectrum gave $\lambda_1=0.010$,
$n_{\mathrm{positive}}=1/192$, $\boldsymbol{D_{\mathrm{KY}}=2.21}$ ---
against the TRUE $L_0{=}16$ system's own $D_{\mathrm{KY}}\approx10.3$
(8 positive exponents, the base-domain search table above) a
substantial under-capture, not a close match. Best-available hypothesis:
\texttt{attn\_window=3} covering $87\%$ of this tiny domain is
effectively a near-global, not a local, mixer, and $d_{\mathrm{lat}}=192$
is plausibly oversized for an attractor whose true dimension is only
$\sim10$. \textbf{Corrected (user-directed): $d_{\mathrm{lat}}=32$
($n_{\mathrm{sites}}{=}4{\times}4$, \texttt{local\_channels=2}, patch
length $4.0$ physical units, not $2.0$) and \texttt{attn\_window=2}.}
This changes the ladder's own fixed-patch-length convention set above
(now $4.0$ physical units and \texttt{local\_channels=2} throughout, not
$2.0$/3) --- $n_{\mathrm{sites}}$ at each rung is correspondingly halved
from the original plan ($4\to8\to16\to32$, not $8\to16\to32\to64$).

\textbf{v2 result: improved but still under-captures.} $\lambda_1=0.019$,
$n_{\mathrm{positive}}=2/32$, $D_{\mathrm{KY}}=3.20$ --- better than v1's
$2.21$ but still far below the true $\approx10.3$. Diagnosed further: at
$n_{\mathrm{sites}}=4$, \texttt{attn\_window=2} with the propagator's own
default \texttt{n\_layers=2} gives an EFFECTIVE receptive field of
$n_{\mathrm{layers}}\times\texttt{attn\_window}=4$ sites --- saturating
a 4-site ring just as completely as v1's nominally-wider window did.
This surfaced a real gap in the v1/v2 discussion above: effective
receptive field is $n_{\mathrm{layers}}\times\texttt{attn\_window}$, not
\texttt{attn\_window} alone (user-directed correction, 2026-09-29).

\textbf{v3 (user-directed): $n_{\mathrm{sites}}=8$,
\texttt{local\_channels=1} ($d_{\mathrm{lat}}=64$), \texttt{n\_layers=1},
\texttt{attn\_window=2} (effective RF$=2$ sites, $62.5\%$ ring coverage,
genuinely partial).} \texttt{val\_kmax\_mse} improved dramatically to
$\mathbf{0.0108}$ --- BETTER than production scale's own $0.03$. But the
Lyapunov diagnostic told a completely different story:
$\lambda_1=-1.405$ (strongly, not marginally, negative),
$n_{\mathrm{positive}}=0/64$, $\boldsymbol{D_{\mathrm{KY}}=0}$ ---
genuine collapse, the best-looking validation loss of the three attempts
corresponding to the WORST dynamics. Not a clean fixed point either
(\texttt{max}$|z|$ drifts $2.3\to4.5$ over the 2000-step rollout, two ICs
stay separated at $25.4$) --- a persistent, non-chaotic drift, not
settled convergence. \textbf{A clean, direct demonstration of why this
project never trusts \texttt{val\_kmax\_mse} alone}: minimizing it
further pushed training toward an easier-to-fit CONTRACTING map, not
richer chaos.

\textbf{v4 (user-directed): hold v3's architecture fixed, test a
different axis --- finer SOLVER grid.} $NX{=}128$ at the SAME
$L_0{=}16$ domain ($dx{=}0.125$, half of v1--v3's $0.25$), patch length
still fixed at $2.0$ physical units so $n_{\mathrm{sites}}$ doubles to
$16$ automatically (not just more solver resolution per site, genuinely
more latent sites too --- effective RF$=2$ sites now covers only
$31.25\%$ of the ring, down from v3's own $62.5\%$). \texttt{attn\_window},
\texttt{n\_layers}, \texttt{local\_channels} all held at v3's own values
for a controlled comparison --- isolating whether v3's collapse was a
receptive-field/coverage issue after all (this change would fix it) or a
different problem, e.g.\ \texttt{n\_layers=1} simply being too shallow a
propagator body regardless of locality (this change would not fix it).
Dataset regenerated (not just resized) at the finer grid, sized from the
start to avoid the L1 truncation episode's own after-the-fact fix
($\sim$0.5GB).

\textbf{v4 result: worse, not better --- and confounded.}
$\lambda_1=0.0021$, $n_{\mathrm{positive}}=1/512$,
$D_{\mathrm{KY}}=1.19$ --- below BOTH v1 ($2.21$) and v2 ($3.20$),
despite restoring \texttt{n\_layers=2} and $8\times$ more latent
dimensions. Caught directly rather than accepted at face value: v4's
own dataset had only $61$ snapshots/trajectory (\texttt{trajectory\_
time=6}, sized down following the L1-truncation lesson) against
v1--v3's $251$ (\texttt{trajectory\_time=25}) --- a real, unintended
FOURTH confounded variable (alongside \texttt{dx}, $n_{\mathrm{sites}}$,
\texttt{n\_layers}) that was never isolated.

\textbf{Revision 2 (2026-09-29, user-directed): step back from the
from-scratch progressive ladder.} The small-$L_0$ architecture search
(v1--v4) accumulated too many simultaneously-varying knobs
(\texttt{attn\_window}, \texttt{n\_layers}, $n_{\mathrm{sites}}$,
\texttt{local\_channels}, \texttt{dx}, dataset length) without ever
isolating one at a time cleanly enough to draw a confident conclusion
--- a real, documented finding in its own right (the \texttt{local\_mlp\_2d}
recipe's own behavior at very small domains is evidently more fragile
and knob-sensitive than at production scale, where Section 230's exact
same architectural family works robustly), but not a productive path to
keep pursuing for the hierarchy experiment itself. \textbf{Pragmatic
pivot}: reuse what is ALREADY validated rather than build a new small
base model from scratch. Level 0 is Section 230's own existing
checkpoint (Section 230/231, $D_{\mathrm{KY}}=71.9$ at its own trained
$44\times44$ size, no retraining). Level 1 is trained directly at G1's
own transfer size ($L_x{=}L_y{=}66$, $NX{=}192$) --- the exact size
where the gap to close is already precisely known (G1: model
$D_{\mathrm{KY}}=43.6$ vs.\ true $192.5$, $77\%$ relative error), rather
than a small-domain rung whose own eventual relevance to that gap was
never established. A real trajectory dataset at this size did not
previously exist (G1 itself only ever computed Lyapunov exponents, no
saved trajectories) --- generated fresh, sized correctly from the start
this time ($n_{\mathrm{train}}{=}24$, $n_{\mathrm{val}}{=}6$,
\texttt{trajectory\_time=10}, $\approx1.8$GB, matching the production
dataset convention).

\subsection{Why the transfer specifically fails: an orientation hypothesis, and a three-part follow-up (2026-09-29)}
\label{sec:orientation_hypothesis}

A chat discussion following the L1 dataset's generation asked directly
why 2D KS's zero-retraining transfer FAILS this badly while 1D KS's own
transfer (Section~\ref{exp:e1}, sub-2\% error at every size) does not
--- both systems gain MORE Fourier modes as $L$ grows, so mode count
alone cannot be the distinguishing factor. \textbf{Sharpened hypothesis}:
what differs is the KIND of new modes. 1D's new modes all sit on the
same axis (finer wavenumber magnitude, same pattern family) --- exactly
what translation-equivariant weight-sharing already generalizes over by
construction. 2D's new modes appear at new ANGLES around the isotropic
critical ring --- genuinely new pattern ORIENTATIONS, not finer
modulation of an already-seen family. An isotropically-SHAPED Conv2d
kernel guarantees the architecture treats every direction symmetrically,
but its LEARNED weights are still free parameters --- nothing guarantees
they generalize correctly across orientation combinations one 44x44
training run never actually realized. Not yet verified directly
(candidate checks: trained-kernel anisotropy; true attractor's own
angular power spectrum at $44\times44$ vs.\ $66\times66$) --- an open
hypothesis, not a settled fact.

\textbf{Three-part follow-up, user-directed, in this order}:
\begin{enumerate}[nosep]
\item \textbf{Hierarchical training at L1, first.} Proceed with the
  frozen-Level-0 + trained-Level-1 plan (\S\ref{sec:residual_architecture})
  using the $L_1{=}66$ dataset just generated, independent of whether
  the orientation hypothesis above is even correct --- tests the
  hierarchy mechanism itself, not the diagnosis.
\item \textbf{Retrain Section 230 on $D_4$-augmented data.} Since
  $L_x{=}L_y$, 90/180/270-degree rotations and reflections (the 8-element
  dihedral group $D_4$) are EXACT, lossless symmetries of both the
  square periodic domain and the isotropic KS equation --- augmenting
  Section 230's own training trajectories with all 8 $D_4$ transforms
  and retraining tests whether the gap is simply insufficient
  orientation diversity in the ORIGINAL training set (fixable with zero
  architecture change), before committing to a bigger architecture
  change.
\item \textbf{Steerable convolutions, same $D_4$-augmented data.} If (2)
  does not fully close the gap, replace the free-parameter Conv2d
  kernels in the encoder/propagator with genuinely rotation-equivariant
  (steerable) kernels --- Weiler \& Cesa's E(2)-equivariant CNN
  framework, or the earlier Cohen--Welling steerable-CNN construction ---
  so orientation generalization is guaranteed by construction rather
  than hoped for from data alone. Physically well-motivated here
  specifically: 2D KS's own governing equation has no preferred
  direction at all, so full rotation equivariance is a genuine symmetry
  of the underlying physics, not merely a convenient inductive bias
  borrowed from vision tasks where it usually is not.
\end{enumerate}

\textbf{System choice: 2D KS first, not Kolmogorov flow.} G1 already
confirmed 2D KS's own physics is genuinely extensive (the
$D_{\mathrm{KY}}/\mathrm{Area}$ match above) --- the cleanest available
test bed for validating the HIERARCHY MECHANISM itself (does
freeze-and-add-a-residual actually work, close gaps, show a sensible
convergent pattern) before applying the identical recipe to Kolmogorov
flow's own still-open extensivity question (G2). Kolmogorov flow remains
the natural next system once the mechanism is validated here, not a
parallel effort --- the same "confirm each rung before extending"
discipline this document has used throughout (Section~\ref{sec:ks}'s own
gate-before-Phase-2 rule, Section~\ref{sec:next_steps}'s system ladder,
G2's own gate on G1 completing first).

\textbf{Base-domain search, run 2026-09-29 (40 directions, \texttt{warmup\_
time=20}, \texttt{total\_time=40}, lighter than G1's own 300-direction
rigor --- adequate for picking $L_0$, not a publication-grade number):}

\begin{table}[h]
\centering
\begin{tabular}{lccccc}
\toprule
$L$ & $NX$ & $\lambda_1$ & $n_{\mathrm{positive}}$/40 & $D_{\mathrm{KY}}$ & $D_{\mathrm{KY}}/\mathrm{Area}$ \\
\midrule
12 & 48 & 0.154 & 3  & 6.07  & 0.0422 \\
16 & 64 & 0.048 & 8  & 10.27 & 0.0401 \\
20 & 80 & 0.124 & 9  & 16.60 & 0.0415 \\
24 & 96 & 0.200 & 14 & 25.83 & 0.0448 \\
\bottomrule
\end{tabular}
\caption{Every candidate is already genuinely chaotic, even $L{=}12$
($\approx1.35\times$ the pattern length $2\sqrt2\pi\approx8.89$) ---
2D KS's isotropic instability ring apparently supports richer coupling
at small domains than 1D's own two discrete critical modes do, and the
$D_{\mathrm{KY}}/\mathrm{Area}$ ratio is already within $\sim10\%$ of the
$L{\in}\{32,66\}$ asymptotic value ($0.044$) at every size tested here
--- a notable finding in its own right: 2D KS appears to reach its
extensive regime much faster, in units of pattern lengths, than 1D KS
does (whose own E1 result showed a much larger, $\sim70\%$, deviation at
a comparable multiple of ITS OWN pattern length, $L{=}25\approx2.8\times$).}
\end{table}

\textbf{$L_0=16$ chosen}: genuinely non-trivial ($D_{\mathrm{KY}}\approx10$,
8 positive exponents) but cheap ($NX{=}64$; a full 40-direction Lyapunov
check takes about a minute). Ladder: $L_0{=}16$ ($NX{=}64$,
$n_{\mathrm{sites}}{=}8$), $L_1{=}32$ ($NX{=}128$, $n_{\mathrm{sites}}{=}16$),
$L_2{=}64$ ($NX{=}256$, $n_{\mathrm{sites}}{=}32$), $L_3{=}128$
($NX{=}512$, $n_{\mathrm{sites}}{=}64$, a stretch goal --- comparable
grid cost to G1's own $NX{=}192$ run) --- note $L_1{=}32$ coincides with
Step 1's own original small-test domain (\S\ref{step:2dks}'s progress
note, $D_{\mathrm{KY}}=45.5$ there), and $L_2{=}64$ is close enough to
G1's own transfer size ($66$) to let the ladder's CUMULATIVE prediction
at $L_2$ eventually be compared directly against G1's own single-model
transfer result ($43.6$ vs.\ true $192.5$) at essentially the same size
--- the natural convergence point validating whether the hierarchical
approach actually beats direct zero-shot transfer, not just whether it
runs.

\subsection{Plan (R1--R6)}
\label{sec:residual_plan}

\begin{description}[style=nextline,leftmargin=1.2em]
\item[R1 --- Ladder datasets, from scratch.] \label{exp:r1}
For each rung $L_k\in\{16,32,64,\ldots\}$: verify chaos directly (done
for $L_0$ above; repeat at each new rung before trusting it, same
discipline as every other gate in this document), then generate a
trajectory dataset via \texttt{generate\_trajectory\_dataset2d}
unchanged, sized for the rung (small at $L_0$, growing with area).
\item[R2 --- Level-0 architecture and training, from scratch, at $L_0$.] \label{exp:r2}
A fresh, small \texttt{KSAutoencoderLocalField2D} + \texttt{local\_mlp\_2d}
propagator at $L_0{=}16$'s own sizing above, Section 224/230's regularizer
recipe unchanged --- NOT a reuse of Section 230's own $L{=}44$ checkpoint;
this is a genuinely new base model, small and cheap to iterate on.
\item[R3 --- Level-1 architecture.] \label{exp:r3}
Build the coarse residual encoder + conditioned propagator described in
\S\ref{sec:residual_architecture}; new, size-independent Conv2d modules,
tested the same way \texttt{KSAutoencoderLocalField2D}/\texttt{local\_mlp\_2d}
were (exact equivariance, gradient flow, shape tests) before any real
training.
\item[R4 --- Train Level 1 at $L_1{=}32$, freezing Level 0.] \label{exp:r4}
Level 0's own zero-retraining prediction at $L_1$ (the existing resize
mechanism) vs.\ R1's true $L_1$ data gives the residual Level 1 must
learn. Record: how much of the gap closes; Level 1's own required
capacity (channels, effective window) and its spatial/spectral pattern
--- the first direct look at ``what does a residual model predicting
large-scale effects actually look like.''
\item[R5 --- Extend the ladder: $L_2{=}64$, then $L_3{=}128$ if
compute allows.] \label{exp:r5}
Add Level 2 (freezing Levels 0+1), then Level 3, repeating R4's own
protocol at each rung. This is where the cascade question
(\S\ref{sec:residual_ladder_revision}) actually gets answered: plot
required Level-$k$ capacity (and closed-gap fraction) against $k$,
alongside the TRUE $D_{\mathrm{KY}}/\mathrm{Area}$ ratio measured
independently at each same rung, and check whether the two converge
together.
\item[R6 --- Kolmogorov flow, once R1--R5 validate the mechanism on
2D KS.] \label{exp:r6}
Repeat the entire ladder for Kolmogorov flow, starting from its own
smallest chaotic multi-forcing-period domain (not yet identified ---
Section~\ref{step:kolmogorov}'s own trained domain has only ONE forcing
period, which \S\ref{sec:2d_transfer_plan}'s own G2 discussion already
flagged as likely too small to have taught genuine local coupling at
all). The key comparison is 2D KS's own ladder (confirmed-extensive
physics) vs.\ Kolmogorov flow's (open question) --- per
\S\ref{sec:residual_hierarchy}'s own prediction, expect 2D KS's required
per-level capacity to shrink/plateau and Kolmogorov flow's to stay
larger/more persistent if its physics is genuinely less local.
\end{description}

\textbf{Status.} Plan revised and dataset generation (R1, $L_0$) started
(2026-09-29). Per this
document's own established discipline, R1's dataset generation begins
immediately below.

%======================================================================
\section{Adding cross-patch coupling (D1--D3), and why it failed}
\label{sec:d1d3}
%======================================================================

\textbf{Question (2026-10-01).} The PDE is local, so a tiled local model
should work. Does the Kolmogorov failure at low drag come from too little
interaction between patches? The full plan is in
\texttt{docs/oct\_1\_exp.md}. One fact from the config changes how to read
every earlier Kolmogorov result. The Section 231/233 propagator has
receptive field $2\ \text{layers}\times\text{window }3 = 6$ sites on an
8-site periodic ring, so it covers 75\% of the training domain. At training
size the ``local'' model is close to global, and it can use periodic
wraparound, a self-interaction that does not exist in a larger domain.

\begin{table}[h]
\centering
\small
\begin{tabular}{lcccc}
\toprule
Experiment & $\alpha$ & native $D_{\mathrm{KY}}$ & $2\times$ error & $3\times$ error \\
\midrule
Sec.\ 233 baseline (window 3) & 0.1 & 12.8 & 4.5\% & 19.4\% \\
D1: window 1 (25\% of ring) & 0.1 & 14.7 & 11.8\% (3 IC) & 7.3\% (3 IC) \\
\midrule
Sec.\ 231 baseline (window 3) & 0.05 & --- & 50.0\% & --- \\
D2: exact Biot--Savart input$^\S$ & 0.05 & 22.3 & 95.4\% (3 IC) & $D_{\mathrm{KY}}\in\{17,18,9\}$ \\
D3: learned regional scale & 0.05 & 19.6 & 98.7\% (3 IC) & $D_{\mathrm{KY}}\in\{2.5,27,5.3\}$ \\
\bottomrule
\end{tabular}
\caption{D1--D3 summary. Truth: $D_{\mathrm{KY}}=15.9/30.3/56.4$ at
$1/2/3\times$ for $\alpha{=}0.1$, and $56.8$ at $2\times$ for
$\alpha{=}0.05$ (no $3\times$ truth exists yet at $\alpha{=}0.05$). D2 and
D3 reuse the Section 231 autoencoder and retrain only the propagator.
$^\S$Computed with a config bug that kept the native box's $L_x$ in the velocity solve. Rerun with the fix, D2 still collapses at $2\times$ ($D_{\mathrm{KY}}$ 0--1.9; \S\ref{sec:c8c9_status}).}
\label{tab:d1d3}
\end{table}

\textbf{What D1 shows.} Narrowing the window to 25\% of the ring makes
$2\times$ worse and $3\times$ better (Table~\ref{tab:d1d3}). This is
non-monotone, so window width is not a reliable control. The structural
reason: the window is counted in sites, and L-transfer adds sites. Any
fixed window covers a shrinking fraction of the domain as $L$ grows,
whatever its trained value.

\textbf{What D2 and D3 show.} Each added a cross-patch coupling: the exact
$1/k^2$ streamfunction velocity (D2), and a trainable regional scale (D3:
$2\times2$ pooling, its own mixer, learned pool and broadcast projections).
Both trained well at native size, and both collapsed under transfer, far
worse than the uncoupled baseline. More expressive coupling made transfer
\emph{worse}. Our working explanation is that every coupling path is
trained only on the input distribution of the 8-site periodic box.
\begin{itemize}[leftmargin=*]
\item D3's regional mixer (window 2) runs on a 4-site regional ring, so it
  saw nothing but whole-ring, wraparound-dominated inputs during training.
\item D2's feedback is exact, but its input amplitude is not: the lowest
  wavenumber drops as $L$ grows, and $|u|\propto 1/k$ for the new modes.
  The learned embedding receives values outside its training range.
\end{itemize}
In both cases the new coupling meets new frequencies and new aggregation
patterns at deployment that it never saw in training. We have not yet
verified this mechanism directly (for example with a
\S\ref{sec:m2_autopsy}-style spectral autopsy), so treat it as a hypothesis.
\S\ref{sec:crop_training} is designed to remove this failure mode by
construction, not to diagnose it.

%======================================================================
\section{Training on crops of a large domain}
\label{sec:crop_training}
%======================================================================

\subsection{The idea}

Simulate one large domain, sample random sub-windows (crops) of it, and
train the weight-shared local model only on crops. Then deploy the same
weights, with circular padding, on a periodic domain of any size, or tile
them across any domain. The setting is also realistic. Real geophysical
data never arrive as a periodic box. They arrive as regional windows of a
larger, open flow, which is exactly the training input this scheme uses.

\textbf{Why this targets the D1--D3 failure mode directly.}
\begin{enumerate}[leftmargin=*]
\item \emph{No wraparound artifact.} A crop has no periodic seam, so the
  model cannot learn the self-interaction that a small box allows
  (\S\ref{sec:d1d3}). This applies to every coupling path, including
  D3's regional mixer.
\item \emph{Training and deployment inputs match.} Each crop is a window
  of the large-$L$ attractor, so it contains the large-scale structures,
  low-wavenumber background and velocity amplitudes the model will meet
  at deployment. In a non-extensive regime the local statistics of a
  small box differ from local statistics inside a large domain. Small-box
  training then fits the wrong local measure, which is a direct
  physical reason for the $\alpha{=}0.05$ failure.
\item \emph{Transfer is no longer extrapolation.} The model never learns
  ``a box of size $L_0$''. It learns a local update rule from samples of
  the target system, and deploying at $2\times$ or $3\times$ applies the
  same rule to more sites.
\end{enumerate}

\subsection{Making it work: the boundary problem}

A crop has no physical boundary condition. Information from outside
enters through its edges. We do not need a different encoder; we only
need to be careful about which sites are scored. All learned layers here
(patchify, finite-radius circular mixers, per-site MLPs, transposed-conv
decoder) are weight-shared local operators. So the output at a site
depends only on inputs inside a known light cone. Two training schemes
follow.

\textbf{Variant A, light-cone masking (exact).} Score only sites whose
full receptive field after $k$ rollout steps lies inside the crop:
$R_k = R_{\mathrm{enc}} + k\,R_{\mathrm{prop}}\,(+R_{\mathrm{dec}})$
sites. Padding at the crop edge then cannot reach the loss. The gradient
equals what the same weights would get on the full periodic domain at
those sites. The cost is the halo. With $R_{\mathrm{prop}}{=}6$ sites
(window 3, 2 layers) and $k{=}12$, $R_k\gtrsim80$ sites per side, roughly
$10\times2\pi$. That is unaffordable for multi-step rollout training, but
it remains usable for one-step and short-horizon checks.

\textbf{Variant B, forced lateral boundaries (the NWP answer).} This is
how a limited-area weather model works. Split the crop into three
regions:
\begin{itemize}[leftmargin=*]
\item an outer ring of $R_{\mathrm{enc}}$ sites, discarded because padding
  corrupts their encoding;
\item a forced halo of $R_{\mathrm{prop}}$ sites, overwritten at every
  rollout step by the encoding of the true field at that time;
\item a free-running interior, the only region scored.
\end{itemize}
The interior then evolves under the model alone, while the halo supplies
the correct outside neighbours at every step. The halo stays at
$R_{\mathrm{enc}}{+}R_{\mathrm{prop}}$ sites for every rollout length.
With the D1 window ($R_{\mathrm{prop}}{=}2$), that is about 8 sites. A
crop of $4\times2\pi$ (32 sites) then leaves half its sites scored. At
deployment on a periodic domain no halo is needed, because the local rule
learned in the interior is the same rule applied everywhere.

The one mismatch is that interior sites next to the halo always see true
neighbours in training but predicted neighbours at deployment. That
mismatch is confined to a strip $R_{\mathrm{prop}}$ wide. Variant A at
short horizon measures it directly, so the plan includes that check.

\textbf{Forcing anisotropy.} Kolmogorov forcing depends on $y$. Crop only
in $x$ and keep full periodic $y$-strips ($L_y{=}2\pi$). This matches the
axis used for every L-transfer test so far, so nothing in $y$ changes.

\textbf{Is a vorticity-only latent local enough?} Velocity at an interior
site depends on vorticity outside the crop through the Biot--Savart
kernel, which decays only as $1/r$ in 2D. Real atmospheric data include
observed winds. So the plan includes an ablation that gives the encoder
$(\omega,u,v)$ instead of $\omega$ alone. With $u$ supplied locally, the
vorticity update becomes local, and only the slow far-field contribution
to $u$'s own evolution remains non-local.

\subsection{Experiments (P0--P5)}

The rules below are pre-registered (written before any run).
\begin{description}[leftmargin=*]
\item[P0 --- data.] \textbf{$\alpha{=}0.1$ COMPLETED} (\texttt{kolmogorov\_wide16\_alpha0.1.h5}; field std $1.531$ vs.\ $1.485$ in the $2\pi$ box, so the wide domain carries slightly more variance per point, consistent with mild non-extensivity). $\alpha{=}0.05$ queued. One wide simulation per $\alpha$:
  $L_x{=}16\times2\pi$, $L_y{=}2\pi$, $N_X{=}2048$, $N_Y{=}128$, same
  $\nu$, forcing and $dt_{\mathrm{snap}}$ as now. Use
  $\alpha\in\{0.1, 0.05\}$ first, then $\{0.2, 0.025\}$. About six
  trajectories already give more independent local samples than the
  current 24-trajectory periodic-box datasets, because one snapshot holds
  roughly 16 box-widths. Generate one dataset at a time, never two
  concurrently.
\item[P1 --- sanity at $\alpha{=}0.1$.] \textbf{COMPLETED: fails the $2\times$ criterion (20.8\%), passes $3\times$ (1.3\%); see \S\ref{sec:p1_results}.} Train with variant B (crops
  $4\times2\pi$, window 1). Evaluate zero-shot at $1/2/3\times$ against
  the existing truth. \emph{Pass:} $2\times$ error $\le10\%$ and $3\times$
  error below Section 233's 19.4\%, over 3 IC seeds.
\item[P2 --- the real test, $\alpha{=}0.05$.] \textbf{NEARLY COMPLETE: PARTIAL by the pre-registered rule ($2\times$ error 20.9\%); see \S\ref{sec:p2_results}.} Same recipe. The
  comparison points are the periodic-box baseline (50\%), D2 (95\%) and
  D3 (99\%). \emph{Strong:} single-digit $2\times$ error.
  \emph{Partial:} below 25\%. \emph{Negative:} 40\% or more, meaning crop
  training does not rescue low-drag transfer. A negative result still
  separates ``the training distribution was wrong'' from ``the physics is
  not local enough''. This needs a $1\times$ and a $3\times$ truth run at
  $\alpha{=}0.05$, using the checkpointed Benettin code.
\item[P3 --- how local can the model be?] At fixed crop size, sweep the
  free-interior one-step and $k$-step error against halo width $h$ and
  against propagator window. The width $h^*(\alpha)$ where the error
  saturates is a data-driven interaction range. \emph{Prediction}
  (\S\ref{sec:transfer_diagnostic}): $h^*$ grows as $\alpha$ falls,
  roughly following $\ell_\alpha\propto\alpha^{-3/2}$. If $h^*$ does not
  saturate within the crop, the regime is non-local at that crop size.
  That is itself a measurement.
\item[P4 --- inputs.] Encoder input $\omega$ versus $(\omega,u,v)$, at
  the $\alpha$ where P2 is weakest.
\item[P5 --- cheap validation and the reverse direction.] Deploy on the
  full $16\times2\pi$ domain and compare energy spectra, spatial
  autocorrelation and dissipation--input PDFs
  (\cite{constanteamores2024}'s statistics), with no Lyapunov truth
  needed. Also deploy zero-shot \emph{down} to the $1\times$ periodic box.
  Any mismatch there measures finite-size effects of the small box, which
  the crop-trained model never saw.
\item[P6 --- same-domain patch training (added 2026-10-03).]
  \textbf{COMPLETED: informative failure, narrow-architecture limited}
  (\S\ref{sec:p6_results}). Train
  only on windows of one $3\times2\pi$ periodic domain at $\alpha{=}0.05$,
  deploy on that full domain, and compare with the $3\times$ truth.
  Design, rationale and rules are in \S\ref{sec:p6}.
\item[P7 --- crop training on 2D Kuramoto--Sivashinsky (added 2026-10-03).]
  \textbf{Step 1 COMPLETED} (native error 13.5\%, transfer 24.4\%); step 2
  planned. \S\ref{sec:p7}.
\item[P8 --- drag sweep, $\alpha{=}0.025$ (and the $\alpha{=}0$ pilot).]
  \textbf{QUEUED} after P6. \S\ref{sec:drag_sweep}.
\end{description}

\subsection{P6: same-domain patch training, the practical setting}
\label{sec:p6}

\textbf{Why this experiment.} For CM1, reanalysis or observations, the
domain cannot be re-simulated larger. One large domain exists, and the
goal is a cheap, local model of it. In that setting L-transfer is not
the question. The question is whether a local model trained only on
\emph{windows} of a domain reproduces the dynamics of the \emph{whole}
domain. The model still sees large-scale structure during training, but
only from a local perspective, as local gradients and background.

\textbf{What it would show if it works.}
\begin{itemize}[leftmargin=*]
\item \emph{Cost:} data and training scale with the window, not the
  domain.
\item \emph{Interpretability:} a small receptive field makes the learned
  rule a local object we can distill into stencil or sparse-PDE form, as
  in the earlier local-closure sections.
\item \emph{DA readiness:} the model is local with a measurable
  localization radius.
\end{itemize}

\textbf{Where the scientific weight is.} The result is only non-trivial
if the windows are \emph{smaller than the coherent structures}. If every
crop contains the largest structures whole, success is expected and
uninformative. The informative case is a model with finite reach
reproducing dynamics of structures larger than its window. Patch-trained
local predictors of large chaotic domains are not new
\cite{pathak2018}. What would be new is Lyapunov-level validation of the
assembled model, together with the dependence of its error on the window
size.

\textbf{Relation to P2.} P2 already gives a preview: trained on crops of
the $16\times2\pi$ domain, it gives $D_{\mathrm{KY}}$ 72.0, 69.6 and 69.9
at $3\times$, to be scored when the $3\times$ truth completes. P6
removes P2's domain mismatch (training at $16\times$, deploying on a
$3\times$ box that is periodic at $6\pi$) by training and deploying on
the same domain.

\textbf{Design (fixed before any run).}
\begin{itemize}[leftmargin=*]
\item \emph{Data:} one $3\times2\pi$ periodic domain at $\alpha{=}0.05$
  (384$\times$128 grid, 24 sites wide), 12 training and 2 validation
  trajectories of 40 time units, with seeds distinct from all other
  datasets.
\item \emph{Narrow architecture.} With the standard encoder (radius
  $2\times3$ layers) and window-3 propagator, each crop side loses 12
  sites to masking and halo, which is as wide as half the domain. P6
  therefore uses an encoder of radius $1\times2$ layers (2 sites) and a
  window-1 propagator (halo 2), so each side loses only 4 sites.
\item \emph{Free-interior sweep:} 6, 8 and 12 sites (0.75, 1 and 1.5
  box widths), i.e.\ crops of 14, 16 and 20 sites.
\item \emph{Upper bound} (\texttt{--periodic-full}): the same
  architecture trained on the whole periodic domain (crop = domain), so
  circular padding is exact and no mask or halo is needed.
\item \emph{Evaluation:} $D_{\mathrm{KY}}$ on the full $3\times$ domain,
  3 seeds per arm, with normalized initial conditions.
\end{itemize}

\textbf{Pre-registered rules.}
\begin{itemize}[leftmargin=*]
\item \emph{Locality length:} the smallest free-interior width at which
  the error has stopped improving, i.e.\ is within 3 points of the
  upper-bound arm. This is reported as the measured locality length at
  $\alpha{=}0.05$.
\item \emph{Success:} at that width, $D_{\mathrm{KY}}$ is within 10\% of
  the truth.
\item \emph{Informative failure:} if even the upper-bound arm misses by
  more than 15\%, the limit is the narrow architecture and not the
  windowing. That would be reported as such, not as a failure of patch
  training.
\end{itemize}



\subsection{Implementation and pre-run deviations (2026-10-02)}
\label{sec:crop_impl}

\textbf{Code.}
\begin{itemize}[leftmargin=*]
\item \texttt{scripts/gen\_kolmogorov\_wide\_dataset.py} (P0). Preemptable,
  one file per trajectory.
\item \texttt{ks\_latent/training/crop.py}, plus the driver
  \texttt{scripts/train\_kolmogorov\_crop.py}.
\item \texttt{scripts/exp\_crop\_eval.py}, for zero-shot evaluation on
  periodic boxes of $1/2/3\times2\pi$.
\item \texttt{tests/unit/test\_crop\_training.py}, with 5 tests. The key
  test checks that, with halo $\ge$ the propagator's per-step receptive
  field, the forced-halo free-interior prediction equals, to $10^{-6}$,
  what the same weights produce on a larger periodic domain. A companion
  test checks that a halo smaller than the receptive field breaks this
  equality.
\end{itemize}
The pipeline runs strictly sequentially:
P0($\alpha{=}0.1$) $\to$ P1 $\to$ P0($\alpha{=}0.05$) $\to$ P2 $\to$ the
$\alpha{=}0.05$, $3\times$ truth run $\to$ C7.

\textbf{Configuration, fixed before any run.}
\begin{itemize}[leftmargin=*]
\item \emph{Data:} wide domain $L_x{=}16\times2\pi$ ($N_X{=}2048$, 128
  sites), $L_y{=}2\pi$. There are 4 training trajectories and 1
  validation trajectory, each 40 time units at $dt_{\mathrm{snap}}{=}0.2$,
  with normalization statistics from the training split.
\item \emph{Stage A (autoencoder):} crops of 48 sites ($6\times2\pi$). The
  reconstruction loss is masked to the interior 24 sites (12 per side,
  the encoder plus decoder receptive field). The latent variance and
  log-determinant regularizers ($w_{\mathrm{var}}{=}0.02$,
  $w_{\mathrm{logdet}}{=}0.0035$, as in Sections 231/233) are applied per
  site to the $c{\times}c$ channel statistics of the interior. Training
  is 8000 steps at batch 32, AdamW at $10^{-3}$, with warmup and cosine
  decay.
\item \emph{Stage B (propagator):} a frozen-AE bank of 32 random crops per
  trajectory. The outer 6 encoder sites per side are dropped, leaving a
  36-site strip, of which 6 sites per side are a forced halo and 24 run
  free. Training is 6000 steps at batch 64, AdamW at
  $3\times10^{-4}$, using Stage 2's $k$ curriculum ($2\to8\to12$) and
  input-noise anneal ($0.10\to0.02$). The best checkpoint is chosen by
  validation $k{=}12$ interior MSE.
\item \emph{Evaluation:} zero-shot on periodic $1/2/3\times$ boxes, 3
  initial-condition seeds each, with initial conditions normalized by the
  wide set's statistics. Lyapunov settings match \texttt{exp\_g3}.
\end{itemize}

\textbf{Deviations from the P1/P2 text above, and why.}
\begin{enumerate}[leftmargin=*]
\item \emph{Propagator window 3, not 1.} This keeps the propagator
  identical to Sections 231/233, so P1 and P2 differ from those baselines
  only in the training data and its masking. Window 1 would have added a
  second change (D1 showed window alone changes $2\times$ error).
\item \emph{Crops are $6\times2\pi$, not $4\times2\pi$.} With window 3,
  the total halo is 12 sites per side, and a 4$\times2\pi$ crop would
  leave only 8 sites free.
\item \emph{No auxiliary prediction loss in Stage A.} This is the
  ``two-stage'' variant ($w_{\mathrm{pred}}{=}0$). The propagator is
  trained cold in Stage B, as D2 and D3 were.
\item \emph{Regularizers are per-site, not on the flattened latent.} The
  flattened $d_{\mathrm{lat}}$ of a crop differs from the native size,
  and the per-site $c{\times}c$ form is the size-independent version. The
  index-band spatial-coherence term ($w_{\mathrm{spatial}}$) is dropped.
\end{enumerate}
Items 3 and 4 mean P1 is not a pure single-variable change from Section
233. If P1 fails, the next control is a periodic-box model trained with
this same trainer (crop = full box, no masking). That would separate
``crop data'' from ``trainer differences.''

\subsection{P1 results (completed 2026-10-03)}
\label{sec:p1_results}

\textbf{Status: COMPLETED} (all 9 evaluations).

\textbf{Training.} Stage A reached an interior validation reconstruction
MSE of 0.0082. Section 233's periodic-box autoencoder reached 0.0045 and
Section 231's 0.0083; the metrics differ (interior only, wide-domain data,
shorter schedule), so the comparison is indicative only. Stage B reached
a best validation $k{=}12$ interior latent MSE of 0.021. Initial
conditions are normalized with the wide set's statistics
($\sigma{=}1.531$).

\begin{table}[h]
\centering
\small
\begin{tabular}{lcccc}
\toprule
size & crop-trained $D_{\mathrm{KY}}$ (seeds 0/1/2) & truth & mean error & Sec.\ 233 error (normalized IC) \\
\midrule
$1\times$ & 15.2 / 17.0 / 18.2 & 15.88 & 8.6\% & 23.4\% \\
$2\times$ & 36.5 / 37.9 / 35.6 & 30.34 & \textbf{20.8\%} & 3.7\% \\
$3\times$ & 54.6 / 55.9 / 56.4 & 56.35 & \textbf{1.3\%} & 19.4\% (raw IC) \\
\bottomrule
\end{tabular}
\caption{P1, crop-trained model at $\alpha{=}0.1$, zero-shot on periodic
boxes.}
\label{tab:p1}
\end{table}

\textbf{Pre-registered rule: P1 fails on the $2\times$ criterion.} The
rule required $2\times$ error $\le10\%$; the result is 20.8\%. The $3\times$
criterion (below 19.4\%) is met decisively: 1.3\% mean error, with seeds at 3.1\%, 0.7\% and 0.0\%. This is the most accurate 2D transfer in this document.

\textbf{What the pattern shows.} $D_{\mathrm{KY}}$ per box width
($2\pi\times2\pi$ area) is:
\begin{itemize}[leftmargin=*]
\item crop-trained model: 16.8, 18.3, 18.5 at $1\times$, $2\times$,
  $3\times$, nearly extensive;
\item truth: 15.9, 15.2, 18.8, with a dip at $2\times$.
\end{itemize}
The crop-trained model has learned roughly the large-domain density: its
$\approx18.5$ per box is close to the $3\times$ truth (18.8). It then
overshoots the small boxes, and most of all the $2\times$ box, where the
truth dips. This is the ``reverse direction'' effect anticipated in P5.
A model trained on large-domain statistics reproduces finite-size
deviations of small periodic boxes only to the extent they are local. The
box-trained Section 233 model matched $2\times$ closely but failed at
$3\times$, which is the opposite trade.

\textbf{Caveat that decides the reading.} The $2\times$ truth is a single
seed. T1 was designed to test whether the $2\times$ dip in
$D_{\mathrm{KY}}/\mathrm{Area}$ is real, and its $2\times$ truth seeds were
never completed. If the dip is within seed-to-seed spread, much of the
20.8\% is truth uncertainty. If it is real, the crop-trained model
misses a genuine finite-size effect. \textbf{Follow-up:} finish T1's
$2\times$ truth seeds (\texttt{scripts/exp\_t1\_error\_bars.py}). \textbf{TO DO} (deprioritized 2026-10-03 in favour of C8--C12; \texttt{scripts/run\_t1\_x2\_truth\_after\_c7.sh} is ready to run).

\subsection{P2 results, interim (2026-10-03)}
\label{sec:p2_results}

\textbf{Status:} all 9 evaluations done. The $\alpha{=}0.05$ $3\times$ truth run, needed to score $3\times$, was moved to after C7 (2026-10-03). The
wide dataset has field std 1.385 (vs.\ 1.394 in the $2\pi$ box). Stage A
reached an interior validation reconstruction MSE of 0.0082; Stage B
reached a best validation $k{=}12$ interior latent MSE of 0.025. The
recipe is otherwise identical to P1.

\begin{table}[h]
\centering
\small
\begin{tabular}{lcccc}
\toprule
size & crop-trained $D_{\mathrm{KY}}$ (seeds 0/1/2) & truth & mean error & prior best at $\alpha{=}0.05$ \\
\midrule
$1\times$ & 24.0 / 20.5 / 24.9 & 23.26 & 7.3\% & --- \\
$2\times$ & 45.1 / 44.8 / 44.9 & 56.84 & \textbf{20.9\%} & 46.8\% (Sec.\ 231, normalized IC) \\
$3\times$ & 72.0 / 69.6 / 69.9 & 75.22 & \textbf{6.3\%} & --- \\
\bottomrule
\end{tabular}
\caption{P2, crop-trained model at $\alpha{=}0.05$, zero-shot on periodic
boxes. For comparison at $2\times$: D2 95\%, D3 98--99\%.}
\label{tab:p2}
\end{table}

\textbf{Pre-registered rule: PARTIAL.} The $2\times$ error is below 25\%
but not single digits. With \emph{no architecture change}, training on
crops of a large domain more than halves the low-drag transfer error of
periodic-box training (46.8\% to 20.9\%). Every architectural attempt
(D2, D3) made it worse.

\textbf{The residual is systematic, and its sign points to physics.} The
three $2\times$ seeds agree to within half a point (20.6--21.1\%), all
\emph{below} the truth. $D_{\mathrm{KY}}$ per box width is:
\begin{itemize}[leftmargin=*]
\item crop-trained model: 23.1, 22.5, 23.5 at $1\times$, $2\times$,
  $3\times$, so it is extensive;
\item truth: 23.3 at $1\times$, then 28.4 at $2\times$, growing faster
  than the area.
\end{itemize}
Compare P1 at $\alpha{=}0.1$, where the residual had the opposite sign:
there the truth dipped at $2\times$ and the model overshot.

The local model therefore captures the extensive, local part of the chaos
and misses an excess that grows with domain size. That is precisely the
collective, large-scale contribution that
\S\ref{sec:transfer_diagnostic}'s TD hypotheses attribute to weak drag,
and that C7's learned consensus is designed to supply. The $3\times$
truth decides whether the gap keeps widening: TD2 predicts it does at
$\alpha{=}0.05$. Note also that the crop-trained model's density (about
23 per box) matches the $1\times$ truth rather than any larger-domain
value. That is consistent with the wide-domain \emph{local} statistics
being close to those of the $2\pi$ box, with the excess residing in
correlations longer than the model's reach.

\subsection{P1--P2 conclusions (2026-10-03)}
\label{sec:p1p2_conclusions}

\begin{enumerate}[leftmargin=*]
\item \textbf{The training distribution was the main problem, not the
  architecture.} With no architectural change, training on crops of a
  large domain cut the low-drag $2\times$ error from 46.8\% to 20.9\%.
  Every attempt to add cross-site coupling while training on the small
  periodic box made transfer worse: D2 (95\%), D3 (98--99\%), and C4
  re-projection (\S\ref{sec:c4_partial}). The D1--D3/C4 failures were
  therefore mainly failures of training only on the 8-site periodic
  ring, not evidence that communication is impossible.
\item \textbf{The remaining low-drag error is systematic and has a
  physical signature.} At $\alpha{=}0.05$ the three $2\times$ seeds miss
  by 20.6--21.1\%, all under-predicting, so the miss is a property of the
  model, not initial-condition luck. The crop-trained model is extensive
  (about 23 per box at every size), while the truth grows faster than the
  area (23.3, then 28.4 per box). A model of finite reach cannot produce
  super-extensive growth. The residual, about 12 of roughly 57
  dimensions at $2\times$, is therefore a direct estimate of the
  \emph{non-local, collective} share of the chaos.
\item \textbf{The sign of the residual tracks the drag regime.} At
  $\alpha{=}0.1$ the truth dips at $2\times$, the model overshoots, and
  $3\times$ is nearly exact (1.3\%). At $\alpha{=}0.05$ the truth grows
  and the model undershoots. In both regimes the crop-trained model
  learns about the large-domain \emph{local} density. Its deviation from
  the truth measures how non-extensive the true dynamics are, which is
  \S\ref{sec:transfer_diagnostic}'s diagnostic working on a learned model
  rather than in principle.
\item \textbf{Reverse transfer works.} Never trained on a periodic
  $2\pi$ box, the crop-trained models reproduce native-size
  $D_{\mathrm{KY}}$ within about 7--9\% at both drags (better than the
  box-trained Section 233 model's own 23\% at $\alpha{=}0.1$).
\item \textbf{Practical takeaway.} Training a local model on windows of
  a large domain and deploying at any size is realistic for observational
  data and cheaper than box training. Here it beat box training exactly
  where box training failed: $3\times$ at $\alpha{=}0.1$ and $2\times$ at
  $\alpha{=}0.05$.
\end{enumerate}

\textbf{Not yet established.}
\begin{itemize}[leftmargin=*]
\item \emph{Truth error bars.} The $2\times$ truths are single seeds (the
  T1 seeds are queued). The $\alpha{=}0.05$ $3\times$ truth is pending:
  a gap larger than 21\% at $3\times$ would strengthen reading 2, and a
  smaller one would weaken it.
\item \emph{Training seeds.} There is one training seed per crop model.
\item \emph{Learnability.} Whether the excess is learnable at all is C7's
  question. If learned consensus closes the gap, the missing ingredient
  was communication. If it does not, the gap stands as a measurement of
  the non-local fraction as a function of $\alpha$ and $L$. Either
  outcome is publishable.
\end{itemize}

\subsection{The $\alpha{=}0.05$ $3\times$ truth, and a revision of the ``super-extensive'' reading (2026-10-03)}
\label{sec:truth3_revision}

\textbf{Truth:} $D_{\mathrm{KY}}=75.22$, $n_{\mathrm{pos}}=36$,
$\lambda_1=0.131$, from 128 directions over 40 time units (seed 0).

\begin{table}[h]
\centering
\small
\begin{tabular}{lcc}
\toprule
model ($\alpha{=}0.05$, $3\times$) & $D_{\mathrm{KY}}$ (seeds) & mean error \\
\midrule
P2 (crop-trained) & 72.0 / 69.5 / 69.9 & \textbf{6.3\%} \\
C7(b) (learned consensus) & 71.4 / 65.2 & 9.2\% \\
C9(i) (exact velocity input) & 64.4 / 68.2 / 67.4 & 11.4\% \\
C9(ii) (box-trained) & 4.9 / 1.6 / 4.0 & collapse \\
\bottomrule
\end{tabular}
\caption{Scores at $3\times$ against the new truth.}
\label{tab:truth3}
\end{table}

\textbf{What changes.}
\begin{enumerate}[leftmargin=*]
\item \emph{Crop training succeeds at $3\times$ at weak drag.} P2 reaches
  6.3\%, in the single-digit range. The added mechanisms are slightly
  worse: consensus 9.2\%, exact velocity 11.4\%.
\item \emph{The truth is not monotonically super-extensive.}
  $D_{\mathrm{KY}}$ per box width is 23.3 at $1\times$, 28.4 at $2\times$
  and 25.1 at $3\times$. The $2\times$ value is the outlier, not part of a
  trend. At $\alpha{=}0.1$ the $2\times$ value was also the outlier, in
  the opposite direction (15.9, 15.2, 18.8). Both $2\times$ truths are
  single seeds with 40-time-unit averaging, and T1 measured about 8\%
  seed-to-seed scatter in truth $D_{\mathrm{KY}}$ at native size.
\item \emph{Consequently the ``systematic 21\% deficit at $2\times$''
  (\S\ref{sec:p2_results}, \S\ref{sec:p1p2_conclusions}) may be largely a
  truth artifact.} The model is extensive at about 23 per box. The truth
  at $1\times$ and $3\times$ (23.3 and 25.1) is close to extensive, with a
  mild growth of about 8\%. The conclusions in
  \S\ref{sec:p1p2_conclusions} that rest on a strongly super-extensive
  truth, items 2--3 and the ``12 of 57 dimensions'' estimate, are
  \emph{suspended} until multi-seed $2\times$ truths exist.
\end{enumerate}

\textbf{Resolution (multi-seed $2\times$ truth, COMPLETED 2026-10-03).}
Seeds 1 and 2 give $D_{\mathrm{KY}}=38.68$ ($n_{\mathrm{pos}}{=}17$) and
53.89, against seed 0's 56.84 ($n_{\mathrm{pos}}{=}28$).

\begin{table}[h]
\centering
\small
\begin{tabular}{lccc}
\toprule
 & $1\times$ & $2\times$ (3 seeds) & $3\times$ \\
\midrule
truth $D_{\mathrm{KY}}$ & 23.26 (1 seed) & \textbf{49.8 $\pm$ 9.8} (56.8 / 38.7 / 53.9) & 75.22 (1 seed) \\
truth per box & 23.3 & 24.9 $\pm$ 4.9 & 25.1 \\
P2 model & 23.1 (7.3\%) & 44.6 (10.4\% vs.\ mean) & 70.5 (6.3\%) \\
C7(b) model & 23.0 & 45.3 (9.0\%) & 68.3 \\
C9(i) model & 22.8 & 44.6 (10.4\%) & 66.7 (11.4\%) \\
\bottomrule
\end{tabular}
\caption{$\alpha{=}0.05$ with the multi-seed $2\times$ truth.}
\label{tab:truth2_multiseed}
\end{table}

\textbf{Consequences.}
\begin{enumerate}[leftmargin=*]
\item \emph{The true system's $D_{\mathrm{KY}}$ depends strongly on the
  initial condition} at $2\times$ with 40-time-unit averaging: a
  $\pm20\%$ spread, with $n_{\mathrm{pos}}$ ranging from 17 to 28. This
  points to long-lived regimes or multistability, consistent with the
  much longer averaging windows of \cite{kolmogorovregimes2026}. A single
  40-unit truth is \emph{not} a reliable reference at this drag.
\item \emph{The ``systematic 21\% deficit at $2\times$'' is not
  statistically significant.} Every window-trained model lies within one
  truth standard deviation (9--10\% from the mean). With the truth
  approximately extensive at $\alpha{=}0.05$ (23--25 per box at all three
  sizes), the window-trained models transfer at $\alpha{=}0.05$ about as
  well as at $\alpha{=}0.1$.
\item \emph{Conclusions resting on single truths are downgraded.} The
  super-extensive reading
  (\S\ref{sec:p2_results}--\S\ref{sec:p1p2_conclusions}) and the
  non-local-excess motivation for C9--C18 no longer stand. The
  $\alpha{=}0.1$ $2\times$ ``dip'' behind P1's $2\times$ miss is likely
  the same effect, and T1's $\alpha{=}0.1$ $2\times$ seeds would confirm
  it.
\item \emph{New standing requirement for the paper:} truth values
  quoted against models must be multi-seed means. Alternatively, each
  truth must be averaged over a window long enough that its running
  estimate has converged (record running estimates as T1 does).
  Model-side errors should be reported against the truth spread, not a
  point value.
\end{enumerate}

\textbf{Action taken.} $\alpha{=}0.05$ $2\times$ truth seeds 1 and 2 (96
directions, 40 time units, checkpointed;
\texttt{scripts/run\_truth\_a005\_x2\_seed.py}) were launched in
parallel, at highest priority. They decide whether the $2\times$ deficit,
and hence the motivation for C9--C18, is real. T1's $\alpha{=}0.1$ $2\times$
seeds address the mirror question at strong drag and are re-promoted
from the deferred list.

\subsection{C8, C16 and C13 results: the chaos lives at the forcing scale (2026-10-03)}
\label{sec:c8_c16_c13_results}

\textbf{C8: COMPLETED. No large-scale deficit at $3\times$.} The unstable
subspaces of the truth ($n_{\mathrm{pos}}{=}36$) and of P2's model (34)
have nearly identical scale content:
\begin{itemize}[leftmargin=*]
\item variance beyond the box ($\lambda_x>2\pi$): 9.1\% vs.\ 10.0\%;
\item variance in non-box modes: 66.8\% vs.\ 66.3\%;
\item variance at $m{=}0$: 2.3\% vs.\ 2.6\%.
\end{itemize}
Only 2 dimensions are missing, and they are not concentrated at large
scales: the beyond-box excess is $-0.15$ dimensions. The coded verdict is
``local.'' The substantive reading is stronger: the model's chaos lives
at the same scales as the truth's.

\textbf{C16: COMPLETED. The chaos lives near the forcing scale, and there
is no cascade-type self-similarity.}
\begin{itemize}[leftmargin=*]
\item \emph{(iii) Unstable dimensions per wavenumber octave} at $3\times$,
  $\alpha{=}0.05$: 0.6 at $|k|\in[1,2)$, 11.4 at $[2,4)$, 20.1 at
  $[4,8)$, 3.7 at $[8,16)$, and essentially zero below $|k|{=}1$. About
  31.5 of 36 unstable dimensions sit around the forcing wavenumber
  ($k{=}4$).
\item \emph{(i) Spectra and structure functions.} At $\alpha{=}0.05$ the
  streamwise energy-spectrum slopes are $+0.83$ for $k\in[0.25,1]$ and
  $-1.07$ for $k\in[1,3]$, with a steep dissipation range ($-8.4$)
  beyond forcing. $S_2$ has slope 0.43 and $S_3$ has mixed sign. These
  are not the Kraichnan inverse-cascade predictions ($-5/3$, $2/3$,
  $S_3>0$). $\alpha{=}0.1$ is similar.
\item \emph{(ii) Band statistics.} Band amplitude falls with scale with
  exponent $-0.55$ ($\alpha{=}0.05$) and $-0.61$ ($\alpha{=}0.1$), near
  the predicted $-2/3$. But the decorrelation time is essentially
  \emph{independent} of scale, at 4--5 time units: exponent $-0.07$ and
  $-0.10$ against the predicted $+2/3$. PDF and autocorrelation
  \emph{shapes} collapse within 15\% over a factor of 16 in scale.
\item \emph{Pre-registered verdict.} The coded rule (shape collapse over
  at least $4\times$) says PROCEED. The physical premise, however, fails:
  at Re $\approx100$ there is no inertial range with
  Kolmogorov--Kraichnan time scaling, and the timescale is set by the
  forcing-scale dynamics at every scale. C17 remains runnable with the
  fitted exponents (no time rescaling between levels), but C18's
  motivation is weak in this regime. We record both rather than
  overriding the rule.
\end{itemize}

\textbf{C13: COMPLETED (linear closure only). The large scales are
eddy-driven, and a local linear closure is not viable.} At every cutoff
and both drags, the exact bookkeeping holds ($T{=}R{+}E$ to better than
$10^{-3}$). At $\alpha{=}0.05$:
\begin{itemize}[leftmargin=*]
\item the eddy forcing's variance is 1.15--1.31$\times$ the large-scale
  tendency's;
\item the exact coarse terms alone explain less than nothing (explained
  fraction $-0.15$ to $-0.31$);
\item a local linear closure in coarse and small-scale-enstrophy features
  fits only 14\% of the eddy forcing (38\% at $\alpha{=}0.1$), leaving
  the total explained fraction near 0 at the $2\pi$ cutoff.
\end{itemize}
\emph{Pre-registered verdict:} not viable with a linear closure, and the
closure would have to be stochastic (remainder above 0.3). The small-CNN
closure named in the pre-registration has not been run, so this verdict
covers linear closures only.

\textbf{What the three results mean together.}
\begin{enumerate}[leftmargin=*]
\item The large-scale, non-local deficit that motivated C9--C18 is
  \emph{not present at $3\times$}. Crop training reproduces the truth's
  dimension (6.3\% error) and its scale distribution.
\item The chaos is concentrated at the forcing scale. That is exactly
  the scale a local model with a few-site reach resolves best, and it
  explains why crop-trained local models transfer once they see
  large-domain training data.
\item The open question narrows to whether the $2\times$ anomaly is real
  (the multi-seed $2\times$ truths are running).
\item The two-scale route (C14--C15) would need a stochastic or
  nonlinear closure for large scales that carry little of the chaos.
  The self-similar route (C17--C18) lacks its physical premise at this
  Re. Both are deprioritized. P6 and the drag sweep become the most
  informative next steps.
\end{enumerate}

\subsection{Drag sweep: $\alpha{=}0.025$ and the $\alpha{=}0$ limit (QUEUED 2026-10-03)}
\label{sec:drag_sweep}

\textbf{Rationale.} The transfer-error-as-diagnostic claim
(\S\ref{sec:transfer_diagnostic}, T2) needs more than two drag values.
\begin{itemize}[leftmargin=*]
\item \emph{$\alpha{=}0.025$} is the next point. It remains stationary
  under the standard 100-unit spinup and runs on the existing pipeline,
  giving a three-point curve of transfer error against $\alpha$.
\item \emph{$\alpha{=}0$} is the end point: the condensate limit, where
  large-scale energy is removed only by viscosity at rate
  $\nu k_{\min}^2$. That is a timescale of roughly 900 time units at
  $3\times$, so the standard spinup would sample a transient. It is
  therefore gated by a pilot.
\item \emph{Forcing note.} In our solver
  $f(y)=-(\nu n^2+\alpha)\,n\cos ny$, so the laminar amplitude, and the
  nominal Re of 100, are the same at every $\alpha$. The solver's own
  notes record chaos in the square box at Re 70--100 without drag, so
  $\alpha{=}0$ is chaotic here.
\end{itemize}

\textbf{Queue} (\texttt{scripts/run\_alpha\_sweep\_queue.sh}, after P6):
\begin{enumerate}[leftmargin=*]
\item $\alpha{=}0.025$ truth runs at $1/2/3\times$ (64/128/176
  directions; \texttt{scripts/run\_truth\_kolmogorov.py}), started in the
  background in parallel. Each warns if $n_{\mathrm{pos}}$ exceeds 80\%
  of the directions, which would mean the estimate is truncated.
\item The $\alpha{=}0$ saturation pilot (\texttt{scripts/exp\_alpha0\_pilot.py}):
  1500 time units at $1\times$ and $2\times$, recording total and
  domain-mode energy.
\item The $\alpha{=}0.025$ wide dataset, P8 training with P2's exact
  recipe, and evaluation at $1/2/3\times$ with 3 seeds each.
\end{enumerate}

\textbf{Pre-registered rules.}
\begin{itemize}[leftmargin=*]
\item \emph{P8 (diagnostic curve).} The TD hypothesis predicts that
  transfer error at fixed size grows as $\alpha$ falls. If P8's $3\times$
  error stays single-digit, locality holds down to $\alpha{=}0.025$, and
  the onset of non-locality lies below it.
\item \emph{$\alpha{=}0$ pilot (gate).} The saturation time is the
  first time at which the energy's 200-unit running mean changes by less
  than 5\% over the previous 200 units. $\alpha{=}0$ datasets and truths
  will use a spinup at least that long. If no saturation is reached
  within 1500 units at $2\times$, $\alpha{=}0$ is restricted to the
  $1\times$ box and framed qualitatively.
\end{itemize}

\subsection{Options after P7: patch training, revisiting P6, and $y$-transfer (2026-10-04)}
\label{sec:post_p7_options}

\textbf{Context from discussion.} Running P2 on a larger domain works
because the deployed domain is periodic, so the model never meets an
edge. Training on crops needs exact targets, which costs a 12-site edge
zone on each side (6 for the encoder, 6 for the decoder; the propagator
halo is 6). P2's crops were 48 sites wide, cut from a 128-site domain;
the $3\times$ domain has only 24 sites, so P2-sized crops do not fit.
Kolmogorov crops have so far been strips spanning the full height, for
two reasons: the domain is only 8 sites tall, and the forcing $\cos4y$
breaks translation invariance in $y$. In $y$, P2 is therefore
effectively box-trained, since it has never seen $y$-structure longer
than $2\pi$, and its behaviour under $y$-transfer is untested. P6's
failure was caused by the narrow architecture, not by windowing, so
windowed training on a single domain is \emph{not} ruled out; the edge
handling needs redesigning.

\textbf{Decision tree (no further launches until P7 reports).}
\begin{description}[leftmargin=1em]
\item[P7 transfers well] (2D-KS $66^2$ error well below step 1's 24\%,
  and native error not much worse).
  \begin{itemize}
  \item \emph{P9: 2D patch training for Kolmogorov.} Use a tall domain
    (e.g.\ $12\times12\cdot2\pi$, a $1536^2$ grid) and $40\times40$-site
    patches with $16\times16$ interiors. Patch offsets in $y$ are
    multiples of 2 sites, one forcing period, which is an exact symmetry,
    so the model needs no $y$-position input. Evaluate at $1\times1$,
    $2\times1$, $1\times2$ and $2\times2$ against multi-seed truths.
    Normalize everything with the training scalar mean and standard
    deviation, and check that the field's standard deviation does not
    drift with domain height. This is the general paradigm: no direction
    is privileged, and no patch contains a full periodic direction.
  \item \emph{P6 revisited (edge handling).} C19 (\S\ref{sec:c19}) tests
    which of these is viable. Candidate fixes:
    (a) window only the propagator, reusing P2's autoencoder on the
    periodic $3\times$ ring, which shrinks the edge loss to 6 sites
    (interiors of 2--8 sites);
    (b) give the encoder true physical context beyond the crop as
    unscored input, so the scored window, not the model's receptive
    field, is what shrinks;
    (c) soft, distance-weighted edge loss instead of hard masking,
    accepting slightly inexact targets near edges;
    (d) keep P2's full-reach encoder and narrow only the propagator,
    since P6 suggests the encoder's reach is the bottleneck.
  \end{itemize}
\item[P7 partially works] (transfer error improves but stays
  double-digit). Diagnose before scaling to Kolmogorov: run P7's native
  versus transfer error split and a C8-style unstable-subspace
  comparison, and test whether larger interiors or more crops help.
  Run P6 option (a) as the cheap Kolmogorov test, since it needs no new
  data.
\item[P7 fails] (no better than step 1). 2D patch training in its
  current form does not generalize, so do not generate tall Kolmogorov
  data. Fall back to cheap tests on existing data: P6(a), and P2's
  $y$-transfer (\texttt{--scale-y}) to establish whether box training
  in $y$ fails the way box training in $x$ did.
\end{description}

\textbf{Running regardless of P7:} the remaining P8 runs
($\alpha{=}0.025$ truths at $2\times$/$3\times$, wide data, training
and evaluation) and the $\alpha{=}0$ chaos check
(\S\ref{sec:alpha0_results}) before any $\alpha{=}0$ dataset.

\subsection{C19: encoder/decoder reach ablation (REDUCED to one arm; results in \S\ref{sec:c19_results})}
\label{sec:c19}

\textbf{Why.} P6 was uninformative because it shrank the encoder
\emph{and} the propagator together (\S\ref{sec:p6_results}). The
discussion that followed separated the three contributions to a crop's
edge loss. The masked margin is the model's receptive field, whatever
the architecture: a ViT with windowed attention of radius $r$ over $n$
layers has reach $nr$ and needs the same margin. In P2 that margin
comes from three parts:
\begin{itemize}[leftmargin=*]
\item \emph{encoder reach} of 6 sites (radius $2\times3$ layers);
\item \emph{decoder reach} of 6 sites, the same convolutions mirrored;
\item \emph{propagator halo} of 6 sites, its window 3 over 2 layers.
\end{itemize}
Only the propagator halo is dynamics. The autoencoder's two reaches are
design choices. If they can shrink without hurting the dynamics, the
autoencoder's edge mask falls from 12 sites to $e{+}d$, and
small-domain crop experiments become feasible: P6 redone properly, and
2D patches (P9). The extreme case $e{=}d{=}0$ is a strictly per-patch
autoencoder. Its latent is then a pure local compression and all
spatial interaction lives in the propagator, so crop training loses only
the 6-site physical halo.

\textbf{Design (fixed before running).} The P2 recipe throughout:
$\alpha{=}0.05$, the $16\times$-wide dataset, 48-site $x$-crops, a
window-3 propagator with halo 6, and the same steps and seed. Only the
autoencoder's site-mixing reach changes. Decoder reach is a new config
option (\texttt{dec\_site\_mix\_radius}, \texttt{dec\_n\_site\_mix\_layers};
it defaults to the encoder's, so earlier checkpoints are unchanged), and
the autoencoder mask is $e{+}d$. A unit test checks that a decoder of
reach $d$ only changes outputs within $d$ sites.
\begin{center}
\small
\begin{tabular}{lccccc}
\toprule
arm & enc.\ reach $e$ & dec.\ reach $d$ & AE mask/side & $c$ & question \\
\midrule
P2 (reference) & 6 & 6 & 12 & 3 & -- \\
\texttt{c19\_e6d0} & 6 & 0 & 6 & 3 & is decoder context needed? \\
\texttt{c19\_e2d0} & 2 & 0 & 2 & 3 & P6's encoder, P2's propagator \\
\texttt{c19\_e0d0} & 0 & 0 & 0 & 3 & per-patch autoencoder \\
\texttt{c19\_e0d0\_c6} & 0 & 0 & 0 & 6 & per-patch, more channels \\
\bottomrule
\end{tabular}
\end{center}
Evaluation is $D_{\mathrm{KY}}$ at $1/2/3\times$, seed 0 only (one seed,
to limit machine load), against the P2 seed-0 reference: 23.95 (2.96\%),
45.14 and 71.96 at $1/2/3\times$. Truths are 23.26 at $1\times$, the
multi-seed $49.8\pm9.8$ at $2\times$ (\S\ref{sec:truth3_revision}; the
built-in 56.84 is a single-seed outlier) and 75.22 at $3\times$.
Autoencoder interior validation error is recorded for every arm.

\textbf{Pre-registered rules.}
\begin{itemize}[leftmargin=*]
\item An arm is \emph{viable} if its $3\times$ error is at most 12\%
  (twice P2's 6.3\% three-seed mean) and its native error at most 8\%.
  These are single-seed thresholds and are deliberately loose.
\item If \texttt{c19\_e0d0} or \texttt{c19\_e0d0\_c6} is viable, the
  autoencoder needs no spatial context. Crop training then costs only
  the propagator halo. Next would be P6 redone on the $3\times$ domain
  with this autoencoder (interiors of 2--12 sites, which now fit) and P9
  with 2D patches.
\item If only \texttt{c19\_e6d0} is viable, the encoder needs context
  but the decoder does not. The edge mask halves, from 12 to 6. P6 would
  still not fit at $3\times$ (6 + 6 propagator halo), but P9 patches can
  shrink.
\item If no reduced arm is viable, P6's failure is attributed to the
  encoder reach. Supporting that, \texttt{c19\_e2d0} isolates P6's
  encoder with P2's propagator. The spatial context the autoencoder sees
  is then part of what makes the dynamics learnable, which fits C8's
  finding that the unstable directions live at forcing scales spanning
  several sites. Small-domain window training then needs a bigger
  domain, P6b at $6\times$.
\item The two per-patch arms separate the explanations: if
  \texttt{c19\_e0d0} fails and \texttt{c19\_e0d0\_c6} works, the issue
  was capacity per site rather than context.
\end{itemize}

\textbf{Queue.} \texttt{scripts/run\_c19\_then\_p8.sh}. It starts
after P7's seed-0 evaluations and runs before P8 training. Every job runs
at background priority (\texttt{taskpolicy -b}, efficiency cores);
training uses MPS.

\subsection{C19 results: decoder reach 0 (COMPLETED 2026-10-04)}
\label{sec:c19_results}

\begin{table}[h]
\centering
\small
\begin{tabular}{lcccccc}
\toprule
 & \multicolumn{2}{c}{$1\times$ (truth 23.26)} & \multicolumn{2}{c}{$2\times$ (truth $49.8\pm9.8$)} & \multicolumn{2}{c}{$3\times$ (truth 75.22)} \\
arm (seed 0) & $D_{\mathrm{KY}}$ ($n_{\mathrm{pos}}$) & err. & $D_{\mathrm{KY}}$ ($n_{\mathrm{pos}}$) & err. & $D_{\mathrm{KY}}$ ($n_{\mathrm{pos}}$) & err. \\
\midrule
P2 ($e{=}6$, $d{=}6$) & 23.95 (11) & 3.0\% & 45.14 (21) & 9.4\% & 71.96 (34) & 4.3\% \\
\texttt{c19\_e6d0} ($e{=}6$, $d{=}0$) & 21.33 (11) & 8.3\% & 44.09 (21) & 11.5\% & 65.20 (31) & 13.3\% \\
\midrule
ratio C19 / P2 & 0.89 & & 0.98 & & 0.91 & \\
\bottomrule
\end{tabular}
\caption{C19, first arm. Autoencoder interior validation error: 0.0705
(C19) against 0.0082 (P2). The $2\times$ errors use the multi-seed truth
mean; the evaluation script's built-in single-seed value, 56.84, gives
22\% for C19 and 21\% for P2. $\lambda_1$ for C19 is 0.110, 0.123 and
0.125.}
\label{tab:c19}
\end{table}

\textbf{Pre-registered verdict: NOT VIABLE, narrowly.} The native
error is 8.3\% (threshold 8\%) and the $3\times$ error 13.3\% (threshold
12\%). Both are single-seed values near their thresholds.

\textbf{The informative part is \emph{how} it fails.}
\begin{itemize}[leftmargin=*]
\item \emph{Transfer is unaffected.} The model's dimension per box is
  21.3, 22.0 and 21.7 at $1/2/3\times$: extensive within 3\%, just like
  P2's 24.0, 22.6 and 24.0. Removing the decoder's context does
  \emph{not} make transfer degrade with size. The ratio to P2 shows no
  trend with size (0.89, 0.98, 0.91).
\item \emph{What drops is the dimension per site,} by about 9\% at every
  size, together with a 9$\times$ loss in reconstruction. A decoder with
  no reach must reconstruct each 16$\times$16-grid patch from its own
  $c{=}3$ channels. The latent that satisfies this keeps less of the
  small-scale chaotic content. That is a capacity or smoothness effect
  on the \emph{local} representation, not a locality effect on the
  dynamics.
\end{itemize}

\textbf{Connections.}
\begin{itemize}[leftmargin=*]
\item This separates two things this document has conflated: the
  \emph{dimension per site} the model achieves, which depends on the
  representation, and the \emph{scaling of dimension with size}, which
  depends on locality and the training distribution and is what
  transfer is about. C19 shows the decoder's reach acts only on the
  first. Box training and coupling (D2, D3, C9) failed on the second.
\item Practically, decoder context is \emph{not} needed for transfer
  but \emph{is} needed for native accuracy at $c{=}3$. The obvious
  repair, a per-patch decoder with more channels (e.g.\ $c{=}6$, as the
  dropped arm \texttt{c19\_e0d0\_c6} intended), is \textbf{QUEUED}
  after P9 and before P8 (\texttt{scripts/run\_c19a\_p9\_p8\_v2.sh}).
  It uses encoder and decoder reach 0, $c{=}6$, and otherwise the P2
  recipe, with seed-0 evaluations at $1/2/3\times$ judged by the same
  pre-registered thresholds. If it
  restores the per-site dimension, the autoencoder margin halves at no
  cost.
\end{itemize}

\textbf{Where this belongs.} It does not support paper 1's headline,
which uses full-reach models on large domains. It is
(i) a \emph{supplementary design ablation} for paper 1's methods: the
crop margin is the receptive field, and here is what each part buys;
and (ii) a \emph{prerequisite for paper 2's practical claim}, training
on the one domain you have (CM1, observations;
\S\ref{sec:corollary_cm1}), which needs small margins. For that claim
the next step is the $c{=}6$ per-patch arm. With it, a P6 redo becomes
feasible on a modestly larger single domain.

\subsection{P9: 2D-patch training for Kolmogorov flow (QUEUED 2026-10-04)}
\label{sec:p9}

\textbf{Why now.} P7 (\S\ref{sec:p7_results}) showed that 2D crops,
cut in both directions from a domain much larger than the crop, give
single-digit native and transfer errors for 2D KS with an unchanged
architecture. Every Kolmogorov model so far used $x$-only strips that
keep a full periodic $y$-direction of 8 sites. In $y$ those models are
therefore box-trained, and box training is exactly what failed under
$x$-transfer. P9 applies P7's recipe to Kolmogorov flow, which also
tests whether the $y$-dependent forcing undermines patch training.

\textbf{Design (fixed before running).}
\begin{itemize}[leftmargin=*]
\item \emph{Data:} $\alpha{=}0.05$ on a $6\times6$-box periodic domain
  ($L_x{=}L_y{=}12\pi$, a $768^2$ grid, $48\times48$ sites, 24 forcing
  periods in $y$). 4 training and 1 validation trajectory of 40 time
  units, seeds 900--904; normalization from the training split.
  (\texttt{gen\_kolmogorov\_wide\_dataset.py} gains \texttt{--y-scale}.)
\item \emph{Model:} the P2 architecture unchanged (encoder radius
  $2\times3$, window-3 propagator, $c{=}3$, full 12-site autoencoder
  margin), trained on $40\times40$-site crops with $16\times16$
  interiors and a 6-site forced halo in both directions.
\item \emph{Forcing:} crop $y$-offsets are multiples of 32 grid points,
  one forcing period ($\pi/2$, 2 sites). Every crop therefore sees
  $\cos4y$ at the same phase relative to its sites, which is an exact
  symmetry of the system, so no forcing input is needed. The deployed
  boxes start at $y{=}0$ and share this alignment. (The new
  \texttt{y\_align} option is unit-tested; default 1 reproduces P7.) An
  arm with arbitrary offsets, testing whether the model infers the
  forcing phase from the flow, is deferred.
\item \emph{Evaluation:} seed 0 at $1\times1$, $2\times1$ and
  $3\times1$ boxes, against truths of 23.26, $49.8\pm9.8$ (multi-seed;
  the eval script's built-in 56.84 is a single-seed outlier) and 75.22.
  This compares directly with P2's seed-0 values: 23.95 (3.0\%), 45.14
  and 71.96 (4.3\%). A $2\times2$ box needs a new truth
  ($D_{\mathrm{KY}}\approx95$), which is deferred.
\end{itemize}

\textbf{Pre-registered rules.}
\begin{itemize}[leftmargin=*]
\item \emph{Success:} $1\times1$ and $3\times1$ errors at or below
  P2's, i.e.\ within 2 points of P2's seed-0 errors. Patch training then
  generalizes to the forced system, and the strip restriction was
  unnecessary.
\item \emph{Partial:} single-digit errors, but worse than P2. The
  model loses something by never seeing a full forcing column in $y$.
  The candidate cause is the large-scale $y$-structure, i.e.\ jets.
\item \emph{Failure:} errors above 15\%, or collapse. The $y$-dependent
  forcing and large-scale jet structure are not recoverable from
  $40\times40$ patches; Kolmogorov needs full-$y$ strips or a
  forcing-aware model.
\item Beyond the error itself: P9 reaching P2's accuracy at
  $1\times1$, even though it was \emph{never} trained on a periodic
  $y$-direction, would itself be the $y$-transfer result discussed in
  \S\ref{sec:post_p7_options}.
\end{itemize}

\textbf{P9 outcome: FAILURE at $3\times$} (\S\ref{sec:p9_results}) --
$42.6\%$ error (seed 0); the mechanism is open (the evaluation boxes are
$2\pi$ tall at every size, so the $y$-forcing alone cannot explain it).

\textbf{Queue} (\texttt{scripts/run\_c19a\_p9\_p8.sh}, all at
background priority): the first C19 arm's evaluations, then P9 (data
generation already started, an estimated 3--4 hours), then P8 training
and evaluations. \textbf{C19 was cut to one arm} (user, 2026-10-04).
P7 showed the full 12-site margin is no obstacle when the training
domain is large, so shrinking the margin matters only for the
small-domain case (P6). The first arm's autoencoder (decoder reach 0)
reached a reconstruction error of 0.0705, against P2's 0.0082, about
9$\times$ worse: decoder context matters for reconstruction at $c{=}3$.
Its dynamics are still being evaluated. The arms with encoder reach 2
and 0 remain available as TO DO.

\subsection{P9 results (seed 0, 2026-10-04): FAILURE at $3\times$, mechanism open}
\label{sec:p9_results}

\begin{table}[h]
\centering
\small
\begin{tabular}{lccccc}
\toprule
size & truth $D_{\mathrm{KY}}$ & P9 $D_{\mathrm{KY}}$ ($n_{\mathrm{pos}}$) & P9 error & P2 error (seed 0) \\
\midrule
$1\times$ & 23.26 & 25.52 (11) & 9.7\% (over) & 3.0\% \\
$2\times$ & $56.84$ / $49.8\pm9.8$ & 53.86 (25) & 5.2\% / 8.1\% & 9.4\% \\
$3\times$ & 75.22 & \textbf{107.28} (52) & \textbf{42.6\% (over)} & 4.3\% \\
\bottomrule
\end{tabular}
\caption{P9: Kolmogorov flow, $\alpha{=}0.05$, trained on $40\times40$-site
2D crops (y-offsets aligned to the forcing period) of a $6\times6$-box
domain, deployed on square periodic boxes. Autoencoder interior
validation error 0.0079 (P2: 0.0082). $\lambda_1$ = 0.192, 0.216, 0.236
at $1/2/3\times$ -- rising with size, unlike P2's roughly flat
$\lambda_1$ across sizes, itself a sign of something size-dependent
entering the dynamics.}
\label{tab:p9_results}
\end{table}

\textbf{Pre-registered verdict: FAILURE.} The $3\times$ error, 42.6\%,
is nearly $3\times$ the 15\% failure threshold, and the model
\emph{overshoots} the true dimension by 32 units -- this is not a
graceful degradation but a qualitative breakdown.

\textbf{The failure is specific to $3\times$, not a smooth trend.} Error
rises non-monotonically: 9.7\% (over) at $1\times$, 5.2--8.1\% at
$2\times$ (on par with or better than P2), then 42.6\% (over) at
$3\times$. $\lambda_1$ also climbs with size (0.192 to 0.236) where
P2's does not. Both point to the same thing: the model's local update
rule becomes dynamically more unstable as more training-pattern copies
are tiled together, and that instability only becomes dominant once the
domain is large enough (between $2\times$ and $3\times$) -- an
emergent, size-triggered failure, not a per-site error that simply
accumulates linearly with area.

\textbf{Correction to a first interpretation (same day).} A first draft
of this section attributed the failure to the $y$-dependent forcing.
That does not follow from the geometry: every Kolmogorov evaluation box
is $2\pi$ tall and periodic in $y$ (\texttt{exp\_crop\_eval.py}); only
$L_x$ grows from $1\times$ to $3\times$. The $y$-structure the deployed
model sees is therefore \emph{identical} at all three sizes, so the
forcing alone cannot explain an error that appears only at $3\times$.

\textbf{What the evaluation actually tests.} P9 was trained on crops of
a $12\pi\times12\pi$ domain and never saw a flow that is periodic at
$2\pi$ in $y$. Deploying it on $2\pi$-tall boxes is therefore a
$y$-\emph{shrinking} transfer combined with $x$-growth. Candidate
explanations, not yet distinguished:
\begin{enumerate}[leftmargin=*]
\item \emph{Training/deployment physics mismatch.} A tall domain
  supports large-scale $y$-structure (jet meandering and rearrangement
  at wavelengths $>2\pi$ in $y$) that a $2\pi$-tall box suppresses. The
  model learned dynamics that include these routes. That the mismatch
  only surfaces at large $L_x$ is plausible but unexplained.
\item \emph{Seed variability.} Confirmed and large: IC seeds 1 and 2
  at $3\times$ give 77.0 (2.4\%) and 64.7 (14.0\%), so the three-seed
  spread is $83\pm22$ (\S\ref{sec:p9_diag}).
\item \emph{Forcing-phase context.} Still possible in a weaker form: the
  model never saw the global $y$-periodicity constraint, and the
  deployed boxes impose one it was not trained on.
\end{enumerate}

\textbf{What is solid.} The pre-registered verdict (FAILURE, $>15\%$
at $3\times$) stands for seed 0, and P9 does not reproduce P2's
accuracy on $2\pi$-tall boxes. The contrast with P7 remains
informative, but its \emph{mechanism} is not yet identified. P7's 2D KS
boxes were square and matched the isotropic training crops, while P9's
boxes are much shorter in $y$ than anything in its training data, so
the two experiments differ in more than the forcing.

\textbf{The fair test P9 has not had.} Deploying on boxes taller in
$y$ ($1\times2$, $2\times2$, which need new truths) would test the
model in the geometry it was trained for. Until then, the honest
summary is: 2D-patch training of Kolmogorov flow from a tall domain does
\emph{not} transfer to the standard $2\pi$-tall boxes, and the cause is
open.

\subsection{P9 diagnosis: seed scatter and rollout statistics (2026-10-04)}
\label{sec:p9_diag}

\textbf{More initial-condition seeds at $3\times$} (same trained model, at
normal priority, about 5 minutes each):
\begin{center}
\small
\begin{tabular}{lcccc}
\toprule
$3\times$ & IC seed 0 & IC seed 1 & IC seed 2 & mean $\pm$ sd \\
\midrule
P9 $D_{\mathrm{KY}}$ ($n_{\mathrm{pos}}$, $\lambda_1$) & 107.3 (52, 0.236) & 77.0 (36, 0.202) & 64.7 (30, 0.155) & $83\pm22$ \\
P2 $D_{\mathrm{KY}}$ ($n_{\mathrm{pos}}$) & 72.0 (34) & 69.6 (34) & 69.9 (34) & $70.5\pm1.3$ \\
truth & 75.22 (seed 0 only) & & & \\
\bottomrule
\end{tabular}
\end{center}
The ``superlinear growth'' reported first (seed 0) is \emph{not} a
consistent property of P9 at $3\times$. Its outcome depends strongly on
the initial condition (spread $\pm26\%$). P2, in contrast, gives the
same answer from every IC to within 2\%. P9's three-seed mean is only
10\% from the truth, but a model whose dimension varies by a factor of
1.7 with the IC is not reliable.

\textbf{Free-rollout statistics} (\texttt{scripts/exp\_rollout\_stats\_kolmogorov.py};
IC seed 0, 1000 model steps = 200 time units, statistics over the second
half; truth from the same IC; kinetic-energy shares):
\begin{center}
\small
\begin{tabular}{llccccc}
\toprule
size & & vort.\ std & KE & beyond-box & $|k|<2$ & forcing $|k|\in[3,5]$ \\
\midrule
$1\times$ & truth & 1.40 & 0.172 & -- & 9.6\% & 54.6\% \\
 & P2 & 1.42 & 0.175 & -- & 9.5\% & 55.1\% \\
 & P9 & \textbf{0.98} & 0.160 & -- & \textbf{49.6\%} & \textbf{26.0\%} \\
\midrule
$2\times$ & truth & 1.41 & 0.172 & 6.5\% & 14.6\% & 53.6\% \\
 & P2 & 1.38 & 0.172 & 7.2\% & 26.2\% & 52.1\% \\
 & P9 & \textbf{0.96} & 0.131 & \textbf{29.0\%} & \textbf{40.6\%} & \textbf{27.1\%} \\
\midrule
$3\times$ & truth & 1.43 & 0.176 & 9.3\% & 21.5\% & 54.4\% \\
 & P2 & 1.42 & 0.180 & 11.5\% & 23.8\% & 54.1\% \\
 & P9 & 1.43 & \textbf{0.207} & \textbf{20.7\%} & 18.6\% & 43.1\% \\
\bottomrule
\end{tabular}
\end{center}

\textbf{Findings.}
\begin{enumerate}[leftmargin=*]
\item \emph{P2's rollouts are statistically faithful at every size.}
  Variance, energy and scale distribution match the truth within a few
  percent at $1/2/3\times$. The only exception is a mild large-scale
  excess at $2\times$ (26\% against 15\%).
\item \emph{P9's rollouts are not, and they are not even consistently
  wrong.} At $1\times$ and $2\times$ P9 sits in a low-variance state
  (std 0.97 against 1.40) with energy moved from the forcing scale into
  large scales: half its energy is at $|k|<2$, against 10--15\% in the
  truth. At $3\times$, from the same IC seed, it sits in a different,
  higher-variance state with 17\% excess kinetic energy and twice the
  truth's beyond-box energy. P9 has at least two dynamical regimes, and
  which one it lands in depends on the domain size and the IC.
\item \emph{The $D_{\mathrm{KY}}$ error tracks the energetic regime, not
  the large-scale pile-up.} The pile-up at $1\times$/$2\times$ coexists
  with roughly correct $D_{\mathrm{KY}}$ (9.7\%, 5--8\%). The excess
  energy of the $3\times$ seed-0 regime coincides with the 42.6\%
  overshoot and the largest $\lambda_1$.
\item \emph{This is a training-distribution effect specific to P9.} P2
  shares P9's architecture and differs only in its training data
  ($x$-strips of a $16\times$-wide, $2\pi$-tall domain, against 2D
  patches of a $6\times6$-box domain), and it shows none of this. The
  tall training domain carries 7\% more field variance (std 1.479
  against 1.39--1.43) and supports large-scale $y$-structure that the
  $2\pi$-tall deployment boxes do not. P9 learned dynamics with energy
  pathways the deployment geometry does not have.
\end{enumerate}

\textbf{Is the truth itself multistable? Not strongly.} Truth and P9
rollouts of 400 time units (10$\times$ the truth Lyapunov runs) from IC
seeds 0 and 1 (\texttt{scripts/run\_multistability\_check.sh}); share of
kinetic energy at $|k|<2$ averaged over each quarter of the run:
\begin{center}
\small
\begin{tabular}{llcccc}
\toprule
 & & seed 0 & & seed 1 & \\
size & & quarters & late std & quarters & late std \\
\midrule
$2\times$ & truth & 13, 14, \textbf{29}, 19\% & 1.40 & 24, 22, 25, 22\% & 1.40 \\
 & P9 & 42, 40, 55, 59\% & 0.86 & 50, 47, 35, 48\% & 0.93 \\
$3\times$ & truth & 18, 21, 19, 19\% & 1.46 & 16, 14, 14, 14\% & 1.45 \\
 & P9 & 34, 19, 16, 21\% & 1.49 & 42, 63, 50, 65\% & 0.89 \\
\bottomrule
\end{tabular}
\end{center}
Seed 2 (same protocol): truth $2\times$ quarters 15, 13, 17, 21\% (std
1.46), $3\times$ 19, 21, 18, 18\% (std 1.48); P9 $2\times$ 48, 57, 40,
45\% (std 0.94), $3\times$ 43, 77, 76, 68\% (std 0.81).
The truth's large-scale share stays within about 14--25\% across all three seeds
and sizes, with one temporary excursion to 29\% (seed 0, $2\times$). That
excursion is consistent with the 40-unit $2\times$ truth scatter, but
it is not a set of distinct coexisting states. P9 instead lands in a
low-variance pile-up (40--65\%) or in an energetic truth-like-variance
regime depending on the IC. Its regimes are not real physics.

\textbf{The decisive check: P9 in its own geometry.} A free P9 rollout on
a $1\times2$ box (taller than $2\pi$, like its training crops), against
the true PDE from the same IC: $|k|<2$ share 60.2\% against 12.2\%,
forcing-scale share 20.4\% against 54.4\%, std 0.97 against 1.43. P9
drifts even in the geometry it was trained for. A second tall box,
$3\times3$ (seed 0), was closer: std 1.40 against 1.49, KE 0.208 against
0.188, $|k|<2$ share 23\% against 17\%, but forcing-scale share still
36\% against 56\%. A matching geometry helps, but does not remove the
drift. The failure is therefore
a \textbf{training-stability problem}, not a train/deploy geometry
mismatch. The tall training data itself has no pile-up (19\% at $|k|<2$,
55\% at the forcing scale, std 1.48), so the error accumulates in free
rollout.

\textbf{Likely cause and fix.} In P9's $40\times40$-site crops the
latent strip is $28\times28$ sites, of which only the $16\times16$
interior (33\%) runs free; the 6-site halo is reset to the truth every
step on all four sides, and rollouts are at most 12 steps. P2's strips
are periodic and entirely free in $y$ (67\% free). P9 was therefore
heavily teacher-forced and never had to stay on the attractor on its own
for long. P7 (2D KS) used the same 33\% geometry and succeeded,
presumably because 2D KS has no large-scale energy pathway comparable
to Kolmogorov's inverse cascade.

\textbf{P9b outcome} (\S\ref{sec:p9b_results}): instability fixed (seed
spread $21.6\%\to2.2\%$ at $3\times$); a smaller, consistent
undershoot (11--15\%) remains, most likely from the stabilizing
training change itself (more end-of-training noise, longer rollouts)
trading some chaotic growth for stability.

\textbf{P9b (COMPLETED 2026-10-05, normal priority).} Same data, same
autoencoder (copied from P9; reconstruction was not the problem), and a
propagator retrained with longer free rollouts: $k_{\max}$ 12$\to$32,
$k_{\mathrm{mid}}$ 8$\to$16, 9000 steps, end noise 0.02$\to$0.05
(\texttt{scripts/run\_p9b\_c19\_p8.sh}). Evaluation: rollout statistics
at $1/2/3\times$, then $D_{\mathrm{KY}}$ at $1/2/3\times$ for IC seeds
0--2. \emph{Rule:} P9b is fixed if its free-rollout statistics stay
within 25\% of the truth's (std, forcing-scale share, $|k|<2$ share) and
its $3\times$ $D_{\mathrm{KY}}$ IC-seed spread is under 10\% with mean
error under 10\%. If longer rollouts are not enough, the next step is
larger crops with a bigger free interior, which needs a larger training
domain.

\subsection{P9b results (COMPLETED 2026-10-05): instability fixed, a smaller systematic bias remains}
\label{sec:p9b_results}

All 9 evaluations (3 sizes $\times$ 3 IC seeds) are in.

\begin{center}
\small
\begin{tabular}{lccccc}
\toprule
size & truth & P9b seeds (0,1,2) & P9b mean $\pm$ sd & mean error & P9 (original) \\
\midrule
$1\times$ & 23.26 & 20.63, 20.32, 21.24 & $20.73\pm0.38$ & 10.9\% & 25.52 (9.7\%, 1 seed) \\
$2\times$ & $49.8\pm9.8$ & 41.59, 41.31, 44.10 & $42.33\pm1.25$ & 15.0\% & 53.86 (8.1\%, 1 seed) \\
$3\times$ & 75.22 & 66.65, 63.64, 63.50 & $64.60\pm1.45$ & 14.1\% & $83.0\pm17.9$ (19.7\%, 3 seeds) \\
\bottomrule
\end{tabular}
\end{center}

\textbf{Against the pre-registered rule} (rollout statistics within 25\%,
$3\times$ seed spread $<10\%$, mean error $<10\%$): \textbf{two of three
criteria pass; the third is close.} The rollout statistics passed
cleanly (\S\ref{sec:p9_diag}: all within 10--15\% of the truth, well
under the 25\% bar). The spread criterion passes decisively: relative
sd at $3\times$ is $1.45/64.6=2.2\%$, against P9's $17.9/83.0=21.6\%$ --
a \textbf{10-fold reduction in seed-to-seed scatter}. The mean-error
criterion (14.1\%) narrowly misses the 10\% threshold.

\textbf{Diagnosis confirmed, cleanly separated into two effects.} P9's
single pathology -- an IC-dependent, size-dependent regime collapse
into a spurious high-energy or low-energy state -- is gone. In its
place is a small, highly consistent systematic undershoot (11--15\%,
not growing monotonically with size: 10.9\%, 15.0\%, 14.1\%). This is
exactly what the training change targeted (longer free rollouts force
the model to be stable under its own dynamics, which it now is) and
exactly the kind of residual error that instability was previously
masking: with P9, no single seed's error was a reliable estimate of
anything, because the model could land in either a 2\% match or a
43\% miss from the same architecture and the same training run,
depending on the IC. P9b replaces that with a boring, reproducible
answer, which is the more useful kind of wrong.

\textbf{What the remaining bias is probably not.} It is not a capacity
limit (P9b uses the same $c{=}3$, same latent size as P2, which does
not undershoot); not the $y$-forcing/geometry mismatch (ruled out in
\S\ref{sec:p9_diag}, and P9b's rollout statistics now match the truth's
energy distribution closely at all three sizes); and not seed noise (sd
of 1.25--1.45 across three different ICs is far too tight for that).
The most likely remaining cause is the training change itself: longer
free rollouts with higher terminal noise (0.05 against the baseline
0.02) bias the propagator toward \emph{damping} slightly, trading some
genuine chaotic growth for stability. That is a tunable side effect,
not a new instability.

\textbf{Verdict and next step.} \textbf{Net improvement, not full
closure.} The informative-failure mode (wild, unusable IC dependence)
is resolved; a milder, consistent, and more tractable bias remains.
\textbf{TO DO}, not run tonight: a second point on the
stability/accuracy trade-off -- an intermediate setting ($k_{\max}=20$,
end noise $0.035$) between P9 (unstable, $k_{\max}{=}12$, noise $0.02$)
and P9b (stable but biased, $k_{\max}{=}32$, noise $0.05$) -- to see
whether a lighter touch keeps the stability gain while recovering some
of the lost magnitude.

\subsection{What C19 and P9b together suggest about improving the
propagator (2026-10-05, discussion)}
\label{sec:propagator_lessons}

C19 (encoder/decoder reach) and P9/P9b (propagator training length)
look at first like the same kind of fix applied twice, but they are
not, and the difference is itself a useful design principle.

\textbf{The autoencoder's reach was a capacity illusion; the
propagator's reach is not, and should not be cut the same way.} A
static snapshot's content is, by construction, entirely contained
within its own patch, so encoder/decoder reach could only ever help
\emph{compress} that content more efficiently (borrow a neighbour's
information to describe your own patch in fewer numbers) -- it could
never supply information that was not already there. That is why
raising the per-site channel count (c19\_e0d0\_c6, $c{=}6$) substitutes
for reach with no loss. The propagator does not have this property:
predicting $z_{i,t+1}$ genuinely needs $z_{i\pm1,t}$, because
information really does propagate between sites between time steps --
that is the light-cone bound (H2, brief Phase 2), not a
compression artefact. No amount of extra per-site latent capacity at
site $i$ substitutes for being unable to look at site $i{+}2$'s current
state when that state causally determines site $i$'s future.
\textbf{C19 therefore does not license shrinking the propagator's
window}, and does not by itself rescue P6's narrow encoder (reach 2):
C16 suggests that reach is below the forcing-scale structure's physical
extent, which would be a genuine light-cone violation, not a
capacity-substitutable shortfall the way the autoencoder's reach was.
That distinction is itself TO DO to confirm: rerun a P6-style narrow
encoder at reach 2 with $c{=}6$-to-8 channels (capacity, not reach) and
see whether it still fails -- if it does, P6's problem really is reach,
not capacity, and sits in a different category from C19's.

\textbf{What does transfer: architecture failures can be training-regime
failures in disguise.} P9's instability looked structural (wild,
IC-dependent collapse) but the identical window-3 architecture worked
cleanly for P7 on 2D KS; what was broken was that training only ever
asked the propagator to run freely for 12 steps inside a crop that was
67\% teacher-forced, so it never had to prove it could stay on the
attractor under its own compounding error. Longer free rollouts
(P9b, $k_{\max}{=}32$) fixed the instability outright. The general
principle, stated once for reuse: \textbf{before changing an
architecture, check whether the failure is a capacity problem (fixed by
more per-unit channels/width) or a training-protocol problem (fixed by
longer free rollouts, not by changing the network) rather than assuming
it is structural.}

\textbf{TO DO, not run tonight -- concrete propagator experiments this
suggests, in priority order:}
\begin{enumerate}[leftmargin=*]
\item \emph{Disentangle P9b's two simultaneous changes.} Going from P9
  to P9b changed both $k_{\max}$ ($12\to32$) and end noise
  ($0.02\to0.05$) at once, so it is not yet known which one fixed the
  instability and which one caused the new 11--15\% undershoot. The
  canonical recipe in \S\ref{sec:consolidation} already holds end noise
  at the 0.02 baseline while keeping $k_{\max}{=}32$, which tests this
  as a byproduct of Phase 1 -- no separate run needed, but call out the
  result explicitly when Phase 1 lands.
\item \emph{A second point on the stability/accuracy trade-off}
  (already flagged in \S\ref{sec:p9b_results}): an intermediate setting
  ($k_{\max}{=}20$, end noise $0.035$) between P9 and P9b, to see
  whether a lighter touch keeps the stability gain while recovering
  some of the lost magnitude.
\item \emph{A targeted rollout loss, rather than relying on
  $k_{\max}$ as a blunt proxy -- IMPLEMENTED 2026-10-05, user-directed.}
  The diagnostic that actually caught both P9's and P8's failures was
  the free-rollout energy-spectrum statistics
  (\texttt{exp\_rollout\_stats\_kolmogorov.py}: variance, and energy
  share at large vs.\ forcing scales); this makes it a differentiable
  auxiliary training loss instead of a post-hoc check.
  \texttt{ks\_latent.training.crop.spectral\_band\_shares(field, Lx, Ly,
  bands)} computes the share of Poisson-solve kinetic energy
  (`$|\widehat{\mathrm{field}}|^2/k^2$', the same quantity the
  diagnostic script uses) in each $|k|$ band, fully differentiably.
  \texttt{CropPropConfig.w\_spectral} (default 0, every existing recipe
  unaffected) adds, every \texttt{spectral\_every} steps, an MSE penalty
  between the rolled-out field's band shares and a target computed once
  from the true training data. \textbf{Exact only for decoder reach 0}:
  the regularizer decodes a training window directly via a weight-shared
  resize of the frozen autoencoder (\texttt{\_resize\_decoder\_for\_\allowbreak
  spectral\_reg}), which would introduce circular-wraparound artefacts at
  that window's own edges for any nonzero decoder reach -- asserted, not
  silently wrong. Wired into \texttt{train\_kolmogorov\_crop.py} as
  \texttt{--w-spectral}/\texttt{--spectral-bands}/\texttt{--spectral-every}
  (works for both Kolmogorov and 2D KS, since the script is dataset-agnostic).
  4 new unit tests (known-mode recovery, differentiability, the
  nonzero-reach assertion, and an end-to-end training smoke test); 402
  related tests pass. \textbf{Not yet evaluated against real data} --
  queued as the 2D KS ablation below.
\item \emph{Spend capacity on per-site expressivity, not reach, at fixed
  window.} The legitimate analogue of C19's finding: keep the window at
  its light-cone-justified value, but test whether a wider/deeper
  per-site update function (larger \texttt{token\_d\_model}, more
  residual blocks, or the existing \texttt{backbone="vit"} option) buys
  stability or accuracy the way extra channels did for the autoencoder.
  Untested; would isolate ``more capacity at fixed reach'' cleanly from
  everything else changed so far.
\item \emph{Confirm or refute the P6/C19 distinction above} (reach-2,
  $c{\ge}6$ encoder): does more per-site capacity rescue a
  physically-too-narrow encoder reach, or not? This is the cleanest test
  of whether ``capacity substitutes for reach'' is specific to the
  autoencoder's snapshot-compression job or is a more general property
  worth exploiting elsewhere.
\end{enumerate}

\subsection{Training-distribution error caught before more compute was spent:
2D-patch training was the wrong default for Kolmogorov (2026-10-05, user-directed)}
\label{sec:kolmogorov_training_redirect}

\textbf{The error.} Phase 1's ``canonical recipe'' (\S\ref{sec:consolidation}) applied 2D
patch training uniformly to Kolmogorov, carrying over P7's successful
recipe for 2D KS without re-checking it against this document's own
already-established P9 verdict (\S\ref{sec:p9_results}: FAILURE -- 2D
patch training does not suit Kolmogorov's anisotropic, forced dynamics,
because cropping in $y$ breaks the thing $x$-strip training preserves).
This was a process error (not re-deriving a result already in the
document), not a new finding.

\textbf{How it was caught.} \texttt{k\_a005} (C19 autoencoder, 2D patch,
long rollout) was compared against the \emph{already-complete}
\texttt{c19\_e0d0\_c6} (C19 autoencoder, $x$-strip, short rollout,
\S\ref{sec:phase0_results}) -- a comparison available as soon as
\texttt{k\_a005} finished, requiring no new runs:
\begin{center}
\small
\begin{tabular}{lccc}
\toprule
recipe (mean error) & $1\times$ & $2\times$ & $3\times$ \\
\midrule
C19, $x$-strip, short rollout (\texttt{c19\_e0d0\_c6}) & 8.7\% & 8.3\% & 2.9\% \\
C19, 2D-patch, long rollout (\texttt{k\_a005}) & 12.9\% & 13.6\% & 15.1\% (1 seed) \\
old AE, 2D-patch, long rollout (P9b, \S\ref{sec:p9b_results}) & 10.9\% & 15.0\% & 14.1\% \\
\bottomrule
\end{tabular}
\end{center}
\textbf{2D-patch training is worse at every size, with two different
autoencoders landing in the same degraded 11--15\% range}: the gap
is attributable to training-data shape, not to which autoencoder sits
on top of it. P9b's long rollouts fixed P9's \emph{instability}
(spread $22\%\to2\%$ at $3\times$) but never closed the accuracy gap to
$x$-strip training -- a fact already on record
(\S\ref{sec:p9b_results}: ``net improvement, not full closure'') that
the Phase 1 setup should have weighed before choosing 2D patches as the
Kolmogorov default.

\textbf{Action taken.} \texttt{k\_a01}, \texttt{k\_b005},
\texttt{k\_a0025}, \texttt{k\_c0025} stopped -- none had produced a
trained checkpoint, so no compute was wasted beyond the training
already in progress for \texttt{k\_a01} (partial, discarded) and
\texttt{k\_b005}/\texttt{k\_c0025} (not yet started). Replaced with the
one combination not yet tested: $x$-strip training (exactly
\texttt{c19\_e0d0\_c6}'s own flags, alpha swapped) with the P9b long
rollout added on top -- isolating whether long rollouts help Kolmogorov
\emph{at all} once the already-known-bad training-distribution choice
is removed. \textbf{$\alpha{=}0.025$} (user-specified): the hardest
case, where the original recipe's native-size collapse was worst
(\S\ref{sec:p8_model_first}: $n_{\mathrm{pos}}$ 4--9 against truth 14).
\begin{description}[leftmargin=*]
\item[\texttt{k\_xs0025}] $x$-strip + long rollout, no regularizer.
\item[\texttt{k\_xsr0025}] $x$-strip + long rollout + spectral
  regularizer (same arm pairing idea as \texttt{k\_b005}/\texttt{k\_c0025},
  now on the correct training distribution).
\end{description}
Evaluated against the existing $\alpha{=}0.025$ truths, 3 IC seeds per
size (\texttt{scripts/run\_k\_xstrip0025.sh}).

\subsection{\texttt{k\_b005}: spectral regularizer tested on $\alpha{=}0.05$
first, with a direct same-dataset baseline (QUEUED 2026-10-05, user-directed,
priority insert)}
\label{sec:k_b005_queued}

User-directed re-prioritization: \texttt{k\_a005}
(\S\ref{sec:k_a005_results}) already showed a confusing non-monotone
pattern across sizes (native 0.9/12.9\% by seed, $2\times$ 29.5\%,
$3\times$ 15.1\%, seed 0), giving a direct, same-dataset baseline to
test the regularizer against -- better motivated than waiting for
$\alpha{=}0.025$ (\S\ref{sec:spectral_redirect}), since we already have
a concrete problem (the $2\times$/$3\times$ degradation) to see whether
it fixes.

\texttt{k\_b005}: identical recipe to \texttt{k\_a005} except
$k_{\max}$ $32\to16$ ($k_{\mathrm{mid}}$ $16\to8$, user: ``only rollout
to 16 steps'') \emph{and} \texttt{--w-spectral 1.0} added. \textbf{Two
variables change at once} (rollout length and the regularizer) -- a
deliberate test of a specific combined hypothesis (can the regularizer
let a shorter, cheaper rollout match or beat the long-rollout result),
not a single-variable ablation of \texttt{k\_a005}; the same
attribution caution as P9$\to$P9b applies; if it helps, a
\texttt{k\_max{=}32}+regularizer arm would be needed to separate the
two effects. Queued with priority ahead of \texttt{k\_a0025}/
\texttt{k\_c0025} in the training lane
(\texttt{scripts/run\_k\_b005.sh}, polls specifically for
\texttt{k\_a01}'s checkpoint so it grabs the GPU immediately once
\texttt{k\_a01} finishes); 2 IC seeds, 3 sizes, against the same truths
as \texttt{k\_a005}.

\subsection{Spectral-regularizer target changed to Kolmogorov $\alpha{=}0.025$
(2026-10-05, user-directed)}
\label{sec:spectral_redirect}

\texttt{p7c\_ks2d} (2D KS, canonical recipe + regularizer) was cancelled
after 750/9000 propagator steps, before any evaluation. Reason: its
baseline (\texttt{p7b\_ks2d}, \S\ref{sec:p7b_results}) already lands
within 1\% of the extensive prediction at all three sizes, leaving
almost no headroom for a regularizer to show a measurable effect.
\textbf{\texttt{k\_c0025}} (Kolmogorov $\alpha{=}0.025$, same canonical
recipe + the regularizer, Kolmogorov-default bands \texttt{0,2,3,5})
replaces it: \texttt{k\_a0025} (canonical recipe, no regularizer,
queued in the same run) is the system with the largest unresolved
failure in this document -- the original recipe collapsed to
$n_{\mathrm{pos}}$ 4--9 against a truth of 14 at native size
(\S\ref{sec:p8_model_first}) -- so there is real headroom here to see
whether the regularizer helps where it is actually needed. Queued to
start once the Phase 1/2 training lane (\S\ref{sec:consolidation})
clears (\texttt{scripts/run\_k\_c0025.sh}); evaluated against the
existing $\alpha{=}0.025$ truths, 2 IC seeds per size.

\subsection{Phase 2 (2D KS) spectral-regularizer ablation
(QUEUED 2026-10-05)}
\label{sec:phase2_spectral_ablation}

User-directed: retrain P7 (2D KS crop training) with the full
consolidation-plan canonical recipe -- the C19 per-patch $c{=}6$
autoencoder (\S\ref{sec:c19_c6_results}) and P9b's long free rollouts
($k_{\max}{=}32$, noise held at the 0.02 baseline per
\S\ref{sec:consolidation}) -- as two arms, with and without the new
spectral-band-share regularizer (\S\ref{sec:propagator_lessons} item 3),
so the regularizer's effect is isolated rather than bundled into one
``everything at once'' retrain.
\begin{description}[leftmargin=*]
\item[\texttt{p7b\_ks2d}] canonical recipe, no regularizer.
\item[\texttt{p7c\_ks2d}] canonical recipe + \texttt{--w-spectral 1.0
  --spectral-bands "0,1,1,999"} (large scale, $|k|<1$, i.e.\ wavelength
  $>2\pi\approx6.28$, against the KS instability wavelength
  $2\sqrt2\pi\approx8.89$, against everything shorter. A first
  reasonable default, not yet tuned.)
\end{description}
Both evaluated at $44^2$/$66^2$/$88^2$ (scales 1.0/1.5/2.0), 3 IC seeds
each -- the $88^2$ seeds are exactly the check
\S\ref{sec:p7_88_check} flagged as needed to tell a genuine
sub-extensive ceiling from an unlucky single-seed draw into a spurious
regime (the same mechanism P9 had). Runs after Phase 0 (C19 verify) and
the Phase 1 datasets clear, one job at a time
(\texttt{scripts/run\_phase2\_2dks\_spectral\_ablation.sh}).

\emph{Pre-registered reading.} If \texttt{p7c\_ks2d}'s three $88^2$
seeds cluster tightly (like P9b's $3\times$, sd/mean $\lesssim3\%$) near
the extensive prediction ($\approx327$--345), the regularizer fixed a
real free-rollout drift and should become part of the canonical recipe.
If \texttt{p7b\_ks2d} (no regularizer) already clusters tightly near
that prediction, the single-seed 215.86 was the P9-style unlucky draw,
not a model problem, and the regularizer is unneeded for 2D KS. If both
arms scatter widely, $88^2$ needs the real truth run before either
reading is possible.

\subsection{C19 \texttt{c19\_e0d0\_c6} results (COMPLETED 2026-10-05):
spatial context in the autoencoder is unnecessary, not merely reducible}
\label{sec:c19_c6_results}

The second, dropped-then-requeued C19 arm (\S\ref{sec:c19}): encoder
reach 0, decoder reach 0, $c{=}6$ channels (against P2's reach 6/6,
$c{=}3$), otherwise the unchanged P2 recipe.

\begin{center}
\small
\begin{tabular}{lcccc}
\toprule
size & truth & \texttt{c19\_e0d0\_c6} (0/0 reach, $c{=}6$) & P2 (6/6 reach, $c{=}3$) & \texttt{c19\_e6d0} (6/0, $c{=}3$) \\
\midrule
$1\times$ & 23.26 & \textbf{23.52 (1.1\%)} & 23.95 (3.0\%) & 21.33 (8.3\%) \\
$2\times$ & $49.8\pm9.8$ & \textbf{50.66 (1.7\%)} & 45.14 (9.4\%) & 44.09 (11.5\%) \\
$3\times$ & 75.22 & \textbf{74.16 (1.4\%)} & 71.96 (4.3\%) & 65.20 (13.3\%) \\
\bottomrule
\end{tabular}
\end{center}

\textbf{Superseded by the three-seed verification}
(\S\ref{sec:phase0_results}): with all IC seeds the per-patch model is
about as accurate as P2 at $1\times$/$2\times$ and better at $3\times$, not
better everywhere; the gate still passes. Original single-seed reading
follows. \textbf{Pre-registered verdict: VIABLE, by a wide margin -- and better
than the reference it was tested against.} The thresholds were 8\%
native / 12\% at $3\times$; this arm clears both by roughly $6\times$,
and beats P2's own accuracy at all three sizes.

\textbf{What this establishes.} A strictly per-patch autoencoder --
each site's latent reconstructs only its own $16\times16$-grid patch,
with \emph{zero} information about any neighbouring patch, at encoding
or decoding time -- matches or exceeds the full-reach encoder/decoder's
accuracy, provided the per-site channel count is raised from 3 to 6.
This is a stronger result than the arm it supersedes
(\texttt{c19\_e6d0}, encoder reach 6 only, which was narrowly not
viable at $c{=}3$, \S\ref{sec:c19_results}): spatial context does not
trade off against channel count in the way that result suggested --
\emph{no} spatial context, with enough channels, is strictly better than
\emph{full} spatial context with too few. The earlier finding (reach 0
costs 9\% native accuracy at $c{=}3$) was a statement about capacity,
not about locality: 3 channels were not enough to hold a
$16\times16$-patch's worth of local chaotic content without borrowing
from neighbours; 6 channels are.

\textbf{Consequence: the autoencoder's edge margin can be eliminated
entirely, not merely halved.} Crop training's edge mask is the sum of
encoder reach and decoder reach (\S\ref{sec:c19}); with both at zero,
training needs no autoencoder margin at all, only the propagator's
6-site forced halo. This removes the obstacle that made P6
(\S\ref{sec:p6_results}) and the general "train a local model on
windows of the one domain you have" claim (paper 2, CM1/Rayleigh--B\'enard,
\S\ref{sec:corollary_cm1}) need either a smaller margin or a larger
domain. \textbf{P6 should be redone} with this autoencoder
(\texttt{c19\_e0d0\_c6}'s architecture) and a propagator window sized
by the light-cone bound (\S\ref{sec:p6_results} proposed $1\times2$,
$2\times2$, or $2\times3$ encoder-reach sweeps; with the autoencoder
margin now zero, the propagator's own window is the only thing left to
size) on the existing $3\times$ domain -- no larger domain is required.

\textbf{One qualification.} This channel increase was tested on a
per-patch (reach-0) autoencoder specifically; it is not yet shown that
raising $c$ alone (reach 0, same recipe otherwise) is what helped, as
opposed to some interaction with reach -- a $c{=}6$, full-reach
(6/6) control would isolate that, and is TO DO but not urgent, since
the practical consequence (zero-margin crops) only needs the reach-0
result in hand.

\subsection{Phase 0 results: C19 verification (2026-10-05) -- gate PASSED}
\label{sec:phase0_results}

Consolidation Phase 0 (\S\ref{sec:consolidation}): more IC seeds for the
per-patch $c{=}6$ autoencoder, a second training seed, and a
full-reach $c{=}6$ control. All on Kolmogorov flow, $\alpha{=}0.05$,
$16\times$-wide $x$-strip training (the P2 recipe, only the
autoencoder changed). $D_{\mathrm{KY}}$ and error against the truths
23.26 / $49.8\pm9.8$ / 75.22; P2's three IC seeds already existed and are
included for a like-for-like comparison.

\begin{center}
\small
\begin{tabular}{llccccc}
\toprule
size & model & IC 0 & IC 1 & IC 2 & mean err. & spread (sd/mean) \\
\midrule
$1\times$ & P2 ($6/6$, $c{=}3$) & 23.95 (3.0) & 20.47 (12.0) & 24.89 (7.0) & 7.3\% & 8.1\% \\
 & C19 ($0/0$, $c{=}6$) & 23.52 (1.1) & 20.22 (13.1) & 26.04 (11.9) & 8.7\% & 10.2\% \\
 & C19, 2nd train seed & 24.17 (3.9) & -- & -- & -- & -- \\
 & control ($6/6$, $c{=}6$) & 25.91 (11.4) & -- & -- & -- & -- \\
\midrule
$2\times$ & P2 & 45.14 (9.4) & 44.84 (10.0) & 44.94 (9.8) & 9.7\% & 0.3\% \\
 & C19 & 50.66 (1.7) & 48.52 (2.6) & 39.56 (20.6) & 8.3\% & 10.4\% \\
 & C19, 2nd train seed & 48.48 (2.7) & -- & -- & -- & -- \\
\midrule
$3\times$ & P2 & 71.96 (4.3) & 69.55 (7.5) & 69.86 (7.1) & 6.3\% & 1.5\% \\
 & C19 & 74.16 (1.4) & 76.07 (1.1) & 70.45 (6.3) & 2.9\% & 3.2\% \\
 & C19, 2nd train seed & 72.99 (3.0) & -- & -- & -- & -- \\
\bottomrule
\end{tabular}
\end{center}
Errors in \% in parentheses. The control's $2\times$/$3\times$ evaluations
are still running.

\textbf{Gate verdict: PASSED.} Against P2's three-seed means: $1\times$
8.7\% against 7.3\% (within the 2-point margin); $2\times$ 8.3\% against
9.7\%; $3\times$ 2.9\% against 6.3\% (and under the 6\% bar); second
training seed 3.0\% at $3\times$ (under 8\%). \textbf{The per-patch
$c{=}6$ autoencoder becomes the canonical autoencoder} for Phases 1--4.
This also corrects the first C19 write-up
(\S\ref{sec:c19_c6_results}), which rested on IC seed 0 only: with all
seeds, per-patch $c{=}6$ is about as accurate as P2 at $1\times$ and
$2\times$ and better at $3\times$, rather than better everywhere. The case
for adopting it is accuracy at least as good, plus a zero crop margin,
plus the DA advantages in \S\ref{sec:consolidation} Phase 4.

\textbf{Finding: the native-size scatter is set by the initial
condition, not by the model.} At $1\times$, P2 and C19 -- two different
autoencoders, trained separately -- give nearly the same
$D_{\mathrm{KY}}$ for the same IC: 23.95 vs 23.52 (IC 0), 20.47 vs 20.22
(IC 1), 24.89 vs 26.04 (IC 2). The low IC-1 value appears in both. Two
independent models agreeing on a per-IC deviation means the deviation
is a property of the dynamics they are both reproducing, not of either
model's error. The evaluation is a finite-time Benettin estimate (2000
steps, 400 time units) from a specific IC on an 8-site box; such
estimates scatter, and the truth's own single-run scatter was already
$\pm20\%$ at $2\times$ (\S\ref{sec:truth3_revision}). The truth runs are also
much shorter than the model evaluations (\texttt{run\_truth\_kolmogorov.py}:
40 time units with 20 of warm-up, against 400 for the model's Benettin
run), so each truth value is the \emph{noisier} of the two estimates. A
fair per-IC comparison should integrate the truth over the same 400
units, or at least several seeds. The single 23.26
native truth is the weak link: part of the ``native error'' is the
comparison of a per-IC finite-time value against one fixed number.

\textbf{Per-IC truths (2026-10-05) -- the measurement explanation is only
partly right.} $1\times$ truths integrated from the \emph{same} three ICs
as the model evaluations, 120 time units, 40 directions:
\begin{center}
\small
\begin{tabular}{lccccc}
\toprule
IC & truth & P2 & C19 & C19, 2nd train seed & control \\
\midrule
0 & 22.46 (12) & 23.95 (+6.6\%) & 23.52 (+4.7\%) & 24.17 (+7.6\%) & 25.91 (+15.3\%) \\
1 & 20.52 (10) & 20.47 ($-$0.3\%) & 20.22 ($-$1.5\%) & -- & -- \\
2 & 21.63 (12) & 24.89 (+15.1\%) & 26.04 (+20.4\%) & -- & -- \\
\bottomrule
\end{tabular}
\end{center}
\begin{itemize}[leftmargin=*]
\item \emph{The IC-1 dip is real dynamics:} the truth dips too, and both
  models match it within 1.5\%.
\item \emph{The old single-run truth was high:} the three-IC truth mean
  is 21.54 against 23.26.
\item \emph{But a genuine native-size model error remains.} Every model
  lies above its per-IC truth except at IC 1, by 15--20\% at IC 2. The
  models' IC scatter ($\pm10\%$) is about $2.5\times$ the truth's
  ($\pm4\%$). Mean per-IC error: P2 7.3\%, C19 8.9\%, essentially
  unchanged.
\end{itemize}
So the models \emph{track} the true IC dependence but \emph{amplify} it
and lean high on the 8-site box. That is more consistent with the
model-side hypothesis (the propagator's 6-site reach nearly wraps an
8-site ring, a regime absent from training) than the first reading
allowed. Item 1 below is revised accordingly; Phase 1's long-rollout
models are the direct test.

\textbf{Speculation and implications.}
\begin{enumerate}[leftmargin=*]
\item \emph{(Revised by the per-IC truths above: only partly a
  measurement problem.) The native-size ``problem'' may be mostly a measurement
  problem.} If the truth is integrated from the same three ICs, the
  natural score is the per-IC error, not the error against a single
  truth. If per-IC errors are small, the native numbers in this whole
  document (P2, P9b, C19, P8) overstate model error at $1\times$, and the
  $1\times$ gates should be restated against a multi-seed truth. This is
  cheap to check: three $1\times$ truths at about an hour each.
  \textbf{TO DO, first priority before any $1\times$ claim.}
\item \emph{This weakens, but does not rule out, the hypothesis that
  propagator capacity or rollout length fixes native-size behaviour.} The
  user-raised mechanism is real in principle: the 8-site native box is
  smaller than every training crop, and the propagator's reach (6 sites)
  nearly wraps the whole ring, an input regime absent from training.
  But a model-side defect would not produce per-IC agreement across two
  different models; most of the scatter is shared. A model-side effect
  can only explain the residual differences (e.g.\ IC 2: 24.89 vs 26.04;
  the control's 25.91 at IC 0 against 23.5--24.2 for the others). Phase 1
  (long rollouts) will test the hypothesis anyway, at no extra cost.
\item \emph{P2's tightness at $2\times$ (0.3\% spread) is suspicious in
  the other direction.} The truth scatters by about $\pm20\%$ at
  $2\times$ across ICs, but P2 gives 44.8--45.1 from all three. A model
  that is \emph{less} IC-sensitive than the true system may be averaging
  over dynamical variability the true flow has, as already suspected in
  \S\ref{sec:p9_diag}. C19's larger spread (including the 39.56 at IC 2)
  could be the more faithful behaviour. This can only be settled with
  per-IC truths at $2\times$.
\item \emph{Extensivity of the C19 model is clean.} Dimension per box
  is 23.3 / 23.1 / 24.5 at $1/2/3\times$ (three-seed means), about
  constant, as an extensive local model should be, against P2's 23.1 /
  22.5 / 23.5.
\item \emph{The control separates channels from reach (completed).}
  Full reach with $c{=}6$ (\texttt{c19\_e6d6\_c6}), IC 0: 25.91 / 46.42 /
  71.84 at $1/2/3\times$, i.e.\ 15.3\% (against the per-IC truth), 6.8\%
  and 4.5\%. Against per-patch $c{=}6$ at the same IC (4.7 / 1.7 / 1.4\%;
  second training seed 7.6 / 2.7 / 3.0\%) and P2 (6.6 / 9.4 / 4.3\%):
  \textbf{at equal channel count, per-patch beats full reach at every
  size}, and doubling channels while keeping full reach lands roughly
  where P2 does (worse at $1\times$). The gain is therefore not from
  channels alone; removing reach is at least neutral and probably
  helps. One IC seed for the control, so a lean rather than a verdict,
  but consistent in direction at all three sizes. A plausible reason:
  with reach, the latent at a site mixes neighbouring content, so the
  propagator's input already contains partly-propagated information it
  must then disentangle; a per-patch latent is a cleaner local state for
  the propagator to act on.
\end{enumerate}

\subsection{P8 truth runs at $\alpha{=}0.025$ (IN PROGRESS)}
\label{sec:p8_truths}

\begin{table}[h]
\centering
\small
\begin{tabular}{lcccc}
\toprule
$\alpha$ & $1\times$ $D_{\mathrm{KY}}$ & $n_{\mathrm{pos}}$ & $\lambda_1$ & source \\
\midrule
0.1 & 15.88 & & & C4 control \\
0.05 & 23.26 & & & C4 / P2 \\
0.025 & 25.23 & 14 & 0.116 & \texttt{truth\_a0p025\_x1\_seed0} (64 dirs) \\
\bottomrule
\end{tabular}
\caption{Native-box ($1\times$) truth $D_{\mathrm{KY}}$ against drag.
The $\alpha{=}0.025$ value is a single seed, so it carries the roughly
$\pm20\%$ single-run scatter of \S\ref{sec:truth3_revision}. The
$2\times$ and $3\times$ runs are still running.}
\label{tab:p8_truth_x1}
\end{table}

\textbf{Reading.} Halving the drag from 0.1 to 0.05 raised the native
dimension by 46\%. Halving it again, to 0.025, adds only about 8\%,
which is within single-seed scatter. At the native box scale the
dimension is therefore close to saturating in $\alpha$. This is
consistent with C8/C16(iii): the chaos lives at the forcing scale
($|k|\in[2,8)$), which drag barely touches once $\alpha\ll\nu n^2=0.16$.
What drag still controls is the beyond-box energy, which grows without
bound as $\alpha\to0$ (\S\ref{sec:alpha0_results}). The TD hypothesis
therefore makes a sharp prediction for P8. If transfer degrades at
$\alpha{=}0.025$, the cause is the \emph{size-dependent large-scale
sector}, not more local chaos, because the native dimension barely
changed. The decisive numbers are the $2\times$ and $3\times$ truths
(extensivity: does $D_{\mathrm{KY}}$ scale like $1\times$ times the area
ratio, as at $\alpha{=}0.05$?) and P8's model error at those sizes.

\subsection{P8 model results at native size (2026-10-05): CONFIRMED
onset of non-locality at $\alpha{=}0.025$ (3 seeds)}
\label{sec:p8_model_first}

Same recipe as P2/P9/P9b, trained on the $\alpha{=}0.025$
$16\times$-wide dataset. Training losses are excellent -- the best of
any model tonight (AE val reconstruction 0.00386 against P2's 0.0082
and P9's 0.0079; propagator val (12-step) 0.00923 against P2's 0.0253)
-- ruling out representation collapse as the explanation for what
follows.

\begin{center}
\small
\begin{tabular}{lcccc}
\toprule
$1\times$ (truth 25.23, $n_{\mathrm{pos}}{=}14$) & seed 0 & seed 1 & seed 2 & mean $\pm$ sd \\
\midrule
$D_{\mathrm{KY}}$ & 9.96 & 16.97 & 18.02 & $14.98\pm4.37$ \\
$n_{\mathrm{pos}}$ & 4 & 8 & 9 & \\
error & 60.5\% & 32.7\% & 28.6\% & 40.6\% (mean) \\
\bottomrule
\end{tabular}
\end{center}

\textbf{Confirmed across all three seeds, not a single-IC artifact.}
Every seed undershoots, by 29--61\%; mean error 40.6\%, by far the
largest native-size error of any model in this document (P2: 3.0\%;
P9: 9.7\%; P9b: 10.9\%). Unlike P9's two-sided, IC-dependent regime
switching (\S\ref{sec:p9_diag}: some seeds too energetic, some too
quiet), every $\alpha{=}0.025$ seed collapses in the \emph{same}
direction -- fewer positive exponents than the truth, degree varying
by IC ($n_{\mathrm{pos}}$ 4, 8, 9 against the truth's 14).

\textbf{Reading.} This matches the large-scale-energy finding already
in hand (\S\ref{sec:p8_largescale}): $\alpha{=}0.025$ puts 24.8\% of
the field's energy beyond the native box, against 15.9\% at
$\alpha{=}0.05$ -- the native-size training crop has no view of that
pathway, and this is consistent with the model settling into an
artificially damped, lower-dimensional native state at this drag. This
is the onset-of-non-locality signature P8 was built to detect
(\S\ref{sec:drag_sweep}), now with seed support behind it rather than
one IC. \textbf{Caveat, honestly stated:} this is still only the
\emph{native}-size result. P8's actual pre-registered question
(\S\ref{sec:status_plan}) is about the \emph{transfer} curve --
whether $2\times$/$3\times$ error grows with falling $\alpha$ -- and
those evaluations are still running. A collapsed native baseline
changes what ``transfer error'' even measures (error relative to an
already-wrong native state is a different question from error
relative to a correct one), so the $2\times$/$3\times$ numbers need to
be read against \emph{both} the truth and this native shortfall, not
truth alone. Queued and running overnight; full P8 write-up to follow
once they land.

\subsection{P8 transfer results and stop (2026-10-05)}
\label{sec:p8_transfer}

\begin{center}
\small
\begin{tabular}{lccccc}
\toprule
size & truth & seed 0 & seed 1 & seed 2 & mean error \\
\midrule
$1\times$ & 25.23 & 9.96 (60.5\%) & 16.97 (32.7\%) & 18.02 (28.6\%) & 40.6\% \\
$2\times$ & 51.51 & 48.28 (6.3\%) & 50.75 (1.5\%) & 30.36 (41.1\%) & 16.3\% \\
$3\times$ & 82.75 & 66.41 (19.7\%) & -- & -- & -- \\
\bottomrule
\end{tabular}
\end{center}
$n_{\mathrm{pos}}$: $1\times$ 4/8/9 (truth 14); $2\times$ 24/24/15;
$3\times$ 32. The $2\times$ and $3\times$ truths are single seeds. The
remaining $3\times$ seeds were stopped by the user as unlikely to change
the picture.

\textbf{Reading.} The P2 recipe at $\alpha{=}0.025$ fails, with the same
symptom family as P9 (\S\ref{sec:p9_diag}): excellent short-horizon
losses, but free-running dynamics that land in IC-dependent,
lower-dimensional states. Every native-size seed undershoots. At
$2\times$ two seeds are accurate and one collapses. A first reading
that ``only the 8-site native box is bad'' did not survive the third
$2\times$ seed. Against the pre-registered P8 rule (\S\ref{sec:drag_sweep})
this is \emph{not} the single-digit $3\times$ error that would have
shown locality holding at $\alpha{=}0.025$. It cannot yet be read as a
clean locality-breakdown curve either: the failure is largest at
native size, and it looks like the training-stability failure that P9b
fixed in P9 more than like transfer error growing with size. The
distinguishing experiment (P8b, the P9b fix plus the C19 autoencoder at
$\alpha{=}0.025$) is listed in \S\ref{sec:status_plan}.

\subsection{P8 wide data: large-scale energy versus drag (2026-10-04)}
\label{sec:p8_largescale}

The $\alpha{=}0.025$ wide dataset ($16\times$, 5 trajectories) is
complete. As planned in \S\ref{sec:alpha0_results}, we measured
directly how much energy sits beyond the native box, i.e.\ in
streamwise modes with wavelength $>2\pi$ ($m<16$ on the $32\pi$ domain).
This uses the training trajectories, every 10th snapshot, kinetic energy
$|\hat\omega|^2/k^2$:
\begin{center}
\small
\begin{tabular}{lccc}
\toprule
$\alpha$ & field std & energy share beyond box ($\lambda>2\pi$) & share at $\lambda\ge 10.7\pi$ ($m\le3$) \\
\midrule
0.1 & 1.53 & 7.7\% & 1.8\% \\
0.05 & 1.39 & 15.9\% & 2.6\% \\
0.025 & 1.19 & 24.8\% & 3.7\% \\
\bottomrule
\end{tabular}
\end{center}
\textbf{Reading.}
\begin{itemize}[leftmargin=*]
\item Each halving of the drag roughly adds 8--9 points to the energy
  beyond the native box, while the native $D_{\mathrm{KY}}$ barely moves
  between 0.05 and 0.025 (\S\ref{sec:p8_truths}). Drag controls the
  \emph{large-scale energy}, not the local chaos, which is exactly the
  separation the TD hypothesis relies on.
\item The very largest scales ($m\le3$, wavelengths of 5--16 boxes)
  still hold only 2--4\%. At $\alpha{=}0.025$ on this domain there is no
  domain-scale condensate yet, unlike $\alpha{=}0$ (86\% in the domain
  mode at $2\times$). The beyond-box energy sits at intermediate scales
  of 1--5 boxes.
\item \emph{Prediction for P8:} a quarter of the energy now lives at
  scales larger than one training box but well inside a 40-site
  crop window. Crop training should still see it locally. If P8's
  $3\times$ error stays near P2's, the extra large-scale energy is
  handled by the local model. If the error grows, the 1--5-box scales
  are where locality starts to fail.
\end{itemize}

\subsection{$\alpha{=}0$ pilot results (COMPLETED 2026-10-04)}
\label{sec:alpha0_results}

\begin{table}[h]
\centering
\small
\begin{tabular}{lcccccccc}
\toprule
 & \multicolumn{7}{c}{energy $E$ (domain-mode share $E_{m=1}/E$), 50-unit means ending at $t$} & \\
size & $t{=}100$ & 200 & 400 & 700 & 1000 & 1200 & 1500 & sat.\ time \\
\midrule
$1\times$ & 0.193 (0.50) & 0.206 (0.78) & 0.224 (0.64) & 0.220 (0.62) & 0.213 (0.77) & 0.221 (0.68) & 0.221 (0.68) & 400 \\
$2\times$ & 0.196 (0.42) & 0.207 (0.75) & 0.249 (0.85) & 0.262 (0.86) & 0.266 (0.86) & 0.267 (0.86) & 0.267 (0.87) & 646 \\
\bottomrule
\end{tabular}
\caption{The $\alpha{=}0$ pilot. The laminar state has $E{=}0.5$; the
initial laminar breakdown happens by $t\approx10$ at both sizes. The
coefficient of variation of $E$ over the last 300 units is 2.7\% at
$1\times$ and 0.47\% at $2\times$.}
\label{tab:alpha0}
\end{table}

\textbf{Gate verdict: PASSED, with a caveat.} Saturation (the
pre-registered 200-unit running-mean criterion) occurs at $t{=}400$ at
$1\times$ and $t{=}646$ at $2\times$. Both are well inside 1500 units
and below the $\nu k_{\min}^2$ estimate (about 900 at $3\times$, about
400 at $2\times$). Extrapolating, $3\times$ should saturate around
900--1000 units. $\alpha{=}0$ datasets and truths therefore need a
spinup of at least 700 units at $2\times$ and an estimated 1000 at
$3\times$. That is 7--10$\times$ the standard 100-unit spinup, so the
current $\alpha{>}0$ pipeline would sample a transient at $\alpha{=}0$.

\textbf{What the saturated state looks like: a domain-scale condensate.}
At $2\times$, 86\% of the energy sits in the domain-scale streamwise mode
($m{=}1$, wavelength $4\pi$, twice the native box). Total energy is
nearly frozen: a coefficient of variation of 0.47\%, against 2.7\% at
$1\times$. This is the inverse-cascade condensate the TD hypothesis
(\S\ref{sec:drag_sweep}) anticipated. Without drag, energy piles up at
the \emph{largest scale the domain allows}, so the dominant structure
\emph{changes with domain size}. That is precisely the non-local,
size-dependent feature that a local, weight-shared model trained on
windows cannot represent, and it goes beyond the C8/C9 picture at
$\alpha{=}0.05$. There, the truth's unstable subspace sat at the forcing
scale and the beyond-box modes carried little of the dimension.

\textbf{Consequences for the program.}
\begin{itemize}[leftmargin=*]
\item \emph{Answer to ``would the system work at $\alpha{=}0$ with
  enough data?''} Probably not for zero-shot domain transfer. The
  equilibrium mean state is set by the domain, so a model trained at one
  size meets a different large-scale flow at another. More training data
  at one size cannot fix this. It remains open whether the
  \emph{fluctuations} riding on the condensate stay local. That question
  is the natural $\alpha{=}0$ experiment: condition the local model on
  the domain-scale mode, which is the C9 global-coupling idea specialized
  to one known mode.
\item \emph{Chaos check needed.} The 0.47\% energy fluctuation at
  $2\times$ is small enough that the condensate could be nearly steady or
  only weakly chaotic. A short Benettin run at $2\times$ after a 700-unit
  spinup should precede any $\alpha{=}0$ dataset generation. If
  $D_{\mathrm{KY}}$ is small, $\alpha{=}0$ is not an interesting
  transfer target in this geometry.
\item \emph{P8 framing.} $\alpha{=}0.025$ sits between the drag-limited
  regime, where transfer works ($\alpha\ge0.05$), and this condensate end
  point. P8's $3\times$ error therefore tests whether the condensate's
  influence has already set in at $\alpha{=}0.025$. Its expected
  signature is a growing domain-mode share in the $\alpha{=}0.025$ wide
  data, which is cheap to check when that dataset exists.
\end{itemize}

\subsection{P6 results (COMPLETED 2026-10-04)}
\label{sec:p6_results}

\begin{table}[h]
\centering
\small
\begin{tabular}{lccccc}
\toprule
arm (narrow arch.) & free interior & $D_{\mathrm{KY}}$ at $3\times$ (seeds 0/1/2) & mean error & AE val & prop val ($k{=}12$) \\
\midrule
\texttt{p6\_i6} & 6 sites (0.75 box) & 8.5 / 7.6 / 7.6 & 89.5\% & 0.0090 & 0.0203 \\
\texttt{p6\_i8} & 8 sites (1 box) & 22.4 / 27.2 / 26.8 & 66.1\% & 0.0092 & 0.0190 \\
\texttt{p6\_i12} & 12 sites (1.5 box) & 0.0 / 0.0 / 0.0 & 100\% & 0.0092 & 0.0187 \\
\texttt{p6\_full} (upper bound) & whole periodic domain & 25.9 / 0.0 / 0.0 & 88.5\% & 0.0090 & 0.0184 \\
\midrule
P2 (standard arch., 16$\times$ crops) & 24 sites & 72.0 / 69.5 / 69.9 & 6.3\% & 0.0082 & 0.025 \\
\bottomrule
\end{tabular}
\caption{P6: same-domain window training at $\alpha{=}0.05$, deployed on
the full $3\times$ domain (truth 75.22). P6 uses the narrow
architecture (encoder radius $1\times2$ layers, propagator window 1).}
\label{tab:p6}
\end{table}

\textbf{Pre-registered verdict: informative failure, the narrow
architecture is the limit.} The rule's ``informative failure'' branch
applies: even the upper-bound arm, trained on the \emph{entire} periodic
domain with no crops, masks or halos, misses by more than 15\%. Two of its
three seeds collapse to a fixed point. The failure is therefore
\emph{not} caused by windowing, and P6 cannot measure a locality length
with this architecture. The window sweep is non-monotone (about 8, about
25, then 0), which is what an unstable architecture produces, not a
smooth dependence on window size.

\textbf{What differs from the configurations that work.}
\begin{itemize}[leftmargin=*]
\item \emph{P2}, the standard architecture trained on crops of the
  $16\times$ domain, reaches 6.3\% at $3\times$.
\item \emph{D1}, with a window-1 propagator but the standard encoder,
  transferred well at $\alpha{=}0.1$.
\end{itemize}
The common factor in the failures is the \emph{encoder's} receptive
field: 2 sites ($\approx1.6$ length units) in P6, against 6 sites
($\approx4.7$) in every working model.

\textbf{Connection to C16(iii): a quantitative design rule.} C16(iii)
found that about 31.5 of the 36 unstable dimensions at $3\times$ sit at
$|k|\in[2,8)$, i.e.\ wavelengths of about 0.8--3.1. A 2-site encoder
reach (about 1.6, roughly one forcing wavelength $\pi/2$) cannot encode
the local structure at the longer end of that band into a single site's
latent, so the dimension-carrying degrees of freedom are not represented
and the dynamics degenerate. The 6-site encoder (about 4.7) spans the
whole band. This suggests a rule to test directly: \emph{the encoder's
receptive field must cover the wavelength band that carries the unstable
dimensions}, measurable beforehand from a truth run's unstable
subspace. It complements the light-cone bound on the propagator's reach
(Phase 2, H2).

\textbf{Connection to the training-loss lesson.} In-domain validation
losses are nearly identical across all four arms (autoencoder 0.009,
propagator 0.018--0.020) and comparable to P2's. Yet the long-time
dynamics range from a fixed point ($D_{\mathrm{KY}}{=}0$) to about 27. As
with C4, C9(ii), D2 and D3, short-horizon, teacher-forced losses do not
certify the long-time dynamics. That is mechanism (b) of
\S\ref{sec:c8_diag}. Every claim in this document rests on Lyapunov-level
evaluation, and this result is a direct demonstration of why that is
necessary.

\textbf{What this means for the practical (CM1-like) claim.} P6 does not
refute it. Training on windows and deploying on the domain still works
with the standard architecture: P2's crops come from the $16\times$
domain and its $3\times$ deployment reaches 6.3\%. What P6 shows is that
the window must be large enough to contain the standard architecture's
12-site margin \emph{plus} a useful interior, so the minimum window is
set by the encoder reach needed for the dimension-carrying scales.
\emph{Follow-up (P6b, not yet queued):} repeat P6 with the standard
architecture on a domain large enough for it (e.g.\ $6\times$, 48 sites:
48-site crops with a 24-site interior, plus a periodic-full arm). Its
truth would need a multi-seed $6\times$ run ($D_{\mathrm{KY}}\approx150$,
about 250 directions, several hours per seed). A cheaper alternative is
an encoder-reach sweep at fixed window, $\{1\times2, 2\times2,
2\times3\}$, which tests the design rule above directly.

\subsection{P7: crop training on 2D Kuramoto--Sivashinsky}
\label{sec:p7}

\textbf{What exists for 2D KS.}
\begin{itemize}[leftmargin=*]
\item \emph{Training:} only a box-trained model, Section 230, on the
  $44\times44$ periodic box.
\item \emph{Native size:} model $D_{\mathrm{KY}}$ 71.9 against an
  \emph{estimated} truth of about 85.6, i.e.\ about 16\% low. That truth
  was area-scaled from the $66^2$ run and never computed directly at
  $44^2$.
\item \emph{Transfer to $66^2$ ($1.5\times$):} the original 77\% was
  mostly the IC-normalization bug. Corrected, it is 23\% from one seed
  (\S\ref{sec:icnorm}).
\item \emph{Physics:} 2D KS is extensive (true $D_{\mathrm{KY}}$/area is
  constant across $44^2$ and $66^2$, G1), which is the regime where local
  models should transfer.
\end{itemize}
For 2D KS the dominant error may therefore be \emph{native-size
accuracy}, not transfer: a 16\% native undershoot plus modest transfer
degradation would explain the 23\%.

\textbf{Step 1 (IN PROGRESS, cheap): is the error native or transfer?}
\begin{enumerate}[label=(\alph*),leftmargin=*]
\item The two deferred normalized-IC transfer seeds at $66^2$:
  \textbf{COMPLETED}, 145.3 and 145.5 (24.5\%, 24.4\%). With seed 0
  (147.4), the transfer error is a tight $24\pm0.5$\%.
\item A direct true-PDE $D_{\mathrm{KY}}$ at the native $44^2$ box
  (150 directions; \texttt{exp\_g1} at scale 1): \textbf{COMPLETED},
  $D_{\mathrm{KY}}=81.77$ with $n_{\mathrm{pos}}=43$ and $\lambda_1=0.165$
  (the area-scaled estimate was 85.6).
\end{enumerate}

\textbf{Step 1 result (COMPLETED 2026-10-03).}
\begin{itemize}[leftmargin=*]
\item \emph{Native, normalized IC:} 71.2, 70.4 and 70.7 (seeds 0--2), a
  13.5\% error against the direct truth.
\item \emph{Transfer to $66^2$:} 24.4\% (above).
\item \emph{Density per unit area:} the truth is close to extensive,
  rising 4.6\% from $44^2$ to $66^2$ (0.0422 $\to$ 0.0442). The model's
  density falls 8\% (0.0366 $\to$ 0.0335).
\end{itemize}
\emph{Pre-registered verdict:} native error is 13.5 of 24.4 points, about
55\%, below the two-thirds threshold. The target is therefore
\emph{transfer itself}, on top of a substantial native undershoot. Step 2
remains well motivated: window training removes the 8-site wraparound in
both directions, which bears on native accuracy, and supplies
large-domain statistics, which bears on transfer.
\emph{Rule:} if the native error, measured against the direct truth,
accounts for at least two-thirds of the $66^2$ transfer error, P7's main
target is native accuracy. Large-domain training is then expected to
help because it removes the $8$-site wraparound in \emph{both}
directions. Otherwise the target is transfer itself.

\textbf{Step 2: COMPLETED, SUCCESS (\S\ref{sec:p7_results})} (after P6, before
P8; \texttt{scripts/run\_p7\_step2\_queue.sh}).
\begin{itemize}[leftmargin=*]
\item \emph{Code.} \texttt{gather\_xy\_crops}. \texttt{train\_crop\_ae},
  \texttt{encode\_crop\_bank} and \texttt{forced\_halo\_rollout} now take a
  $y$-crop and $y$-halo; with neither set they reproduce the $x$-only
  behaviour exactly. The driver gains \texttt{--crop-sites-y}, and
  \texttt{scripts/exp\_crop\_eval\_ks2d.py} evaluates on square periodic
  boxes. \texttt{scripts/gen\_ks2d\_wide\_dataset.py} generates the wide
  data and is preemptable.
\item \emph{Tests.} 5 new tests (20 crop tests in total; 385 related
  tests pass):
  \begin{itemize}[leftmargin=*]
  \item 2D gather wraps both axes;
  \item the 2D forced-halo free interior equals, to $10^{-5}$, the same
    weights on a larger 2D periodic domain;
  \item with an $x$-halo only, the sites within reach of the crop's
    $y$-edges are corrupted while deeper sites remain exact, which shows
    the $y$-halo is necessary;
  \item a 2D crop-autoencoder and crop-bank smoke test.
  \end{itemize}
  An end-to-end smoke run (training, then resize to a deployed square
  box) passes.
\item \emph{Configuration.} The Section 230 architecture is unchanged
  (encoder radius $2\times3$, window-3 propagator, $c{=}3$), so P7
  differs from the box-trained model only in its training data. The
  wide domain is $6\times44$ per axis ($768^2$ grid, $48\times48$ sites,
  4 training and 1 validation trajectory of 25 time units, spinup 200;
  generation about 16 minutes per trajectory). Crops are $40\times40$
  sites, which leaves a 12-site margin per side and a $16\times16$-site
  free interior.
\item \emph{Evaluation.} At $44^2$ (truth 81.77) and $66^2$ (truth
  192.5), 3 seeds each.
\end{itemize}
\textbf{Original step-2 plan (for reference):}
\begin{itemize}[leftmargin=*]
\item \emph{Isotropy changes the crop geometry.} Kolmogorov flow allowed
  $x$-only crops with full periodic $y$-strips because its forcing
  depends on $y$. 2D KS is isotropic, so $x$-only crops would keep the
  $y$-direction wraparound problem: an encoder receptive field of 6 sites
  on an 8-site ring. P7 therefore needs \emph{2D crops}, with loss
  masking and forced halos in both directions.
\item \emph{Implementation.} \texttt{gather\_x\_crops},
  \texttt{train\_crop\_ae} and \texttt{forced\_halo\_rollout} are
  generalized to 2D windows. The forced-halo exactness test is extended
  to 2D, comparing against the same weights on a larger 2D periodic
  domain.
\item \emph{Data.} A wide square 2D KS domain, e.g.\ $4\times44$ by
  $4\times44$ (a $512^2$ grid, $32\times32$ sites), with normalization
  from its training split.
\item \emph{Protocol.} As P2: crop autoencoder, then a forced-halo
  propagator with the Section 230 architecture.
\item \emph{Evaluation.} At $44^2$ (against the step-1 direct truth) and
  $66^2$ (truth 192.5), 3 seeds each. Evaluation at $88^2$ needs a new
  truth with $D_{\mathrm{KY}}\approx340$, i.e.\ at least 500 directions
  and many CPU hours, so it is optional.
\end{itemize}
\emph{Rules:}
\begin{itemize}[leftmargin=*]
\item \emph{success:} native and $66^2$ errors both single-digit;
\item \emph{partial:} both below the box-trained model's (16\% native,
  23\% at $66^2$);
\item \emph{negative:} no better than the box-trained model, which would
  mean the 2D KS error lies in the local rule, not in the training
  distribution.
\end{itemize}

\subsubsection{\texttt{p7b\_ks2d} (canonical recipe) results (COMPLETED
2026-10-05): the $88^2$ sub-extensivity was the training recipe, not
the system}
\label{sec:p7b_results}

Consolidation Phase 2 (\S\ref{sec:consolidation}): 2D KS retrained with
the canonical recipe (C19 per-patch $c{=}6$ autoencoder, P9b long free
rollouts, $k_{\max}{=}32$). Autoencoder interior validation error 0.0118
(original P7: 0.0089, about 30\% worse); propagator validation error at
full rollout 0.0196 (original: 0.0221 at only 12 steps, despite scoring
over $32$). 2 IC seeds, $44^2$/$66^2$/$88^2$:
\begin{center}
\small
\begin{tabular}{lccccccc}
\toprule
 & \multicolumn{2}{c}{$44^2$ (truth 81.77)} & \multicolumn{2}{c}{$66^2$ (truth 192.5)} & \multicolumn{2}{c}{$88^2$ (no truth)} \\
 & seed 0 & seed 1 & seed 0 & seed 1 & seed 0 & seed 1 \\
\midrule
$D_{\mathrm{KY}}$ & 83.19 & 84.03 & 187.97 & 187.74 & 331.98 & 334.57 \\
$n_{\mathrm{pos}}$ & 42 & 43 & 97 & 95 & 170 & 168 \\
error & 1.7\% & 2.8\% & 2.4\% & 2.5\% & -- & -- \\
density ($D_{\mathrm{KY}}/\mathrm{area}$) & 0.0430 & 0.0434 & 0.0432 & 0.0432 & 0.0429 & 0.0433 \\
\bottomrule
\end{tabular}
\end{center}

\textbf{Against the original P7} (\S\ref{sec:p7_results},
\S\ref{sec:p7_88_check}): native and $66^2$ errors both improve
(1.2/5.7\% $\to$ 1.7(2.8)/2.4(2.5)\%), and \textbf{the $88^2$ result is
transformed}: 215.86 (density 0.0279, 64\% of the extensive prediction,
flagged as possibly a spurious single-IC regime) becomes 331.98/334.57
(density 0.043, matching the $44^2$/$66^2$ density to within 1\% and
landing inside the pre-registered 327--345 bracket) -- \textbf{from two
independent IC seeds that agree with each other to within 1\%.}

\textbf{This resolves the open question from \S\ref{sec:p7_88_check}.}
That section listed three possibilities for the 215.86 reading: genuine
sub-extensive scaling, a rare real branch, or a model-specific spurious
regime with no physical counterpart -- explicitly the same mechanism
category as P9's failure. Two IC seeds of a model that differs from the
original \emph{only} in autoencoder design and rollout length, now
agreeing tightly at the extensive prediction, make (c) -- a training
artefact -- by far the best-supported explanation. The original P7 was
not probing a real sub-extensive ceiling; it was exhibiting the same
free-rollout instability P9 had, just manifesting as an $88^2$-specific
energy deficit rather than $3\times$-specific excess.

\textbf{What is and is not yet shown.} Still no direct $88^2$ truth
(\S\ref{sec:p7_88_check}'s point (b), ``a real but rare branch,'' is not
fully excluded without one) -- but two independent seeds agreeing to
within 1\% on both $D_{\mathrm{KY}}$ and density is strong evidence
against instability being the source of the earlier number, since an
unstable model would not be expected to reproduce itself this tightly
across ICs (contrast P9's $83\pm22$ at a single size from three seeds of
the \emph{old}-recipe model). A direct truth remains the only way to
fully close this and is recommended before this becomes the paper's
flagship transfer number; see \S\ref{sec:p7_88_check} for the proposed
cheap calibration probe before committing to the full run.

\textbf{Connection to the $1\times$/native-size discussion}
(\S\ref{sec:phase0_results}): this is the first direct test, on a
different system, of whether the canonical recipe's long rollouts
improve behaviour beyond fixing P9's specific $3\times$ instability.
Native-size 2D KS error did improve here (1.2\%$\to$1.7/2.8\%, not
better, but $66^2$ and especially the $88^2$ scaling improved sharply),
which is at least consistent with -- though does not on its own prove --
the hypothesis that long rollouts help a propagator cope with
deployment regimes absent from training (here, larger domains; for
Kolmogorov, \S\ref{sec:phase0_results}'s concern was smaller domains).
\texttt{k\_a005} (Kolmogorov, same canonical recipe, running now) is the
direct test on the native-size question itself.

\subsubsection{P7 step 2 results (COMPLETED 2026-10-04): SUCCESS}
\label{sec:p7_results}

Run notes. The first launch crashed because the queue opened the
dataset file while it was still being written (an HDF5 lock). The
queue now waits for the writer process to exit. At the user's request
the evaluation used seed 0 only, at background priority. Training took
8 minutes (autoencoder) and 19 minutes (propagator) on MPS. Crop
validation was 0.0089 (autoencoder interior) and 0.022 (propagator at
$k{=}12$).

\begin{table}[h]
\centering
\small
\begin{tabular}{lcccccc}
\toprule
 & \multicolumn{3}{c}{$44^2$ (native, truth 81.77, $n_{\mathrm{pos}}$ 43)} & \multicolumn{3}{c}{$66^2$ (truth 192.5)} \\
model & $D_{\mathrm{KY}}$ & $n_{\mathrm{pos}}$ & error & $D_{\mathrm{KY}}$ & $n_{\mathrm{pos}}$ & error \\
\midrule
box-trained, step 1 (3 seeds) & 70.4--71.2 & 34--36 & 13.5\% & 145.3--145.5 & 69 & 24.4\% \\
\textbf{P7 crop-trained} (seed 0) & \textbf{82.78} & 41 & \textbf{1.2\%} & \textbf{181.61} & 89 & \textbf{5.7\%} \\
\bottomrule
\end{tabular}
\caption{P7: 2D KS, trained on $40\times40$-site 2D crops of a
$6\times$-per-axis periodic domain, deployed on square periodic boxes.
The architecture is identical to the box-trained model (Section 230);
only the training data differ. $\lambda_1$ = 0.237 at $44^2$ and 0.246
at $66^2$.}
\label{tab:p7_results}
\end{table}

\textbf{Pre-registered verdict: SUCCESS.} Native and $66^2$ errors are
both single-digit. This is one seed; the step-1 seed spread was
$\pm0.5$ points, so the effect is about 20 times larger than the seed
noise.
\begin{itemize}[leftmargin=*]
\item native error falls 11-fold, from 13.5\% to 1.2\%;
\item transfer error falls more than 4-fold, from 24.4\% to 5.7\%.
\end{itemize}

\textbf{Connections to the document's hypotheses.}
\begin{enumerate}[leftmargin=*]
\item \emph{The training distribution decides transfer, now in a
  genuinely 2D, isotropic system.} The same architecture, with no other
  change, goes from 24\% to 6\% at $66^2$ and from 13.5\% to 1.2\% at
  native size. This is the 2D counterpart of the Kolmogorov
  box-versus-window result (P1/P2). There, though, crops were strips
  that kept a full periodic $y$-direction. Here they are cropped in both
  directions, so \emph{no training sample contains a whole periodic
  direction}. That is the strongest version of the patch-training
  paradigm tested so far.
\item \emph{The native undershoot was a training-distribution artefact,
  not missing capacity.} Step 1 attributed 13.5 of the 24.4 points to
  native error and raised the capacity question of C0. P7 removes it
  with the same $c{=}3$ and receptive field. The box-trained model's
  deficit came from training on a box whose periodicity the encoder
  could see: its 6-site receptive field covered most of an 8-site ring
  in each direction (\S\ref{sec:p7}). Capacity per site was not the
  problem for 2D KS, which reframes C0's capacity concern for this
  system.
\item \emph{The residual transfer error is mostly the truth's
  finite-size super-extensivity.} The truth grows by $192.5/81.77 =
  2.35$ from $44^2$ to $66^2$, against an area ratio of 2.25. The model
  grows by $181.6/82.8 = 2.19$, slightly \emph{below} strict
  extensivity. A local, weight-shared model with a fixed per-site rule
  is extensive by construction, so it cannot reproduce the extra
  $\approx$4\% density that the true system shows at this size.
  Kolmogorov showed the same pattern at $\alpha{=}0.05$ (truth density
  rising with size, model near-extensive). If the truth's density
  saturates at larger boxes, as extensivity theory predicts, the error
  should \emph{fall} with size. An $88^2$ evaluation would test this,
  but needs an expensive truth ($D_{\mathrm{KY}}\approx340$).
\item \emph{Locality (H1) and the light-cone bound.} The deployed model
  has a fixed physical reach: encoder 6 sites, propagator 6 per step.
  It reproduces a 192-dimensional attractor at 1.5$\times$ the training
  box to within 6\%, so 2D KS dynamics are local at that reach. This
  matches the 1D KS conclusion and the extensivity picture underlying
  the whole program.
\end{enumerate}

\textbf{Consequences for next steps} (\S\ref{sec:post_p7_options},
branch ``P7 transfers well''):
\begin{itemize}[leftmargin=*]
\item 2D patch training works when the training domain is large
  relative to the crop and the deployed boxes. That is the basis for
  P9, 2D patches for Kolmogorov. Its extra difficulties are the
  $y$-dependent forcing (handled by 2-site-aligned offsets) and possibly
  the large-scale sector at small $\alpha$ (P8, \S\ref{sec:alpha0_results}).
\item C19 (\S\ref{sec:c19}), now running, decides whether the 12-site
  autoencoder margin can shrink. If so, P9's crops and its training
  domain can both be much smaller, and the P6-style single-domain test
  becomes feasible.
\item Seeds 1 and 2 at both sizes, and an $88^2$ truth, would firm up
  the result for publication. They are deferred because of machine load.
\item \emph{QUEUED (2026-10-04): model-only $88^2$ evaluation} (seed 0;
  512 of 768 directions; \texttt{scripts/run\_p7\_x2\_model\_only.sh},
  after P9). It runs without a truth. Extensivity of the two existing
  truths (0.042 and 0.044 per unit area) brackets the expectation at
  $D_{\mathrm{KY}}\approx330$--345.
  \emph{Rule:} if the model is stable and inside or near that bracket,
  extensive transfer to $2\times$ per side is supported. If the model
  sits at $\approx 4.0\times$ its native value ($\approx 330$) while the
  truth density keeps rising, the finite-size reading of the 5.7\% is
  supported but not proven; that needs the full $88^2$ truth, deferred
  for cost. A collapse or a value well below 300 would mean the
  1.5$\times$ success does not extend to $2\times$.
\end{itemize}

\subsubsection{$88^2$ model-only check (2026-10-04): density BELOW extensivity, unconfirmed (1 seed)}
\label{sec:p7_88_check}

Queued follow-up (\S\ref{sec:p7_results}) now returned: seed 0, 512
directions, no truth run (see rule there). Result:
$D_{\mathrm{KY}}=215.86$ ($n_{\mathrm{pos}}=101$ of 768,
$\lambda_1=0.178$), wall time 79 min.

\textbf{Against the pre-registered bracket: below it, not within it.}
The two real truths give a density per unit area of $81.77/44^2=0.0422$
and $192.5/66^2=0.0442$; extrapolating to $88^2$ predicts
$D_{\mathrm{KY}}\approx327$--345. The model gives 215.86, density
0.0279 -- about 64\% of the predicted density, not ``stable and near
the bracket.'' This is the \emph{opposite} direction from P9's failure
mode (\S\ref{sec:p9_diag}): there the model's chaos grew faster than
area; here it may be growing slower, i.e.\ saturating rather than
extensive.

\textbf{UPDATE (2026-10-05): resolved, see \S\ref{sec:p7b_results}.} The
canonical-recipe retrain (\texttt{p7b\_ks2d}) gives 331.98/334.57 from
two IC seeds, matching the $44^2$/$66^2$ density and the predicted
bracket -- consistent with (c) below (a training-recipe artefact in the
original P7), not a real sub-extensive ceiling. Original reasoning
follows for the record.

\textbf{Not yet a verdict -- one seed only, and this could be the exact
P9 mechanism, not just noise.} P9 showed a crop-trained propagator can
have \emph{multiple free-running regimes at one fixed size}, selected
by the IC alone -- not because the true system is multistable (checked
directly and it is not, \S\ref{sec:p9_diag}), but because short forced
rollouts in training left the model under-constrained once deployed
free-running at a size larger than any training crop. The symptom was a
shifted energy spectrum (too much or too little large-scale energy
depending on branch), and it moved $D_{\mathrm{KY}}$ by a factor of
$\approx1.7$ (65 to 107) from IC alone, at one fixed trained model. This
$88^2$ result carries every risk factor that made P9's first read
wrong, and none of the checks that would rule it out: one IC seed; no
rollout-statistics diagnostic run (so there is no way yet to see
whether this trajectory's energy spectrum looks like the
$44^2$/$66^2$ trajectories' or has drifted); and $88^2$ is the largest
relative size jump tested ($4\times$ the native area, against
$66^2$'s $2.25\times$), giving a crop-trained model's free dynamics the
most room to leave its training-distribution support. 215.86 could
therefore be (a) genuine sub-extensive scaling, (b) a real but rare
low-dimensional branch if the true system has some multistability at
this size, or (c) a model-specific spurious regime with no physical
counterpart, the same category as P9's failure. There is currently no
way to distinguish these three. \textbf{TO DO}: 2--3 more IC seeds at
$88^2$, and the same free-rollout energy-spectrum diagnostic used for
P9 (\texttt{exp\_rollout\_stats\_kolmogorov.py}'s
2D-KS analogue does not yet exist -- would need porting), to see
whether the model's rollout drifts to a less energetic state at this
size the way P9's drifted to a more energetic one.



\subsection{\texttt{k\_a005} (Kolmogorov, canonical recipe) results
(IN PROGRESS 2026-10-05)}
\label{sec:k_a005_results}

Kolmogorov $\alpha{=}0.05$ retrained with the canonical recipe (C19
per-patch $c{=}6$, 2D patch training on the $6\times6$-box domain, P9b
long rollouts). Direct test of the hypothesis raised in
\S\ref{sec:phase0_results}: does fixing the propagator's free-rollout
training fix the native-size ($1\times$) behaviour.

\textbf{Native size ($1\times$, seed 0): the hypothesis's strongest
result yet.} $D_{\mathrm{KY}}=22.65$, $n_{\mathrm{pos}}{=}12$, against
the per-IC truth 22.46, $n_{\mathrm{pos}}{=}12$ -- \textbf{0.9\% error
and an exact match on the number of positive exponents.} Against every
other model at this exact IC (\S\ref{sec:phase0_results}): P2 6.6\%,
C19 (per-patch $c{=}6$, short rollout) 4.7\%, C19 2nd train seed 7.6\%,
full-reach $c{=}6$ control 15.3\%. This is the first model to match
both the value and the exponent count, and it isolates long rollouts
(not channel count or reach) as the thing that fixed it, since C19
itself (same autoencoder, short rollout) still had 4.7\% error at this
IC.

\textbf{$2\times$ (seed 0): a sharp reversal.} $D_{\mathrm{KY}}=40.04$,
$n_{\mathrm{pos}}{=}19$, against the multi-seed truth mean $49.8\pm9.8$
-- \textbf{29.5\% error}, worse than P2 (9.4\%) or C19 (1.7\% at this
IC). So the native-size fix did not transfer cleanly to $2\times$; if
anything the model now undershoots badly at a size where the
short-rollout C19 model was accurate. This is the opposite pattern from
\texttt{p7b\_ks2d} (\S\ref{sec:p7b_results}), where long rollouts helped
at every size tested, largest size included.

\textbf{Reading, provisional (seed 1 and $3\times$ still running).} The
clean story from \texttt{p7b\_ks2d} -- ``long rollouts fix free-running
drift, uniformly'' -- does not hold for Kolmogorov. A plausible
reconciliation: Kolmogorov's forced, anisotropic dynamics may respond to
longer forced rollouts differently at different sizes, in a way 2D KS's
unforced, isotropic dynamics does not. Alternatively this is seed 0
being unrepresentative at $2\times$ the way it was at $3\times$ for the
original P9 (\S\ref{sec:p9_diag}) -- seed 1's $2\times$ result is the
direct check. $3\times$ (seed 0) is now in too: $D_{\mathrm{KY}}=63.89$,
$n_{\mathrm{pos}}{=}31$, truth 75.22 -- \textbf{15.1\% error}. So for
seed 0 across all three sizes: $1\times$ excellent (0.9\%), $2\times$
poor (29.5\%), $3\times$ poor but milder (15.1\%) -- a non-monotonic
pattern, worst at $2\times$, not a steady trend with size. \textbf{Not
yet a verdict.} Remaining: x1/x2/x3 seed 1, and the rollout-statistics
diagnostic (queued, \texttt{scripts/run\_consolidation\_p1p2.sh}), which
will show directly whether the $2\times$/$3\times$ free rollouts have
drifted in energy content the way P9's did. If seed 1 reproduces the
same non-monotonic shape, it is a property of this recipe on
Kolmogorov, not seed noise.

\subsection{If it works: implications}

\textbf{Data assimilation.} The learned propagator's Jacobian is banded by
construction, and $h^*$ from P3 is a measured localization radius. That
is exactly the quantity the localization assertion of brief \S15 needs,
and here it comes from data rather than tuning. One model then serves
regional and global DA at any domain size, and the required ensemble size
should track the local dimension per site, not the global attractor
dimension (H4). Variant B is literally the limited-area training setup,
so a model trained on regional windows can be used inside a regional
filter with lateral boundaries from a parent analysis.

\textbf{Meteorology.} Training cost scales with crop size, not domain
size, and training data can be any collection of regional windows of
reanalysis or high-resolution simulation. Real data are not homogeneous,
though. Latitude-dependent Coriolis, orography and land--sea contrast
mean the local rule must be conditioned on local static fields, as the
anisotropic phase channels did for $f(y)$. Teleconnections are the
atmospheric version of the Biot--Savart far field. P3 is the tool that
says, regime by regime, whether a crop-local model can see enough.

\textbf{Chaotic systems.} Crop-trainability is an operational test of
extensivity: a system can be learned from crops only if its local
statistics do not depend on the domain beyond $h^*$. $h^*(\alpha)$ is
measured from rollout error alone, with no Lyapunov truth runs. It is
therefore a cheaper companion to \S\ref{sec:transfer_diagnostic}'s
transfer-error diagnostic, and divergence of $h^*$ would flag the onset of
large-scale organization (the self-aggregation application).

\textbf{Machine learning.} This is the physics version of training
fully-convolutional networks on image crops and applying them to whole
images. The light-cone argument gives an exactness guarantee (variant A)
that vision crop training lacks. One large simulation provides many
nearly independent samples through spatial ergodicity, so data efficiency
improves. Small receptive fields also make the learned rule easier to
distill into a stencil or sparse PDE (the earlier local-closure
sections). The broader lesson from D1--D3 is that any coupling mechanism
added for transfer must be trained on the input distribution it will see
at deployment. Crops satisfy that by construction.

\subsection{Prior art: what is not new}

Training local predictors with an overlap halo on one large-domain
simulation is not new. Parallel local reservoir computers did this for
large-$L$ KS \cite{pathak2018}. Patch-decomposed autoencoder/NODE models
do related things on KS and Kolmogorov flow \cite{constanteamores2024},
and translation-equivariant operators deployed zero-shot on larger
domains are recent \cite{zeroshotlargerdomains2026}. Limited-area ML
weather models are also an active area; we have not yet surveyed them.
Our candidate contributions are narrower:
\begin{enumerate}[leftmargin=*]
\item a learned \emph{latent} local field trained on crops with
  forced-halo rollouts;
\item Lyapunov-spectrum validation of the \emph{assembled} model across
  domain sizes;
\item $h^*(\alpha)$ as a locality measurement in drag-damped 2D
  turbulence;
\item use of the resulting banded latent in ensemble DA.
\end{enumerate}
Before any novelty claim, a targeted literature search like the one in
\S\ref{sec:td_literature} is required, with a focus on patch-trained PDE
surrogates and ML limited-area models.

%======================================================================
\section{Does latent capacity keep up with domain size in 2D?}
\label{sec:capacity}
%======================================================================

\subsection{The concern, and a correction to its premise}

\textbf{The concern (2026-10-02).} In 1D KS the ratio of latent size to
full-state size stays fixed under L-transfer. The worry is that in 2D,
with Section-231-style models, the latent grows more slowly than the
full state. If so, information theory says a growing domain must
eventually exceed the latent's capacity. A proposed test was to grow only
one dimension.

\textbf{The arithmetic says otherwise for every model here.} Patch size is
held fixed under transfer (in both $x$ and $y$), and so is the number of
channels per site. Then $d_{\mathrm{lat}} = c\,N_{\mathrm{sites}}$ and the
full state $N_{\mathrm{full}} = N_{\mathrm{sites}}\times(\text{patch
area})$. Both grow in proportion to the domain area, in 1D and in 2D,
whether one dimension grows or both. Table~\ref{tab:capacity} checks this
row by row.

\begin{table}[h]
\centering
\small
\begin{tabular}{lcccccc}
\toprule
System, size & $N_{\mathrm{full}}$ & $d_{\mathrm{lat}}$ &
$d_{\mathrm{lat}}/N_{\mathrm{full}}$ & true $D_{\mathrm{KY}}$ &
$d_{\mathrm{lat}}/D_{\mathrm{KY}}$ & transfer \\
\midrule
1D KS, $L{=}100$ & 256 & 48 & 0.188 & 22.1 & 2.17 & (native) \\
1D KS, $L{=}200$ & 512 & 96 & 0.188 & 44.6 & 2.15 & $<2\%$ \\
1D KS, $L{=}400$ & 1024 & 192 & 0.188 & 89.1 & 2.15 & $<2\%$ \\
\midrule
2D KS, $44^2$ & 16384 & 192 & 0.0117 & $\approx85.6^\dagger$ & 2.24 & (native) \\
2D KS, $66^2$ & 36864 & 432 & 0.0117 & 192.5 & 2.24 & 77\%$^\ddagger$ \\
\midrule
Kolm.\ $\alpha{=}0.1$, $1\times$ & 16384 & 192 & 0.0117 & 15.9 & 12.1 & (native) \\
Kolm.\ $\alpha{=}0.1$, $2\times$ & 32768 & 384 & 0.0117 & 30.3 & 12.7 & 4.5\% \\
Kolm.\ $\alpha{=}0.1$, $3\times$ & 49152 & 576 & 0.0117 & 56.4 & 10.2 & 19.4\% \\
Kolm.\ $\alpha{=}0.05$, $1\times$ & 16384 & 192 & 0.0117 & 23.3 & 8.25 & (native) \\
Kolm.\ $\alpha{=}0.05$, $2\times$ & 32768 & 384 & 0.0117 & 56.8 & 6.76 & 50\% \\
\bottomrule
\end{tabular}
\caption{Latent capacity across every transfer in this document.
$^\dagger$Area-scaled from the $66^2$ truth, using the constant
$D_{\mathrm{KY}}/\mathrm{Area}$ measured in G1. $^\ddagger$Raw (not normalized)
initial condition; 23\% with a normalized one (\S\ref{sec:icnorm}). 1D values are the
transferred model's, which matched truth within 2\% (E1). Every
Kolmogorov transfer grew $x$ only; 2D KS grew both directions.}
\label{tab:capacity}
\end{table}

Four things follow from Table~\ref{tab:capacity}.
\begin{enumerate}[leftmargin=*]
\item \emph{The latent-to-state ratio is constant in 2D too}, at 0.0117
  in every 2D row. 2D compresses the raw state about $16\times$ harder
  than 1D. That compression ratio is fixed by patch area and does not
  change with $L$.
\item \emph{The proposed one-dimension test has in effect already been
  run.} Every Kolmogorov transfer grew $x$ only, and it still failed at
  $\alpha{=}0.05$. Because the ratio is constant whichever directions
  grow, a one-axis test cannot isolate the ratio anyway. What it can
  isolate is geometry (\S\ref{sec:capacity_experiments}, C3).
\item \emph{Raw capacity does not separate success from failure.} The
  ratio that matters is latent size over attractor dimension,
  $d_{\mathrm{lat}}/D_{\mathrm{KY}}$:
  \begin{itemize}[leftmargin=*]
  \item about 2.2 for 1D KS, which transferred cleanly;
  \item about 2.2 for 2D KS, which failed;
  \item 7--13 for Kolmogorov flow, which failed at $\alpha{=}0.05$
    despite the most headroom of all.
  \end{itemize}
  A generic embedding bound gives a useful reference line: a smooth
  generic map into $\mathbb{R}^m$ is injective on a set of dimension $d$
  when $m>2d$ \cite{embedology1991}. Both KS models sit just above it
  ($48 > 44.3$; $192 > 171$). Kolmogorov sits 3--6$\times$ above it.
\item \emph{The failure signature is collapse, not saturation.} A
  capacity-limited model should reproduce $D_{\mathrm{KY}}$ up to a
  ceiling near $d_{\mathrm{lat}}$ and then plateau. Instead the failing
  models land far below both truth and capacity: 43.6 of a possible 432
  (2D KS), 28.4 of 384 (Section 231), and 1--3 of 384 (D2, D3). The model
  is losing dynamics, not running out of room for them.
\end{enumerate}

\subsection{The right form of the concern: information density, not state size}

The intuition is still right in principle once it is stated in the right
quantities. A forecasting latent must carry the attractor's information.
Roughly, it needs $d_{\mathrm{lat}}\gtrsim D_{\mathrm{KY}}$ dimensions
(with $2D_{\mathrm{KY}}{+}1$ generically sufficient), plus enough
precision to absorb new information at the KS-entropy rate $h_{\mathrm{KS}}$
(the sum of the positive exponents) over the forecast horizon. Define the
information density per site
$\rho_{\mathrm{site}} = D_{\mathrm{KY}}/N_{\mathrm{sites}}$.

\textbf{Extensive regime: capacity never runs out.} If chaos is
extensive, then $D_{\mathrm{KY}}\propto$ area and
$h_{\mathrm{KS}}\propto$ area, so $\rho_{\mathrm{site}}$ is independent
of $L$. A local latent with fixed $c$ grows at exactly the rate the
information does. A fixed per-site budget $c\ge\kappa\rho_{\mathrm{site}}$
($\kappa\approx2$ for the embedding margin) then holds for every $L$. This
is the central reason local latents are the right architecture for
extensive chaos.

\textbf{Non-extensive regime: capacity eventually runs out.} If
$D_{\mathrm{KY}}$ grows faster than area, $\rho_{\mathrm{site}}(L)$ rises
and fixed $c$ is exhausted at some $L_{\mathrm{cap}}$ where
$\kappa\rho_{\mathrm{site}}(L_{\mathrm{cap}}) = c$. Measured values of
$\rho_{\mathrm{site}}$:
\begin{itemize}[leftmargin=*]
\item 1D KS and 2D KS: about 1.35, against $c{=}3$. Both are tight, and
  1D still worked.
\item Kolmogorov $\alpha{=}0.1$: 0.25, 0.24, 0.29 at $1/2/3\times$.
\item Kolmogorov $\alpha{=}0.05$: 0.36, then 0.44 at $2\times$.
\end{itemize}
The non-extensive growth is real, but at these sizes it is far from the
$c{=}3$ ceiling. So capacity is a future limit, not the current one.

\textbf{Where the information lives matters as much as how much there is.}
In the non-extensive regime the extra degrees of freedom are
\emph{collective}. They are large-scale modes (the beginning of
condensate-like organization) spanning many sites. A local latent can
represent such a mode easily: every site's gauge channel holds the local
value of the large-scale field, which is a distributed and highly
redundant encoding. A second, subtler 1D/2D difference is also real: the
dimension of a 2D patch of side $\ell$ should go as
$d_{\mathrm{local}}(\ell)\approx\rho\ell^2+\beta\ell+\gamma$. The
perimeter term $\beta\ell$, the information shared with neighbours, has
no 1D analogue beyond a constant. For patches about one correlation
length wide, which is what we use, it may not be small. This changes the
correct per-site budget, not its scaling with $L$.

\subsection{An alternative hypothesis: consensus, not capacity}

A distributed representation of a collective mode only stays a single
mode if all the sites carrying it update consistently. In our models each
site updates from a neighbourhood of radius $R$ sites per step. Nothing
forces two distant sites' copies of a domain-scale mode to agree, except
information relayed across $\sim L/R$ steps. In the PDE that agreement is
enforced instantly by the pressure (Biot--Savart) solve.

So the copies drift apart under chaos. If training has taught the model
to damp disagreement, the mode is suppressed instead. Suppression predicts
exactly the collapse-toward-low-$D_{\mathrm{KY}}$ signature above, and it
worsens with $L$ because more copies must agree over longer distances. 1D
KS at these sizes has short correlations and no collective modes, so no
consensus is needed. This hypothesis fits Kolmogorov at low drag better
than 2D KS. For 2D KS, the strong initial-condition dependence of the
transferred spectrum (43.6 vs.\ $\sim$144) suggests multistability, which
could be a symptom of failed consensus but is not yet shown to be one.

The consensus hypothesis makes a different prediction from the capacity
hypothesis. More channels should not help, but a mechanism that forces
overlapping or redundant estimates to agree should help. That is the
discriminating experiment (C2 versus C4/C5 below).

\subsection{Ways to grow the latent with the information, and keep transfer}

One design rule carries over from D1--D3: added capacity must come from
weight sharing driven by physical size. At deployment there must be no
new untrained parameters, and every new slot must see an input
distribution it saw in training.
\begin{enumerate}[leftmargin=*]
\item \textbf{Budget $c$ from measurement.} Keep the current
  fixed-patch, fixed-$c$ design (latent $\propto$ area) and set
  $c\ge\kappa\,\rho_{\mathrm{site}}$ from the measured 2D local-dimension
  curve (C1), including the perimeter term. This is sufficient whenever
  the regime is extensive.
\item \textbf{Redundant overlapping measurements with fusion.} Encode with
  stride $s$ smaller than patch $p$, so each physical point is seen by
  $r=(p/s)^2$ sites in 2D. That is $r$ redundant measurements. Before
  each propagator step, apply a local fusion operator $\Phi$. Two
  choices:
  \begin{itemize}[leftmargin=*]
  \item \emph{Re-projection:} decode, overlap-add with a partition-of-unity
    window, and re-encode. This is the map $z\mapsto E(D(z))$ onto
    mutually consistent latents.
  \item \emph{Learned weighted consensus:} each site also predicts a
    variance, and overlapping estimates are combined inverse-variance
    weighted. That is the same operation LETKF uses to combine
    overlapping local analyses.
  \end{itemize}
  Latent size grows by the constant factor $r$, so it stays $\propto$
  area and transfer is unaffected, since $\Phi$ is local and
  weight-shared. What redundancy buys is not asymptotic capacity. It buys
  \emph{error correction and consensus}: independent local errors average
  down, and copies of a shared collective mode are pulled back together
  every step. That is exactly what the consensus hypothesis says is
  missing. It also tests the earlier conjecture that overlapping patches
  would expose cross-patch interactions.
\item \textbf{Multiresolution latent with scale-shared weights.} Add
  coarse levels (pooling factor 2 per axis), each with $c$ channels per
  coarse site, and add levels at deployment as $\log_2(L/\ell_0)$ grows.
  Total size is $c\cdot N_{\mathrm{sites}}(1+\tfrac14+\tfrac1{16}+\dots) =
  \tfrac43\,c\,N_{\mathrm{sites}}$ in 2D, still $\propto$ area. That buys
  dedicated slots for collective modes at every scale for a 33\%
  overhead. Two conditions decide whether added levels transfer:
  \begin{itemize}[leftmargin=*]
  \item Share weights across levels as well as positions, so a new level
    adds no parameters. Scale-sharing is physically justified only inside
    a self-similar cascade range. It should fail at the forcing scale and
    at the arrest scale $\ell_\alpha$, and where it fails is itself a
    measurement.
  \item Train on crops of a large domain (\S\ref{sec:crop_training}), so
    coarse levels never see wraparound. D3 broke this rule.
  \end{itemize}
\item \textbf{Redundancy in time instead of space.} Give each site a
  delay embedding of its last $k$ latents (Takens). Capacity then grows
  per site without changing spatial structure, and it addresses
  Mori--Zwanzig memory directly. It is cheap, but it does not fix
  consensus.
\item \textbf{Adaptive capacity (speculative).} Use a set of
  permutation-equivariant ``slots'' that attach to coherent structures
  (object-centric attention), with the number of slots tied to area, i.e.\
  to the expected number of vortices $\sim A/\ell_\alpha^2$. Capacity then
  follows the structures rather than the grid. This is the most novel
  option and the least proven.
\end{enumerate}

\subsection{Experiments (C0--C6)}
\label{sec:capacity_experiments}

The rules below are pre-registered (written before any run). C0, C1 and
C4 need little or no training, so they run first. Order: C0, C1, C4, then
C2 and C3, then C5 and C6 on the crop-training infrastructure of
\S\ref{sec:crop_training}.
\begin{description}[leftmargin=*]
\item[C0 --- capacity audit, no training.] \textbf{COMPLETED, see \S\ref{sec:c0_results}.} For every existing checkpoint
  and transfer size, encode held-out data and compute two quantities:
  \begin{itemize}[leftmargin=*]
  \item the per-site effective rank, i.e.\ the participation ratio of the
    $c{\times}c$ per-site latent covariance pooled over sites and time;
  \item the global participation ratio over all $d_{\mathrm{lat}}$
    channels.
  \end{itemize}
  \emph{Capacity-bound prediction:} per-site rank approaches $c$ as $L$
  grows, coinciding with transfer failure. \emph{Rule:} if per-site rank
  is below $c/2$ at every failing transfer, capacity is ruled out as the
  primary cause.
\item[C1 --- 2D local dimension (brief D5 in 2D).] Run two-NN on
  independent 2D patches of side $\ell$ for 2D KS and for Kolmogorov at
  $\alpha\in\{0.05,0.1\}$. Fit $\rho\ell^2+\beta\ell+\gamma$, and compare
  the implied per-site budget, perimeter term included, with $c{=}3$.
\item[C2 --- capacity sweep.] Retrain the $\alpha{=}0.05$ model with
  $c\in\{2,3,6,12\}$ (crop-trained once P0 exists, otherwise the Section
  231 recipe). Measure $2\times$ transfer error against $c$.
  \emph{Capacity-bound:} error at least halves from $c{=}3$ to $c{=}12$.
  \emph{Not capacity-bound:} roughly flat.
\item[C3 --- geometry, the sharpened one-axis test.] Transfer 2D KS from
  $44^2$ to $88\times44$ ($x$ only, $2\times$ area) and to $66^2$ (both
  axes, $2.25\times$ area), with 3 IC seeds each. Truth comes from area
  scaling plus one direct truth run at $88\times44$. Because
  $d_{\mathrm{lat}}/N_{\mathrm{full}}$ is the same in both cases, this is
  a test of geometry, not of the ratio. Is growing both directions, which
  admits genuinely 2D large-scale modes, worse than growing one at matched
  area? The multistability seen in G1 is a confound, which is why 3 seeds
  are mandatory.
\item[C4 --- zero-training fusion, the cheap consensus test.] \textbf{COMPLETED, see \S\ref{sec:c4_partial}.} During
  rollout, apply re-projection $z\leftarrow E(D(z))$ every
  $m\in\{1,4,16\}$ steps. Use the existing Section 231 ($\alpha{=}0.05$)
  and 233 ($\alpha{=}0.1$) models at $1\times$ and $2\times$, and measure
  $D_{\mathrm{KY}}$ and transfer error. \emph{Control:} if re-projection
  changes native-size $D_{\mathrm{KY}}$ by more than 10\%, it is
  perturbing the dynamics rather than enforcing consistency, and the
  result is uninterpretable. \emph{Consensus-supporting:} $2\times$ error
  improves while native stays within 10\%.
\item[C5 --- redundancy versus matched capacity.] Train an overlapping encoder ($s{=}p/2$,
  $r{=}4$) with inverse-variance fusion against a non-overlapping control
  with $4c$ channels, so total latent size is equal. This is the key
  control: it separates redundancy and consensus from raw capacity.
  \emph{Rule:} redundancy helps beyond capacity if the overlapping model
  beats the control at $2\times$ by at least 10 percentage points.
\item[C6 --- multiresolution, scale-shared, crop-trained.] Levels are
  added at deployment. Use the P2 decision rule. Also record
  per-level reconstruction error against scale, to locate where scale
  sharing breaks relative to $\ell_\alpha$.
\item[C7 --- learned consensus, crop-trained, rank-protected.] Added
  2026-10-02 after C0 and the partial C4 result. Design and rules are in
  \S\ref{sec:c7}. C7 merges C5 with \S\ref{sec:crop_training} and
  supersedes C5 as the main consensus test. \textbf{STOPPED 2026-10-03:
  learned consensus does not appear to help.} Arm (b) is complete at
  $1\times$/$2\times$ (2 of 3 seeds at $3\times$), with no improvement
  over P2. Arms (c) and (d) are \textbf{TO DO}
  (\S\ref{sec:c7_results}). Effort moved to C8--C12.
\item[C8 --- where the missing dimensions live (diagnostic).] Fourier content
  of the true and modelled unstable subspaces at $\alpha{=}0.05$.
  \S\ref{sec:c8c12}.
\item[C9 --- learn the local, solve the non-local.] Crop-trained local
  rule plus an exact, globally computed velocity input that is identical at
  training and deployment. \S\ref{sec:c8c12}.
\item[C10 --- physics-anchored hybrid propagator.] An exact coarse
  Kolmogorov step on the anchor channel plus a learned local closure.
  \S\ref{sec:c8c12}.
\item[C11 --- release the halo, train on long-horizon statistics.] Includes
  the free-rollout upper-bound control on the full wide domain.
  \S\ref{sec:c8c12}.
\item[C12 --- reach sweep (P3 realized).] Error against propagator reach,
  with and without the C9 velocity input. \S\ref{sec:c8c12}.
\item[C13 --- large-scale closure diagnostic (no training).] How much of
  the true large-scale tendency the exact coarse PDE explains, how much a
  local closure adds, and how much is noise. \S\ref{sec:two_scale}.
\item[C14 --- large-scale-only model.] Exact coarse dynamics plus a
  learned eddy-forcing closure, trained a posteriori on the whole wide
  domain. \S\ref{sec:two_scale}.
\item[C15 --- coupled two-scale model.] C14 coupled two-way to the
  crop-trained local fine model. \S\ref{sec:two_scale}.
\item[C16 --- self-similarity diagnostic (no training).] Spectra,
  structure functions, rescaled-statistics collapse under
  coarse-graining, and unstable dimensions per octave.
  \S\ref{sec:self_similar}.
\item[C17 --- zero-shot scale transfer.] P2's model applied to
  coarse-grained, rescaled fields: L-transfer's analogue across scale.
  \S\ref{sec:self_similar}.
\item[C18 --- recursive scale-shared model.] One local rule reused at
  levels $\ell_0, 2\ell_0, 4\ell_0,\dots$ with Kolmogorov--Kraichnan
  rescaling and two-way inter-level coupling. \S\ref{sec:self_similar}.
\end{description}

\textbf{What would settle it.} The two outcomes are distinguished as
follows.
\begin{itemize}[leftmargin=*]
\item \emph{Capacity-limited:} C2 improves with $c$ while C4 and C5 do
  not, and C0 shows per-site rank saturating.
\item \emph{Consensus-limited, the current favourite:} C2 is flat while
  C4 and/or C5 improve over their controls.
\end{itemize}
Both outcomes are publishable and point to different architectures. The
first calls for budgeting channels from measured information density. The
second calls for building fusion into the propagator, which also aligns
the model directly with how ensemble DA already combines overlapping
local analyses.

\subsection{Completed: initial-condition normalization check (2026-10-02)}
\label{sec:icnorm}

{\sloppy\textbf{Status: COMPLETED.} \emph{Found while preparing C0, not
pre-registered.} Every training dataset stores
$(u-\bar u)/\sigma$, with $\sigma$ computed on the training split. All
transfer-evaluation scripts (\texttt{exp\_g1/g2/g3},
\texttt{run\_ltransfer\_test}) encoded \emph{raw} solver output instead,
so the encoder's input was $\sigma$ times too large: $\sigma$ is 1.32 for
1D KS, 1.39 and 1.49 for Kolmogorov ($\alpha{=}0.05/0.1$), and 3.82 for
2D KS. Lyapunov exponents are measured after a 200-step warmup, so the
initial condition mainly selects which basin of the model's own latent
dynamics the rollout settles in.
\texttt{scripts/exp\_ic\_normalization\_check.py} reruns each case twice,
changing nothing but the initial-condition scaling (seed 0, settings
identical to \texttt{exp\_g3}).\par}

\begin{table}[h]
\centering
\small
\begin{tabular}{lcccc}
\toprule
Case & truth & raw IC $D_{\mathrm{KY}}$ (err) & normalized IC $D_{\mathrm{KY}}$ (err) & $\sigma$ \\
\midrule
Sec.\ 231, Kolm.\ $\alpha{=}0.05$, $2\times$ & 56.8 & 28.4 (50.0\%) & 30.2 (46.8\%) & 1.39 \\
Sec.\ 233, Kolm.\ $\alpha{=}0.1$, $2\times$ & 30.3 & 31.7 (4.5\%) & 31.5 (3.7\%) & 1.49 \\
D3, Kolm.\ $\alpha{=}0.05$, $2\times$ & 56.8 & 1.02 (98\%) & 1.33 (98\%) & 1.39 \\
Sec.\ 230, 2D KS, $1.5\times$ & 192.5 & 43.6 (77\%) & \textbf{147.4 (23\%)} & 3.82 \\
\bottomrule
\end{tabular}
\caption{Raw versus normalized transfer initial condition (seed 0). Each
raw run reproduces its originally reported number exactly, which confirms
the check isolates the scaling alone.}
\label{tab:icnorm}
\end{table}

\textbf{What changes.}
\begin{itemize}[leftmargin=*]
\item \emph{Kolmogorov is unaffected} ($\sigma\approx1.4$). The
  $\alpha{=}0.1$ headline result stands (3.7\% vs.\ 4.5\%). The
  $\alpha{=}0.05$ failure is real (47\% vs.\ 50\%), and so is D3's collapse.
\item \emph{2D KS's headline failure was mostly an artifact.} At the
  largest mis-scale ($3.8\times$) the rollout landed in a low-dimensional
  basin. The earlier ``43.6 vs.\ $\sim$144 depending on initial
  condition'' observation was this same effect. With a correct initial
  condition the transfer error is 23\%: not clean, but a different
  finding from 77\%.
\item \emph{Seeds 1--2 (added 2026-10-03, normalized IC only):} 145.3 and
  145.5, i.e.\ 24.5\% and 24.4\%. The corrected 2D KS transfer error is
  therefore $\approx$24\% and highly reproducible across 3 seeds.
\item \emph{Consequence for \S\ref{sec:capacity}.} The 2D KS row of
  Table~\ref{tab:capacity} no longer describes a near-total failure.
  Every evaluation from here on (C0, C4) uses normalized initial
  conditions.
\end{itemize}

\subsection{Completed: C0 capacity audit (2026-10-02)}
\label{sec:c0_results}

\textbf{Status: COMPLETED.} Script: \texttt{scripts/exp\_c0\_capacity\_audit.py}.
For each model and size, the script integrates 3 independent true
trajectories (100 snapshots each), normalizes them with the training
split's statistics, and encodes them. It also rolls the model forward 600
steps from each encoded initial condition, discarding the first 200.
Statistics are computed on both latent sets:
\begin{itemize}[leftmargin=*]
\item the per-site participation ratio of the $c{\times}c$ channel
  covariance (PR\textsubscript{cov}, pooled over sites and time);
\item the same on the channel correlation matrix (PR\textsubscript{cor},
  which is scale-free);
\item the global participation ratio over all $d_{\mathrm{lat}}$ channels
  (gPR);
\item the rollout-to-true ratio of total latent variance.
\end{itemize}
$c{=}3$ throughout. A per-site two-NN intrinsic dimension was also
computed but is \emph{not reported}. It returned values up to 8.3 for
3-dimensional vectors, which is impossible, so the estimator is invalid
at this ambient dimension.

\begin{table}[h]
\centering
\footnotesize
\setlength{\tabcolsep}{4pt}
\begin{tabular}{lccc|ccc|cccc}
\toprule
 & & & & \multicolumn{3}{c|}{true, encoded} & \multicolumn{4}{c}{model rollout} \\
Model & size & $d_{\mathrm{lat}}$ & err & PR\textsubscript{cov} & PR\textsubscript{cor} & gPR & PR\textsubscript{cov} & PR\textsubscript{cor} & gPR & var \\
\midrule
1D KS (224) & $1\times$ & 48 & --- & 1.68 & 2.21 & 11.7 & 1.79 & 2.31 & 14.6 & 0.96 \\
 & $2\times$ & 96 & $<$2\% & 1.63 & 2.12 & 14.8 & 1.80 & 2.32 & 27.4 & 0.86 \\
 & $4\times$ & 192 & $<$2\% & 1.68 & 2.21 & 22.3 & 1.79 & 2.32 & 45.8 & 0.90 \\
\midrule
2D KS (230) & $1\times$ & 192 & --- & 1.59 & 1.69 & 5.2 & 1.48 & 1.54 & 8.9 & 1.28 \\
 & $1.5\times$ & 432 & 23\%$^*$ & 1.48 & 1.52 & 3.8 & 1.73 & 1.84 & 13.8 & 0.83 \\
\midrule
Kolm.\ $\alpha{=}0.1$ (233) & $1\times$ & 192 & --- & 1.56 & 1.80 & 3.3 & 1.54 & 1.78 & 4.3 & 1.10 \\
 & $2\times$ & 384 & 4.5\% & 1.48 & 1.74 & 2.4 & 1.61 & 1.84 & 8.6 & 1.01 \\
 & $3\times$ & 576 & 19\% & 1.58 & 1.87 & 4.2 & 1.62 & 1.86 & 9.9 & 1.11 \\
\midrule
Kolm.\ $\alpha{=}0.05$ (231) & $1\times$ & 192 & --- & 1.74 & 1.93 & 6.0 & 1.62 & 1.85 & 8.0 & 1.21 \\
 & $2\times$ & 384 & 47\% & 1.62 & 1.90 & 5.6 & 1.55 & 1.82 & 8.7 & 1.04 \\
D2 (biot--savart) & $1\times$ & 192 & --- & \multicolumn{3}{c|}{(same AE as 231)} & 1.67 & 1.88 & 9.0 & 1.04 \\
 & $2\times$ & 384 & 95\% & & & & \textbf{1.14} & \textbf{1.29} & \textbf{2.9} & 1.44 \\
D3 (multiscale) & $1\times$ & 192 & --- & & & & 1.57 & 1.79 & 5.6 & 1.29 \\
 & $2\times$ & 384 & 98\% & & & & \textbf{1.25} & \textbf{1.40} & \textbf{3.0} & 1.06 \\
\bottomrule
\end{tabular}
\caption{C0 capacity audit (maximum per-site participation ratio is
$c{=}3$). $^*$Normalized-IC transfer error (\S\ref{sec:icnorm}). gPR uses
300 (true) or 1200 (rollout) samples and is a coarse, low-sample
statistic.}
\label{tab:c0}
\end{table}

\textbf{Pre-registered rule outcome.} The rule rules capacity out only if
per-site rank is below $c/2=1.5$ at every failing transfer. That condition
is \emph{not met}: the values are 1.48 to 1.62, right at the threshold. So
C0 does not formally rule capacity out. In hindsight the threshold was
poorly placed, because the gauge-anchored channel 0 has about 0.45--0.85
of the learned channels' standard deviation, which caps PR\textsubscript{cov}
well below $c$ even when all three channels are in use. We record that
rather than move the threshold after seeing the data.

\textbf{What the data do show.}
\begin{enumerate}[leftmargin=*]
\item \emph{There is no capacity-saturation signature anywhere.} The
  capacity hypothesis predicted per-site rank rising toward $c$ as $L$
  grows, coinciding with failure. Instead per-site rank is flat or
  slightly falling with $L$ in every system (1D: 1.68, 1.63, 1.68; 2D KS:
  1.59, 1.48; Kolm.\ $\alpha{=}0.1$: 1.56, 1.48, 1.58; $\alpha{=}0.05$:
  1.74, 1.62).
\item \emph{Channel usage is the same in the success and the failures.}
  The 1D KS model, which transfers to $4\times$ within 2\%, uses its
  channels as fully as any failing 2D model, or more fully
  (PR\textsubscript{cor} 2.1--2.2 vs.\ 1.5--1.9). Nothing per-site
  separates success from failure. Together with
  Table~\ref{tab:capacity}, C0 gives \emph{no support} for capacity as the
  cause. C2 (the channel sweep) remains the decisive test.
\item \emph{No failure is an amplitude collapse.} The rollout-to-true
  variance ratio is 0.83--1.44 in every case, including D2 and D3 at
  $2\times$, whose $D_{\mathrm{KY}}$ is 1--3. The failing models keep
  their energy and lose \emph{dimensions}.
\item \emph{New: D2 and D3 fail by rank collapse, i.e.\
  over-synchronization.} At $2\times$ their rollouts drop to per-site
  PR\textsubscript{cov} 1.14 and 1.25 and global PR about 3, down from
  1.57--1.67 and 5.6--9 at native size, while total variance holds or
  grows. Each site's channels become nearly collinear, and the whole field
  is dominated by about three global directions. The added cross-site
  coupling has \emph{locked} the sites into a low-dimensional coherent
  state. That is the opposite of the ``copies drift apart'' branch of the
  consensus hypothesis: it is excess consensus, not missing consensus.
  The plain Section 231 model at $\alpha{=}0.05$ shows no such
  signature (1.55, 8.7) yet still underestimates $D_{\mathrm{KY}}$ by
  47\%. So the plain model's failure and the D2/D3 failures are
  different mechanisms.
\end{enumerate}

\textbf{Implication for C4.} Re-projection is a consensus-enforcing
operator. Given finding 4, it could \emph{reduce} dimension further, as
it did for D2/D3. C4's native-size control (native $D_{\mathrm{KY}}$
within 10\%) is therefore essential, and a drop in transferred
$D_{\mathrm{KY}}$ under re-projection would count against the
missing-consensus branch, not for it.

\subsection{Completed: C4 re-projection (2026-10-02)}
\label{sec:c4_partial}

\textbf{Status: COMPLETED.} The $\alpha{=}0.1$ $m{=}16$ runs were skipped by decision, since the result was already clear. Script: \texttt{scripts/exp\_c4\_reprojection.py}.
All initial conditions are normalized. Every run covers the same 2000
propagator steps, with the composite step counts divided by $m$. Results
so far, Section 231 model at $\alpha{=}0.05$:

\begin{table}[h]
\centering
\small
\begin{tabular}{lcccc}
\toprule
$m$ & native $D_{\mathrm{KY}}$ (shift vs.\ $m{=}0$) & $2\times$ $D_{\mathrm{KY}}$ & $2\times$ error & rule \\
\midrule
0 (none) & 21.62 (---) & 30.21 & 46.8\% & baseline \\
1 & 24.97 (+15.5\%) & \textbf{2.02} & 96.4\% & uninterpretable (native $>$10\%) \\
4 & 23.91 (+10.6\%) & \textbf{12.11} & 78.7\% & uninterpretable (native $>$10\%) \\
16 & 21.34 ($-$1.3\%) & \textbf{18.90} & 66.7\% & \textbf{no improvement} \\
\bottomrule
\end{tabular}
\caption{C4, Section 231 half complete (truth 23.26 native, 56.84 at
$2\times$). Transferred $D_{\mathrm{KY}}$ falls monotonically as
re-projection becomes more frequent: 30.2, 18.9, 12.1, 2.0 for no
re-projection, then $m{=}16, 4, 1$.}
\label{tab:c4_partial}
\end{table}

\textbf{Control: Section 233 at $\alpha{=}0.1$, a model that already transfers.}

\begin{table}[h]
\centering
\small
\begin{tabular}{lcccc}
\toprule
$m$ & native $D_{\mathrm{KY}}$ (shift) & $2\times$ $D_{\mathrm{KY}}$ & $2\times$ error & rule \\
\midrule
0 (none) & 12.16 (---) & 31.46 & 3.7\% & baseline \\
1 & 0.00 ($-$100\%) & 2.00 & 93.4\% & uninterpretable \\
4 & 12.37 (+1.7\%) & \textbf{1.86} & \textbf{93.9\%} & \textbf{no improvement (destroys transfer)} \\
\bottomrule
\end{tabular}
\caption{C4 control (truth 15.88 native, 30.34 at $2\times$).}
\label{tab:c4_control}
\end{table}

The control is the sharpest result in C4. At $m{=}4$, re-projection leaves
the working model's native dynamics essentially unchanged (+1.7\%), yet
collapses its $2\times$ transfer from 3.7\% error to 93.9\%. So the
damage is \emph{size-specific}, not a generic perturbation of the
dynamics. A plausible reading is that the autoencoder was fit only at
native size: at $2\times$ the propagator's outputs, including modes at
wavenumbers absent from the training box, sit further from what $E\circ D$
reproduces, and the projection removes them. This is the same
train-only-at-native-size mismatch behind D2 and D3, now exposed in the
autoencoder itself. It argues directly for crop training, where the
autoencoder sees large-domain inputs (\S\ref{sec:crop_training}).

\textbf{Pre-registered verdict, Section 231.} $m{=}16$ is the only
interpretable setting, with native $D_{\mathrm{KY}}$ within 1.3\%. It
gives \emph{no improvement}: $2\times$ error rises from 46.8\% to
66.7\%. Every setting is worse than the baseline at $2\times$, and the
damage grows with re-projection frequency. This is a clean dose--response
against an untrained projection used as consensus.

\textbf{Reading.} Under the pre-registered rule, $m{=}1$ is
uninterpretable, because re-projection moved native $D_{\mathrm{KY}}$ by
more than 10\%. Substantively, though, it reproduces the D2/D3 failure
exactly (\S\ref{sec:c0_results}, finding 4): a mild effect at native size
and a collapse to $D_{\mathrm{KY}}\approx2$ at $2\times$. This time it is
caused by an operator with \emph{no new learned weights at all}.

The operator explains why. $z\mapsto E(D(z))$ projects the latent onto
the set of latents the autoencoder can reproduce, followed by the
encoder's pooling. It is contractive by construction: any component the
propagator places outside that set is discarded at every application.
At native size the propagator's outputs stay close to that set. At
$2\times$ they apparently drift further, so the projection removes much
more, including the components that carry the chaotic degrees of
freedom.

So C4 does \emph{not} show that consensus is unnecessary. It shows that an
\emph{untrained, contractive} consensus collapses dimension at new sizes.
Communication itself is not optional. Outside the pressure solve the PDE
propagates large-scale information by local hops, and a local model must
carry large-scale dynamics somehow. What the evidence constrains is the
\emph{form} that communication can take.

\subsection{C7: a learned consensus mechanism that cannot over-synchronize}
\label{sec:c7}

\textbf{The common failure.} D2 (exact Biot--Savart input), D3 (learned
regional scale) and C4 (untrained re-projection) share one failure mode.
At a new domain size, the added cross-site communication
\emph{over-synchronizes} the sites into a low-rank state while the energy
is preserved: per-site rank falls toward 1, global rank to about 3, and
variance holds (Table~\ref{tab:c0}). A mechanism meant to carry
large-scale dynamics has to be built so that this collapse is neither the
default nor learnable.

\textbf{Design: four requirements, each aimed at one observed failure.}
\begin{enumerate}[leftmargin=*]
\item \emph{Starts as a no-op (against C4's forced averaging).} The
  consensus step is a residual, gated message-passing update between
  neighbouring sites,
  \[
  z_j \leftarrow z_j + g_j \odot W_o\,\textstyle\sum_{k\in\mathcal N(j)} a_{jk}\,m_{jk},
  \]
  with $W_o$ zero-initialized. At initialization the model is exactly
  the existing local propagator. Communication appears only where
  training finds it reduces rollout error, and it is never imposed.
\item \emph{Learns what to trust (against blind averaging).} Each site
  outputs a mean and a log-variance for its estimate of shared quantities.
  Messages are combined inverse-variance weighted
  ($a_{jk}\propto\sigma_k^{-2}$). This is the same rule LETKF uses to
  merge overlapping local analyses, so confident local detail is kept
  rather than averaged away. With overlapping patches (stride $p/2$), each
  physical point has $r{=}4$ redundant estimates to fuse. The fusion is
  the trained replacement for C4's $E\circ D$.
\item \emph{Trained on the deployment distribution (against D2/D3's
  distribution mismatch).} Train on crops of a large domain with
  forced-halo rollouts (\S\ref{sec:crop_training}, variant B). Also run
  multi-size rollouts: crop widths drawn from
  $\{2,3,4\}\times2\pi$ in each batch. The mechanism then never learns
  behaviour that exists only on an 8-site periodic ring.
\item \emph{Rank collapse is penalized directly (against the C0/D2/D3
  signature).} Add a rank-preservation term on rollout latents,
  \[
  \mathcal L_{\mathrm{rank}} = -\lambda\,\log\det\!\big(\mathrm{Cov}_{\mathrm{sites},t}[z^{\mathrm{roll}}_j] + \epsilon I\big),
  \]
  computed per site ($c{\times}c$) and, optionally, on a low-pass global
  projection. This is the same idea as the \texttt{w\_logdet} Stage-1
  regularizer that already fixed this project's latent-channel collapse,
  now applied to the \emph{dynamics} instead of the encoder. Its target is
  the true-encoded statistics (Table~\ref{tab:c0}), so it does not push
  rank above what the data support.
\end{enumerate}

\textbf{Experiment.} Kolmogorov flow at $\alpha{=}0.05$, using P0's wide
dataset. Four arms, all crop-trained:
\begin{enumerate}[label=(\alph*),leftmargin=*]
\item local propagator only;
\item plus learned consensus (requirements 1--3);
\item plus learned consensus and $\mathcal L_{\mathrm{rank}}$
  (requirements 1--4);
\item non-overlapping $4c$-channel control at matched latent size, the C5
  capacity control.
\end{enumerate}
Evaluate at $1/2/3\times$ with 3 initial-condition seeds each, and report
both $D_{\mathrm{KY}}$ and the C0 rank statistics of the rollouts.

\textbf{Pre-registered rules.}
\begin{itemize}[leftmargin=*]
\item \emph{Consensus helps:} (b) or (c) beats (a) and (d) at $2\times$ by
  at least 10 percentage points, with native $D_{\mathrm{KY}}$ within
  10\%.
\item \emph{Rank protection is what matters:} (c) succeeds where (b)
  shows the C0 rank-collapse signature (per-site PR\textsubscript{cov}
  below 1.3).
\item \emph{Crop training alone suffices:} (a) already reaches single
  digits. Then consensus is unnecessary in this regime, and P2 is the
  headline instead.
\item \emph{Negative:} all four arms stay above 25\%. The remaining gap
  is then non-local physics that neither local relay nor learned
  consensus captures, which is evidence for an explicit, properly
  normalized long-range term (a corrected D2).
\end{itemize}

\textbf{Implementation and pre-run choices (2026-10-02, before any C7
run).} \textbf{Status: STOPPED after arm (b); arms (c), (d) TO DO} (\texttt{scripts/exp\_c7\_pipeline.py} resumes from saved outputs).
\begin{itemize}[leftmargin=*]
\item \emph{Consensus block}
  (\texttt{backbone="local\_mlp\_2d\_consensus"},
  \texttt{\_ConsensusRelax2D}). Before each of the 2 local mixer blocks,
  each site emits a 16-dimensional message and a log-precision. Then 4
  rounds of precision-weighted averaging run over the radius-1 (Chebyshev)
  neighbourhood. The result enters only through
  $\sigma(\mathrm{gate}(h))\odot W_o m$ with $W_o$ zero-initialized. This
  is the learned analogue of an iterative (Jacobi-like) elliptic
  relaxation: a local, size-agnostic way to approximate long-range
  coupling. It \emph{adds} context; it does not overwrite the state the
  way $E\circ D$ did.
\item \emph{Tests} (\texttt{tests/unit/test\_crop\_training.py}). The
  block is the identity at initialization, exactly reduces to plain
  \texttt{local\_mlp\_2d}, and is exactly shift-equivariant. It resizes
  through the standard loader. Its forced-halo interior is exact when the
  halo covers its reach.
\item \emph{Halo grows with consensus reach.} The per-step receptive field
  is $n_{\mathrm{layers}}(\text{window}+n_{\mathrm{iter}}r) = 2(3+4)=14$
  sites. Arms (b) and (c) therefore use 64-site crops (halo 14, free
  interior 24, the same as P2).
\item \emph{Rank penalty} (arm c). The penalty is
  $\lambda\,\mathrm{ReLU}(\log\det\Sigma_{\mathrm{true}} -
  \log\det\Sigma_{\mathrm{pred}})$ on the per-site $c{\times}c$ channel
  covariance of predicted interior rollouts, with $\lambda{=}0.02$ fixed
  in advance. It is one-sided: it only resists collapse below the true
  encoded statistics.
\item \emph{Arms (b) and (c) reuse P2's autoencoder,} so only the
  propagator differs from arm (a). Arm (d) trains its own $c{=}12$
  autoencoder.
\item \emph{Two parts of the requirement list are dropped, with reasons.}
  \begin{itemize}[leftmargin=*]
  \item Mixed crop widths (requirement 3) are redundant. The exactness test
    shows the free-interior update is independent of crop width once the
    halo covers the reach.
  \item A separate overlapping-patch encoder (requirement 2's stride
    $p/2$) is not built. The encoder's 6-site receptive field already makes
    neighbouring latents overlapping, redundant measurements of the same
    physical region.
  \end{itemize}
\end{itemize}

\textbf{Dependencies and order.} C7 needs P0 (the wide dataset) and the
crop-training data loader with forced halos. It should run after P1
confirms crop training works at $\alpha{=}0.1$. If P2 alone succeeds, C7
becomes an ablation of a working system rather than a rescue.

\subsection{C7 results so far (2026-10-03)}
\label{sec:c7_results}

\textbf{Arm (b), learned consensus, no rank penalty: COMPLETED at
$1\times$/$2\times$.} At $3\times$, seeds 0 and 1 give $D_{\mathrm{KY}}$
71.4 and 65.2 (P2: 72.0 and 69.6). These will be scored once the
$\alpha{=}0.05$ $3\times$ truth run (now running, as C8's prerequisite)
completes. \textbf{C7 was stopped on 2026-10-03} because learned
consensus does not appear to help. Arm (b)'s third $3\times$ seed, arm
(c) (rank penalty) and arm (d) (matched-capacity control) remain
\textbf{TO DO}. If C9 or C10 changes the picture, arm (c) becomes the
natural ablation.

\begin{table}[h]
\centering
\small
\begin{tabular}{lccc}
\toprule
size & C7(b) $D_{\mathrm{KY}}$ (seeds 0/1/2) & mean error & P2 (arm a) mean error \\
\midrule
$1\times$ & 22.5 / 22.6 / 24.0 & 3.1\% & 7.3\% \\
$2\times$ & 43.9 / 46.6 / 45.4 & 20.3\% & 20.9\% \\
\bottomrule
\end{tabular}
\caption{C7 arm (b) against arm (a), which is P2 (truth: 23.26 native,
56.84 at $2\times$).}
\label{tab:c7b}
\end{table}

\textbf{Pre-registered rule: consensus does not help.} The rule required
(b) to beat (a) at $2\times$ by at least 10 points; the result is
20.3\% vs.\ 20.9\%, within seed spread. The consensus block, with a
14-site per-step reach and inverse-variance fusion, slightly tightens
native-size accuracy but leaves the $2\times$ deficit untouched. The model
remains extensive (about 22.7 per box at $2\times$).

\subsection{Why cross-patch communication has not helped at $\alpha{=}0.05$}
\label{sec:c8_diag}

The evidence now spans six interventions:
\begin{itemize}[leftmargin=*]
\item box training, 47\% error at $2\times$;
\item D2 (exact Biot--Savart input), 95\%;
\item D3 (learned regional scale), 99\%;
\item C4 (untrained re-projection), destructive;
\item P2 (crop training), 21\%;
\item C7(b) (learned consensus), 20\%.
\end{itemize}
Every model that works is \emph{extensive}, at about 23 $D_{\mathrm{KY}}$
per box, while the truth grows faster than the area (28.4 per box at
$2\times$). Four mechanisms can explain why adding local communication
does not close that gap. They are listed in the order we now consider
most likely.

\begin{enumerate}[label=(\alph*),leftmargin=*]
\item \textbf{Forced halos hide the large-scale problem from the loss.}
  In variant-B training, everything larger than the free interior
  ($3\times2\pi$) is \emph{prescribed} through the halo at every step. The
  model is never asked to generate or sustain the dynamics of structures
  larger than its interior. At deployment on a periodic box there is no
  halo, so those dynamics must arise internally, and nothing in training
  rewarded that. C7 is subject to exactly this: its consensus block could
  only learn to relay information the halo already supplied.
\item \textbf{The loss horizon and scale weighting ignore slow collective
  modes.} Rollouts are at most 12 steps of $dt_{\mathrm{snap}}{=}0.2$,
  i.e.\ 2.4 time units. That is shorter than a Lyapunov time and far
  shorter than the turnover time of domain-scale structures. The latent
  MSE is also dominated by small scales. The slow, large-scale modes that
  carry the super-extensive excess are almost unconstrained.
\item \textbf{The instantaneous non-local coupling is missing.} Velocity
  at every point is set by the vorticity everywhere through the $1/k^2$
  solve. Relay through local neighbourhoods, however many rounds, is not
  equivalent: the coupling acts every step, at all distances, while
  domain-scale vortices advect and merge within a few steps. D2 tried to
  supply this coupling but was trained on the box. Its input was
  therefore out of distribution at deployment, and on a crop the solve
  is not even well defined, because the velocity depends on vorticity
  outside the crop.
\item \textbf{Local states at large $L$ are off-distribution.} Larger
  domains host stronger large-scale shear and velocity than small boxes.
  Crops from the wide domain partly fix this (P2), but a crop
  ($6\times2\pi$) is still smaller than the largest structures of a
  $16\times2\pi$ domain.
\end{enumerate}

\textbf{Design principle.} The common lesson is the one numerical models
already follow: \emph{learn the local, solve the non-local, and train
where the loss can see the large scales}. Local tendencies are learned
from patches. The elliptic solve is done exactly, by the same operator at
training and deployment, so its inputs are always in distribution. And
the free part of the training rollout must be large and long enough that
large-scale dynamics are generated by the model, not imposed on it.
C8--C12 implement this, in order of cost and diagnostic value.

\subsection{Paused (2026-10-03): exact resume point}
\label{sec:paused_state}

All background runs were stopped on user request at $\sim$18:50. Nothing
was lost -- every pipeline here is preemptable and skips completed
outputs. To resume, simply relaunch the three queue scripts below in any
order (they wait on each other / on file existence as already coded);
each will pick up exactly where it left off.

\begin{description}[leftmargin=*]
\item[P6 (\texttt{scripts/run\_p6\_queue.sh}).] C16 and C13 diagnostics:
  done. $6$-site free-interior arm (\texttt{p6\_i6}): done, collapsed
  (\S\ref{sec:p6} table pending write-up). $8$-site arm
  (\texttt{p6\_i8}): $3\times$ seeds 0,1 done; seed 2 was mid-run when
  killed (no partial file written, will simply rerun). $12$-site arm and
  the full-periodic upper-bound arm: not started.
\item[P7 step 2 (\texttt{scripts/run\_p7\_step2\_queue.sh}).] Wide
  $6\times44$ 2D-KS dataset generation: trajectory 0/5 done ($\sim$32
  min), saved at
  \texttt{artifacts/datasets/ks2d\_wide6\_parts/traj0.npy}; trajectories
  1--4 and the final HDF5 assembly remain. Training and evaluation have
  not started (they wait on the full dataset and on the P6 queue
  finishing).
\item[P8 / drag sweep (\texttt{scripts/run\_alpha\_sweep\_queue.sh}).]
  Not started (waits on P6 and P7 step 2).
\end{description}

\textbf{To resume:} re-launch all three with \texttt{setopt NO\_BG\_NICE}
first (zsh's default background-job niceness caused the earlier
slowdowns, \S\ref{sec:todo_deferred}):
\begin{verbatim}
setopt NO_BG_NICE
nohup scripts/run_p6_queue.sh > artifacts/logs/p6_queue.log 2>&1 &
nohup scripts/run_p7_step2_queue.sh > artifacts/logs/p7_step2_queue.log 2>&1 &
nohup scripts/run_alpha_sweep_queue.sh > artifacts/logs/alpha_sweep_queue.log 2>&1 &
\end{verbatim}

\subsection{Current staging: highest likelihood of a useful result first (2026-10-03)}
\label{sec:staging}

Pending experiments are ordered by the likelihood that they produce a
decisive or positive result, weighed against cost. Diagnostics come
first because they always produce an answer and they gate the expensive
builds.
\begin{enumerate}[leftmargin=*]
\item \textbf{C8 analysis} (queued; waits for the $\alpha{=}0.05$
  $3\times$ truth run). Decides whether the missing dimensions are
  large-scale.
\item \textbf{C16, self-similarity diagnostic} (queued;
  \texttt{scripts/exp\_c16\_self\_similarity.py}; no training, about 20
  minutes). Gates C17/C18. Part (iii) reuses C8's truth basis.
\item \textbf{C13, closure diagnostic} (queued;
  \texttt{scripts/exp\_c13\_closure\_diagnostic.py}; no training). Gates
  C14/C15. It computes the exact eddy forcing from the full field. A
  bookkeeping check on one $\alpha{=}0.1$ trajectory confirms that
  $T\approx R+E$ to better than $10^{-3}$ of $\mathrm{var}(T)$. The same
  check already hints that the resolved terms alone explain little: the
  eddy forcing's variance is 4--12$\times$ that of the tendency, so the
  large scales sit in a strong balance between resolved advection and
  eddy forcing.
\item \textbf{P6, same-domain patch training} (queued). Highest
  likelihood of a positive result, since native-size models consistently
  land within 3--8\%. It measures a locality length at
  $\alpha{=}0.05$.
\item \textbf{C11, free rollouts, halo release and long-horizon
  statistics} (code to be written while P6 runs). The most direct test
  of the leading explanation, that training hides the large scales.
\item \textbf{C17} if C16 passes; \textbf{C10} if C13 shows the closure
  is tractable.
\item C14, C15, C18, C12, then the deferred items
  (\S\ref{sec:todo_deferred}).
\end{enumerate}

\subsection{Deferred experiments (TO DO, as of 2026-10-03)}
\label{sec:todo_deferred}

Stopped or deprioritized in favour of C8--C12. All can be resumed from
their saved outputs.
\begin{itemize}[leftmargin=*]
\item \textbf{C7 arms (c) and (d)}, plus arm (b)'s third $3\times$ seed
  (\texttt{scripts/exp\_c7\_pipeline.py}). Consensus did not help in arm
  (b).
\item \textbf{T1 $2\times$ truth seeds 0--2} at $\alpha{=}0.1$
  (\texttt{scripts/run\_t1\_x2\_truth\_after\_c7.sh}). These decide
  whether the $2\times$ dip behind P1's $2\times$ miss is real.
\item \textbf{2D KS normalized-IC seeds 1--2}
  (\texttt{scripts/run\_icnorm\_ks2d\_seeds.sh}). These confirm the 77\%
  to 23\% correction of \S\ref{sec:icnorm}.
\end{itemize}

\textbf{Operational note.} The slow background runs of 2026-10-02/03 had
one root cause. zsh's default \texttt{BG\_NICE} option starts every
background job at nice 5, which on Apple Silicon pushes long jobs onto
the efficiency cores, a 5--15$\times$ slowdown. All jobs from C8 onward
are launched with \texttt{setopt NO\_BG\_NICE}, so they run at nice 0.

\subsection{C8--C12: making local training capture emergent large-scale chaos}
\label{sec:c8c12}

The rules below are pre-registered (2026-10-03, before any of these
runs). All runs are at $\alpha{=}0.05$ unless stated. Success is measured
against P2's $2\times$ error (20.9\%) and, once available, the
$\alpha{=}0.05$ $3\times$ truth.

\begin{description}[leftmargin=*]
\item[C8 --- where do the missing dimensions live? (diagnostic, cheap).]
  The Benettin checkpoint stores the orthonormal tangent basis, and the
  span of its leading $n$ vectors is the leading Oseledets subspace. The
  queued $\alpha{=}0.05$ $3\times$ truth run therefore gives the unstable
  subspace at no extra cost. Compute the fraction of its variance at
  streamwise wavelengths $\lambda_x$ in three bands:
  \begin{itemize}[leftmargin=*]
  \item longer than the training box, $\lambda_x>2\pi$, i.e.\ modes that
    do not exist at native size;
  \item longer than the crop interior, $\lambda_x>6\pi$;
  \item shorter than these, the local band.
  \end{itemize}
  Do the same for the P2 model at $3\times$, with latent tangent vectors
  mapped to physical space through the decoder Jacobian. Also compare the
  temporal autocorrelation of the lowest $k_x$ modes between truth and
  model rollouts, which tests mechanism (b).
  \emph{Prediction:} the truth's unstable subspace carries substantially
  more variance at $\lambda_x>2\pi$ than the model's.
  \emph{Rule:} if the truth-minus-model excess at $\lambda_x>2\pi$ accounts
  for at least half the missing dimensions' variance, mechanisms (a)/(c)
  are confirmed and C9/C10 take priority. If the excess is mostly local,
  the deficit is a local-rule accuracy problem and C11's long-horizon
  statistics come first.
\item[C9 --- learn the local, solve the non-local.] Give the propagator,
  at every site, the velocity $(u,v)$ obtained from the \emph{exact}
  Biot--Savart solve of the site-resolution anchor field (channel 0, the
  local mean vorticity) over the \emph{whole periodic domain}. This is
  D2's physics, fixed in the three ways D2 lacked:
  \begin{itemize}[leftmargin=*]
  \item \emph{Training on crops, with a consistent solve.} At each
    rollout step, build the full wide-domain anchor field: truth
    everywhere except the crop's free interior, where the model's
    predicted channel 0 is used. Solve on the wide domain (a $128\times8$
    FFT, cheap) and crop. Inside the interior this equals the velocity the
    deployed model would compute from the same anchor values, so the input
    is consistent at training and deployment by construction. This needs
    a unit test analogous to the forced-halo exactness test.
  \item \emph{Deployment.} Solve on the deployed periodic domain from the
    model's own anchor field, using the same operator on the same
    quantity.
  \item \emph{Normalization and safety.} Normalize velocity by the
    wide-set statistics. This fixes D2's out-of-range amplitudes, because
    the training data contain large-domain velocities. Enter it through a
    zero-initialized gated embedding, so the model is the identity at
    initialization.
  \end{itemize}
  Site resolution resolves exactly the large scales this targets
  ($\lambda\ge2$ sites); small-scale velocity stays with the local rule.
  \emph{Arms:} (i) crops of the wide domain; (ii) trained only on the
  $2\pi$ periodic box, the ``train small, transfer large'' goal.
  \emph{Rules:} strong if $2\times$ error is single-digit; partial if
  below 15\%, a clear gain over P2's 20.9\%. The $\alpha{=}0.1$ control
  must stay within P1's accuracy.
\item[C10 --- physics-anchored hybrid propagator.] Advance channel 0 with
  an exact coarse Kolmogorov step at site resolution: advection by C9's
  velocity, $\nu\nabla^2$, $-\alpha\omega$ and $f(y)$, with sub-stepping
  over $dt_{\mathrm{snap}}$. Learn only a local correction to channel 0
  (the closure) and the evolution of the hidden channels $1..c{-}1$,
  which carry sub-patch information. Large-scale dynamics, where the
  non-local excess lives if C8 confirms it, are then handled by exact
  physics; learning is confined to the local closure. This is the
  latent-field analogue of learned corrections to coarse solvers
  \cite{barsinai2019,kochkov2021}. \emph{Resolution caveat:} at 8 sites
  per $2\pi$ the site-grid Nyquist wavenumber is 4, exactly the forcing
  wavenumber. The primary C10 configuration therefore uses $8\times8$
  patches (16 sites per $2\pi$, $4\times$ the latent size). The 16-patch
  version, with forcing-scale dynamics left to the closure, is the
  secondary arm. \emph{Arms and rules:} as C9.
\item[C11 --- release the halo; train on long-horizon statistics.]
  Applied to the best of P2, C9 and C10.
  \begin{itemize}[leftmargin=*]
  \item \emph{Halo-release curriculum:} the halo is forced only for the
    first $m$ steps of each rollout, with $m$ annealed from $k$ to 0, so
    the interior must sustain its own large scales (addresses (a)).
  \item \emph{Long rollouts with a statistics loss:} $K$ up to 100--200
    steps, at least 4 Lyapunov times. Pointwise error is meaningless at
    that horizon, so the loss compares statistics instead: the low-$k_x$
    energy spectrum, the autocorrelation of the lowest modes, and the PDF
    of large-scale velocity, between free-running crops and truth
    (addresses (b)). Monitor the C0 rank statistics throughout.
  \item \emph{Upper-bound control:} free rollouts on the entire
    $16\times2\pi$ periodic wide domain. These are exact, with no halo
    and no crop. This shows what the architecture can do when allowed to
    see large scales in training. If even this fails at $2\times$, the
    limit is the form of the local rule, not the training scheme.
  \end{itemize}
  \emph{Rule:} as C9, plus the control's error defines the achievable
  floor for every crop-trained variant.
\item[C12 --- reach sweep (P3 realized).] Run consensus rounds
  $\{0,4,8,16\}$ with crops sized to match, crossed with the C9 velocity
  input on or off, at $\alpha\in\{0.05,0.1\}$. Report $2\times$ error and
  the saturation reach $h^*(\alpha)$. This measures the effective
  interaction length directly and tests the $\ell_\alpha$ scaling of
  \S\ref{sec:transfer_diagnostic}.
\end{description}

\textbf{Order.}
\begin{enumerate}[leftmargin=*]
\item C8, which needs only the queued truth run and decides the priority
  of everything else.
\item C9, the cheapest change aimed at the most likely mechanism.
\item C11's upper-bound control, which sets the floor.
\item C10.
\item The remaining C11 variants.
\item C12.
\end{enumerate}
Every arm reports $D_{\mathrm{KY}}$ at $1/2/3\times$ with 3 seeds and the
C0 rank statistics of its rollouts.

\subsection{Status of C8/C9, and a bug that invalidates D2's transfer tests (2026-10-03)}
\label{sec:c8c9_status}

\textbf{Bug found while building C9.} \texttt{aux\_cfg\_to\_propagator\_cfg}
converts the auxiliary propagator config used by every training script
into the propagator config. It copied a \emph{hand-maintained} list of
fields and silently dropped every field added after it was written:
\texttt{biot\_savart\_Lx/Ly}, \texttt{consensus\_*} and
\texttt{multiscale\_pool\_*}.
\begin{itemize}[leftmargin=*]
\item \emph{C7 and D3 are unaffected.} They were trained with exactly the
  defaults the dropped fields fall back to.
\item \emph{D2 is not.} Its resize loader set the new domain length, but
  the conversion discarded it. Every D2 transfer evaluation therefore
  solved for velocity on the \emph{native $2\pi$ $k$-grid} instead of the
  actual domain. Large-scale velocities were wrong at $2\times$ and
  $3\times$, so D2's 95\% ``failure'' (\S\ref{sec:d1d3}), and its C0
  rows, were never a valid test of an exact velocity input. Native-size
  D2 numbers are correct.
\item \emph{Fix.} The conversion now copies every shared field
  automatically. Every dropped field had identical defaults on both
  sides, so the fix changes nothing for any other checkpoint.
\item \emph{Regression tests.} Two new tests cover it: every shared field
  is copied, and a resized Biot--Savart checkpoint gets the new $L_x$ in
  the model itself, not only in its config. 410 related tests pass.
\end{itemize}
\textbf{D2 re-evaluated with the fix: COMPLETED.} At $2\times$, with
normalized initial conditions and the correct $L_x$, D2 gives
$D_{\mathrm{KY}}=1.6$, 1.9 and 0.0 (seeds 0/1/2; 97--100\% error). It
still collapses. The bug was real but was \emph{not} the cause of D2's
failure: a box-trained model with an unnormalized exact-velocity input
collapses under transfer either way. C9 changes three things that D2
lacked: training data from a large domain, a velocity solve consistent
between crop and deployment, and normalization. C9 arm (ii) isolates the
normalization effect for a box-trained model.

\textbf{C8 implementation} (\texttt{scripts/exp\_c8\_unstable\_subspace.py}).
\textbf{Status: IN PROGRESS.}
\begin{itemize}[leftmargin=*]
\item \emph{Truth side.} The $\alpha{=}0.05$ $3\times$ truth run is
  running (128 directions, total time 40). Its final checkpoint stores
  the orthonormal Benettin basis in physical space.
\item \emph{Model side.} The same Benettin run on P2's model at $3\times$
  saves its latent basis, and the leading vectors are mapped to physical
  space through the decoder Jacobian.
\item \emph{Band fractions.} Both subspaces (the leading $n_{\mathrm{pos}}$
  vectors, orthonormalized) are split into streamwise bands:
  beyond-box ($m\in\{1,2\}$, $\lambda_x>2\pi$) and non-box
  ($m\not\equiv0 \bmod 3$). A check on synthetic fields returns 1.00 for a
  pure $m{=}1$ field and the expected 4/384 for noise.
\item \emph{Mechanism (b).} An autocorrelation part compares the mean
  energy and e-folding time of the $m{=}1,2,3$ modes along a true
  trajectory and a decoded model rollout.
\end{itemize}

\textbf{C8 interim results.}
\begin{itemize}[leftmargin=*]
\item \emph{Model side: complete.} P2 at $3\times$ gives
  $D_{\mathrm{KY}}=71.96$ with 34 positive exponents, reproducing the
  evaluation run. Its decoded unstable subspace is saved.
\item \emph{Mechanism (b): complete.} One true trajectory and one decoded
  model rollout at $3\times$, 100 time units each.
\end{itemize}

\begin{table}[h]
\centering
\small
\begin{tabular}{lccc}
\toprule
mode (wavelength) & energy, model / truth & e-folding time, truth & e-folding time, model \\
\midrule
$m{=}1$ ($6\pi$) & 0.82 & 4.8 & 4.2 \\
$m{=}2$ ($3\pi$) & 1.30 & 4.6 & 3.2 \\
$m{=}3$ ($2\pi$, box fundamental) & 1.69 & 5.4 & 4.6 \\
\bottomrule
\end{tabular}
\caption{C8 mechanism-(b) check, $\alpha{=}0.05$, $3\times$ (single
realizations, so indicative only).}
\label{tab:c8_autocorr}
\end{table}

The crop-trained model over-energizes the $2\pi$ (box-fundamental)
scale, under-energizes the domain scale, and its large-scale modes
decorrelate 15--30\% faster than the truth's. This is consistent with
mechanism (b): large-scale dynamics that are too local and too short-lived.
The truth-side band analysis waits for the $3\times$ truth run.

\textbf{C9 implementation.} \textbf{Status: arm (i) COMPLETED; arm (ii)
IN PROGRESS.} Results are in \S\ref{sec:c9_results}.
\begin{itemize}[leftmargin=*]
\item \emph{Velocity normalization.} The D2 backbone gains a velocity
  normalization (\texttt{biot\_savart\_vel\_scale}, set to the training
  set's own Biot--Savart velocity std) and a zero-initialized velocity
  embedding. The defaults reproduce D2.
\item \emph{Training-only velocity override.} In crop training
  (\texttt{BSContext}, patch-aligned crop bank), each rollout step places
  the model's predicted channel 0 into the free-interior sites of the
  truth anchor field of the \emph{whole $16\times2\pi$ domain}. It then
  solves Biot--Savart there and crops the velocity to the strip. At
  deployment the model solves on its own periodic domain from its own
  channel 0.
\item \emph{Exactness test.} One forced-halo strip step equals, at the
  free interior, the same weights' deployment step on the full periodic
  domain. A companion test confirms that a crop-only solve (what D2 would
  have done on a crop) gives a different velocity.
\item \emph{Arms.} (i) crop-trained on the wide $\alpha{=}0.05$ set,
  reusing P2's autoencoder, so only the propagator differs from P2;
  (ii) box-trained on the $2\pi$ set, reusing Section 231's autoencoder.
  Arm (ii) is D2 plus velocity normalization and zero-initialized
  embedding.
\item \emph{Evaluation.} Both arms are evaluated at $1/2/3\times$ with 3
  seeds and normalized initial conditions.
\end{itemize}

\subsection{C9 results (2026-10-03)}
\label{sec:c9_results}

\begin{table}[h]
\centering
\small
\begin{tabular}{lccc}
\toprule
size & C9(i) $D_{\mathrm{KY}}$ (seeds 0/1/2) & C9(i) mean error & P2 mean error \\
\midrule
$1\times$ & 21.3 / 23.3 / 23.7 & 3.5\% & 7.3\% \\
$2\times$ & 43.3 / 47.5 / 43.0 & \textbf{21.5\%} & 20.9\% \\
$3\times$ & 64.4 / 68.2 / 67.4 & 11.4\% & 6.3\% \\
\bottomrule
\end{tabular}
\caption{C9(i): crop-trained, with the exact global Biot--Savart velocity
input (truth 23.26 native, 56.84 at $2\times$).}
\label{tab:c9i}
\end{table}

\textbf{C9(i), pre-registered verdict: NO IMPROVEMENT.} The $2\times$
error is 21.5\%, against a partial-success threshold of 15\% and P2's
20.9\%. The exact non-local velocity, computed consistently between crop
training and deployment, sharpens native accuracy slightly but adds no
dimensions at $2\times$: the model remains extensive (about 22 per box).
Mechanism (c), the missing instantaneous coupling, is therefore
\emph{not} the main cause of the deficit, at least when the coupling
enters as an input to a learned local rule.

\textbf{C9(ii), box-trained with normalized velocity: FAILS under
transfer} ($1\times$ and $2\times$ complete; $3\times$ seed 0 done).
Native size is accurate (22.7, 22.2, 21.5; 2--7\% error). At $2\times$
all three seeds collapse to $D_{\mathrm{KY}}=2.9$, 3.0 and 3.1 (about
95\% error), and $3\times$ seed 0 collapses to 4.9. This is the same
failure as D2 and D3. Velocity normalization alone does not rescue a
box-trained model.

\textbf{Reading.} Across D2, D3, C9(i) and C9(ii), the dominant factor is
\emph{the training data}. Box-trained models collapse under transfer
whatever coupling they are given. Crop-trained models are stable and
extensive whatever coupling they are given. Neither learned consensus
(C7) nor exact non-local physics as an input (C9) closes the
super-extensive gap. That leaves two directions:
\begin{itemize}[leftmargin=*]
\item how training treats the large scales: forced halos and short
  horizons, i.e.\ mechanisms (a) and (b), tested by C11;
\item explicitly modelling the large-scale dynamics: the two-scale
  (C13--C15) and self-similar (C16--C18) routes.
\end{itemize}
C8's truth-side band analysis, due once the $3\times$ truth run
completes, decides between them.

\subsection{A two-scale model: predicting the large scales separately (C13--C15)}
\label{sec:two_scale}

\textbf{Motivation (2026-10-03).} Three lines of evidence point to the
missing dimensions being large-scale:
\begin{itemize}[leftmargin=*]
\item every working model is extensive, while the truth grows faster than
  the area;
\item C8's mode statistics (Table~\ref{tab:c8_autocorr}) show the
  crop-trained model under-energizes the domain scale, over-energizes
  the $2\pi$ box scale, and decorrelates its large-scale modes 15--30\%
  too fast;
\item C9(i)'s first $2\times$ seed shows that giving a learned local rule
  the exact global velocity as an \emph{input} does not recover them
  (23.7\% vs.\ P2's 20.6\% on the same seed; provisional, 1 of 3 seeds).
\end{itemize}
This motivates modelling the large scales as their own component, a
physically grounded version of the multiscale idea. It is not D3 again;
the differences are listed below.

\textbf{The physics obstacle: large scales are not closed in 2D.}
Kolmogorov flow is forced at $k{=}4$, and energy cascades
\emph{upward}, so the large scales are fed by small-scale eddies. The
coarse-grained vorticity $\bar\omega = G_\ell*\omega$ obeys
\[
\partial_t\bar\omega + \bar{\mathbf u}\cdot\nabla\bar\omega
 = \nu\nabla^2\bar\omega - \alpha\bar\omega + \bar f
   - \nabla\cdot\boldsymbol\tau,\qquad
 \boldsymbol\tau = \overline{\mathbf u\,\omega}-\bar{\mathbf u}\,\bar\omega ,
\]
where $\bar{\mathbf u}$ is the Biot--Savart velocity of $\bar\omega$.
Every term except the eddy forcing $-\nabla\cdot\boldsymbol\tau$ is
\emph{exact and cheap} on a coarse grid. The eddy forcing depends on the
unresolved scales and is partly random as seen from the large scales,
including backscatter. A large-scale-only model therefore needs a
closure for it, ideally with a stochastic part. This is the classical
subgrid-parameterization (LES closure) problem.

\textbf{What C9 and the earlier experiments teach, and how each lesson
enters the design.}
\begin{enumerate}[leftmargin=*]
\item \emph{The non-local part at large scales is known physics.}
  Advection by the Biot--Savart velocity, viscosity, drag and the
  projected forcing are all exact on a coarse grid. C9 suggests that
  supplying the exact velocity as an \emph{input} to a learned rule is
  not enough, because the rule must still learn large-scale advection.
  Supplying the exact \emph{dynamics} (C10, applied to the coarse field)
  is the stronger form, leaving only the eddy forcing to be learned.
\item \emph{Coarse dynamics transfer in $L$ for free.} The large-scale
  modes are exactly the ones that appear as the domain grows. A coarse
  solver in physical units handles new low-$k$ modes correctly at any
  size; only the local closure must transfer.
\item \emph{A-priori vs.\ a-posteriori training.} Our failures share
  forced halos and 12-step horizons, which is effectively \emph{a-priori}
  training. The data-driven closure literature for 2D and
  quasi-geostrophic turbulence reports that a-priori-trained closures can
  be unstable or inaccurate when run freely, and that \emph{a-posteriori}
  training through long free rollouts fixes much of this
  \cite{guan2022,frezat2022}. That is our mechanism (b), already known in
  this community.
\item \emph{A coarse model is cheap enough to train globally.} Cropping
  is needed only because the fine model cannot be rolled out over the
  whole $16\times2\pi$ domain. A coarse $\bar\omega$ model on that domain
  is tiny, so it can be trained with free rollouts of thousands of steps
  on the \emph{full} wide periodic domain: no halos, no crops, and
  statistics-based losses on the slow modes. This removes mechanisms (a)
  and (b) for exactly the part of the system where they bite. The fine
  model stays crop-trained and local, the regime where it already works
  (1--8\% native error in P1/P2).
\item \emph{Lessons from D3, C4 and D2 are respected.} The coarse
  dynamics are not learned from wraparound-dominated small-box data
  (D3's failure). The scales are not forced to agree by averaging (C4's
  failure). Inputs are normalized by large-domain statistics (D2's
  failure).
\end{enumerate}

\textbf{Design: a two-scale model.}
\begin{itemize}[leftmargin=*]
\item \emph{Scale split.} A fixed filter $G_\ell$ at a physical cutoff
  $\ell_c$ (primary choice: the $2\pi$ box scale) gives
  $\omega=\bar\omega+\omega'$.
\item \emph{Large-scale component.} An exact coarse Kolmogorov step for
  $\bar\omega$, plus a learned eddy-forcing closure. The closure is a
  local function of the coarse state and of pooled small-scale latents,
  with a learned noise term. It is trained a posteriori with long free
  rollouts on the full wide domain.
\item \emph{Small-scale component.} P2's crop-trained local latent model,
  additionally given the large-scale velocity and strain from the coarse
  component (C9's input, but taken from the model's own coarse state).
\item \emph{Two-way coupling.} Small to large through the eddy-forcing
  closure; large to small through sweeping and strain.
\end{itemize}
For data assimilation the state is (coarse field, fine latents). The
coarse block is small and globally coupled, and the fine block is local
and localizable, so a multiscale-localization scheme is the natural
filter (\S\ref{sec:c9_da}).

\textbf{Experiments (pre-registered 2026-10-03, before any run).}
\begin{description}[leftmargin=*]
\item[C13 --- closure diagnostic (no training, cheap).] \textbf{COMPLETED (linear closure): not viable} (\S\ref{sec:c8_c16_c13_results}). On the wide
  $\alpha{=}0.05$ (and $0.1$) datasets, filter at $\ell_c$ and compute the
  true large-scale tendency $\partial_t\bar\omega$. Spectral time
  derivatives are used where available; otherwise centred differences at
  $dt_{\mathrm{snap}}$, which is adequate for slow large scales but must
  be checked against a halved-$dt$ subset. Report the fraction of its
  variance explained by:
  \begin{enumerate}[label=(\roman*),leftmargin=*]
  \item the exact coarse PDE terms alone;
  \item (i) plus a local fitted closure (linear, then small CNN) in the
    coarse and small-scale state;
  \item the unexplained remainder, which is the noise the closure must
    carry.
  \end{enumerate}
  Repeat for $\ell_c\in\{\pi, 2\pi, 4\pi\}$.
  \emph{Rule:} a large-scale model is viable if (ii) explains at least
  80\% of the tendency variance at $\ell_c{=}2\pi$. A remainder above
  30\% means the closure must be stochastic.
\item[C14 --- large-scale-only model.] Exact coarse dynamics plus the
  learned closure (deterministic, then stochastic), trained a posteriori
  with free rollouts on the full $16\times2\pi$ domain. The loss
  combines short-horizon trajectory error with long-horizon statistics:
  the low-$k$ spectrum, mode autocorrelation and the PDF of the
  large-scale velocity. Evaluate against the truth at $1\times$,
  $3\times$ and $16\times$ on:
  \begin{itemize}[leftmargin=*]
  \item the low-$k$ energy spectrum;
  \item mode e-folding times (Table~\ref{tab:c8_autocorr});
  \item the large-scale share of the unstable subspace from C8.
  \end{itemize}
  \emph{Rule:} success if band energies are within 15\% and e-folding
  times within 15\% of the truth at every size.
\item[C15 --- coupled two-scale model.] C14's large-scale component
  coupled two-way to the crop-trained fine model. Report the full
  $D_{\mathrm{KY}}$ at $1/2/3\times$ with 3 seeds.
  \emph{Rule:} as C9, strong if the $2\times$ error is single-digit and
  partial if below 15\%.
\end{description}

\textbf{Order and gating.} C13 runs first, since it decides viability and
the cutoff. C14 and C15 proceed only if C13 passes. The whole direction
is conditional on C8: if C8 finds the missing dimensions are mostly
\emph{not} large-scale, its motivation disappears, and local-rule
accuracy (C11's statistics losses) takes priority. C9's remaining seeds
may also revise motivation point 3.

\subsection{Self-similarity across scales: one local rule at every level (C16--C18)}
\label{sec:self_similar}

\textbf{Hypothesis (raised 2026-10-03).} The chaotic dynamics may be
approximately self-similar across a range of spatial scales. If so,
large-scale dynamics could be generated by a coarser copy of the
\emph{same} local model, coupled to the finer copies below it, instead
of by a separately learned large-scale component.

\textbf{A clarification first.} The Kaplan--Yorke dimension measures the
attractor's size in state space. A fractal attractor is not by itself
evidence of self-similarity across \emph{spatial} scales. The two can be
linked quantitatively, though, by counting how many unstable dimensions
fall in each octave of wavenumber. C8's subspace analysis can measure
exactly that (C16).

\textbf{Why the hypothesis is physically motivated.}
\begin{enumerate}[leftmargin=*]
\item \emph{2D turbulence has a scale-invariant range.} Between the
  forcing scale ($k{=}4$ here) and the drag arrest scale $\ell_\alpha$,
  energy cascades upward in an inverse cascade with approximate scale
  invariance \cite{kraichnan1967,boffettaecke2012}. With constant energy
  flux $\varepsilon$, an eddy of size $\ell$ has velocity
  $u_\ell\sim(\varepsilon\ell)^{1/3}$, turnover time
  $\tau_\ell\sim\varepsilon^{-1/3}\ell^{2/3}$ and vorticity
  $\omega_\ell\sim\varepsilon^{1/3}\ell^{-2/3}$. Coarse-graining by a
  factor of 2 should therefore leave the statistics invariant after
  rescaling time by $2^{2/3}\approx1.59$ and vorticity by $2^{-2/3}$. A
  ``larger version of the same model'' is legitimate inside this range,
  \emph{provided} each level steps in rescaled time and amplitude. That
  rescaling is the ingredient C6 lacked: it shared weights across levels
  without it.
\item \emph{It fits the drag dependence we observe.}
  \begin{itemize}[leftmargin=*]
  \item At $\alpha{=}0.1$, $\ell_\alpha$ is small. Our domain sizes lie
    beyond the arrest scale, where large scales are drag-dominated and
    chaos is extensive, and a single-scale local model transfers (P1).
  \item At $\alpha{=}0.05$, the truth's $D_{\mathrm{KY}}$ grows faster
    than the area, which indicates our sizes are \emph{inside} the
    cascade range, where structures keep growing with $L$. That is the
    regime where scale-shared levels should help and a single-scale
    local model cannot, which is exactly where P2, C7 and C9 fall short.
  \end{itemize}
\item \emph{It is the only principled route for ``train small, capture
  emergent large-scale chaos''.} A model trained on a small domain has
  never seen scales larger than that domain. Any local model, C9 included,
  can extrapolate to new scales only if their dynamics follow rules it
  has already learned. Scale-sharing makes that assumption explicit and
  testable. The rule is learned from scale ratios available
  \emph{within} the small domain, and levels are added as $L$ grows, with
  no new parameters.
\item \emph{It has well-understood analogues.} Structurally, each level
  is a learned renormalization-group step (compare H3). Numerically it
  is multigrid or a fast multipole method with a learned local operator
  (\texttt{docs/spatially\_structured\_latent\_DA\_FMM\_research\_idea.md}).
  It is also the recursive, weight-shared form of the C15 two-scale
  model: coarse levels pass sweeping and strain down, and finer levels
  pass eddy forcing up through pooled latents.
\end{enumerate}

\textbf{Where it should fail, by construction.} Self-similarity breaks at
every characteristic scale: the forcing scale, the arrest scale
$\ell_\alpha$, and the box scale when a domain-scale condensate forms.
At Re $\approx100$, Kolmogorov flow may have a short inertial range,
perhaps a factor of 2--4 in scale. If C16 finds the range too short, the
direction reduces to C15's two-scale model with separate weights per
level. That is still useful, but it is not a scale-transfer result.
Where scale-sharing fails is itself a measurement of the characteristic
scales.

\textbf{Experiments (pre-registered 2026-10-03, before any run).}
\begin{description}[leftmargin=*]
\item[C16 --- self-similarity diagnostic (no training, cheap).] \textbf{COMPLETED: shape collapse yes, cascade time scaling no} (\S\ref{sec:c8_c16_c13_results}). On the
  wide $\alpha{=}0.05$ and $0.1$ datasets:
  \begin{enumerate}[label=(\roman*),leftmargin=*]
  \item the energy spectrum and its slope (the inverse-cascade
    prediction is $k^{-5/3}$) and the second- and third-order structure
    functions;
  \item a collapse test: block-average the vorticity by $\times2$ and
    $\times4$, rescale amplitude by $2^{-2/3}$ per level and time by
    $2^{2/3}$, and compare one-point PDFs, spatial correlation functions
    and temporal autocorrelations across levels;
  \item unstable dimensions per octave of $k$, from C8's truth basis,
    compared with the self-similar prediction.
  \end{enumerate}
  \emph{Output:} the scale range $[\ell_{\min},\ell_{\max}]$ over which
  rescaled statistics collapse within 15\%, at each $\alpha$.
  \emph{Rule:} C17--C18 proceed only if this range spans at least a
  factor of 4 at $\alpha{=}0.05$. If the fitted exponents differ from the
  Kolmogorov--Kraichnan values, the fitted ones are used.
\item[C17 --- zero-shot scale transfer.] The direct analogue of
  L-transfer, across scale instead of size. Feed P2's model with
  coarse-grained, amplitude-rescaled fields: block-average by 2, then
  re-grid to the native patch size. Compare its one-step and $k$-step
  predictions with truth sampled at the rescaled interval
  $0.2\times2^{2/3}\approx0.32$, against a model trained directly at that
  level as the reference. \emph{Rule:} scale transfer works if the
  zero-shot $k$-step error is within $1.5\times$ that of the
  level-trained reference.
\item[C18 --- recursive scale-shared model.] Levels $\ell_0, 2\ell_0,
  4\ell_0,\dots$, with the number of levels growing as
  $\log_2(L/\ell_0)$. One shared local rule plus per-level rescaling,
  either physics-set (Kolmogorov--Kraichnan) or learned as one scalar per
  level, with two-way inter-level coupling as in C15.
  \begin{itemize}[leftmargin=*]
  \item \emph{Key arm:} trained only on the $2\pi$ box, which is the
    train-small goal.
  \item \emph{Comparison arm:} trained on crops of the wide domain.
  \end{itemize}
  Evaluate $D_{\mathrm{KY}}$ at $1/2/3\times$ (3 seeds) and statistics at
  $16\times$. \emph{Rule:} as C9 (strong if the $2\times$ error is
  single-digit, partial if below 15\%). In addition, the box-only arm
  must come within 5 points of the crop arm for the train-small claim to
  hold.
\end{description}

\textbf{Order.} C16 first, since it is cheap and decides viability. Then
C17, which tests the core assumption with no new architecture. C18 only
if both pass. C16(iii) reuses C8's truth basis, so it follows the
$3\times$ truth run.

\subsection{Implications of C9 for latent data assimilation and L-transfer}
\label{sec:c9_da}

\textbf{The DA state is unchanged.} C9's velocity input is a
deterministic, differentiable function of the latent field's channel 0
(the coarse vorticity), not additional state. The DA state is still $z$,
and the forecast model is still $z\mapsto z$, so the existing machinery
applies unchanged:
\begin{itemize}[leftmargin=*]
\item the particle-flow filter and LETKF;
\item observation operators applied through the decoder;
\item the Lyapunov tools.
\end{itemize}
The added cost is one FFT on the coarse site grid per step, which is
negligible next to the network.

\textbf{What changes: Jacobian structure.} Without C9 the propagator is
strictly local, so $\partial z_{n+1}/\partial z_n$ is banded, and that
was the clean argument for localization. With C9, every site depends
weakly on every other site's coarse vorticity through the $1/k^2$ solve.
The velocity induced by a distant vortex decays as $1/r$, and only the
large scales carry it. The Jacobian becomes \emph{banded plus a smooth,
large-scale non-local part}. Operational DA already handles exactly this,
because pressure is global in atmospheric models too. The standard
remedies apply:
\begin{itemize}[leftmargin=*]
\item localize the analysis update rather than the dynamics;
\item use multiscale localization: a short radius for small scales, a
  longer one for the coarse anchor field.
\end{itemize}
The accurate description of the model is therefore ``local learned
dynamics plus one exact global linear solve,'' not ``strictly local.''

\textbf{Savings.}
\begin{itemize}[leftmargin=*]
\item \emph{State reduction} is a constant $\approx85\times$ at every
  domain size: $16{,}384\to192$ at $1\times$ and $32{,}768\to384$ at
  $2\times$ (Table~\ref{tab:capacity}).
\item \emph{Ensemble size} is the larger saving. Global latent DA without
  localization needs roughly $D_{\mathrm{KY}}$ members to span the
  unstable subspace, which grows with the domain (about 23 at $1\times$
  and 57 at $2\times$ for $\alpha{=}0.05$). With localization the
  requirement should follow the local dimension within one localization
  region, which does not grow with the domain (H4). A transferable local
  latent with a measured localization radius (P3/C12) is what makes this
  testable.
\item \emph{Forecast cost.} One surrogate step covers 0.2 time units,
  where the solver needs 20 steps at $128^2$.
\end{itemize}
\emph{Caveat:} 192 latents is about $8\times$ the attractor dimension at
$1\times$. The latent is sized for locality, not minimality, and that
trade pays off specifically in localized DA on large domains.

\textbf{L-transfer under C9.} The zero-shot mechanism is unchanged, with
one component now exact.
\begin{itemize}[leftmargin=*]
\item \emph{Learned parts.} The local mixer and the latent and velocity
  embeddings are per-site and shared. Deploying on a larger domain adds
  sites.
\item \emph{Non-local part.} The velocity is solved by FFT on the
  deployed domain, with $k$ in physical units. The new low-$k$ modes of a
  larger domain therefore enter automatically with correct physics. This
  is exactly what D2 got wrong through the config bug
  (\S\ref{sec:c8c9_status}).
\item \emph{Fixed constant.} The velocity normalization is fixed from the
  training set.
\end{itemize}
The remaining risk is \emph{off-distribution local states}: larger
domains grow larger structures with stronger large-scale flow than the
training data contain. C9's two arms test this directly. Arm (i) is
trained on crops of a large domain, whose data already contain such
velocities; arm (ii) is trained on the $2\pi$ box only and must
extrapolate.

For systems that cannot be re-simulated larger, such as CM1 or
observations, the arm (i) recipe applies to the largest affordable
domain. Transfer beyond it then rests on two things: the local rule,
monitored by the residual diagnostic, and the non-local solve, which is
exact by construction.

\textbf{DA after transfer.} Deploy the transferred model, localize with
the radius measured in physical units (P3/C12), and let the forecast
carry the large-scale coupling exactly. If C9 succeeds, this combination
is the payoff the program was designed for: a learned latent forecast
model whose filter cost per localization region does not grow with the
domain, while still representing domain-scale coupling.

\subsection{Corollary: patch-trained global models, and the path to RBC and CM1}
\label{sec:corollary_cm1}

\textbf{What is already shown.} At $\alpha{=}0.1$, a local model trained
only on $6\times2\pi$ crops of one large simulation reproduces the
$D_{\mathrm{KY}}$ of a $3\times2\pi$ periodic domain within 1.3\%, and the
native box within 9\% (\S\ref{sec:p1_results}). The global model is
assembled purely from a patch-trained local rule, and it captures the
large-domain chaos where a box-trained model missed by 19\%. At
$\alpha{=}0.05$ the same construction misses by 21\%, and that
remainder is the non-local part C9--C11 target. If C9/C10 close it, the
claim becomes: \emph{a local rule learned from windows, combined with an
exact solve of the one non-local term, reproduces the emergent
large-domain chaos.}

\textbf{Why this matters beyond Kolmogorov flow.}
\begin{itemize}[leftmargin=*]
\item Training cost and data requirements scale with the patch, not the
  domain.
\item Observational and reanalysis data arrive as windows.
\item A patch-trained local model is the natural object for localized
  ensemble data assimilation, since its Jacobian is banded and its
  localization radius is measured (P3/C12).
\end{itemize}

\textbf{Rayleigh--B\'enard convection (next system).} Boussinesq,
incompressible, with an elliptic pressure solve: structurally the same
problem as Kolmogorov flow. The C9/C10 template carries over directly:
exact pressure solve plus a learned local closure. Crop horizontally and
keep full vertical columns between the walls, as we keep full periodic
$y$-strips here. A Rayleigh--B\'enard demo dataset already exists in
\texttt{artifacts/datasets/}. The emergent large-domain phenomenon to
target is the growth of large-scale circulation and superstructures with
aspect ratio.

\textbf{CM1 \cite{bryanfritsch2002}.} CM1 integrates the compressible
non-hydrostatic equations. Pressure is prognostic, so information
travels at finite (acoustic) speed. The dynamics are therefore strictly
local, with a physical light cone of $c_s\Delta t$ per step; in that
sense CM1 is more favourable than an incompressible flow. The flagship
emergent target is convective self-aggregation, which appears only above
a domain-size threshold \cite{mullerheld2012,jeevanjeeromps2013}. The
test would be:
\begin{enumerate}[leftmargin=*]
\item train on column patches of a large radiative--convective
  equilibrium simulation;
\item deploy to larger domains;
\item check whether aggregation emerges at the right size.
\end{enumerate}
This is the \S\ref{sec:transfer_diagnostic} diagnostic and the corollary
in one experiment.

\textbf{CM1 caveats.}
\begin{itemize}[leftmargin=*]
\item \emph{Vertical structure.} It is non-periodic and must stay whole
  in each column.
\item \emph{Inhomogeneous lower boundaries.} The local rule must be
  conditioned on static fields.
\item \emph{Moist thresholds.} Convective triggering is intermittent and
  heavy-tailed.
\item \emph{Radiation and gravity waves.} Gravity waves are fast and long
  range, though still finite speed.
\item \emph{Data volume.}
\end{itemize}

\textbf{Conditions for the corollary to transfer to a new system.}
\begin{enumerate}[leftmargin=*]
\item local dynamics, either finite propagation speed or an exactly
  solvable non-local term;
\item statistical homogeneity within the training windows, or explicit
  conditioning on what is not homogeneous;
\item crops larger than the interaction reach (P3/C12 measure this);
\item at least approximate extensivity at the crop scale.
\end{enumerate}
When condition 4 fails, the size of the transfer error is itself the
measurement of how non-local the emergent dynamics are.

%======================================================================
\section{Status, publication plan and decision tree (updated 2026-10-04)}
\label{sec:status_plan}
%======================================================================

\subsection{What is established}
\begin{enumerate}[leftmargin=*]
\item \textbf{1D KS:} clean zero-shot L-transfer (under 2\% in
  $D_{\mathrm{KY}}$ to $4\times$), plus localized latent DA.
\item \textbf{2D Kolmogorov, $\alpha{=}0.1$:} a model trained on windows
  of a large domain transfers to $3\times$ with 1.3\% error (P1).
\item \textbf{2D Kolmogorov, $\alpha{=}0.05$:} the same recipe gives
  6.3\% at $3\times$ and about 7\% native (P2). Box-trained models
  collapse at $2\times$ (47--99\%).
\item \textbf{2D KS (P7, 2026-10-04):} training on 2D crops, cut in
  \emph{both} directions from a $6\times$-per-axis domain, gives 1.2\%
  native and 5.7\% at $66^2$ (seed 0), against 13.5\% and 24.4\% for
  the identical box-trained architecture (\S\ref{sec:p7_results}). This
  is the first test in which no training sample contains a whole
  periodic direction. The remaining error is consistent with the truth's
  finite-size super-extensivity.
\item \textbf{The training distribution decides success, not the
  architecture}, now shown in two 2D systems and for \emph{native} as
  well as transfer error. Every box-trained model collapses under transfer,
  whatever coupling it has (D2, D3, C9(ii)). Every window-trained model is
  stable. Learned consensus (C7) and the exact non-local velocity (C9)
  never helped.
\item \textbf{Scale content is right.} The chaos lives at the forcing
  scale, and the window-trained model reproduces both its dimension and
  its scale distribution (C8, C16).
\item \textbf{The capacity concern (\S\ref{sec:capacity}) is not the
  bottleneck for 2D KS.} P7 removed the native undershoot that motivated
  it with the same $c{=}3$; the undershoot came from the training box,
  not from too few channels.
\item \textbf{Negative and limiting results.} P6's narrow architecture
  (encoder reach 2) fails even on the full periodic domain, which
  implicates encoder reach, not windowing (\S\ref{sec:p6_results}). Drag-free
  Kolmogorov flow saturates into a domain-scale condensate
  (\S\ref{sec:alpha0_results}), a size-dependent mean state that a local
  model cannot be expected to transfer.
\item \textbf{Training stability is a first-class failure mode
  (P9/P9b, 2026-10-05).} Short-horizon teacher-forced losses do not
  certify free-running dynamics: P9 drifted into spurious, IC-dependent
  regimes even in its own geometry while the true flow is not
  multistable. Longer free rollouts in training (P9b) cut the $3\times$
  seed spread from 22\% to 2\%, leaving a consistent 11--15\%
  undershoot. Free-rollout statistics (variance, energy split by scale)
  are a cheap diagnostic that catches this before Lyapunov evaluation.
\item \textbf{The autoencoder needs capacity, not spatial context
  (C19, verified 2026-10-05).} A per-patch autoencoder (encoder and
  decoder reach 0) with $c{=}6$ matches P2's full-reach $c{=}3$ model at
  $1\times$/$2\times$ and beats it at $3\times$ (three-IC-seed mean errors
  8.7/8.3/2.9\% against 7.3/9.7/6.3\%; a second training seed gives
  3.9/2.7/3.0\%), and removes the autoencoder's crop margin entirely. It
  is now the canonical autoencoder (\S\ref{sec:phase0_results}).
\item \textbf{Native-size ($1\times$) behaviour.} Per-IC truths show the
  models track the true IC dependence but amplify it ($\pm10\%$ against
  $\pm4\%$) and overshoot on average (per-IC mean error 7--9\%). The old
  single-run native truth (23.26) was about 8\% high; the three-IC mean
  is 21.54. Long-rollout training (Phase 1) is the test of whether this
  is a propagator effect on the 8-site box.
\item \textbf{Integrity.} The decision rules were pre-registered. Two
  bugs that changed conclusions were found and corrected: unnormalized
  transfer ICs (\S\ref{sec:icnorm}) and dropped config fields
  (\S\ref{sec:c8c9_status}).
\end{enumerate}

\subsection{Assessment}
\textbf{Good news.} The core thesis holds more broadly than expected. A
local latent field trained on windows of a large domain transfers
zero-shot, even at weak drag, where box training fails completely. P7
upgrades ``windows'' from strips to genuinely 2D patches in an isotropic
system, and the contribution can be stated as a general training
paradigm (patch training on a large domain) for that class of system.
P9 shows the paradigm does not automatically carry over: 2D patches of
a tall Kolmogorov domain fail on the standard $2\pi$-tall boxes at
$3\times$. Whether that is the forcing, a training/deployment geometry
mismatch, or seed scatter is open (\S\ref{sec:p9_results}). That
is a clean and practically relevant lesson, and it fits the setting
where re-simulation is impossible (CM1, observations).

\textbf{Weaker than hoped.} The ``measurable non-local remainder'' story
rested largely on a single $2\times$ truth value that may be an outlier
(\S\ref{sec:truth3_revision}). If that is confirmed, the
physics-diagnostic angle needs a regime where locality genuinely breaks,
which the drag sweep (P8, $\alpha{=}0$) is designed to find.

\subsection{Publication plan}
\textbf{Paper 1} (core result, close to complete):
\emph{Local latent fields for spatiotemporal chaos: zero-shot domain
transfer, training on windows, and localized data assimilation.}
\begin{itemize}[leftmargin=*]
\item 1D KS: clean transfer plus DA.
\item 2D Kolmogorov: transfer works when chaos is extensive, and window
  training is essential while box training collapses.
\item 2D KS: 2D patch training fixes both native and transfer error
  (P7), giving a second, isotropic 2D system. P9 (\S\ref{sec:p9_results})
  did not carry over to Kolmogorov when deployed on standard $2\pi$-tall
  boxes ($42.6\%$ at $3\times$, seed 0); the cause is not yet
  identified.
\item The negative results: consensus and exact non-local inputs do not
  help.
\item Lyapunov-level validation throughout, which is unusual in this
  literature.
\end{itemize}
Still needed for paper 1:
\begin{itemize}[leftmargin=*]
\item \textbf{canonical model decided} (per-patch $c{=}6$,
  \S\ref{sec:phase0_results}); the full-reach $c{=}6$ control's
  $2\times$/$3\times$ evaluations will finish the attribution;
\item multi-seed truths at $1\times$ (and ideally $2\times$) from the
  \emph{same} ICs as the model evaluations, so native and $2\times$ errors
  can be scored per IC;
\item a second training seed for P1 and P2, which have one each;
\item P7 evaluation seeds 1 and 2 (only seed 0 was run, to limit machine
  load); optionally an $88^2$ 2D-KS truth to test whether the
  finite-size error shrinks with size;
\item truth error bars ($\alpha{=}0.05$ $2\times$ seeds running; T1 at
  $\alpha{=}0.1$ deferred);
\item a literature search on patch-trained surrogates;
\item a \textbf{2D DA demonstration} with a transferred model. DA was the
  original goal and is so far demonstrated only in 1D.
\end{itemize}

\textbf{Paper 2:} the drag sweep and transfer-error diagnostic, and the
application to Rayleigh--B\'enard and CM1
(\S\ref{sec:corollary_cm1}). Its practical claim, training on the one
domain you have, needs crops much smaller than the domain, and so needs
small autoencoder margins. C19 (\S\ref{sec:c19_results}) is the first
evidence: decoder reach 0 leaves transfer intact but lowers the
dimension per site by about 9\% at $c{=}3$. The per-patch $c{=}6$ arm
and a P6 redo are prerequisites for that claim.

\textbf{Supplementary material for paper 1:} C19 as a design ablation.
The crop margin equals the model's receptive field, whatever the
architecture; decoder context buys native accuracy (3.0\% against
8.3\%) but not transferability.

\subsection{Why a C19-style model should suit data assimilation:
hypotheses and tests (2026-10-06, speculative)}
\label{sec:c19_da_hypotheses}

User-requested speculation on how the per-patch (C19) architecture should
help DA, especially in 2D and on large domains. \textbf{Nothing in this
subsection is a result.} Each item is a hypothesis with a mechanism, a
prediction, and a test that could refute it. The only established fact
used is the block-diagonal decoder Jacobian (\S\ref{sec:da_2d}, exact,
unit-tested).

\paragraph{Reference numbers for 2D KS (\texttt{p7b\_ks2d}).} Patch $16$
pixels, site length $h{=}44/40{=}1.1$, $c{=}6$, so
$d_{\mathrm{latent}}{=}6P$ with $P$ sites ($9600$ at $44^2$, $21{,}600$ at
$66^2$, $38{,}400$ at $88^2$). True $D_{\mathrm{KY}}$ is $81.8$ at $44^2$
and $192.5$ at $66^2$, a density of about $0.043$ per unit area, or
$0.05$ per site. The default observation grid (stride 32 pixels, i.e.\
one observation per $2{\times}2$ block of sites) observes one site in
four directly.

\subsubsection*{The core argument}
In physical-space LETKF the observation operator is local by
construction: a gauge reads one grid point. In latent DA it is
$h(D(z))$, and for a global or full-reach decoder one physical
observation depends on many latent sites. Localization in latent space
then approximates a non-local operator, and the approximation worsens as
decoder reach grows. With decoder reach 0, $h\circ D$ is exactly as local
as in physical space: each observation depends on one site's $c$
channels. A C19 model therefore lets the latent filter inherit, without
approximation, the structure that makes operational local filters
scale: independent local analyses, cost linear in area, and an ensemble
size set by local rather than global dimension. The propagator remains
non-local over a forecast window, but that dependence enters only
through the ensemble covariance, which is where localization is
designed to act.

\subsubsection*{Hypotheses}
\begin{description}[leftmargin=*]
\item[H1. A true local (Hunt et al.\ 2007) LETKF is exact for this
  architecture.] Each site's analysis needs only the observations inside
  its localization disk, and the observation-space quantities for those
  observations depend only on sites inside the disk. \emph{Prediction:}
  at native size, a per-site LETKF and the current dense
  covariance-localized EnKF with the same Gaspari--Cohn radius give
  analysis RMSE within a few percent of each other. For a full-reach
  decoder (reach 6), truncating each local problem to the disk discards
  part of $h\circ D$, so the gap should grow with decoder reach.

\item[H2. Analysis cost scales linearly with area, and the dense
  implementation does not scale at all.] The current filter
  (\texttt{ks\_latent/da/letkf.py}) forms a dense $d\times d$ forecast
  covariance and an equally large taper. In float64 that is $0.74$\,GB
  per matrix at $44^2$, $3.7$\,GB at $66^2$, and $11.8$\,GB at $88^2$,
  before temporaries. \textbf{The transfer DA gate in
  \S\ref{sec:consolidation} Phase 4 therefore probably cannot be run
  at $66^2$ or above with the current code.} A per-site LETKF costs
  $P$ independent problems of size $N_{\mathrm{ens}}\times
  N_{\mathrm{ens}}$, each with $O(10)$ observations, so cost is
  $O(P\,N_{\mathrm{ens}}^3)$ and parallel across sites.
  \emph{Prediction:} wall time and memory proportional to $P$ from
  $44^2$ to $88^2$ and beyond.

\item[H3. The observation Jacobian costs $c$ forward passes, whatever
  the domain size.] Because the decoder Jacobian is block-diagonal and
  identical in structure at every site, one forward-mode JVP with the
  tangent ``unit vector in channel $k$ at every site at once'' returns
  column $k$ of every site's block simultaneously (Jacobian colouring,
  Curtis--Powell--Reid). The full $(m,d)$ Jacobian then costs $c{=}6$
  JVPs, against $m$ reverse passes now ($m{=}400$ at $44^2$, growing
  with area). The propagator's Jacobian is banded and colours the same
  way with $(2r_p{+}1)^2c$ colours, which matters for PFF, 4D-Var and
  Lyapunov analysis. \emph{Prediction:} the coloured Jacobian equals
  \texttt{jacrev} to round-off and its cost is flat in domain size.
  \textbf{Correction after measurement (\S\ref{sec:da_2d}'s T1
  experiment, 2026-10-06):} $m$ does NOT grow with area under the
  deployment convention actually used (observation stride scales with
  the domain to hold physical density fixed, so $m{=}400$ at every
  size tried). The predicted speedup is confirmed regardless (it grows
  with size, as stated), but for a different mechanism: \texttt{jacrev}'s
  fixed-$m$ reverse passes each get more expensive as the decoder
  itself grows, while the coloured method's fixed 6 passes absorb that
  same growth far more slowly. Right prediction, wrong stated reason --
  see the T1 experiment write-up for the numbers.

\item[H4. The ensemble size needed is set by local, not global,
  dimension.] In extensive chaos the number of unstable directions grows
  with area (82, 193, about 330 at $44^2$, $66^2$, $88^2$). An
  unlocalized filter needs $N_{\mathrm{ens}}$ of that order. A
  localized one needs $N_{\mathrm{ens}}$ of order the dimension inside one
  localization disk. With $c_{\mathrm{GC}}{=}3$ sites (support radius
  $6.6$ units, area about $137$), that is about $0.043\times137\approx6$.
  \emph{Prediction:} with localization, the $N_{\mathrm{ens}}$ at which
  analysis RMSE saturates is between 16 and 32 and the same at every
  size. Without localization it grows roughly in proportion to area,
  and at fixed $N_{\mathrm{ens}}$ the unlocalized filter degrades as the
  domain grows. This is the 2D version of brief \S15's first headline
  experiment.

\item[H5. Transfer costs the filter almost nothing.] Model, observation
  operator and localization are all local, so a filter at $2\times$ is
  the native filter applied to more sites. \emph{Prediction:} with the
  same $N_{\mathrm{ens}}$, $c_{\mathrm{GC}}$ and observation density,
  analysis RMSE per grid point at $66^2$ and $88^2$ is within the 25\%
  gate of native, and closer than that if the model's own transfer
  error is small (P7: 5.7\% in $D_{\mathrm{KY}}$ at $66^2$).

\item[H6. Risk: seams at patch edges.] A per-patch decoder has no
  mechanism that smooths across patch boundaries. Analysis increments
  that leave the decoder's data manifold may decode to fields that are
  discontinuous at patch edges, putting spurious energy at high
  wavenumbers that the propagator then has to absorb. A full-reach
  decoder would hide this. \emph{Prediction (as a risk):} the
  gradient jump across patch edges, relative to the gradient inside
  patches, is larger in analyses than in truth or in free forecasts, and
  grows as $N_{\mathrm{ens}}$ falls. If it appears, a re-encoding step
  $z\leftarrow E(D(z))$ after each analysis should remove it without
  costing accuracy.

\item[H7. Unobserved sites are corrected only through the ensemble
  covariance, and that is an advantage.] Three sites in four have no
  observation in their own patch, so their update comes entirely from
  flow-dependent cross-site covariances. A full-reach decoder also
  spreads an observation through $h\circ D$ itself, a fixed,
  climatological spreading. \emph{Prediction:} per-patch matches or beats
  full-reach at moderate $N_{\mathrm{ens}}$ ($\ge32$); full-reach may win
  at very small $N_{\mathrm{ens}}$, where its built-in spreading acts as a
  static prior. Unobserved-site RMSE should improve as $c_{\mathrm{GC}}$
  grows until it covers the observation spacing, then degrade as
  sampling noise enters.

\item[H8. The localization radius is set by the forecast window, not
  the autoencoder.] With encoder reach 0, the brief's bound
  $r_{\mathrm{loc}}\ge\max(\text{encoder receptive field},
  v_*\Delta t)$ reduces to its physical light-cone term over the
  assimilation window. \emph{Prediction:} the optimal
  $c_{\mathrm{GC}}$ grows with \texttt{n\_prop\_steps}, and at fixed window
  is the same for every model with zero autoencoder reach.

\item[H9. The doubled channel count is a moderate cost.] $c{=}6$ has
  about 100 latent variables per unstable direction locally
  ($0.05$ per site), so the forecast ensemble occupies a thin subspace of
  each local latent block. Ensemble methods handle that, but sampling
  noise in the within-site channel covariances grows with $c$.
  \emph{Prediction:} at matched localization, the $c{=}6$ per-patch
  model's RMSE is within about 10\% of a $c{=}3$ full-reach model at
  large $N_{\mathrm{ens}}$, and any gap widens at small $N_{\mathrm{ens}}$.
  This is the Phase 4 comparison already planned in
  \S\ref{sec:consolidation}.

\item[H10. Kolmogorov flow may want anisotropic localization.] Its
  small scales are locally close to isotropic, but the forced jets
  correlate along $x$ over long distances. \emph{Prediction:} an
  elliptical taper ($c_x>c_y$) beats the isotropic one, by an amount
  that grows as drag falls.
\end{description}

\subsubsection*{Tests}
In priority order. T1 and T2 are code; the rest are runs.
\begin{enumerate}[leftmargin=*,label=T\arabic*.]
\item \textbf{Coloured observation Jacobian} (H3). Implement it, assert
  equality with \texttt{jacrev} on \texttt{p7b\_ks2d}, and time both at
  $44^2/66^2/88^2$. \emph{Refuted} if the two disagree beyond round-off,
  or if cost is not flat in size. Small; also speeds up the existing
  dense filter.
\item \textbf{Per-site local LETKF} (H1, H2). A new analysis class with
  the same \texttt{analyze} interface: for each site, gather observations
  within $2c_{\mathrm{GC}}$, weight them by Gaspari--Cohn ($R$
  localization), solve in ensemble space, and update that site's $c$
  channels. Validate against the dense filter at native size. \emph{H1
  refuted} if RMSE differs by more than 5\% at $N_{\mathrm{ens}}{\ge}32$.
  \emph{H2 refuted} if wall time is clearly superlinear in $P$. This is
  required before any transfer DA run in 2D, because of the memory
  numbers under H2.
\item \textbf{$N_{\mathrm{ens}}\times$ size $\times$ localization sweep}
  (H4, H5). $N_{\mathrm{ens}}\in\{8,16,32,64,128\}$ at $44^2$, $66^2$,
  $88^2$, with and without localization. Truth spinup paid once per
  size. Report analysis RMSE, spread/RMSE, and free-run RMSE. \emph{H4
  refuted} if, with localization, the saturating $N_{\mathrm{ens}}$
  grows with size, or if without localization it does not. \emph{H5
  refuted} if transferred RMSE exceeds native by more than 25\% at the
  same settings.
\item \textbf{Seam diagnostic} (H6). For analysis, forecast and truth
  fields, compute the ratio of mean $|\nabla u|$ across patch edges to
  that inside patches, and the energy above the patch wavenumber
  $2\pi/h$. Repeat with re-encoding. \emph{H6 refuted} (good news) if
  analyses match truth to within the forecast's own deviation.
\item \textbf{Per-patch versus full-reach decoder} (H1, H7, H9). Same
  filter, same $c_{\mathrm{GC}}$, \texttt{p7b\_ks2d} against a full-reach
  2D-KS model, split RMSE into observed and unobserved sites, sweep
  $N_{\mathrm{ens}}$ and $c_{\mathrm{GC}}$. Needs a full-reach 2D-KS
  model trained on the same crops; the existing box-trained one would
  confound the comparison.
\item \textbf{Localization radius versus window} (H8).
  \texttt{n\_prop\_steps}$\in\{2,5,10\}$ $\times$
  $c_{\mathrm{GC}}\in\{1.5,2,3,4,6\}$ at native size, $N_{\mathrm{ens}}{=}32$.
  \emph{Refuted} if the optimal radius does not grow with the window.
\item \textbf{Kolmogorov, anisotropic taper} (H10). After T2, a
  separable taper with $c_x\ne c_y$ on the $x$-strip model at
  $\alpha{=}0.05$ and $0.025$. Lowest priority; only if Kolmogorov DA is
  in the paper.
\end{enumerate}

\subsubsection*{T1 and T2: implemented and unit-tested (2026-10-06)}
User-directed (``write code for T1 and T2''). Both are new,
reviewer-facing capabilities, not just speculation anymore -- though the
real $N_{\mathrm{ens}}\times$size sweep (T3) that would actually confirm
or refute H1/H2/H4/H5 is still TO DO, for the same reason as
\S\ref{sec:da_2d}: it is a real cost to commit to, not a unit test.

\textbf{T1 (\texttt{ks\_latent/da/jacobian\_coloring.py}).}
\texttt{coloured\_obs\_jacobian\_2d\_compact} reconstructs the exact
decoder Jacobian from $c$ forward-mode \texttt{jvp} passes (one per
latent channel, the tangent a one-hot vector in that channel at
\emph{every} site at once) instead of $m$ reverse-mode passes. This
works because the colouring argument is airtight, not approximate: each
output row is owned by exactly one site (block-diagonality, already
proved by the existing unit test in \S\ref{sec:da_2d}), so a tangent
that is nonzero at every site's channel $k$ produces, at that row, only
its own site's contribution -- every other site's component of the
tangent multiplies an exactly-zero Jacobian entry. 10 unit tests: the
coloured Jacobian agrees with \texttt{jacrev} to $<10^{-9}$ at two
observation densities (\texttt{stride}$\in\{2,4\}$); the compact and
dense forms agree and the dense form's off-block entries are exactly
\texttt{0.0}, not merely small; \texttt{assert\_decoder\_zero\_reach}
passes for zero radius, passes for zero layers at nonzero radius, and
raises for genuine nonzero reach; a direct call-counting test wraps the
observation operator and confirms it is invoked exactly \texttt{channels}
times (not $m$) at \emph{both} a dense and a sparse observation grid --
the literal content of H3's ``flat cost'' claim, demonstrated rather than
timed. All 10 pass.

\textbf{T2 (\texttt{ks\_latent/da/letkf\_local.py}).}
\texttt{LETKFLocal2D} is a genuine per-gridpoint local filter (Hunt et
al.\ 2007's own formulation, R-localization rather than the dense
filter's Schur-product covariance localization -- the naming distinction
\texttt{ks\_latent/da/letkf.py}'s own docstring already flags as
missing). Per site: gather the (small, compact-support-bounded) set of
observations with nonzero Gaspari-Cohn weight, build the LOCAL slice of
the ensemble covariance by slicing columns of the perturbation matrix
(never forming a $(d,d)$ object), get the Jacobian from T1's compact
form, solve a small linear system, and update only that site's $c$
channels. 4 unit tests: matches the analytic Kalman posterior (mean/cov
to 5\%, $N{=}8000$) with a wide taper on a synthetic block-diagonal
decoder; shape/finiteness; a loud \texttt{ValueError} on a dimension
mismatch; and a compact-support test stronger than the dense filter's
own localization test (\S\ref{sec:p6_results}'s sibling,
\texttt{test\_letkf\_localization\_shrinks\_far\_correlations}, only
shows a \emph{smaller} shift) -- with a tight taper a site at maximum
torus distance from the only observation sees \emph{exactly} \texttt{0.0}
shift, and a real nonzero shift with a wide one. All 4 pass. Wired into
\texttt{ks\_latent.da.cycling.run\_da\_experiment} as
\texttt{CycleConfig.method="letkf\_local2d"} (2 new dispatch tests:
beats the free run end to end on a block-diagonal linear toy system, and
raises loudly if any of the five required geometry fields is missing)
and into \texttt{scripts/run\_da\_pff\_2d\_ltransfer.py} as
\texttt{--method letkf\_local2d} (asserts
\texttt{assert\_decoder\_zero\_reach} on the loaded checkpoint before
running, so a full-reach checkpoint fails loudly rather than silently
producing a wrong analysis).

\textbf{Smoke-validated end to end (2026-10-06), same caveat as
\S\ref{sec:da_2d}'s dense-filter smoke test.} \texttt{p7b\_ks2d} at
native size ($d_{\mathrm{latent}}{=}9600$, $40^2$ sites), 5-unit spinup,
8 ensemble members, 2 cycles of 2 steps, \texttt{--method
letkf\_local2d}: 400 observations identified correctly from the obs
grid, ran to completion with no crashes or NaNs, valid JSON. \textbf{The
reported numbers are meaningless} (same reason as before: far too short
a spinup to be on the attractor) and were discarded.

\textbf{What T1/T2 did NOT yet establish, as of the previous entry.}
Neither test suite measured wall time or memory against the dense
filter (H2's actual claim), nor ran the per-site filter against the
dense one on real $p7b\_ks2d$-scale data (H1's actual claim) -- both
unit suites only validated correctness against analytic ground truth on
small synthetic systems, necessary but not sufficient for H1/H2. Closed
below (user-directed, same day: ``set up and start the T1 and T2
experiments'').

\textbf{CORRECTION (2026-10-07), applies to every 2D DA run below (T1--T5,
and E1/E4 in \texttt{docs/paper1\_tasks.md}).}
\texttt{scripts/run\_da\_pff\_2d\_ltransfer.py} takes its native model
configuration from the stored crop checkpoint ($N_X{=}640$, $40\times40$
sites) and pairs it with $L{=}44$, so each deployed site spans $1.1$ units.
The model was trained on 16-point patches of $5.5$ units (wide domain:
264 units on 768 points). \texttt{exp\_crop\_eval\_ks2d.py} deploys
correctly ($L{=}44$ on $128^2$, $8\times8$ sites). Every 2D DA run below
therefore applied the model to fields resolved $5\times$ finer than its
training data. The structural results stand (exact block-diagonal Jacobian,
colouring speed-up as a function of $d$, dense-filter memory failure at
$d{=}38{,}400$), but every RMSE, skill and calibration value below is
\emph{not} a valid measurement of the trained model and must be re-run at
a correct deployment ($L{=}220/330/440$ for $40^2/60^2/80^2$ sites, which
keeps $d$ unchanged; task E11).
\textbf{Second correction (same day).} Both ensemble filters returned analyses
still attached to the autograd graph of the observation-Jacobian computation, and
\texttt{run\_da\_experiment} stored them, so every cycle's graph was retained
($\approx0.24$\,GB per cycle at $d{=}384$). Fixed by detaching (regression tests
added). Computed values are unaffected, but every memory and timing figure below
is contaminated, including the dense filter's $88^2$ failure, which should be
stated as memory accounting ($11.8$\,GB per $d\times d$ float64 matrix at
$d{=}38{,}400$) rather than as an observed kill. The leak also explains the
E6/E1c memory blow-ups.

\subsubsection*{T1 experiment: real timing on \texttt{p7b\_ks2d} (2026-10-06)}
\texttt{scripts/exp\_t1\_jacobian\_timing.py}: at native size and two
zero-retraining transfer sizes ($66^2$, $88^2$), build the observation
Jacobian both ways and compare.

\begin{center}
\begin{tabular}{rrrrrr}
\toprule
$L$ & $d_{\mathrm{latent}}$ & $m$ & \texttt{jacrev} & coloured & speedup \\
\midrule
44 & 9{,}600  & 400 & 0.150 s & 0.023 s & 6.4$\times$ \\
66 & 21{,}600 & 400 & 0.477 s & 0.044 s & 10.8$\times$ \\
88 & 38{,}400 & 400 & 0.994 s & 0.065 s & 15.3$\times$ \\
\bottomrule
\end{tabular}
\end{center}
Exact agreement at every size (max abs.\ error $2$--$4\times10^{-16}$,
float64 roundoff).

\textbf{A real finding, and a correction to H3's own stated mechanism.}
H3 predicted the speedup comes from avoiding $m$ passes that grow with
domain area. \textbf{$m$ does NOT grow here} -- the deployment
convention already built (\texttt{run\_da\_pff\_2d\_ltransfer.py} scales
\texttt{obs-stride} with the domain so physical observation DENSITY
stays fixed under transfer) holds $m=400$ exactly constant at every
size. What actually drives the growing speedup is a different, equally
real mechanism: \texttt{jacrev}'s $m$ reverse passes each cost more as
the decoder itself grows (a \texttt{ConvTranspose2d} over a bigger
physical grid), so its wall time grows with $d_{\mathrm{latent}}$ even
at fixed $m$ (0.15s$\to$0.99s, a $6.6\times$ growth from $44^2$ to
$88^2$); the coloured method pays that same per-pass growth, but only
over a FIXED 6 passes, so its own time grows far more slowly (0.023s
$\to$0.065s, $2.8\times$) and the ratio between the two widens with
size. \textbf{H3 is confirmed as stated (speedup grows with transfer
size), for a mechanism the hypothesis did not correctly name} -- the
pre-registered prediction and the measurement agree; the stated reason
did not, and is corrected here rather than quietly fixed.

\subsubsection*{T2 experiment: dense vs.\ per-site local LETKF on real data (2026-10-06)}
Same checkpoint (\texttt{p7b\_ks2d}), same seed, same everything except
the analysis method: \texttt{scripts/run\_da\_pff\_2d\_ltransfer.py
--method letkf} vs.\ \texttt{--method letkf\_local2d}, native size
($44^2$, $d_{\mathrm{latent}}{=}9600$), a real 200-unit spinup,
$N_{\mathrm{ens}}{=}64$, 20 cycles.

\textbf{A timing artifact caught and corrected before it was written up
wrong.} The first (dense) run's wall-clock time, 1282s, included the
machine going to sleep mid-run -- caught by the user, not discovered
independently. \textbf{The draft that treated 1282s as evidence for
H2's dense-matrix cost concern was wrong and is retracted here}, not
quietly removed. The second (local) run was protected with
\texttt{caffeinate -i -w <pid>} (held awake for exactly as long as that
process runs, not a guessed timeout) and its wall time, 480s, is clean.
Comparing CPU time (\texttt{user+sys}, immune to the sleep interval
either way) rather than wall time: dense $412{+}121{=}533$s, local
$408{+}108{=}516$s -- \textbf{statistically indistinguishable at this
size.} The real finding is therefore the opposite of first appearances:
\textbf{at native size, the dense filter's extra per-cycle
matrix-formation cost is a small fraction of the (shared, $\sim$8-minute)
spinup cost, not yet a bottleneck.} H2's actual memory/cost concern was
always about TRANSFER (3.7 GB at $66^2$, 11.8 GB at $88^2$ per dense
matrix) -- this comparison, at native size only, cannot speak to that;
T3 (native-vs-transfer, not yet run) is still the test that matters for
H2.

\textbf{Accuracy comparison (H1).}
\begin{center}
\begin{tabular}{rrrr}
\toprule
 & dense (\texttt{letkf}) & local (\texttt{letkf\_local2d}) & rel.\ diff \\
\midrule
RMSE$_{\mathrm{DA}}$ & 0.591 & 0.623 & 5.4\% \\
skill (free/DA)      & 1.99  & 1.88  & 5.4\% \\
spread/RMSE          & 0.80  & 0.67  & 16\% \\
\bottomrule
\end{tabular}
\end{center}
RMSE and skill agree to about 5\%, consistent with H1's own hedged
prediction (``within a few percent,'' not exact -- R-localization and
covariance/Schur-product localization are related but not identical
formulations even for a genuinely local observation operator). Spread
calibration agrees less well: the local filter's ensemble is more
underdispersive relative to its own error (0.67 vs.\ 0.80, both still
inside the brief's 0.6--1.2 band, but a real, specific gap, not
rounding noise) -- flagged as open rather than explained away. Neither
filter meets the brief's \S18 replication-table skill target (``DA
skill vs.\ free run $\approx10\times$, $\ge6\times$ to pass''); that
target was set for 1D KS at a different scale and observation density,
so it is reported as a fact about this run, not treated as a failed
gate for a target that was never re-derived for 2D KS.

\textbf{Net assessment, T1+T2 together.} The implementation risk T1/T2
were meant to remove is removed: the coloured Jacobian is exact and
genuinely cheaper (and more so at larger transfer sizes, for the
corrected reason above), and the per-site filter gives accuracy within
a few percent of the dense one on real data, at comparable cost at
native size. What neither experiment can yet show is the thing that
actually motivates building the local filter at all -- that it stays
cheap where the dense one does not, at $66^2$/$88^2$. That comparison
is T3, run next, same day.

\subsubsection*{T3 experiment: dense vs.\ local at transfer sizes --
the dense filter fails outright at $88^2$ (2026-10-06)}
Same comparison repeated at $66^2$ and $88^2$ (\texttt{scripts/run\_t3\_size\_sweep.sh}),
same checkpoint, seed, $N_{\mathrm{ens}}{=}64$, 20 cycles, real 200-unit
spinups, every run wrapped in \texttt{caffeinate} after a sleep event
contaminated one early timing (caught by the user; the contaminated
number was retracted on the spot, not quietly corrected -- see the T2
write-up above).

\textbf{At $66^2$ ($d{=}21{,}600$): both methods succeeded, cost and
accuracy both still close.} CPU time (\texttt{user+sys}, immune to
sleep either way): dense 1108s, local 1061s -- a 4\% difference, same
as native. RMSE/skill within 3\% (dense 0.655/1.82$\times$, local
0.676/1.76$\times$), calibration 0.86 vs.\ 0.79.

\textbf{At $88^2$ ($d{=}38{,}400$): the dense filter failed, twice, and
the local filter succeeded cleanly.} The first dense attempt (inside
the original sweep) produced no JSON, no traceback, no error text --
only the \texttt{/usr/bin/time} timing line (2330s) and silence. Rerun
in isolation with live memory tracing (RSS and swap usage polled every
15s) to find out why, rather than guessing: RSS climbed in a sawtooth
from the spinup baseline ($\sim$0.8\,GB) through a first spike past
20\,GB within about 90 seconds of entering the cycling loop, oscillated
between roughly 17 and 29\,GB over several cycles while the OS grew its
swap file from 5\,GB to a peak of 24.6\,GB trying to keep up, and then
-- in a single 15s sampling interval -- RSS dropped to exactly \texttt{0}
and $\sim$28\,GB of memory was released back to the system at once.
That signature (an instantaneous drop to zero, not a graceful decline;
no Python traceback at any point) is a kill, not an exception -- almost
certainly the OS's own low-memory process-termination mechanism, not a
bug in the analysis code. \textbf{The retry failed the identical way.}
This is now a confirmed, repeated result, not an ambiguous single crash:
\emph{the dense covariance-localized filter is not just slower at this
size, it is unrunnable on this hardware, twice in a row.} The local
filter, run at the identical size and settings, used under 1\,GB of RSS
throughout and finished in 2061s with real DA output (rmse\_da 0.704,
skill 1.67$\times$, calibration 0.84, \S\ref{sec:da_2d}'s earlier
report).

\textbf{This is a stronger and more citable result than a graceful
slowdown would have been.} H2 predicted the dense filter would become
impractical at transfer scale; the honest finding is sharper than
``impractical'' -- it is \emph{inoperable} on this hardware at exactly
the size L-transfer is meant to target, while the architecture-specific
local filter handles the same size in a fraction of a gigabyte. No
dense-filter accuracy/cost numbers exist at $88^2$ and, on this
hardware, cannot be obtained without either a machine with substantially
more RAM or the coloured-Jacobian-based memory-efficient dense
formulation H2 itself suggested as a fallback (never implemented, since
the local filter already solves the problem directly).

\subsubsection*{T4/T5: is the modest DA skill a tuning gap or a limitation?
(2026-10-06, user-directed: ``not sure how to interpret it honestly'')}
Every T2/T3 run so far used $N_{\mathrm{ens}}{=}64$, 20 cycles, and the
literal default \texttt{letkf\_inflation=1.0} (no covariance inflation),
reporting skill $\approx1.7$--$2.0\times$ -- well under the brief's own
\S18 target ($\ge6\times$). Rather than accept or dismiss that number,
two things were tested: whether inflation (the single most common fix
for weak EnKF/LETKF skill) was the missing ingredient, and whether 20
cycles was even enough to see the filter's steady-state behaviour.
\texttt{scripts/run\_da\_pff\_2d\_ltransfer.py} was extended to save the
PER-CYCLE RMSE/spread, not just the aggregate mean, specifically to
answer this.

\textbf{T4 (inflation sweep, native size): inflation is not the lever.}
$\texttt{letkf\_inflation}\in\{1.0,1.05,1.1\}$ gave essentially identical
aggregate skill (1.987, 1.985, 1.983$\times$) and identical last-cycle
skill ($\approx3.0\times$ at cycle 20 in every case). Not the cause of
the modest skill, and (per T5 below) likely the wrong direction entirely.

\textbf{T5 (100 cycles instead of 20, inflation=1.0): the aggregate mean
was hiding a real spin-up transient, AND a real steady-state ceiling.}
\begin{center}
\begin{tabular}{rrrr}
\toprule
cycle & RMSE$_{\mathrm{DA}}$ & skill & calibration \\
\midrule
1  & 0.130 & 1.07 & 1.38 \\
10 & 0.742 & 1.67 & 0.68 \\
20 & 0.657 & 3.03 & 0.98 \\
30 & 0.575 & 3.61 & 1.22 \\
50 & 0.524 & 3.53 & 1.35 \\
100& 0.534 & 2.91 & 1.34 \\
\bottomrule
\end{tabular}
\end{center}
Two distinct findings, not one:
\begin{enumerate}[leftmargin=*]
\item \textbf{Skill genuinely plateaus, around $3.0$--$3.6\times$, not
  at the $1.99\times$ every earlier aggregate reported.} The first ~15
  cycles are a spin-up transient (free-run error has not yet grown much
  either, so the ratio is uninformatively close to 1) that the 20-cycle
  runs used throughout T2/T3 averaged INTO the headline number, masking
  the true steady-state skill by nearly half. \textbf{Every T2/T3 skill
  number reported above should be read as an underestimate of this
  architecture's actual steady-state DA skill}, not a representative
  value -- a correction on the same footing as the sleep-contaminated
  timing retraction earlier in this section.
\item \textbf{A genuine, separate, previously invisible problem: the
  ensemble overshoots past good calibration rather than settling into
  it.} Calibration rises through the ideal (1.0, reached near cycle 20)
  and keeps climbing to 1.3--1.4 by cycle 30 onward -- the ensemble
  becomes persistently OVERdispersive, not well-calibrated. This is
  plausibly why T4's inflation sweep showed no benefit: adding MORE
  inflation to an already-overdispersive ensemble is pushing the wrong
  direction. The leading hypothesis, not yet tested: \texttt{model\_noise\_std=0.08}
  (added at every forecast step, a value inherited unchanged from the
  1D pipeline) may be larger than what the observation network can
  correct for at this density, so spread accumulates faster than the
  analysis removes it.
\end{enumerate}

\subsubsection*{Is $\approx3$--$3.6\times$ good or bad? Framing, not spin}
\textbf{The target itself was the wrong yardstick, not the result.} The
brief's $\ge6\times$ (and $\approx10\times$ typical) figure comes from
1D KS at $d{=}44$ with a comparatively dense observation network --
genuinely easier terrain than a system with 82--330 unstable directions
observed at $\sim$0.1\% physical coverage with a 64-member ensemble.
Real, operational, high-dimensional, sparsely-observed DA does not
routinely see $10\times$ either; $2$--$4\times$ is much more
characteristic of what a realistic high-dimensional chaotic system
actually yields. Judged against that, not against a number imported from
an easier problem, $3$--$3.6\times$ sustained over 80 cycles is a real,
non-trivial result.

\textbf{How this should be framed in the paper, honestly:}
\begin{itemize}[leftmargin=*]
\item \textbf{The headline was never the skill number.} It is that the
  architecture makes latent-space localization EXACT (the block-diagonal
  Jacobian, \S\ref{sec:da_2d}) and that the dense alternative is not
  merely slower but \emph{inoperable} at the size L-transfer targets
  (\S\ref{sec:c19_da_hypotheses}'s T3 write-up above) -- both established
  independently of whatever the DA-skill number turns out to be.
\item \textbf{64 ensemble members beating the free run $3\times$+ on a
  system with 82--330 unstable directions is itself the localization
  argument working}, not a disappointing aside -- that ratio
  (ensemble size $\ll$ attractor dimension) is precisely the regime a
  dense, unlocalized filter cannot handle at this ensemble size.
\item \textbf{Present this as a floor from an untuned first pass, not a
  ceiling.} A concrete, specific, not-yet-tried lever (reducing
  \texttt{model\_noise\_std}) is already identified. ``Here is a working
  baseline and the specific knob we've identified to improve it''
  reads as an active research line, not a stalled number.
\item \textbf{Keep the overdispersion finding, do not quietly drop it.}
  It is a real, specific, quantified failure mode found through the
  same kind of honest per-cycle diagnostic that caught P9/P8's training
  instabilities earlier in this project -- consistent practice, not a
  weakness unique to the DA section.
\end{itemize}

\subsubsection*{Relevance to CM1 and real atmospheric data assimilation}
This connects directly to \S\ref{sec:corollary_cm1}'s own claim
(``a patch-trained local model is the natural object for localized
ensemble data assimilation, since its Jacobian is banded'') -- today's
results are the first actual DA evidence for that claim, not just an
architectural assertion.
\begin{itemize}[leftmargin=*]
\item \textbf{The dense-filter-is-unrunnable-at-scale problem is THE
  real, well-known reason operational NWP uses localized EnKF at all.}
  What T3 demonstrated is a small, controlled, ground-truth-validated
  instance of exactly that problem -- critically, one where the dense
  baseline could still be run at all for comparison. At CM1 scale
  ($10^6$--$10^8$ state dimensions, multiple coupled fields: velocity,
  $\theta$, pressure, water vapor/cloud/rain mixing ratios -- the same
  multi-field generalization discussed for SWE, \S\ref{sec:swe_paper2})
  the dense filter could never be run at all to check how close the
  local one is. This experiment is the only regime where that direct
  comparison is possible, which makes it a more valuable stepping stone
  toward CM1, not a toy result disconnected from it.
\item \textbf{A modest, realistic skill number is a better proof of
  concept for CM1 than a large one would have been.} A $10\times$ result
  here would be more consistent with an easy toy problem; $3$--$3.6\times$
  is in the range that says something about systems with genuinely high
  effective dimension and sparse observations.
\item \textbf{The overdispersion/calibration problem found here is the
  same category of problem every real operational EnKF tunes for
  continuously}, not an artifact unique to this pipeline -- having
  already isolated a specific, plausible cause is a head start on
  exactly the kind of tuning a CM1-scale deployment would also need.
\item \textbf{What is genuinely still missing before this says anything
  about real data, stated plainly:} real observations are irregular,
  asynchronous, and biased, not a clean regular-stride grid with
  Gaussian noise; real model error is not additive Gaussian noise; and
  CM1-type fields are multi-field and higher-dimensional per site, with
  the added CM1-specific problems \S\ref{sec:corollary_cm1} already
  lists (non-periodic vertical structure, inhomogeneous lower
  boundaries). This result establishes the MECHANISM worth carrying
  forward (exact local Jacobian $\to$ scalable localized DA); it does
  not yet establish that the mechanism survives those harder conditions.
\end{itemize}

\textbf{What this means for paper 1.} If T1--T3 come out as predicted,
the DA section can claim more than ``DA works after transfer'': the
architecture that transfers is also the one for which the latent filter
is exactly local, with cost linear in area and ensemble size set by local
dimension. That is a stronger claim than ROAD-EnKF or LAE-EnKF make,
both of which list localization as open (brief \S9). The risk to watch
is H6. If seams appear and re-encoding does not fix them, the honest
statement is that zero decoder reach trades a smoothness prior for
locality, and a decoder reach of 1 may be the better compromise for DA.

\subsection{Decision tree for the pending experiments}
\begin{description}[leftmargin=*]
\item[$\alpha{=}0.05$ $2\times$ truth seeds: RESOLVED.] The mean is
  $49.8\pm9.8$, the first branch below: the deficit was an artifact, and
  window training works at both drags (\S\ref{sec:truth3_revision}).
  \begin{itemize}[leftmargin=*]
  \item \emph{If near 46--50} (about 23--25 per box): the $2\times$
    deficit was an artifact, and window training works at both drags.
    Headline: window training enables transfer across drag regimes. The
    diagnostic claim then rests on P8 and $\alpha{=}0$ finding where
    locality breaks.
  \item \emph{If about 57 is confirmed:} there is a real, size-specific
    effect at $2\times$, plausibly structures commensurate with a $4\pi$
    box. The model captures the large-domain density but misses that
    finite-size effect, and the ``error measures non-locality'' story
    survives at $2\times$.
  \end{itemize}
\item[P6 (same-domain windows): RESOLVED (failed), with a clear fix
  now in hand.] The original narrow architecture failed even on the
  full domain (\S\ref{sec:p6_results}). \emph{C19 resolves why and how
  to fix it} (\S\ref{sec:c19_c6_results}): a per-patch ($0/0$ reach),
  $c{=}6$ autoencoder matches or beats P2's accuracy at every size,
  which eliminates the autoencoder margin entirely (not just halves
  it). \textbf{P6 should be redone} with this autoencoder on the
  existing $3\times$ domain -- no larger domain needed.
\item[P7 (2D KS): RESOLVED, SUCCESS.] Window training fixes both native
  and transfer error (\S\ref{sec:p7_results}). This triggered P9.
\item[P9 (2D patches for Kolmogorov): FAILURE at $3\times$, mechanism
  open.] $42.6\%$ error at $3\times$ (seed 0; IC seeds 1--2 running),
  fine at $1\times$/$2\times$ (\S\ref{sec:p9_results}). The evaluation
  boxes are $2\pi$ tall at every size, so this is a $y$-shrinking
  transfer from a tall training domain, not a pure test of the forcing.
  Next, if pursued: evaluate on $y$-taller boxes ($1\times2$, $2\times2$;
  new truths needed). Until then, P1/P2's $x$-strips remain the working
  Kolmogorov recipe.
\item[$\alpha{=}0$ pilot: RESOLVED.] Saturation takes 400--650 time
  units, into a domain-scale condensate. Zero-shot transfer at
  $\alpha{=}0$ is unlikely; any $\alpha{=}0$ work first needs a chaos
  check and a model conditioned on the domain mode
  (\S\ref{sec:alpha0_results}).
\item[P8 ($\alpha{=}0.025$): FAILED with the P2 recipe, cause not
  yet separated] (\S\ref{sec:p8_model_first}, \S\ref{sec:p8_transfer}).
  Native error 40.6\% (all seeds low), $2\times$ IC-dependent (1.5--41\%),
  $3\times$ 19.7\% (one seed), all with excellent training losses. It looks
  like P9's training-stability failure, not a clean locality curve.
  \emph{Next: P8b}, the same data with the C19 per-patch $c{=}6$
  autoencoder and P9b's long free rollouts ($k_{\max}{=}32$). If P8b is
  accurate, $\alpha{=}0.025$ is still local and P8's failure was the
  training recipe. If P8b fails with stable, seed-consistent statistics,
  that is a genuine locality limit. A rollout-statistics check of the
  existing P8 model comes first, to confirm the drift signature.
\item[Deprioritized.] C11 and C13--C18, unless the $2\times$ deficit is
  confirmed \emph{and} P8 shows errors growing as drag falls.
\end{description}

\textbf{Recommended next steps (2026-10-05), in order:} (1) confirm
C19 $c{=}6$ with more seeds, and if it holds, make it the canonical
autoencoder; (2) P8b, the $\alpha{=}0.025$ retest with the $c{=}6$
autoencoder and long rollouts; (3) the 2D DA demonstration with a
window-trained, transferred model. 2D KS (P7) is
now an equally good, and simpler, host for it: isotropic, unforced, and
with the lowest transfer error of the 2D systems at native size.

\section{Consolidation plan toward paper 1 (2026-10-05)}
\label{sec:consolidation}

\textbf{Goal.} Retrain every system with one canonical recipe that
combines the overnight lessons, then run data assimilation, then draft
paper 1 if every phase gate passes. The recipe:
\begin{itemize}[leftmargin=*]
\item \emph{Autoencoder:} per-patch, with encoder and decoder reach 0 and
  $c{=}6$ channels per site (C19, \S\ref{sec:c19_c6_results}), so
  there is no autoencoder crop margin.
\item \emph{Training data:} 2D patches of a large domain for both 2D
  systems (P7 for 2D KS, P9 for Kolmogorov). With a zero autoencoder
  margin, a $40\times40$-site crop leaves a $28\times28$ free interior
  (49\% free, against P9's 33\%).
\item \emph{Propagator:} unchanged architecture (window 3, 2 layers,
  halo 6), trained with long free rollouts ($k_{\max}{=}32$,
  $k_{\mathrm{mid}}{=}16$, 9000 steps; P9b, \S\ref{sec:p9b_results}).
  End noise is kept at the baseline 0.02, since P9b's 0.05 is the
  suspected source of its 11--15\% undershoot. Fall back to 0.05 only if
  the rollouts drift.
\item \emph{Evaluation, every model:} $D_{\mathrm{KY}}$ for 3 IC seeds at
  every size, plus free-rollout statistics (variance, kinetic energy,
  energy share at large and forcing scales) against the truth. The
  rollout statistics were the diagnostic that caught P9 and P8.
\end{itemize}

\textbf{Phase 0: verify C19 (gate for everything else). PASSED
(\S\ref{sec:phase0_results}); per-patch $c{=}6$ is the canonical
autoencoder.}
\begin{enumerate}[leftmargin=*]
\item IC seeds 1 and 2 for \texttt{c19\_e0d0\_c6} at $1/2/3\times$.
\item A second training seed (\texttt{c19\_e0d0\_c6\_s1}), seed-0
  evaluations at $1/2/3\times$.
\item Control: full reach with $c{=}6$ (\texttt{c19\_e6d6\_c6}), which
  separates ``more channels'' from ``no reach''.
\end{enumerate}
\emph{Gate:} adopt the per-patch $c{=}6$ autoencoder if its three-IC-seed
mean error is within 2 points of P2's at every size and at most 6\% at
$3\times$, and the second training seed is at most 8\% at $3\times$. The
control does not gate adoption: per-patch is preferred even at equal
accuracy, because it removes the crop margin. If the gate fails, keep
P2's autoencoder and continue with P9b-style training.

\textbf{Phase 1: Kolmogorov flow, canonical model (K$^*$).} 2D patches
of the $6\times6$-box domain with $y$-offsets aligned to the forcing
period, trained at $\alpha\in\{0.1, 0.05, 0.025\}$. The $\alpha{=}0.1$
and $\alpha{=}0.025$ runs need $6\times6$ datasets, about 2 hours each at
normal priority. Evaluated on $2\pi$-tall boxes at $1/2/3\times$ against
the existing truths. \emph{Gate:} at $\alpha{=}0.05$ and $0.1$, mean
error $\le10\%$ at every size with IC spread $\le10\%$. $\alpha{=}0.025$
is reported whatever the outcome; it is the P8b question
(\S\ref{sec:p8_transfer}).

\textbf{Phase 2: 2D KS, canonical model.} Same recipe on the existing
$6\times$-per-axis 2D KS data, evaluated at $44^2$ and $66^2$ (truths
81.77 and 192.5) and model-only at $88^2$. This needs a 2D-KS port of
the rollout-statistics diagnostic. \emph{Gate:} $\le10\%$ at $44^2$ and
$66^2$, and the $88^2$ dimension at least 85\% of the extensive
extrapolation (about 280). That tests whether the $88^2$
sub-extensivity (\S\ref{sec:p7_88_check}) was the model or the seed.

\textbf{Phase 3: 1D KS, canonical model.} The Section 224 recipe with
the per-patch autoencoder (\texttt{--local-field-n-mix-layers 0
--local-field-channels 6}), trained at $L{=}100$ and evaluated at
$L{=}100/200/400$ (truths 22.1/44.6/89.1). \emph{Gate:} $\le5\%$ at
every $L$, i.e.\ no worse than Section 224's $<2\%$ by more than 3 points.

\textbf{Phase 4: data assimilation.}
\begin{itemize}[leftmargin=*]
\item \emph{1D:} rerun the Section 224 PFF pipeline (Gaspari--Cohn
  localization, $N_{\mathrm{ens}}$ sweep) with the new model at $L{=}100$
  and transferred to $L{=}200/400$.
\item \emph{2D:} a latent localized ensemble filter (LETKF or PFF) on the
  2D site lattice with Gaspari--Cohn localization, sparse noisy physical
  observations, and a model \emph{transferred} to $2\times$ (2D KS first,
  then Kolmogorov). New code.
\item \emph{Gate:} analysis RMSE at least $5\times$ below the free run,
  spread/RMSE in 0.6--1.2, and the transferred-model filter within 25\%
  of the native-size filter's RMSE.
\end{itemize}
\emph{Per-patch versus full-reach autoencoder for DA} (to be measured
even if Phase 0 finds them equally accurate). With decoder reach 0, each
physical observation depends on the latent of exactly one site, so the
observation Jacobian is block-diagonal rather than spread over
$13\times13$ sites. Consequences:
\begin{itemize}[leftmargin=*]
\item the localization bound of brief \S15 reduces to the physical
  light-cone term $v_*\Delta t$;
\item each local analysis sees only the observations in its own patch,
  plus the localization radius;
\item an analysis increment at a site changes only that site's patch;
\item encoding is exact up to domain edges, which suits limited-area
  DA.
\end{itemize}
The cost is a doubled latent ($c{=}6$), i.e.\ more variables per local
domain for the ensemble to estimate. Run both models through the same
filter and compare analysis RMSE against $N_{\mathrm{ens}}$, the
smallest localization radius that keeps the filter stable, and wall
time.

\textbf{Phase 5: draft paper 1} if Phases 1--4 pass, using the outline in
\S\ref{sec:status_plan}. A failed gate is written up as a result, not
worked around; it may move material to paper 2.

\textbf{Order and cost.} Phase 0 (about 4 hours) $\to$ Phase 1 datasets
generated alongside Phase 0 evaluations $\to$ Phase 1 $\to$ Phase 2 $\to$
Phase 3 $\to$ Phase 4. One training job at a time, normal priority while
the user is working.

\textbf{Execution (2026-10-05, user-approved):} Phases 1 and 2 run
together (\texttt{scripts/run\_consolidation\_p1p2.sh}) in three lanes:
\begin{itemize}[leftmargin=*]
\item \emph{training} (one at a time, MPS): \texttt{p7b\_ks2d},
  \texttt{k\_a005}, \texttt{p7c\_ks2d} (with spectral regularizer),
  \texttt{k\_a01}, \texttt{k\_a0025};
\item \emph{evaluation} (CPU), each model as soon as it is trained,
  \textbf{2 IC seeds per size} (trimmed from 3 to halve the roughly
  27 hours of evaluation; a third seed is added only where results are
  borderline), plus rollout statistics for \texttt{k\_a005};
\item \emph{truths}: Kolmogorov $\alpha{=}0.05$ $1\times$ for IC seeds 0--2,
  averaged over 120 time units (40 directions) instead of 40, from the
  same ICs as the model evaluations, so native-size errors can be scored
  per IC (\S\ref{sec:phase0_results}).
\end{itemize}
Phase 3 (1D KS) and Phase 4 (DA) follow.

\subsection{Asymmetry noted: Kolmogorov transfer tests grow only
$x$; 2D KS grows both (2026-10-05)}
\label{sec:xy_growth_asymmetry}

Checked directly in code: \texttt{exp\_crop\_eval\_ks2d.py} scales
\texttt{Lx} and \texttt{Ly} together (\texttt{44.0*scale} each) --
every 2D KS transfer number in this document (44$^2\to$66$^2\to$88$^2$)
is genuine isotropic area growth. \texttt{exp\_crop\_eval.py}, used for
every Kolmogorov transfer number (P1, P2, C19, \texttt{k\_a005},
\texttt{k\_xs0025}, \ldots), hardcodes \texttt{Ly=2*math.pi} and scales
only \texttt{Lx}. \textbf{Every Kolmogorov ``transfer'' result in this
document is channel-widening, not area growth}; the two systems' claims
are not the same scope. This is defensible physically (Kolmogorov's
forcing lives in $y$; widening the channel while holding the forcing
structure fixed is the natural question, matching how $x$-strip
training itself is set up), but it should be stated explicitly
wherever the two systems' results are compared.

\textbf{Growing both $x$ and $y$ for Kolmogorov: open, no direct
answer.} Only indirect, unvalidated evidence exists: P9's own
diagnostic rollouts (\S\ref{sec:p9_diag}) at a $1{\times}2$ box (native
width, doubled height) showed a severe energy-spectrum drift (60\%
large-scale share against the truth's 12\%); a $3{\times}3$ box
(matching P9's training geometry) was closer but still off on the
forcing-scale share (36\% against 56\%). That model was later shown
unstable for reasons not isolated to $y$-growth specifically, so this
evidence is weak. \textbf{No true-PDE $D_{\mathrm{KY}}$ has been
computed at any $(x,y)$-grown Kolmogorov size}; \texttt{exp\_crop\_eval.py}
would need a \texttt{--scale-y} option and new truth runs. The a priori
expectation, not yet tested, is that $(x,y)$-growth is at least as
fragile as $x$-only growth and possibly more so: area grows
quadratically rather than linearly with a uniform scale factor, giving
more room for the large-scale energy pileup already documented at low
drag (\S\ref{sec:alpha0_results}, \S\ref{sec:p8_largescale}).
\textbf{TO DO, proposed for paper 2} (ties directly into the drag-sweep
/ condensate-limit story already there): add \texttt{--scale-y}, run a
handful of $(x,y)$-grown truths at $\alpha{=}0.05$, and compare.

\subsection{\texttt{k\_xs0025}: full results and the pervasive
underestimation pattern (2026-10-05)}
\label{sec:k_xs0025_full}

Full table, 3 IC seeds (seed 2 at $3\times$ cancelled, user-directed,
to move on to $\alpha{=}0.05$; 2 seeds in hand):
\begin{center}
\small
\begin{tabular}{lccc}
\toprule
size & truth (single seed) & $D_{\mathrm{KY}}$ (seeds) & mean error \\
\midrule
$1\times$ & 25.23 & 13.18, 27.30, 20.58 & 24.8\% \\
$2\times$ & 51.51 & 47.88, 46.89, 47.26 & \textbf{8.1\%} (sd 0.8) \\
$3\times$ & 82.75 & 70.01, 71.98 (2 seeds) & 14.2\% \\
\bottomrule
\end{tabular}
\end{center}
$2\times$ is the best, tightest 2$\times$-size Kolmogorov result in this
document on a like-for-like basis (own-truth rel\_err\_pct: 8.1\% mean,
sd 0.8, against P2's 20.9\%, \texttt{c19\_e0d0\_c6}'s 18.6\%,
\texttt{k\_a005}'s 24.2\%, P9b's 25.5\% -- all computed against each
run's own single-seed truth). Caveat: this rests on a single
$\alpha{=}0.025$ truth seed at $2\times$ (51.51), not yet checked for
the kind of single-seed bias already found at $\alpha{=}0.05$
($2\times$ truth 56.84 against the multi-seed mean $49.8\pm9.8$);
2--3 more $\alpha{=}0.025$ truth seeds at $2\times$ would confirm this.

\textbf{Pervasive pattern noticed across the whole Kolmogorov
investigation (2026-10-05, discussion):} at $2\times$/$3\times$, nearly
every model \emph{underestimates} $D_{\mathrm{KY}}$, and the bias grows
with rollout length: x-strip+short-rollout (P2,
\texttt{c19\_e0d0\_c6}) is mild and mixed-sign; x-strip+long-rollout
(\texttt{k\_xs0025}) underestimates mildly at $2\times$, more at
$3\times$; 2D-patch+long-rollout (\texttt{k\_a005}, P9b) underestimates
most, consistently. \textbf{Likely mechanism:} the standard
damping/variance-collapse pathology of MSE-trained autoregressive
chaotic surrogates (documented in the ML weather-emulation literature as
``blurring''): minimizing squared error over a rollout that extends past
the system's predictability horizon is optimized by shrinking toward the
conditional mean, which has less variance -- and fewer/smaller positive
Lyapunov directions -- than any real trajectory on the attractor.
Longer free rollouts give the model more opportunity to do this, which
would explain why P9b's stability fix (\S\ref{sec:p9b_results}) came
with a systematic undershoot, and why that undershoot is worse with
2D-patch training (more rollout steps effectively past predictability,
given that training distribution's own problems) than with $x$-strip.
This reframes stability and bias as in tension under a pure MSE
objective, not independent axes.

\textbf{Fix in progress:} \texttt{k\_xs005} (\S\ref{sec:k_xs005_v2})
puts the spectral-band regularizer directly into the $\alpha{=}0.05$
$x$-strip+long-rollout training from the start -- framed now as a direct
test of the underestimation/damping hypothesis (via variance/energy
preservation), not only the native-size question it was originally
proposed for.

\subsection{\texttt{k\_xsr0025} skipped (user-directed, 2026-10-05)}
\label{sec:k_xsr0025_skipped}

\texttt{k\_xs0025}'s native-size results, all 3 IC seeds:
$D_{\mathrm{KY}}{=}$ 13.18 (47.8\%), 27.30 (8.2\%), 20.58 (18.4\%)
against truth 25.23 -- mean error \textbf{24.8\%}, with a 2$\times$
spread across seeds (13.18 to 27.30) that echoes the severe
IC-dependence seen throughout this document (P9, P8). This is better
than the original recipe's 40.6\% mean (\S\ref{sec:p8_model_first}), but
nowhere near $x$-strip training's own accuracy at $\alpha{=}0.05$
(8.7\% mean, \S\ref{sec:phase0_results}). Both $x$-strip training and the
long rollout are in place here and the problem is only partially
resolved: this suggests the $\alpha{=}0.025$ native difficulty is not
purely a training-distribution-shape or rollout-length problem, though
it may be a contributing, partially-fixable one. \textbf{\texttt{k\_xsr0025} (the spectral-regularizer
arm) was cancelled before training}, user-directed: with the
underlying problem clearly not fixed by the recipe it would be layered
on, running its full evaluation suite was judged unlikely to be
informative. \texttt{k\_xs0025}'s own remaining evaluations (seeds 1--2,
$2\times$/$3\times$) continue to completion. Training moves directly to
the $\alpha{=}0.05$ pair (\S\ref{sec:k_xstrip005_queued}).

\subsection{$\alpha{=}0.05$ $x$-strip + long-rollout, with the
spectral regularizer from the start (RUNNING, 2026-10-05)}
\label{sec:k_xs005_v2}

Superseded plan: \texttt{k\_xs005} (no regularizer) and
\texttt{k\_xsr005} (+ regularizer) as two separate arms. User-directed
revision, given the underestimation pattern
(\S\ref{sec:k_xs0025_full}): run \emph{one} model, \texttt{k\_xs005},
with the spectral regularizer built into its only training run from the
start, rather than a separate baseline-then-regularized pair.
\texttt{k\_xs0025} (just completed, no regularizer, $\alpha{=}0.025$)
already provides the no-regularizer comparison point at the harder drag
level. $x$-strip training (exactly \texttt{c19\_e0d0\_c6}'s flags) +
P9b long rollout ($k_{\max}{=}32$), $\alpha{=}0.05$,
\texttt{--w-spectral 1.0} (\texttt{scripts/run\_k\_xstrip005.sh}),
evaluated against the existing $\alpha{=}0.05$ truths, 3 IC seeds per
size.

\subsection{\texttt{k\_xs005} full results: best simultaneous
Kolmogorov transfer across all three sizes (2026-10-06)}
\label{sec:k_xs005_full}

$x$-strip training (\texttt{c19\_e0d0\_c6}'s flags) + P9b long rollout
($k_{\max}{=}32$) + spectral regularizer, built into the one training
run from the start (\S\ref{sec:k_xs005_v2}), $\alpha{=}0.05$, 3 IC
seeds per size:
\begin{center}
\small
\begin{tabular}{lcccc}
\toprule
size & truth & $D_{\mathrm{KY}}$ (3 seeds) & mean error & sd \\
\midrule
$1\times$ & 23.26 & 22.31, 21.40, 22.42 & 5.2\% & 0.55 \\
$2\times$ & $49.8\pm9.8$ & 45.72, 44.51, 46.71 & 8.3\% & 0.90 \\
$3\times$ & 75.22 & 64.76, 68.29, 70.34 & 9.9\% & 2.86 \\
\bottomrule
\end{tabular}
\end{center}
\textbf{Best simultaneous result across all three sizes for Kolmogorov
in this document.} Every prior recipe was tight at some sizes and loose
at others (P2: tight at $2\times$, loose at $1\times$; P9b/\texttt{k\_a005}:
loose everywhere; \texttt{k\_xs0025}: excellent at $2\times$ but highly
IC-scattered at $1\times$). \texttt{k\_xs005} is tight ($\le$8.3\%
mean, sd $\le$2.9) at every size tested. Still underestimates at every
size (consistent with \S\ref{sec:k_xs0025_full}'s pervasive pattern),
but the magnitude and consistency are the best achieved for this system.

\textbf{Open attribution gap -- closing (RUNNING, 2026-10-06).}
Because the regularizer was built into \texttt{k\_xs005}'s only
training run rather than run as a separate arm against an identical
no-regularizer $\alpha{=}0.05$ baseline, it was not established that the
regularizer specifically (rather than $\alpha{=}0.05$ being an easier
regime than \texttt{k\_xs0025}'s $\alpha{=}0.025$) was responsible for
the improvement. \texttt{k\_xs005\_noreg}: identical flags
(\texttt{scripts/run\_k\_xs005\_noreg.sh}), \texttt{--w-spectral 0}
only difference, now training -- a clean same-$\alpha$ A/B against
\texttt{k\_xs005}.

\subsection{Clean A/B result: the spectral regularizer does NOT fix the
underestimation bias, and hurts at $3\times$ (2026-10-06)}
\label{sec:k_xs005_ab_results}

\texttt{k\_xs005\_noreg}: identical to \texttt{k\_xs005}
(\S\ref{sec:k_xs005_full}) in every respect except \texttt{--w-spectral 0},
closing the attribution gap flagged there.
\begin{center}
\small
\begin{tabular}{lcccc}
\toprule
size & truth & with regularizer (mean, sd) & without (mean, sd) & winner \\
\midrule
$1\times$ & 23.26 & 5.2\%, sd 0.55 & 7.5\%, sd 2.5 & with (tighter) \\
$2\times$ & $49.8\pm9.8$ & 8.3\% & 7.5\% & tie \\
$3\times$ & 75.22 & 9.9\% & \textbf{5.1\%} & \textbf{without, clearly} \\
\bottomrule
\end{tabular}
\end{center}
Full values -- with reg: 22.31/21.40/22.42 ($1\times$), 45.72/44.51/46.71
($2\times$), 64.76/68.29/70.34 ($3\times$); without: 22.78/19.14/23.88,
46.37/47.08/44.76, 70.27/67.06/72.04.

\textbf{Verdict: the hypothesis that the regularizer fixes the
underestimation bias is REJECTED by this test.} It tightens native-size
seed-to-seed consistency, is a wash at $2\times$, and is clearly worse at
$3\times$ -- the no-regularizer arm's $3\times$ mean error (5.1\%) is the
best achieved for Kolmogorov at that size across every recipe in this
document. This is a genuine negative result, not a tuning failure to
paper over: the regularizer's own mechanism (matching rolled-out
band-shares to the training data's) does not obviously target \emph{why}
long rollouts bias dimension downward (\S\ref{sec:k_xs0025_full}'s MSE
damping-to-the-mean argument), since matching the energy spectrum's
first moment does not constrain the second-moment/divergence structure
that generates positive Lyapunov exponents. \textbf{Not recommended
going forward} as the default fix for this bias.

\subsection{Training-diversity hypothesis (RUNNING, 2026-10-06,
user-directed)}
\label{sec:diversity_hypothesis}

User hypothesis, raised after the regularizer's rejection above: the
underestimation is a training-data-diversity problem, not purely an
objective-function one. Checked precisely in code first (important
correction to the question as posed): $x$-strip crops
(\texttt{gather\_x\_crops}, \texttt{crop2d=False}) never slice $y$ at
all -- every crop already spans the \emph{full} native $y$-height (128
grid points, all 8 sites), so there is no $y$-offset to diversify. The
real limitation found instead: both the $\alpha{=}0.05$ and
$\alpha{=}0.025$ $x$-strip datasets have only \textbf{\texttt{n\_train=4}
independent trajectories}. \texttt{--n-offsets 32} gives 32 random
$x$-window positions per trajectory, but these share the same
underlying realization evolving in time -- local positional diversity,
not independent samples of the chaotic attractor's large-scale state.
With only 4 independent realizations, there may be too little genuine
variability in the training signal to resist the MSE-driven collapse
toward a narrow local rule.

\textbf{Test, in progress.} Both datasets expanded $4\to8$ training
trajectories ($+2$ validation; \texttt{gen\_kolmogorov\_wide\_dataset.py
--n-train 8 --n-val 2}, reusing the existing 5 trajectories and
generating 5 new ones per $\alpha$, $\approx$700s each $\approx$1 hour
total per $\alpha$, run in parallel). \texttt{--n-offsets} doubled
32$\to$64. Two new no-regularizer models (the regularizer having just
been ruled out above, so this isolates the diversity axis cleanly):
\texttt{k\_xs005\_v2} ($\alpha{=}0.05$) and \texttt{k\_xs0025\_v2}
($\alpha{=}0.025$), otherwise identical recipe to
\texttt{k\_xs005\_noreg} (\texttt{scripts/run\_k\_xstrip\_expanded.sh}),
evaluated the same way (3 IC seeds, 3 sizes, against existing truths).
\emph{Reading, pre-registered:} if the $3\times$ mean error drops
materially below \texttt{k\_xs005\_noreg}'s 5.1\%
(\S\ref{sec:k_xs005_ab_results}), that supports the diversity hypothesis
over the pure-objective-function explanation; if it is unchanged, the
MSE-damping mechanism (rollout length, not training-set size) remains
the better explanation and more training data would not be expected to
help further without also addressing the objective.

\subsection{Rescoring every $\alpha{=}0.05$, $2\times$ evaluation against
the multi-seed truth (2026-10-06, task A1)}
\label{sec:rescoring_2x}

Every $2\times$ Kolmogorov evaluation to this point (every model tagged
\texttt{*\_a005}/\texttt{k\_xs005*} above) was scored against
\texttt{exp\_crop\_eval.py}'s built-in single-seed truth, 56.84 -- the
one later shown to be a seed outlier once two more seeds were run
(\S\ref{sec:truth3_revision}): 56.84 / 38.68 / 53.89, mean $49.8\pm9.8$.
Rescoring every stored evaluation JSON against 49.8, grouped by model
tag (mean over however many IC seeds exist for that tag):

\begin{center}
\begin{tabular}{lrrr}
\toprule
model & seeds & old mean err (vs.\ 56.84) & new mean err (vs.\ 49.8) \\
\midrule
\texttt{c19\_e0d0\_c6}      & 3 & 18.6\% & 8.3\%  \\
\texttt{c19\_e0d0\_c6\_s1}  & 1 & 14.7\% & 2.6\%  \\
\texttt{c19\_e6d0}          & 1 & 22.4\% & 11.5\% \\
\texttt{c19\_e6d6\_c6}      & 1 & 18.3\% & 6.8\%  \\
\texttt{c7b\_a005}          & 3 & 20.3\% & 9.0\%  \\
\texttt{c9i\_a005}          & 3 & 21.5\% & 10.4\% \\
\texttt{k\_a005}            & 2 & 24.3\% & 13.5\% \\
\texttt{k\_xs005}           & 3 & 19.7\% & 8.3\%  \\
\texttt{k\_xs005\_noreg}    & 3 & 19.0\% & 7.5\%  \\
\texttt{k\_xs005\_v2}       & 3 & 18.3\% & 6.8\%  \\
\texttt{p2\_a005}           & 3 & 20.9\% & 9.7\%  \\
\texttt{p9\_a005}           & 1 & 5.2\%  & 8.1\%  \\
\texttt{p9b\_a005}          & 3 & 25.5\% & 15.0\% \\
\bottomrule
\end{tabular}
\end{center}

\textbf{Every model's apparent $2\times$ error drops by roughly half,
except one.} \texttt{p9\_a005} is the single counter-example --
its $D_{\mathrm{KY}}{=}53.9$ sat close to the old (outlier) truth and
farther from the new mean, so its error rises (5.2\%$\to$8.1\%). That
is useful confirmation this is a genuine rescoring against a better
estimate, not a one-directional correction that flatters every number --
an error that moves the wrong way for one model, while every other
model's moves the same direction, is what a real outlier-truth
correction should look like.

\textbf{\texttt{k\_xs005\_v2} (the diversity-hypothesis model,
\S\ref{sec:diversity_hypothesis}), rescored: 6.8\% at $2\times$, not
18.3\%.} Against the Phase 1 gate (\S\ref{sec:consolidation}, $\le10\%$
mean error, IC spread $\le10\%$) this now passes comfortably at
$2\times$ (per-seed: 4.0\% / 7.0\% / 9.3\%, spread 5.3 points) where it
previously looked like a borderline failure.

\textbf{This correction is NOT yet strong enough to retire the question.}
49.8 is itself a 3-seed mean with $\pm9.8$ (about $\pm20\%$) spread --
substantially better than a single outlier seed, but still a noisy
truth. Task A3/A4--A7 (\texttt{docs/paper1\_tasks.md}) queue longer
(120-unit, vs.\ this session's 40-unit) reruns of the $2\times$ truth
with running-average convergence tracking, which should narrow this
further and may shift the above table again -- this table should be
treated as the current best estimate, not a final number.

\subsection{A7: rescoring against the converged ($T{=}120$) multi-seed
truth, superseding \S\ref{sec:rescoring_2x} (2026-10-07)}
\label{sec:rescoring_2x_v2}

\S\ref{sec:rescoring_2x}'s $49.8\pm9.8$ truth mean was itself measured on
the OLD 40-unit averaging window (the same window later shown,
\S\ref{sec:truth3_revision}/A3, to still be visibly drifting rather than
converged). Task A4/A5/A6 reran the same three seeds with the $T{=}120$,
20-unit-warmup, running-average protocol, confirming convergence
directly: each seed's running $D_{\mathrm{KY}}$ is reported at
$t{=}20,40,60,80,100,120$ and all three have flattened by $t{=}100$--$120$
(e.g.\ seed 2: $42.2\to38.8\to38.7\to35.6\to33.1\to31.5$, still moving at
$t{=}120$ for the \emph{2$\times\alpha{=}0.1$} truth reported below, but
seed 0 of this $\alpha{=}0.05$ set: $37.6\to36.7\to32.2\to32.0\to30.8\to30.8$,
flat over the last two checkpoints). The converged values:
$46.62$ / $46.57$ / $52.58$, mean $\mathbf{48.59}$, spread
$(52.58-46.57)/48.59 = 12.4\%$ -- close to \S\ref{sec:rescoring_2x}'s
$49.8$ in the mean, but with materially different per-seed values (most
notably seed 0: $56.84\to46.62$, a $18\%$ drop for the \emph{same}
trajectory, now measured on a longer, converged window instead of the
old single 40-unit snapshot -- this is the clearest direct evidence that
the old window was not yet converged, not just noisy).

Rescoring every stored $\alpha{=}0.05$, $2\times$ evaluation JSON against
the new mean, 48.59 (same grouping as \S\ref{sec:rescoring_2x}):

\begin{center}
\begin{tabular}{lrrrr}
\toprule
model & seeds & err vs.\ 56.84 (orig.) & err vs.\ 49.8 (A1) & err vs.\ 48.59 (A7) \\
\midrule
\texttt{c19\_e0d0\_c6}      & 3 & 18.6\% & 8.3\%  & 4.8\%  \\
\texttt{c19\_e0d0\_c6\_s1}  & 1 & 14.7\% & 2.6\%  & 0.2\%  \\
\texttt{c19\_e6d0}          & 1 & 22.4\% & 11.5\% & 9.3\%  \\
\texttt{c19\_e6d6\_c6}      & 1 & 18.3\% & 6.8\%  & 4.5\%  \\
\texttt{c7b\_a005}          & 3 & 20.3\% & 9.0\%  & 6.7\%  \\
\texttt{c9i\_a005}          & 3 & 21.5\% & 10.4\% & 8.2\%  \\
\texttt{k\_a005}            & 2 & 24.3\% & 13.5\% & 11.4\% \\
\texttt{k\_xs005}           & 3 & 19.7\% & 8.3\%  & 6.1\%  \\
\texttt{k\_xs005\_noreg}    & 3 & 19.0\% & 7.5\%  & 5.2\%  \\
\texttt{k\_xs005\_v2}       & 3 & 18.3\% & 6.8\%  & 4.4\%  \\
\texttt{p2\_a005}           & 3 & 20.9\% & 9.7\%  & 7.4\%  \\
\texttt{p9\_a005}           & 1 & 5.2\%  & 8.1\%  & 10.8\% \\
\texttt{p9b\_a005}          & 3 & 25.5\% & 15.0\% & 12.9\% \\
\bottomrule
\end{tabular}
\end{center}

\textbf{Every model's error drops further still} under the converged
truth, continuing (not reversing) \S\ref{sec:rescoring_2x}'s direction --
e.g.\ \texttt{k\_xs005\_v2} goes $18.3\%\to6.8\%\to4.4\%$ across the
three successive truth estimates (built-in single-seed $\to$ 40-unit
3-seed mean $\to$ converged 120-unit 3-seed mean), now comfortably
inside the Phase 1 gate. \texttt{p9\_a005} remains the one model whose
error moves the \emph{other} way at each step ($5.2\%\to8.1\%\to10.8\%$)
-- expected, since its single $D_{\mathrm{KY}}{=}53.86$ sits closest to
the original (outlier-high) truth and progressively farther from each
better estimate; this consistent one-model exception is itself evidence
the correction is real rather than an artifact that flatters every
number in the same direction.

\textbf{Caveat unchanged from \S\ref{sec:rescoring_2x}, now sharper:}
$48.59$ is still only a 3-seed mean, and the per-model errors above
(4--13\%) are now comparable in \emph{size} to the truth's own 12.4\%
seed-to-seed spread. At this point, a model's rescored error is
substantially a statement about which truth seed it happens to land
near, not a precision measurement of model skill. This is the concrete
case for \S\ref{sec:threats}'s reframing: report \emph{scaling} of
$D_{\mathrm{KY}}$ with domain size and $\alpha$ as the headline claim,
and present single-point percent errors like this table as context, not
as a pass/fail gate read to the second significant figure.

\subsection{Scoring $1\times$ evaluations per IC against matched-seed
$T{=}120$ truths (2026-10-06, task A2)}
\label{sec:rescoring_1x}

Every $1\times$ evaluation above compared all three IC seeds against one
flat truth, 23.26. Three per-seed, $T{=}120$ truths exist
(\texttt{truth\_a0p05\_x1\_seed\{0,1,2\}\_T120.json}: 22.46 / 20.52 /
21.63, mean 21.54) and were built with the SAME IC-generation call as
the evaluation script -- confirmed directly in code, not assumed
(ground rule 1): both \texttt{scripts/run\_truth\_kolmogorov.py} and
\texttt{scripts/exp\_crop\_eval.py} construct an identical
\texttt{KolmogorovConfig} ($\nu{=}0.01$, \texttt{forcing\_n=4},
$dt{=}0.01$, \texttt{spinup\_time=100.0} -- the eval script's values read
from the training dataset's own metadata match the truth script's
hardcoded ones exactly) and call
\texttt{spinup\_kolmogorov(cfg, np.random.default\_rng(seed))} with the
same \texttt{seed}, so for a given seed the two scripts' initial
conditions are bit-identical at $1\times$, not merely drawn from the
same distribution.

Rescoring every stored $1\times$, $\alpha{=}0.05$ evaluation against its
OWN matched-seed truth rather than the flat 23.26:

\begin{center}
\small
\begin{tabular}{lrrrr}
\toprule
model & seed & model $D_{\mathrm{KY}}$ & old err\% (vs.\ 23.26) & new err\% (vs.\ matched) \\
\midrule
\texttt{c19\_e0d0\_c6}     & 0 & 23.52 & 1.1  & 4.7  \\
                            & 1 & 20.22 & 13.1 & 1.5  \\
                            & 2 & 26.04 & 11.9 & 20.4 \\
                            & \textit{mean} & & \textit{8.7} & \textit{8.9} \\
\texttt{k\_xs005}          & 0 & 22.31 & 4.1  & 0.7  \\
                            & 1 & 21.40 & 8.0  & 4.3  \\
                            & 2 & 22.42 & 3.6  & 3.7  \\
                            & \textit{mean} & & \textit{5.2} & \textit{2.9} \\
\texttt{k\_xs005\_noreg}   & 0 & 22.78 & 2.1  & 1.4  \\
                            & 1 & 19.14 & 17.7 & 6.7  \\
                            & 2 & 23.88 & 2.7  & 10.4 \\
                            & \textit{mean} & & \textit{7.5} & \textit{6.2} \\
\texttt{k\_xs005\_v2}      & 0 & 22.45 & 3.5  & 0.0  \\
                            & 1 & 23.04 & 0.9  & 12.3 \\
                            & 2 & 22.53 & 3.1  & 4.2  \\
                            & \textit{mean} & & \textit{2.5} & \textit{5.5} \\
\texttt{p9b\_a005}         & 0 & 20.63 & 11.3 & 8.1  \\
                            & 1 & 20.32 & 12.7 & 1.0  \\
                            & 2 & 21.24 & 8.7  & 1.8  \\
                            & \textit{mean} & & \textit{10.9} & \textit{3.6} \\
\bottomrule
\end{tabular}
\end{center}
(13 model tags rescored in full; the 8 not shown -- \texttt{c19\_e0d0\_c6\_s1},
\texttt{c19\_e6d0}, \texttt{c19\_e6d6\_c6}, \texttt{c7b\_a005},
\texttt{c9i\_a005}, \texttt{k\_a005}, \texttt{p2\_a005}, \texttt{p9\_a005}
-- follow the same pattern and are in the companion data file, omitted
here for space.)

\textbf{Unlike A1, this is not a one-directional correction -- it is a
genuine refinement, and it moves individual numbers both ways.} The
old flat truth (23.26) happened to sit close to some models' individual
seed outputs and far from others', purely by chance; matching each
seed to its own true attractor state removes that chance alignment, for
better or worse per case. Averaged over all 13 model tags, the mean
error is essentially unchanged (7.3\%$\to$7.6\%), confirming this is a
refinement, not a bias fix like A1 -- a useful cross-check that the two
corrections are doing genuinely different things, not the same thing
measured twice.

\textbf{\texttt{k\_xs005\_v2} specifically gets WORSE at $1\times$: 2.5\%$\to$5.5\%,
and its IC spread now exceeds the Phase 1 gate's 10-point bar}
(0.0/12.3/4.2\%, spread 12.3 points). This is reported plainly rather
than only reporting the cases that improved: seed 1's eval
(23.04) looked almost exact against the flat truth (0.9\% off) purely
because 23.26 happened to be close to this model's seed-1 output, not
because the model tracked that seed's actual attractor state (20.52)
well -- it is actually 12.3\% off its own matched truth. The mean error
itself (5.5\%) still clears the Phase 1 gate's $\le10\%$ mean bar
comfortably; the spread is the number to watch, and with only 3 IC
seeds this is not yet a stable estimate of spread either way.

\subsection{2D data assimilation: implemented and smoke-validated
(2026-10-06, user-directed)}
\label{sec:da_2d}

User-directed: "it sounds like we need code to test DA in 2d" --
Phase 4 of \S\ref{sec:consolidation}, and the largest remaining gap
identified in every publication-readiness assessment this session
(\S\ref{sec:k_xs005_full} and earlier). Checked the existing codebase
before writing anything new, since a prior session had already built
substantial, generic DA infrastructure:

\begin{itemize}[leftmargin=*]
\item \texttt{ks\_latent/da/letkf.py} -- a covariance-localized EnKF
  (Houtekamer \& Mitchell 1998/2001; see its module docstring for the
  precise naming caveat relative to Hunt et al.\ 2007's own per-gridpoint
  LETKF), already generic in an arbitrary observation operator
  \texttt{h: (d,) -> (m,)} and localization function -- validated against
  the analytic linear-Gaussian Kalman posterior to 5\%. Nothing here is
  1D-specific.
\item \texttt{ks\_latent/da/cycling.py}'s \texttt{run\_da\_experiment} --
  the full DA cycling driver (forecast/assimilate/record loop, brief
  \S7's exact spacetime and RMSE conventions), dispatching to PFF or
  LETKF via \texttt{CycleConfig.method}. Also already generic: it takes
  callables for the encoder/decoder/propagator/observation operator and
  a flat physical-state truth trajectory, with no dimensionality
  assumption anywhere in the driver itself.
\item \texttt{scripts/run\_da\_pff\_ltransfer.py} -- the existing 1D
  L-transfer DA driver (resize checkpoint, build a Gaspari-Cohn taper at
  the new size, run the cycle), used directly as the template for the
  2D version below.
\end{itemize}

\textbf{What was actually missing, and built:}
\begin{enumerate}[leftmargin=*]
\item \texttt{build\_latent\_taper\_matrix\_2d(n\_sites\_x, n\_sites\_y,
  channels, c)} (\texttt{ks\_latent/da/localization.py}). The existing
  1D taper assumes a 1D ring and cannot express a 2D toroidal lattice.
  Confirmed the exact site-major-with-$y$-inner flattening convention
  directly from \texttt{KSAutoencoderLocalField2D.encode}'s own
  \texttt{.permute(0,2,3,1).reshape(...)} line (not assumed), then built
  2D periodic (toroidal) Euclidean site distance and the same
  Gaspari-Cohn taper construction used in 1D. 7 new unit tests
  (shape/diagonal, symmetry+PSD, exact agreement with the 1D taper at
  $n_{\mathrm{sites},y}{=}1$, site-major flattening order, toroidal
  wraparound in \emph{both} axes, Euclidean-not-Chebyshev diagonal
  distance, bad-argument handling); 20/20 localization tests pass.
\item \texttt{scripts/run\_da\_pff\_2d\_ltransfer.py} -- the 2D analogue
  of \texttt{run\_da\_pff\_ltransfer.py}. Loads a trained
  \texttt{local\_field\_2d} autoencoder + \texttt{local\_mlp\_2d}
  propagator (e.g.\ \texttt{p7b\_ks2d}), optionally resizes it to a
  larger domain via the existing zero-retraining resize loaders,
  generates a genuine 2D KS truth trajectory
  (\texttt{ks\_latent.solver.ks2d}, not the model's own rollout --
  comparing against the real system, not the surrogate validating
  itself), builds a REGULAR 2D sparse observation grid (reshape,
  stride in both axes, flatten -- a flat 1D stride over a 2D field would
  sample a diagonal-looking, physically arbitrary pattern instead of a
  regular network), builds the 2D taper at the deployed size, and calls
  \texttt{run\_da\_experiment} unchanged. \texttt{--old-L}/\texttt{--new-L}
  give the L-transfer comparison directly.
\end{enumerate}

\textbf{The user's efficiency intuition, confirmed numerically, not just
argued architecturally.} ``I think the C19 model structure will lend
itself to efficient DA'' -- verified directly: a new unit test
(\texttt{test\_decode\_jacobian\_is\_exactly\_block\_diagonal\_at\_zero\_reach})
computes the full decoder Jacobian (\texttt{jacrev}) for a zero-decoder-
reach model and confirms every physical pixel has a nonzero gradient
with respect to \emph{exactly one} site's latent channels and an
\emph{exactly zero} (not merely small) gradient to every other site.
This is a sparsity guarantee, not an approximation -- a reviewer-facing,
testable instantiation of \S\ref{sec:consolidation}'s Phase 4 DA
discussion (observation Jacobian is block-diagonal, so each local
analysis only needs its own patch's observations plus the localization
radius).

\textbf{Smoke-validated end to end (2026-10-06).} A degenerate but
genuine run (5-time-unit spinup, 8 ensemble members, 2 cycles of 2
steps, \texttt{p7b\_ks2d} at native size, LETKF + the new 2D taper)
completed with no crashes, no NaNs, and a correctly-shaped report --
confirming the full pipeline (resize-transfer loading, 2D truth
generation, 2D sparse observation, 2D taper, LETKF analysis at
$d_{\mathrm{latent}}{=}9600$) is wired correctly. \textbf{The reported
numbers from that run are meaningless} (5 units is far too short to be
on the attractor) and were discarded; only the real run below counts.

\textbf{Real run: TO DO, not yet launched (cost estimate given, not
run tonight).} Phase 13's headline experiments
(\S\ref{sec:consolidation} Phase 4): $N_{\mathrm{ens}}$ sweep
($\{8,16,32,64,128,256\}$) at native size and at a transferred size
(e.g.\ $66^2$), comparing analysis RMSE against the free run and against
each other (gate: transferred-filter RMSE within 25\% of native).
Dominant costs: the 200-time-unit spinup at $640^2$ (plausibly
30--90 minutes based on this system's own 2D KS data-generation timing
elsewhere in this document) is paid \emph{once} per size and reused
across the whole $N_{\mathrm{ens}}$ sweep; each LETKF analysis call
computes a \texttt{jacrev} decoder Jacobian of shape
$(m, d_{\mathrm{latent}})$ with $d_{\mathrm{latent}}{=}9600$ at native
size (quadrupling at $2\times$) -- untimed at real settings
($n_{\mathrm{cycles}}{\approx}40$, the brief's own convention elsewhere
in this document) and should be timed at one ensemble size before
committing to the full sweep, per this project's own established
discipline for expensive runs. \textbf{Caveat added later the same day
(\S\ref{sec:c19_da_hypotheses}, H2):} the current filter forms dense
$d\times d$ matrices, about $3.7$\,GB each at $66^2$, so the transferred
run probably needs the per-site local LETKF (test T2 there) first.

\subsection{Other code gaps toward a minimal viable paper (2026-10-06)}
\label{sec:mvp_gaps}

Assessed directly in response to ``what else would be needed'':
\begin{itemize}[leftmargin=*]
\item \emph{$N_{\mathrm{ens}}$-sweep and native-vs-transfer comparison
  orchestration.} Not a missing capability -- \texttt{run\_da\_pff\_2d\_ltransfer.py}
  already takes \texttt{--n-ensemble}/\texttt{--old-L}/\texttt{--new-L}; a sweep is
  multiple invocations plus a small results-collection script, not new
  DA machinery. Cheap, to be written alongside the real run above.
\item \emph{A physical-space (non-latent) DA baseline} (plain LETKF or
  SEC directly on the PDE state, with an ensemble of true-PDE
  integrations as the forecast model, instead of the learned
  propagator) -- genuinely new code (a real forecast model replacing the
  learned propagator in the cycling loop) and a real cost driver
  (ensemble-of-PDE-integrations is far more expensive per forecast step
  than the latent propagator). Worth having for the paper (reviewers
  will ask ``how does this compare to not using a learned latent
  space''), but larger than the 2D DA validation above and not required
  to show the \emph{transfer} result, which only needs the learned-model
  filter at two sizes. Lower priority; flagged, not started.
\item \emph{Figures} (Lyapunov spectra overlaid on truth per size, energy
  spectra, DA RMSE-vs-$N_{\mathrm{ens}}$ curves) -- no plotting code
  written this session; straightforward once final numbers are in,
  not a research risk.
\item \emph{1D KS Phase 3 retrain} -- no new code, an existing-script
  training run (\S\ref{sec:consolidation}).
\end{itemize}
\textbf{Net assessment: DA is achievable at reasonable cost, and is
recommended to keep rather than drop}, now that the smoke test confirms
the pipeline works and the block-diagonal Jacobian result gives a
concrete, citable efficiency argument specific to this project's
architecture -- strengthening, not just restoring, the original
DA-focused framing of the whole research program.

\subsection{Kolmogorov $\alpha{=}0.1$ (\texttt{k\_xs01}): the 3$\times$
underestimation, six training variants, and a resolved two-way
hyperparameter interaction (2026-10-07)}
\label{sec:k_xs01_ablation}

\texttt{k\_xs01} (B1, canonical $x$-strip + long-rollout recipe,
$\alpha{=}0.1$) trained cleanly but FAILED the $3\times$ gate on a
consistent, low-variance signature: all three IC seeds
\emph{underestimate} $D_{\mathrm{KY}}$ (model 47.5/46.9/48.5 vs.\ true
56.35), mean 15.4\%, spread only 2.9 points -- a systematic bias, not
noise. 1$\times$ passed its mean bar narrowly (8.9\%, spread 10.9, just
over the 10-point gate); $2\times$ passed cleanly (4.8\%). This matches
the same pervasive-underestimation signature already documented for
$\alpha{=}0.05$ (\S\ref{sec:p9b_results}).

\textbf{Mechanism hypothesis.} Long-horizon pointwise MSE forces the
propagator to match one specific chaotic realization at every lead time
in the training rollout; beyond the predictability horizon the
MSE-optimal prediction is the conditional mean, which has structurally
lower variance (hence lower $D_{\mathrm{KY}}$) than any real trajectory
-- the same variance-collapse/``blurring'' failure documented in
ML-weather-emulation, not fixed by more capacity (a bigger model trained
under the same objective converges to the same optimum). Two pieces of
circumstantial support: the canonical recipe's input-noise reduction
(the document's original hypothesis for the undershoot) was already
applied in \texttt{k\_xs01} and did not fix it; and the long-rollout fix
for P9b's instability ($k_{\max}$ $12\to32$) is the same change that
introduced/worsened the undershoot earlier in the project's timeline.
This motivated a direct test: does reducing $k_{\max}$ (the rollout
length Stage B trains against) shrink the $3\times$ undershoot?

\textbf{\texttt{k\_xs01\_v2} ($k_{\max}$ $32\to16$, $w_{\mathrm{logdet}}$
$0.0035\to0.005$): mechanism confirmed, new side effect.}

\begin{center}
\begin{tabular}{lrr}
\toprule
size & \texttt{k\_xs01} ($k_{\max}{=}32$) & \texttt{k\_xs01\_v2} ($k_{\max}{=}16$) \\
\midrule
$1\times$ & mean 8.9\%, spread 11.0 & mean \textbf{0.9\%}, spread \textbf{2.1} \\
$2\times$ & mean 4.8\%, spread 7.6  & mean 7.7\%, spread \textbf{16.1} \\
$3\times$ & mean 15.4\%, spread 2.9 & mean \textbf{7.6\%}, spread 4.7 \\
\bottomrule
\end{tabular}
\end{center}

$1\times$/$3\times$ (the sizes the mechanism predicted would improve)
both did, substantially. But $2\times$ acquired a new, genuinely
different problem: one seed now \emph{over}estimates by 18\% (the
systematic sign flips) and the cross-seed spread more than doubled --
reducing $k_{\max}$ relieved MSE-damping but reintroduced some of the
instability the long rollout was originally meant to fix.

\textbf{\texttt{k\_xs01\_v3} ($k_{\max}$ $16\to20$, $w_{\mathrm{var}}$
ablation): $k_{\max}$ looked dominant, and a real scale-anchoring bug was
caught.} Splitting the difference on $k_{\max}$ (20) and raising
$w_{\mathrm{logdet}}$ further (0.006) were tested together with removing
$w_{\mathrm{var}}$ (the marginal per-channel variance floor) entirely, to
test whether it is redundant with $w_{\mathrm{logdet}}$ (the full-
covariance log-det barrier, which catches correlation-driven collapse a
marginal floor provably misses). \textbf{$w_{\mathrm{var}}{=}0$ produced a
severe, distinct failure}, not a non-issue: its own diagnostic ($l_{\mathrm{var}}$,
unweighted) went from 831 at step 250 to $2.48\times10^8$ by step 8000 --
six orders of magnitude -- while reconstruction loss stayed flat
(0.014--0.018). This is an encoder/decoder joint rescaling degeneracy:
the decoder's weights shrink to exactly compensate the encoder's latent
blowing up, so reconstruction looks fine while the representation's
absolute scale is meaningless. $w_{\mathrm{var}}$ is load-bearing after
all -- not for collapse prevention (that diagnosis was correct,
$w_{\mathrm{logdet}}$ does subsume it there), but for \emph{scale
anchoring}, a separate job $w_{\mathrm{logdet}}$ has no opinion on (a
log-det barrier only penalizes a degenerate/singular covariance; it has
no preference against all channels growing together). Killed immediately
once confirmed; relaunched with $w_{\mathrm{var}}{=}0.005$ (a quarter of
the original 0.02 -- enough to anchor scale, not enough to dominate
collapse prevention again).

\begin{center}
\begin{tabular}{lrrr}
\toprule
size & \texttt{v2} ($k_{\max}{=}16$) & \texttt{v3} ($k_{\max}{=}20$, fixed $w_{\mathrm{var}}$) \\
\midrule
$1\times$ & 0.9\%, spread 2.1  & 5.2\%, spread 8.3 \\
$2\times$ & 7.7\%, spread 16.1 & 7.2\%, spread 7.0 (fixed) \\
$3\times$ & \textbf{7.6\%}, spread 4.7 & 14.4\%, spread 4.6 (reverted) \\
\bottomrule
\end{tabular}
\end{center}

$w_{\mathrm{var}}{=}0.005$ fixed $2\times$'s blown-out spread exactly as
intended, with no effect on $3\times$. But $k_{\max}$ $16\to20$ reverted
almost all of \texttt{v2}'s $3\times$ gain (7.6\%$\to$14.4\%, nearly back
to the original \texttt{k\_xs01}'s 15.4\%) -- at this point the working
hypothesis was that \textbf{$k_{\max}$, not the $w_{\mathrm{var}}$/$w_{\mathrm{logdet}}$
collapse fixes, is the dominant driver of the $3\times$ underestimation},
since a 4-step rollout-length change ($16\to20$) undid most of a fix
that $w_{\mathrm{var}}$/$w_{\mathrm{logdet}}$ clearly did NOT undo at
$2\times$.

\textbf{\texttt{k\_xs01\_v4} (third $k_{\max}{=}20$ point,
$\mathrm{logdet\_eps}{=}0$): the $k_{\max}{=}20$ finding reproduced, with
a new confound introduced unnoticed.} Before touching
$\mathrm{logdet\_eps}$ (the additive stabilizer in the log-det barrier,
$\mathrm{Cov}(z)+\epsilon I$), $\mathrm{Cov}(z)$'s smallest eigenvalue at
initialization was measured empirically across 10 seeds of this exact
architecture: range $0.0008$--$0.0044$ (mean $0.003$), every seed
positive -- confirming $\epsilon{=}0$ (the bare, unstabilized barrier) is
safe at init for this architecture, and that a genuinely protective hard
floor (e.g.\ $\epsilon{=}-0.01$) would be infeasible at step 1 for every
seed tested, a systematic scale mismatch rather than bad luck on one
seed. \texttt{v4} repeated \texttt{v3}'s $k_{\max}{=}20$/$w_{\mathrm{logdet}}{=}0.006$/
$w_{\mathrm{var}}{=}0.005$ recipe with $\mathrm{logdet\_eps}{=}0$ instead
of the script default ($10^{-3}$) -- intended as a third, independent
$k_{\max}{=}20$ data point to rule out seed noise.

\begin{center}
\begin{tabular}{lr}
\toprule
size & \texttt{v4} ($k_{\max}{=}20$, $\epsilon{=}0$) \\
\midrule
$1\times$ & 3.9\%, spread 5.6 (best 1$\times$ of v2--v4) \\
$2\times$ & 2.9\%, spread 4.5 (best 2$\times$ of v2--v4) \\
$3\times$ & \textbf{11.6\%}, spread 12.2 \\
\bottomrule
\end{tabular}
\end{center}

$3\times$ again landed near \texttt{v3}'s bad result, seemingly
confirming $k_{\max}{=}20$ as the culprit -- but the within-model spread
was wide (seed0 4.3\%, seeds 1/2 16.5\%/13.9\%), and critically,
\textbf{\texttt{v4} changed $\mathrm{logdet\_eps}$ from \texttt{v3} at
the same time it was being read as confirming $k_{\max}$} -- an
unacknowledged second variable riding along with the one being tested.

\textbf{\texttt{k\_xs01\_v5} (first-ever test of $w_{\mathrm{rank}}$,
C7, on $k_{\max}{=}16$): looked like a dramatic win, with a flipped
bias.} $w_{\mathrm{rank}}$ (\texttt{CropPropConfig.w\_rank}, planned
since the C7 experiment but never previously run on any system) is a
one-sided hinge, $w_{\mathrm{rank}}\cdot\mathrm{relu}(\mathrm{ld\_target} -
\mathrm{per\_site\_logdet}(\mathrm{pred}))$, penalizing the free-running
rollout only if it becomes \emph{less} information-rich than the true
encoded data's own per-site log-det. Launched reverting to the
better-known $k_{\max}{=}16$ (per user correction -- $k_{\max}{=}16$, not
20, was always the intended best-$3\times$ baseline to build on) plus
$w_{\mathrm{rank}}{=}0.02$ (the never-run C7-planned value), keeping
\texttt{v3/v4}'s $w_{\mathrm{logdet}}{=}0.006$/$w_{\mathrm{var}}{=}0.005$.

\begin{center}
\begin{tabular}{lrr}
\toprule
size & \texttt{v4} ($k_{\max}{=}20$) & \texttt{v5} ($k_{\max}{=}16$, $w_{\mathrm{rank}}{=}0.02$) \\
\midrule
$1\times$ & \textbf{3.9\%}, spread 5.6  & 8.9\%, spread 7.2 \\
$2\times$ & \textbf{2.9\%}, spread 4.5  & 11.9\%, spread 13.8 \\
$3\times$ & 11.6\%, spread 12.2 & \textbf{2.3\%}, spread \textbf{1.2} (best of any variant) \\
\bottomrule
\end{tabular}
\end{center}

The initial reading: a genuine bias flip, not a clean win -- \texttt{v5}
excellent at $3\times$ (tightest spread of the whole sequence) but now
\emph{over}estimating at $1\times$/$2\times$, where \texttt{v4} was
better. The two obvious candidate explanations were $k_{\max}$ (16 vs.\
20, the variable the whole v2--v4 sequence had been attributing the
$3\times$ effect to) and $w_{\mathrm{rank}}$ itself (the new ingredient).

\textbf{Ablation step 1 -- \texttt{k\_xs01\_v6a} (``v5 minus
$w_{\mathrm{rank}}$''): $w_{\mathrm{rank}}$ never fired, and is
conclusively ruled out.} A single-variable control, holding every other
field at \texttt{v5}'s exact value and setting $w_{\mathrm{rank}}{=}0$,
was launched to isolate $w_{\mathrm{rank}}$'s contribution. (Two
reconstruction-command bugs were caught and fixed before this control
ran cleanly: the first attempt omitted \texttt{v5}'s zero-reach encoder/
decoder architecture --\allowbreak \texttt{site\_mix\_radius=0},
\texttt{n\_site\_mix\_layers=0} on both sides, confirmed directly from
\texttt{v5}'s checkpoint \texttt{.meta.json}, not assumed -- producing an
immediate NaN; the second passed an explicit \texttt{--crop-sites-y 8}
that, combined with the script's \texttt{halo\_sites\_y = halo if
args.crop\_sites\_y else 0} line, forced a nonzero $y$-halo and a
negative free interior, also NaN. Both are documented here as the kind
of reconstruction error that a ``replicate it with the same flags''
ablation is exposed to when the original command isn't saved verbatim --
the fix in both cases was reading the stored checkpoint's resolved
config directly rather than re-deriving it from the training log's
prose.) Once correctly configured, \textbf{every single logged
checkpoint of \texttt{v6a}'s training was bit-identical to \texttt{v5}'s},
not just the final step -- direct proof that $w_{\mathrm{rank}}{=}0.02$'s
hinge term was exactly zero at every step of \texttt{v5}'s entire run:
the free-rollout's per-site log-det never fell below the true data's, so
the penalty never activated and contributed zero gradient throughout.
\textbf{$w_{\mathrm{rank}}$ explains none of \texttt{v5}'s behavior
relative to \texttt{v4}} -- it was inert. This also means
\texttt{v6a}'s evaluation was skipped (it would only reproduce
\texttt{v5}'s numbers exactly) and the real remaining confound,
$\mathrm{logdet\_eps}$ (\texttt{v4} used $\epsilon{=}0$; \texttt{v5}
silently used the script default $\epsilon{=}10^{-3}$, with no flag
carrying that choice forward), still needed isolating.

\textbf{Ablation step 2 -- \texttt{k\_xs01\_v7} ($k_{\max}{=}16$,
$\mathrm{logdet\_eps}{=}0$): the two-way interaction, resolved.} Holding
$k_{\max}{=}16$ fixed (matching \texttt{v5}/\texttt{v6a}) and setting
$\mathrm{logdet\_eps}{=}0$ (matching \texttt{v4}'s tested value, not
\texttt{v5}'s default) isolates $\mathrm{logdet\_eps}$ as the sole
remaining difference from \texttt{v5}.

\begin{center}
\begin{tabular}{lrrr}
\toprule
size & \texttt{v4} ($k_{\max}{=}20$, $\epsilon{=}0$) & \texttt{v5}/\texttt{v6a} ($k_{\max}{=}16$, $\epsilon{\approx}10^{-3}$) & \texttt{v7} ($k_{\max}{=}16$, $\epsilon{=}0$) \\
\midrule
$1\times$ & 3.9\%, spread 5.6  & 8.9\%, spread 7.2  & \textbf{2.7\%}, spread 3.3 (best) \\
$2\times$ & \textbf{2.9\%}, spread 4.5  & 11.9\%, spread 13.8 & 8.5\%, spread 10.9 \\
$3\times$ & 11.6\%, spread 12.2 & \textbf{2.3\%}, spread \textbf{1.2} & 9.8\%, spread 1.5 (tight, high mean) \\
\bottomrule
\end{tabular}
\end{center}

\textbf{This resolves the ablation, and overturns the v2--v4 narrative
that $k_{\max}$ alone is the dominant $3\times$ lever.} Two independent
effects are both real, and interact rather than one dominating:
\begin{itemize}[leftmargin=*]
\item \textbf{$k_{\max}$ (16 vs.\ 20) matters at $3\times$}: both
  $k_{\max}{=}16$ variants (\texttt{v5}, \texttt{v7}) clearly beat
  $k_{\max}{=}20$ (\texttt{v4}) there (2.3\%/9.8\% vs.\ 11.6\%),
  consistent with the original mechanism hypothesis.
\item \textbf{$\mathrm{logdet\_eps}$ matters just as much, with the
  opposite scale-dependence.} Holding $k_{\max}{=}16$ fixed, moving
  $\epsilon$ from its script default ($\approx10^{-3}$, \texttt{v5}) to
  $0$ (\texttt{v7}) \emph{quadruples} the $3\times$ error (2.3\%$\to$9.8\%)
  while roughly \emph{halving} the $1\times$ error (8.9\%$\to$2.7\%) and
  improving $2\times$ (11.9\%$\to$8.5\%).
\item \textbf{No single configuration wins at all three scales.}
  \texttt{v4} wins $2\times$; \texttt{v5} wins $3\times$ (and has the
  tightest spread of any variant); \texttt{v7} wins $1\times$. The
  apparent ``bias flip'' between \texttt{v4} and \texttt{v5} was a
  conflation of two independently-acting, partially-offsetting
  hyperparameters, not a single lever -- a caution directly relevant to
  \S\ref{sec:threats}'s broader point that single-point percent-error
  comparisons this close to the truth's own measurement noise invite
  over-interpreting small configuration changes.
\end{itemize}

\textbf{Status: open, not yet resolved into a new canonical recipe.} No
configuration tested so far wins cleanly at all three scales
simultaneously; which trade-off to accept (if any) is a design decision,
not a further ablation question, and is left for \S\ref{sec:threats}'s
``report scaling, not single-point precision'' framing to absorb rather
than forcing a single-number verdict.

\section{A second paper: multi-field, two-timescale local chaos
(shallow water equations)}
\label{sec:swe_paper2}

User-directed (2026-10-06): a solver for the rotating shallow water
equations (SWE) was built and validated as a candidate third 2D system
(\texttt{ks\_latent/solver/swe.py}, \texttt{SWEConfig}) -- see its own
module/class docstrings for the physics and the ETDRK4-with-matrix-$L$
("integrating factor") numerics. Asked directly whether this is worth
pursuing now or belongs in a later paper, the assessment below was
written before any training on it, following this project's own
pre-registration discipline (\S\ref{sec:consolidation}).

\subsection{What SWE is, and why it is not a small addition to paper 1}

Unlike both existing 2D systems, SWE has (a) THREE coupled prognostic
fields (\texttt{u,v,eta}, not one) and (b) two well-separated
propagation speeds -- the advective speed and the gravity-wave speed
$\sqrt{gH}$, typically much faster. Both are genuinely new axes, not
restatements of anything paper 1 already tests:

\begin{itemize}[leftmargin=*]
\item \textbf{Multi-field capacity.} C19's finding (encoder/decoder
  reach is a capacity substitute, not something locality itself needs)
  should carry over -- feed \texttt{(u,v,eta)} into the same per-patch
  encoder, same zero reach, more channels -- but how many more is an
  empirical question, not a 3x guess: geostrophic balance correlates the
  three fields locally, so the true local degrees of freedom are
  probably less than three independent copies. Answerable by rerunning
  the existing D6/C0 protocols on SWE, not new code.
\item \textbf{Two light-cones, not one.} The localization-radius rule
  (brief \S15: $r_{\mathrm{loc}} \ge \max(\text{encoder field},
  v_*\Delta t)$) now has to satisfy whichever of the two speeds is
  faster. If that is the gravity wave, the required propagator window
  could be materially larger than Kolmogorov flow's own
  (\texttt{window=3}, 2 layers), with knock-on effects on the crop size
  needed for training. This is exactly what the D7 spreading-velocity
  measurement (\texttt{ks\_latent/analysis/spreading.py}, already built
  for KS/Kolmogorov) is for -- unrun on SWE so far.
\item \textbf{A genuinely new DA angle, not yet tested anywhere.} The
  propagator could instead apply the KNOWN linear SWE operator exactly
  (already diagonalized in closed form by the solver itself,
  \S\ref{sec:swe_paper2}'s own module) and let a learned residual handle
  only the nonlinear coupling -- the same idea as the
  \texttt{local\_mlp\_2d\_biotsavart} ablation that gave Kolmogorov's
  propagator the exact Biot-Savart velocity instead of making it relearn
  the relationship. Untested; a clean, specific hypothesis for this
  system alone.
\end{itemize}

\textbf{The honest cost comparison with Kolmogorov flow's own history.}
Finding a genuinely chaotic, well-resolved Kolmogorov-flow config took a
real search (\S\ref{sec:consolidation}'s own citation of
\texttt{KolmogorovConfig}'s docstring: an early naive sweep found the
system apparently stable at every Reynolds number tried, a Benettin-
warmup methodology bug, not physics). \textbf{SWE is at the identical
starting point right now.} The wind-stress forcing added this session
(\texttt{tau0=0.1}, \texttt{SWEConfig}'s own docstring) is an explicitly
UNVALIDATED starting guess, not a confirmed chaotic attractor -- no
Lyapunov spectrum has been computed for this system at all. Before any
L-transfer or architecture question is even askable, SWE needs the full
Kolmogorov-flow arc repeated: a chaos/parameter search, its own D6/D7
measurements, a capacity audit for the multi-field encoder, and a full
crop-training and evaluation pipeline -- on a harder system (three
fields, two timescales) than either existing one.

\subsection{Verdict: paper 2 (or its own paper), not paper 1}

\textbf{Recommendation: park SWE for now.} The code is written, unit-
tested (19 tests, \texttt{tests/unit/test\_solver\_swe.py}, including an
exact match to the known inertia-gravity-wave dispersion relation and a
tangent-linear step validated against finite differences), and reusable
later at no sunk cost. Paper 1 is close to complete without it: three
systems already in hand (1D KS, 2D KS, Kolmogorov flow), the core
thesis (window training transfers, box training collapses) demonstrated
on all of them, a genuine architectural finding (C19) with a DA
efficiency consequence stronger than either direct competitor (ROAD-
EnKF, LAE-EnKF list localization as future work; this project has a
concrete, testable mechanism for it). Adding a fourth system that first
needs its own from-scratch chaos search would delay paper 1 to chase a
result that is not yet known to exist, for a question (`does the
paradigm survive multiple coupled fields and two timescales') that is
better asked, and better reads, as its own paper once paper 1 has
established the baseline claim.

\textbf{What paper 1 needs instead, all cheaper than SWE and already
itemized} (\S\ref{sec:status_plan}, \S\ref{sec:da_2d}): the real 2D DA
run, now actually unblocked by T1/T2
(\S\ref{sec:c19_da_hypotheses}) rather than blocked by them; the
diversity-hypothesis $2\times$ anomaly (\texttt{k\_xs005\_v2}: native
2.5\%, $2\times$ 18.3\%, $3\times$ 8-9\% -- non-monotonic and
unexplained, needs understanding before it can be written up either
way); and the smaller itemized gaps (multi-seed truths, a second
training seed, a few eval seeds).

\section{Threats to validity}
\label{sec:threats}

\begin{itemize}[nosep]
\item \textbf{Single architecture, single seed} for every number
  reported above except where noted (Section~\ref{exp:e4} is the
  proposed fix).
\item \textbf{DA sweep numbers carry $\sim5$--$6\%$ run-to-run noise
  even at a fixed seed} (Section~\ref{exp:e3}'s $L=400$ result
  discovered and quantified this directly), consistent with this
  project's own documented MPS non-determinism. Every qualitative
  conclusion drawn from the DA sweeps has margins well outside this
  band, but individual reported skill/RMSE values should be read as
  indicative to $\sim1$ significant figure of precision beyond the
  first, not exact; a fully rigorous version would average multiple
  seeds per DA point.
\item \textbf{One system.} Every quantitative result in this paper is
  KS. The literature review underlying this project argues a
  multi-system ladder (KS $\to$ Lorenz-96 $\to$ Rayleigh--B\'enard) is
  the stronger publishable structure than a single-system result; this
  paper reports the first, most-developed rung of that ladder.
\item \textbf{The negative result (Section~\ref{sec:negative}) is real
  and should not be minimized}: at the size it was trained on, the
  local field's own localization is not yet shown to beat a strong,
  realistic baseline. The paper's claim is scoped around transfer
  specifically because that is where no comparator can even compete,
  not because the fixed-size comparison was avoided.
\item \textbf{Prior art on spatial localization of KS is pre-neural and
  should be cited, not rediscovered as novel}: Wittenberg \& Holmes
  \cite{wittenberg1999} showed KS dynamics are well localized in both
  real and Fourier space using wavelet projections 25 years before this
  work, and built low-dimensional local models on large domains from
  that observation --- essentially this paper's L-transfer claim, by
  different means. The defensible novelty claim is narrower and more
  recent: a \emph{learned} local latent field, trained end-to-end,
  transfers across domain size and makes a specific, quantitative
  ensemble-DA localization argument concrete. A closely related recent
  finding is worth flagging directly: AROMA \cite{aroma2024} reports KS
  as an explicit \emph{failure case} for local neural-field latents,
  attributing it to decoder reconstruction quality; our reconstruction
  results should be checked against their stated MSE bar
  ($\sim 10^{-10}$--$10^{-12}$ for accurate spectrum reconstruction)
  before any spectral claim in this paper is treated as settled.
\item \textbf{Kolmogorov-flow $D_{\mathrm{KY}}$ numbers carry real,
  measured truth-side uncertainty that is comparable to, or larger than,
  many of the model errors this document reports -- this should
  materially limit how precisely any Kolmogorov $D_{\mathrm{KY}}$-matching
  claim is stated in the paper.} Directly measured, not assumed: three
  independent $2\times$ truth seeds, same protocol (200-unit spinup,
  120-unit averaging, running-average convergence tracked), give
  $46.6/46.6/52.6$ -- a 13\% range from seed alone, with at least one of
  the three (seed 0) still visibly trending at $t{=}120$, not plateaued
  (\S\ref{sec:rescoring_2x}, task A4--A6). The ORIGINAL (40-unit
  averaging) truths spanned $38.7$--$56.8$, a 47\% range at the same
  physical configuration. \textbf{No multi-seed truth exists at all at
  $1\times$ or $3\times$} -- every $1\times$/$3\times$ model evaluation
  in this document is scored against a SINGLE truth realization, with no
  estimate of that single number's own seed-to-seed spread. Given the
  measured $2\times$ spread, a model's 10--15\% deviation from a single
  $1\times$/$3\times$ truth seed cannot currently be distinguished from
  the truth estimate's own finite-time sampling noise -- calling such a
  deviation a clean gate failure (as earlier passes of this document
  did for \texttt{k\_xs01}'s $3\times$ result) overstates the precision
  the measurement actually has. A model's SEED-TO-SEED CONSISTENCY
  (e.g.\ a tight spread across its own 3 IC seeds) remains a genuine,
  trustworthy signal regardless of this issue; whether its central value
  is a real deviation from the true attractor or an artifact of which
  single truth realization it was compared against is not yet
  established without multi-seed truths at every size, not just
  $2\times$. \textbf{By contrast, the 1D KS $D_{\mathrm{KY}}$ numbers in
  this document (Section~\ref{sec:p7_results} and earlier) are the
  reliable ones}: 1D KS's much lower dimensionality and this project's
  own 3-seed cross-check there showed seed-to-seed agreement to
  $0.2214\pm0.0001$ at $L{=}100$ and $0.2222\pm0.0007$ at $L{=}200$ in
  $D_{\mathrm{KY}}/L$ (independent retraining, not just independent
  truth seeds) -- a sub-percent range, not the double-digit-percent
  range seen here. \textbf{Recommendation for the
  paper: lead the quantitative $D_{\mathrm{KY}}$-matching claim with 1D
  KS, where the numbers are trustworthy to the precision stated; present
  Kolmogorov-flow $D_{\mathrm{KY}}$ results with explicit truth-side
  error bars (once multi-seed $1\times$/$3\times$ truths exist) rather
  than as single-point comparisons, and do not claim a specific
  percent-level match there without those error bars in hand.}
\textbf{The more defensible claim for Kolmogorov flow (user-directed,
2026-10-07): match the SCALING, not the point.} Given the truth-side
noise above, asking the model to hit a single $D_{\mathrm{KY}}$ number
to within a few percent is asking more precision of the comparison than
the measurement supports. The claim this project can actually stand
behind is that the model is PREDICTIVE of how $D_{\mathrm{KY}}$ moves as
domain size and drag ($\alpha$) change -- the shape of the scaling
curve, and how that curve shifts between drag regimes -- which is a
real, checkable, and more interesting property than point-matching, and
survives exactly the measurement noise documented above (a consistent
model-side trend across seeds is informative even when any single point
sits inside the truth's own uncertainty band). Testing this needs
truth-side scaling data this project does not yet have at the needed
density: multi-seed truths at $1\times$/$3\times$ (not just $2\times$)
and across $\alpha\in\{0.025,0.05,0.1\}$ (not just $0.05$), all at the
120-unit protocol. Not yet run; see \texttt{docs/paper1\_tasks.md} for
the queued plan once decided.
\end{itemize}

\section{Related work}
\label{sec:related}

Learned finite-difference/local closure models on coarse grids trace to
Bar-Sinai et al.\ \cite{barsinai2019}; the local latent field here is
the same lineage transplanted into a learned bottleneck rather than
physical variables directly. Constante-Amores, Linot \& Graham's
patch-decomposed autoencoder + neural-ODE work
\cite{constanteamores2024,constanteamores2024b} is the closest
architectural relative in the learned-latent literature, with an
important regime contrast: their minimal-flow-unit domains are
non-extensive by construction, which is exactly why a global latent
vector is an adequate choice there and is not here. On the
data-assimilation side, ROAD-EnKF \cite{roadenkf} and recent
learning-enhanced ensemble filters \cite{learningensemblefilters}
report strong latent-EnKF results on KS at both small and large
ensemble sizes; LAE-EnKF \cite{laeenkf} explicitly lists localization
and inflation of a learned latent as future work, which both supports
this paper's novelty claim and dates it --- this is a competitive,
active area.

\begin{thebibliography}{99}

\bibitem{kassam2005} A.-K. Kassam and L. N. Trefethen, ``Fourth-order
time-stepping for stiff PDEs,'' \emph{SIAM J. Sci. Comput.}, 26(4):1214--1233, 2005.

\bibitem{edson2019} R. A. Edson, J. E. Bunder, T. W. Mattner, and A. J.
Roberts, ``Lyapunov exponents of the Kuramoto--Sivashinsky PDE,''
\emph{ANZIAM J.}, arXiv:1902.09651, 2019.

\bibitem{koopman1909} Koopman-operator analysis of the
Kuramoto--Sivashinsky equation reporting $D_{\mathrm{KY}}\approx 23.2$
at $L=100$, arXiv:1909.00076.

\bibitem{anderson2012} J. L. Anderson, ``Localization and
sampling error correction in ensemble Kalman filter data
assimilation,'' \emph{Mon. Wea. Rev.}, 140(7):2359--2371, 2012.

\bibitem{wittenberg1999} A. J. Wittenberg and P. Holmes, ``Scale and
space localization in the Kuramoto--Sivashinsky equation,''
\emph{Chaos}, 9(2):452--465, 1999.

\bibitem{guan2022} Y. Guan, A. Chattopadhyay, A. Subel, and P. Hassanzadeh, ``Stable a posteriori LES of 2D turbulence using convolutional neural networks: backscattering analysis and generalization to higher Re via transfer learning,'' \emph{J. Comput. Phys.} \textbf{458}, 111090 (2022). [To verify before citing.]

\bibitem{frezat2022} H. Frezat, J. Le Sommer, R. Fablet, G. Balarac, and R. Lguensat, ``A posteriori learning for quasi-geostrophic turbulence parametrization,'' \emph{J. Adv. Model. Earth Syst.} \textbf{14}, e2022MS003124 (2022). [To verify before citing.]

\bibitem{kochkov2021} D. Kochkov, J. A. Smith, A. Alieva, Q. Wang, M. P. Brenner, and S. Hoyer, ``Machine learning--accelerated computational fluid dynamics,'' \emph{Proc.\ Natl.\ Acad.\ Sci.\ USA} \textbf{118}(21), e2101784118 (2021). [Learned local corrections to a coarse solver; its generalization claims, including to larger domains, should be checked against the paper before being cited in support of C10.]

\bibitem{barsinai2019} Y. Bar-Sinai, S. Hoyer, J. Hickey, and M. P.
Brenner, ``Learning data-driven discretizations for partial
differential equations,'' \emph{PNAS}, 116(31):15344--15349, 2019.

\bibitem{aroma2024} AROMA: local neural fields for reduced-order
modeling, arXiv:2406.02176.

\bibitem{constanteamores2024} R. Constante-Amores, A. J. Linot, and
M. D. Graham, ``Data-driven prediction of large-scale spatiotemporal
chaos with distributed low-dimensional models,'' arXiv:2410.01238.

\bibitem{constanteamores2024b} R. Constante-Amores, A. J. Linot, and
M. D. Graham, ``Dynamics of a data-driven low-dimensional model of
turbulent minimal pipe flow,'' arXiv:2408.03135.

\bibitem{roadenkf} ROAD-EnKF: a latent ensemble Kalman filter with a
KS test case, arXiv:2301.11961.

\bibitem{learningensemblefilters} Learning Enhanced Ensemble Filters,
arXiv:2504.17836.

\bibitem{laeenkf} LAE-EnKF: latent autoencoder ensemble Kalman filter,
arXiv:2603.06752.

\bibitem{karimipaul2010} M. Karimi and M. R. Paul, ``Extensive chaos in
Rayleigh--B\'enard convection,'' \emph{Phys. Rev. E}, 85(4):046201 (Rapid
Communication), 2010.

\bibitem{egolf2000} D. A. Egolf, I. V. Melnikov, W. Pesch, and R. E.
Ecke, ``Importance of local pattern properties in spiral defect
chaos,'' \emph{Nature}, 404:733--736, 2000.

\bibitem{crossnewell1984} M. C. Cross and A. C. Newell, ``Convection
patterns in large aspect ratio systems,'' \emph{Physica D},
10(3):299--328, 1984.

\bibitem{kennedyohagan2001} M. C. Kennedy and A. O'Hagan, ``Bayesian
calibration of computer models,'' \emph{J. R. Stat. Soc. B},
63(3):425--464, 2001.

\bibitem{mehtaschwab2014} P. Mehta and D. J. Schwab, ``An exact mapping
between the Variational Renormalization Group and Deep Learning,''
arXiv:1410.3831, 2014.

\bibitem{mlrgstatphys2023} Machine Learning Renormalization Group for
Statistical Physics, arXiv:2306.11054, 2023.

\bibitem{mlclosurereview2024} Scientific machine learning for closure
models in multiscale problems: a review, arXiv:2403.02913, 2024.

\bibitem{coarsegrainneuralop2024} Coarse Graining with Neural Operators
for Simulating Chaotic Systems, arXiv:2408.05177, 2024.

\bibitem{zeroshotlargerdomains2026} Zero-shot generalization of
transformer neural operators to larger domains, arXiv:2606.14597, 2026.

\bibitem{falsepromisezeroshot2026} The False Promise of Zero-Shot
Super-Resolution in Machine-Learned Operators, arXiv:2510.06646, 2026.

\bibitem{discrepancyframework2022} Discrepancy Modeling Framework:
Learning missing physics, modeling systematic residuals, and
disambiguating between deterministic and random effects,
arXiv:2203.05164, 2022.

\bibitem{kraichnan1967} R. H. Kraichnan, ``Inertial ranges in
two-dimensional turbulence,'' Phys.\ Fluids 10, 1417 (1967).

\bibitem{smithyakhot1993} L. M. Smith and V. Yakhot, ``Bose condensation
and small-scale structure generation in a random force driven 2D
turbulence,'' Phys.\ Rev.\ Lett.\ 71, 352 (1993).

\bibitem{chertkov2007} M. Chertkov, C. Connaughton, I. Kolokolov, and
V. Lebedev, ``Dynamics of energy condensation in two-dimensional
turbulence,'' Phys.\ Rev.\ Lett.\ 99, 084501 (2007).

\bibitem{boffettaecke2012} G. Boffetta and R. E. Ecke,
``Two-dimensional turbulence,'' Annu.\ Rev.\ Fluid Mech.\ 44, 427--451
(2012), doi:10.1146/annurev-fluid-120710-101240.

\bibitem{danilovgurarie2000} S. D. Danilov and D. Gurarie,
``Quasi-two-dimensional turbulence,'' Phys.-Usp.\ 43, 863--900 (2000),
doi:10.1070/PU2000v043n09ABEH000782.

\bibitem{mishra2015} P. K. Mishra et al., ``Dynamics of reversals and
condensates in 2D Kolmogorov flows,'' arXiv:1412.4959.

\bibitem{ruelle1982} D. Ruelle, ``Large volume limit of the distribution
of characteristic exponents in turbulence,'' Commun.\ Math.\ Phys.\ 87,
287--302 (1982).

\bibitem{egolfgreenside1994} D. A. Egolf and H. S. Greenside, ``Relation
between fractal dimension and spatial correlation length for extensive
chaos,'' Nature 369, 129 (1994).

\bibitem{paul2007} M. R. Paul, M. I. Einarsson, P. F. Fischer, and M. C.
Cross, ``Extensive chaos in Rayleigh--B\'enard convection,'' Phys.\ Rev.\
E 75, 045203(R) (2007), doi:10.1103/PhysRevE.75.045203.

\bibitem{grappinleorat1991} R. Grappin and J. L\'eorat, ``Lyapunov
exponents and the dimension of periodic incompressible Navier--Stokes
flows: numerical measurements,'' J.\ Fluid Mech.\ 222, 61--94 (1991).

\bibitem{plattsirovich1991} N. Platt, L. Sirovich, and N. Fitzmaurice,
``An investigation of chaotic Kolmogorov flows,'' Phys.\ Fluids A 3,
681--696 (1991).

\bibitem{chandlerkerswell2013} G. J. Chandler and R. R. Kerswell,
``Invariant recurrent solutions embedded in a turbulent two-dimensional
Kolmogorov flow,'' J.\ Fluid Mech.\ 722, 554--595 (2013).

\bibitem{kolmogorovregimes2026} M. Y. Vinograd, J. Cullen, and P. Clark
Di Leoni, ``Dimensional regimes in Kolmogorov flow,'' arXiv:2602.08960
(2026).

\bibitem{clearypage2025} A. Cleary and J. Page, ``Estimating the
Reynolds number dependence of the chaotic attractor dimension in
two-dimensional Kolmogorov flow,'' arXiv:2505.21361 (2025).

\bibitem{yamadaohkitani1988} M. Yamada and K. Ohkitani, ``The inertial
subrange and non-positive Lyapunov exponents in fully-developed
turbulence,'' Prog.\ Theor.\ Phys.\ 79, 1265--1268 (1988).

\bibitem{embedology1991} T. Sauer, J. A. Yorke, and M. Casdagli, ``Embedology,'' \emph{J. Stat. Phys.} \textbf{65}, 579--616 (1991).

\bibitem{pathak2018} J. Pathak, B. Hunt, M. Girvan, Z. Lu, and E. Ott,
``Model-free prediction of large spatiotemporally chaotic systems from
data: a reservoir computing approach,'' Phys.\ Rev.\ Lett.\ 120, 024102
(2018).

\bibitem{ozalpmagri2025} E. \"Ozalp and L. Magri, ``Stability analysis
of chaotic systems in latent spaces,'' Nonlinear Dyn.\ 113, 13791 (2025),
arXiv:2410.00480.

\bibitem{vannieuwenburg2017} E. P. L. van Nieuwenburg, Y.-H. Liu, and
S. D. Huber, ``Learning phase transitions by confusion,'' Nature Phys.\
13, 435 (2017), arXiv:1610.02048.

\bibitem{mullerheld2012} C. J. Muller and I. M. Held, ``Detailed
investigation of the self-aggregation of convection in cloud-resolving
simulations,'' J.\ Atmos.\ Sci.\ 69, doi:10.1175/JAS-D-11-0257.1 (2012).

\bibitem{jeevanjeeromps2013} N. Jeevanjee and D. M. Romps, ``Convective
self-aggregation, cold pools, and domain size,'' Geophys.\ Res.\ Lett.\
40, doi:10.1002/grl.50204 (2013).

\bibitem{wing2017review} A. A. Wing, K. Emanuel, C. E. Holloway, and C.
Muller, ``Convective self-aggregation in numerical simulations: a
review,'' Surv.\ Geophys.\ 38, 1173--1197 (2017),
doi:10.1007/s10712-017-9408-4.

\bibitem{patriziorandall2019} C. R. Patrizio and D. A. Randall,
``Sensitivity of convective self-aggregation to domain size,'' J.\ Adv.\
Model.\ Earth Syst.\ 11, 1995--2019 (2019), doi:10.1029/2019MS001672.

\bibitem{wing2020rcemip} A. A. Wing et al., ``Clouds and convective
self-aggregation in a multimodel ensemble of radiative-convective
equilibrium simulations,'' J.\ Adv.\ Model.\ Earth Syst.\ (2020),
doi:10.1029/2020MS002138.

\bibitem{brenowitzbretherton2019} N. D. Brenowitz and C. S. Bretherton,
``Spatially extended tests of a neural network parametrization trained
by coarse-graining,'' J.\ Adv.\ Model.\ Earth Syst.\ 11, 2728--2744
(2019), doi:10.1029/2019MS001711.

\bibitem{wangyuvalogorman2022} P. Wang, J. Yuval, and P. A. O'Gorman,
``Non-local parameterization of atmospheric subgrid processes with
neural networks,'' J.\ Adv.\ Model.\ Earth Syst.\ 14, e2022MS002984
(2022), arXiv:2201.00417.

\bibitem{bryanfritsch2002} G. H. Bryan and J. M. Fritsch, ``A benchmark
simulation for moist nonhydrostatic numerical models,'' Mon.\ Wea.\
Rev.\ 130, 2917--2928 (2002).

\end{thebibliography}

\end{document}
